% !TEX program = pdflatex
\documentclass[UTF8,12pt,a4paper,openany,fontset=none,zihao=false]{ctexrep}
\usepackage[utf8]{inputenc}
\usepackage[T1,OT1]{fontenc}
\usepackage[scaled=0.84]{beramono}
\AddToHook{cmd/ttfamily/before}{\fontencoding{T1}}
\renewcommand{\CJKrmdefault}{gkai}
\renewcommand{\CJKsfdefault}{gkai}
\renewcommand{\CJKttdefault}{gkai}
\AddToHook{file/c70gkai.fd/after}{%
  \DeclareFontShape{C70}{gkai}{m}{it}{<->CJK*gkaiu}{\CJKnormal}%
  \DeclareFontShape{C70}{gkai}{bx}{it}{<->CJKb*gkaiu}{\CJKbold}%
}
\AtBeginDocument{\def\UrlFont{\fontencoding{T1}\ttfamily\small}}
\usepackage{amsmath,amssymb,mathtools,amsthm}
\usepackage{geometry}
\geometry{top=25mm,bottom=25mm,left=22mm,right=22mm,headheight=15pt}
\usepackage{xcolor,booktabs,longtable,array,enumitem}
\usepackage{xurl}
\usepackage[unicode,colorlinks=true,linkcolor=blue,citecolor=blue,urlcolor=blue]{hyperref}
\usepackage{bookmark}
\usepackage{tocloft}
\usepackage{fancyhdr}
\usepackage{microtype}
\usepackage{needspace}
\hypersetup{
  pdftitle={可验证自动化的用途治理、劳动收入与服务市场},
  pdfauthor={Gao Zheng（高政）},
  pdfsubject={精确签发、劳动权利、公共探索、标准化服务与成本回收的条件性经济模型},
  pdfkeywords={可验证自动化;用途治理;劳动收入;探索机会;语义微分原子;任务商;标准化服务;完整成本}
}
\theoremstyle{definition}
\newtheorem{definition}{定义}[chapter]
\newtheorem{theorem}{定理}[chapter]
\newtheorem{lemma}{引理}[chapter]
\newtheorem{proposition}{命题}[chapter]
\newtheorem{corollary}{推论}[chapter]
\newtheorem{remark}{备注}[chapter]
\newenvironment{certificate}{\par\noindent\textbf{证书。}\ignorespaces}{\par}
\renewcommand{\proofname}{证明}
\renewcommand{\bibname}{参考文献}
\renewcommand{\contentsname}{目录}
\renewcommand{\cftchapleader}{\cftdotfill{\cftdotsep}}
\DeclareMathOperator{\im}{im}
\DeclareMathOperator{\rank}{rank}
\DeclareMathOperator{\supp}{supp}
\DeclareMathOperator{\id}{id}
\DeclareMathOperator*{\argmin}{arg\,min}
\DeclareMathOperator*{\argmax}{arg\,max}
\newcommand{\NN}{\mathbb N}
\newcommand{\QQ}{\mathbb Q}
\newcommand{\RR}{\mathbb R}
\newcommand{\EE}{\mathbb E}
\newcommand{\1}{\mathbf 1}
\setlength{\parindent}{0pt}
\setlength{\parskip}{0.45em}
\setlength{\emergencystretch}{3em}
\setcounter{tocdepth}{1}
\setcounter{secnumdepth}{2}
\setlist{topsep=0.35em,itemsep=0.15em,parsep=0pt,leftmargin=2em}
\ctexset{chapter={format=\Huge\bfseries,beforeskip=10pt,afterskip=24pt},
  section={format=\Large\bfseries},subsection={format=\large\bfseries}}
\pagestyle{fancy}
\fancyhf{}
\fancyhead[L]{\small\nouppercase{\leftmark}}
\fancyfoot[C]{\thepage}
\renewcommand{\headrulewidth}{0.3pt}
\renewcommand{\chaptermark}[1]{\markboth{\thechapter\quad #1}{}}
\fancypagestyle{plain}{\fancyhf{}\fancyfoot[C]{\thepage}\renewcommand{\headrulewidth}{0pt}}
\Urlmuskip=0mu plus 1mu

\newlist{semanticsteps}{enumerate}{1}
\setlist[semanticsteps]{label=\textnormal{(\arabic*)},leftmargin=2.4em,
  topsep=0.25em,itemsep=0.15em,parsep=0pt}
\allowdisplaybreaks[2]

\begin{document}
\title{可验证自动化的用途治理\\与劳动收入及服务市场\\
\large 任务充分表示、公共探索与可执行制度修复\\
\large 条件性经济模型与形式证明}
\author{Gao Zheng（高政）}
\date{2026-11-09}
\pagenumbering{roman}
\maketitle
\setcounter{page}{2}

\chapter*{摘要}
\addcontentsline{toc}{chapter}{摘要}

本文在明确的任务域、规则版本、核算期限及资源边界内，研究可验证自动化的用途治理、劳动收入保障、
按需探索机会与标准化服务市场的联合可行性。分析对象包括技术能力、部署权限、劳动请求权、公共服务资格、
实际交付、逐人结算和企业完整成本。制度目标、经验条件与逻辑结论分别登记，持续安全、持续支付和持续经营均由证书推导。

首先，本文区分候选构造、独立执行、验证与资金维持，给出技术阶段转换的充分条件。
实际用途、资格、证据、权限和资源预留构成签发条件；公共申请、拒绝、等待、退款与复核进入同一完整任务状态。
全执行覆盖及单步保持证书保证有限运行安全，实际选择的下降策略和有界调度保证声明服务的按期完成。
直接替代量的约束与总体劳动净效应使用各自的同口径证书。

其次，本文用真实履约、可比较贡献和已落实资金构造正基本报酬与差异报酬，
把自愿进入退出、既得请求权、及时支付和可达复核落实为逐步保持的制度关系。
探索服务按经核验需求在实际预算和容量内分配；劳动评分与私人财富的变化，只有在完整准入与执行条件相应保持时，
才获得精确的非干扰结论。访问采购与服务商收入使用同一交易标识，货币转移与真实资源耗费分别核算。

在表示层，本文构造保留结果与真实类型的最粗静态细化，以及保持全部规定有限后继的最粗动态任务商。
语义微分原子明确每次作用的读取、改变、保持字段和逻辑规则；相同合同下的标准服务实现通过动作响应覆盖及链交换证书，
形成共同观察链商。逐人分配缺陷则由独立的资金修复复形处理：对偶证书识别旧障碍，新增合法方向给出非负修复，
链收缩证明指定缺陷的消除可以保持合法参与的阶段循环。

最后，本文以新增可达任务刻画生产可能集的扩张，以语义保持的复用和完整成本刻画服务市场的可行价格区间。
固定维护合同与真实任务差异分别给出成本回收构造；对不增加服务价值且不再影响收入的投入，证明其最优水平为零。
差异化定价模型提供可核对的正交易与正利润均衡。全部结论在同一资源、资金、时间和权利见证下合取；
制度净优势和最小兼容性分别由共同口径比较及下界证书确认。

\textbf{关键词：}可验证自动化；用途许可；劳动收入；探索机会；标准化服务；任务充分商；语义微分原子；完整成本；有效竞争。

\chapter*{研究范围与阅读结构}
\addcontentsline{toc}{chapter}{研究范围与阅读结构}

本文研究技术能力、实际部署、收入权利和服务供给能否在明确经济目标下共同成立。
公共身份决定申请入口的适用范围，实际用途与作用权限决定执行资格；两者分别登记。
高保障任务与受保护常规任务按照用途和实际控制行为划分，服务申请者的组织类型不替代该划分。

全文采用定义、引理、命题、定理、推论和备注。每个派生结论附证书与证明；
能够由初态、单步经济作用和有限执行得到的性质，均通过数学归纳确认为定理。
证明中的语义原子相对于登记的观察精度、合法作用接口和允许分解定义。
每一步明确前件、代入、逻辑规则、后继谓词及保持字段；
逻辑蕴涵与现实因果识别使用不同的证书。

正文按以下相互衔接的层次展开：
\begin{enumerate}
\item 技术承接与治理：构造来源、运行独立、维护资金、用途许可、公共签发及有限履约。
\item 劳动与探索：直接替代和净劳动效应、真实参与、逐人请求权、按需资源分配及合并核算。
\item 表示与标准：任务类型、动态充分商、观察链商、标准接口和实际状态读回。
\item 修复与生产：逐人侵占障碍、合法追回与补付、参与循环及高保障任务的严格能力扩张。
\item 市场与制度：复用成本、真实需求、成本回收、有效竞争、条件净优势及共同资源总定理。
\end{enumerate}

限定内部资料提供市场参与、真实资格、资金覆盖、兼容扩展、资源修复及任务保持的理论接口
\cite{InternalEconomics,InternalTheory}。
公开研究提供劳动替代与新增任务、行为激励、独立证明检查、监管责任及标准同调结构的文献参照
\cite{AcemogluRestrepo2019,Autor2015,Becker1968,Necula1997,IAEACNS,Hatcher2002}。
标准化服务的接口参照采用明确版本的公开企业平台规范，竞争模型采用公开经济学文献中的可比方法
\cite{JakartaPlatform11,KatzShapiro1994,Hotelling1929}。

最粗性相对于已登记的任务观察及全部后继；最小性相对于已登记的目标、基线和候选族。
有限核算期限不等于状态空间有限：动态分区算法的终止需要有限验证域或独立稳定性证书，
连续需求和价格模型则在各自声明的实数域内证明。
不同模块共享实际资源时，必须提交一个共同可行见证。
完整成本覆盖、有限期限偿付和真实交付分别核验；净经济利润已经扣除正常资本成本。


\clearpage
\tableofcontents
\clearpage
\pagenumbering{arabic}

\chapter{形式语言、经济对象与证书语义}
\label{base:chapter}

\section{声明域与不同类型的前提}

\begin{definition}[共同声明域]
固定
\[
\mathfrak D=(\nu,\mathcal I,\mathcal K,T,\mathcal U,\mathcal B),
\]
其中 \(\nu\) 为规则版本，\(\mathcal I\) 为主体集合，\(\mathcal K\) 为任务集合，\(T<\infty\) 为核算与执行期限，
\(\mathcal U\) 为计量单位及折算规则，\(\mathcal B\) 为资源与风险边界。
同一命题中的资格、权利、资金、转移和结果均使用该声明域。
\end{definition}

\begin{definition}[制度目标、经验条件与派生结论]
制度目标规定被保护的任务与权利、允许的用途、风险容限和比较准则。
经验条件规定需求、能力、预算、核验性能和执行覆盖的实际取值。
派生结论是在这些定义和前件下，经逻辑规则、代数运算或数学归纳获得的命题。
\end{definition}

例如，保护一类真人履约资格是制度目标；有效需求有正下界是经验条件；在需求、份额及生产率上界共同成立时，
人类劳动量具有正下界是派生结论。用途限制同样属于登记的制度情景，不能从某项技术的可核验性直接推导。

\begin{definition}[证书及其适用性]
命题 \(P\) 的证书写为
\[
\mathfrak C_P=(\mathfrak D,\mathcal E_P,\mathcal M_P,\mathcal W_P,\mathcal V_P),
\]
其中 \(\mathcal E_P\) 是证据来源，\(\mathcal M_P\) 是数学表示，\(\mathcal W_P\) 是解、分离函数、策略或执行回执等见证，
\(\mathcal V_P\) 是检查规则。形式检查确认表示内的证明义务；表示有效、数据真实及实际实现相符分别承担现实适用义务。
\end{definition}

\begin{proposition}[证书适用的合取结构]
\label{base:certificate}
设 \(H\) 是命题 \(P\) 的前件，形式检查器在声明语义内满足
\[
\operatorname{Check}(\mathfrak C_P)=1\Longrightarrow (H\Rightarrow P).
\]
若 \(\operatorname{Check}(\mathfrak C_P)=1\)，且 \(\operatorname{ModelValid}\)、\(\operatorname{EvidenceValid}\)、\(\operatorname{ImplementationConforms}\) 及 \(H\) 的实际取值均通过核验，
则可以确认该声明域内的实际命题 \(P\)。
\end{proposition}
\begin{certificate}
模型与实际对象的解释映射、各前件的来源记录、实现符合性证明、形式推导和最终结果回执。
\end{certificate}
\begin{proof}
\begin{semanticsteps}
\item 解释映射和证据有效性将实际对象送入 \(H\) 所描述的模型对象；形式检查给出 \(H\Rightarrow P\)，
以肯定前件规则得到模型中的 \(P\)。
\item 实现符合性和结果解释将该命题传回声明的实际观察。
\item 谓词链为
\[
\begin{aligned}
&\operatorname{EvidenceValid}\land\operatorname{ModelValid}\land H_{\rm actual}
\Longrightarrow H_{\rm model},\\
&H_{\rm model}\land\operatorname{Check}(\mathfrak C_P)=1
\Longrightarrow P_{\rm model},\\
&P_{\rm model}\land\operatorname{ImplementationConforms}
\Longrightarrow P_{\rm actual}.
\end{aligned}
\]
每个解释箭头均是证书中的单独义务，故结论成立。
\end{semanticsteps}
\end{proof}

\begin{remark}[形式正确性与现实适用性]
形式检查器无需承担无法观察的事实判断。需求测量、司法执行、资金到位与环境假设的适用性由各自证据支持。
一个满足代数恒等式的账本可以仍然缺乏实际付款资金；一个有证明的程序也可以因执行环境偏离前提而失去部署许可。
\end{remark}

\section{语义微分原子与类型化谓词推演}
\label{atom:section}

经济制度的细粒度证明需要同时登记事实的可观察分辨率、实际作用的提交边界和逻辑推导的合法规则。
本节把这些对象组织为任务相对的语义微分原子；后文的签发、预留、劳动支付、探索分配和服务结算均采用该证明粒度。
一个原子可以同时改变多个必须一致提交的字段。其原子性相对于登记的任务、执行接口和允许分解，
并不要求任何物理变化都存在与观察无关的绝对最小单位。

\begin{definition}[带类型的经济观察]
\label{atom:typed}
在共同声明域 \(\mathfrak D\) 内，给定观察族
\[
\mathcal O=\{O_k:X\to Y_k:k\in K\}.
\]
每个 \(Y_k\) 同时登记值域、计量单位、主体及合同索引、规则版本、时点和证据精度。
资格观察可取 \(Y_k=\{0,1\}\)，知识状态可取
\(Y_k=\{\mathrm{true},\mathrm{false},\mathrm{unknown}\}\)，
余额、时间和容量分别取各自有单位的数值域。
不同单位之间的运算须具有声明的换算映射；主体不同的请求权使用不同坐标。
\end{definition}

\begin{definition}[经济作用的语义微分]
\label{atom:differential}
对合法作用 \(a\)、实际响应 \(\xi\) 和后继
\(x'=\delta(x,a,\xi)\)，若 \(Y_k\) 嵌入一个声明的加法群，则定义
\[
\partial_a O_k(x,\xi)=O_k(x')-O_k(x).
\]
对一般类型，保存有序变化对
\[
\partial_a^{\rm pair}O_k(x,\xi)=(O_k(x),O_k(x')).
\]
资格从未知到已确认、许可从有效到失效，分别保留其类型化变化对；
不把这些变化与货币增量相加。
作用的观察支集为
\[
S_{a,x,\xi}=\{k\in K:O_k(x')\ne O_k(x)\}.
\]
\end{definition}

\begin{definition}[语义微分原子的证书记录]
\label{atom:record}
一个登记原子记为
\[
\mathfrak a=(\nu,r,a,\mathcal P_a,\mathcal X_a,
\delta_a,\mathcal Q_a,S_a,\mathcal F_a,\mathcal W_a).
\]
这里 \(r\) 为作用及合同标识，\(\mathcal P_a\) 为前置谓词，\(\mathcal X_a\) 为完整响应域，
\(\mathcal Q_a\) 为后置谓词，\(S_a\) 为允许改变的观察字段，
\(\mathcal F_a\) 为其余字段保持的条件，\(\mathcal W_a\) 为权限、资源、时间和实现见证。
它必须交付
\[
\begin{aligned}
&\mathcal P_a(x)\land\xi\in\mathcal X_a(x)
  \Longrightarrow x'=\delta_a(x,\xi),\\
&\mathcal P_a(x)\land\xi\in\mathcal X_a(x)
  \Longrightarrow\mathcal Q_a(x,x'),\\
&k\notin S_a
  \Longrightarrow O_k(x')=O_k(x).
\end{aligned}
\]
定义~\ref{tra:realization} 的实现符合性将该记录连接到真实作用。
对同步更新的余额、义务和最终到账，原子边界位于一致提交点；
未最终结算的外部行为另设待确认状态，不先行签发最终履约事实。
\end{definition}

\begin{definition}[任务相对的不可再分条件]
\label{atom:minimal}
给定有限、完整登记的允许细化候选表 \(\mathcal D_a\)。
若其中没有一个长度至少为二的合法序列，可以同时保持原子的全部响应、观察、累计费用、期限和中间权利约束，
且序列每一步都改变至少一个必保留观察，则称 \(a\) 相对于 \(\mathcal D_a\) 为不可再分原子。
语义不变的内部读取可以作为实现步骤；需要保留的等待、费用或权利变化须进入 \(\mathcal D_a\) 的核验。
\end{definition}

\begin{proposition}[有限候选下的原子性验收]
\label{atom:minimal-check}
若 \(\mathcal D_a\) 已被证明覆盖声明范围内全部允许细化，
且每个候选至少违反不可再分条件中的一项，则 \(a\) 在该范围内不可再分。
\end{proposition}
\begin{certificate}
完整候选表、覆盖范围、每个候选的具体失败谓词与见证。
候选表的空缺不视为排除证据；时间、费用和实际响应也属于候选比较的观察。
\end{certificate}
\begin{proof}
\begin{semanticsteps}
\item 反设存在合格细化 \(d_\star\)。由候选覆盖，
\[
\operatorname{Refines}(d_\star,a)\Longrightarrow d_\star\in\mathcal D_a.
\]
\item 逐候选证书给出
\[
\forall d\in\mathcal D_a,\quad
\neg\operatorname{Refines}(d,a).
\]
\item 对 \(d_\star\) 作全称实例化，得到
\(\neg\operatorname{Refines}(d_\star,a)\)，与第一步矛盾。
\item 排除所有合格细化后，按定义~\ref{atom:minimal} 得到结论。
\end{semanticsteps}
\end{proof}

\begin{definition}[单步逻辑证书]
\label{atom:inference}
固定一个证明的明确前件集 \(\mathcal H\)。第 \(j\) 步逻辑证书记为
\[
\mathfrak p_j=(J_j,\mathsf R_j,P_j,w_j),
\qquad J_j\subseteq\{1,\ldots,j-1\},
\]
其中 \(\mathsf R_j\) 是允许推理规则，\(P_j\) 为结论谓词，\(w_j\) 为相应代数、索引或见证记录。
该步只使用 \(\mathcal H\) 和已完成步骤中的 \(P_k\ (k\in J_j)\)。
允许规则包括定义展开、合取引入及消去、肯定前件、等式替换、域内全称实例化、
携带实际见证的存在引入、相同类型内的代数运算，以及具有基例和归纳步的数学归纳。
存在量词中的见证若用于实际执行，另须保存选择函数与其执行记录。
\end{definition}

\begin{theorem}[有限谓词推演链的健全组合]
\label{atom:proof-induction}
设 \(\mathcal H\) 成立，且每个单步逻辑证书均满足
\[
\mathcal H\land\bigwedge_{k\in J_j}P_k
\Longrightarrow P_j.
\]
则对任意有限 \(m\)，
\[
\boxed{\mathcal H\Longrightarrow\bigwedge_{j=1}^{m}P_j.}
\]
若 \(P_m\) 是最终经济命题，该命题由同一前件集成立。
\end{theorem}
\begin{certificate}
完整前件表、每一步的前序索引、规则名称、类型检查及见证。
所有引用定理同时记录其前件在本证明中的代入，不只引用结论。
\end{certificate}
\begin{proof}
\begin{semanticsteps}
\item 对 \(j=1\)，\(J_1=\varnothing\)，单步证书给出
\(\mathcal H\Rightarrow P_1\)。
\item 假定已经得到 \(P_1,\ldots,P_j\)。由于
\(J_{j+1}\subseteq\{1,\ldots,j\}\)，合取消去得到所有所需前件
\(\bigwedge_{k\in J_{j+1}}P_k\)。
\item 将这些前件与 \(\mathcal H\) 合取，再使用第 \(j+1\) 步证书：
\[
\mathcal H\land\bigwedge_{k\in J_{j+1}}P_k
\Longrightarrow P_{j+1}.
\]
\item 合取引入得到 \(\bigwedge_{k=1}^{j+1}P_k\)。以 \(j\) 作数学归纳，有限终点 \(m\) 的结论成立。
\end{semanticsteps}
\end{proof}

\begin{theorem}[类型化经济增量的有限累计]
\label{atom:telescope}
对同一加法型观察 \(O_k\)，任意已给定的长度 \(n\) 的合法执行满足
\[
\boxed{
O_k(x_n)-O_k(x_0)
=\sum_{j=0}^{n-1}\partial_{a_j}O_k(x_j,\xi_j).
}
\]
\end{theorem}
\begin{certificate}
实际相接的状态序列、每次作用与响应、相同观察单位和每一步微分值。
请求标识或主体索引改变时，先提交相应映射的保持证书。
\end{certificate}
\begin{proof}
\begin{semanticsteps}
\item 长度零时，左侧为零，右侧为空和，等式成立。
\item 假定长度 \(n\) 成立。后继定义给出
\[
\partial_{a_n}O_k(x_n,\xi_n)=O_k(x_{n+1})-O_k(x_n).
\]
\item 将该式加到归纳等式两侧。相同单位内的加法允许合并，且中间项抵销：
\[
[O_k(x_n)-O_k(x_0)]+[O_k(x_{n+1})-O_k(x_n)]
=O_k(x_{n+1})-O_k(x_0).
\]
\item 得到长度 \(n+1\) 的等式，数学归纳完成。
\end{semanticsteps}
\end{proof}

\begin{proposition}[一致支付原子的逐字段保持]
\label{atom:payment}
固定真实权利人 \(i\) 与合同 \(r\)。令 \(B\) 为可履约资产，
\(L=\sum_{j,k}L_{jk}\) 为真实未付义务，\(Y_{ir}\) 为累计最终到账，
\(\Gamma_{ir}=L_{ir}+Y_{ir}\) 为累计成立权利。
在
\[
B\ge L\ge0,\qquad L_{jk}\ge0\ (\forall j,k),\qquad 0\le u\le L_{ir}
\]
且最终性支付证书成立时，登记一致原子
\[
(B,L_{ir},Y_{ir})\mapsto(B-u,L_{ir}-u,Y_{ir}+u).
\]
该原子保持资产覆盖、逐人权利和其他人的既有义务。
\end{proposition}
\begin{certificate}
匹配的主体、合同、金额、原债权、最终到账回执和一次性结算标识；
其余 \(L_{jk},Y_{jk},\Gamma_{jk}\) 的保持表。
\end{certificate}
\begin{proof}
\begin{semanticsteps}
\item \(u\le L_{ir}\le L\le B\) 分别给出
\(B-u\ge0\) 与 \(L_{ir}-u\ge0\)。
\item 只有匹配义务改变，因此 \(L'=L-u\)。相减得到
\[
B'-L'=(B-u)-(L-u)=B-L\ge0.
\]
\item 同一权利的两个坐标同额反向改变：
\[
L'_{ir}+Y'_{ir}=(L_{ir}-u)+(Y_{ir}+u)=\Gamma_{ir}.
\]
\item 对 \((j,k)\ne(i,r)\)，字段保持证书给出
\(L'_{jk}=L_{jk}\)、\(Y'_{jk}=Y_{jk}\)。逐人权利没有被总额相抵删除。
\item 将各字段后置谓词合取，得到所声明的一致支付结论。
\end{semanticsteps}
\end{proof}

\begin{remark}[逻辑蕴涵、真实作用与政策因果识别]
本节的谓词箭头表示在明确前件和结构方程内可核验的逻辑蕴涵。
实际作用箭头还要求时间顺序、实施权限和响应覆盖。
比较制度变动的现实因果效应时，须另外固定对照方案、干预变量和识别条件；
账本恒等式或一个有效证明不自行识别未观察到的反事实。
文中所有持续性结论均由初态、单步规则和进展或资金证书推导为定理。
\end{remark}


\section{任务状态、允许动作与全部响应}

\begin{definition}[完整经济任务]
一个任务为
\[
\mathcal M=(X,X_0,A,\Xi,\delta,\mathcal O,G,U),
\]
其中 \(X\) 是状态集合，\(X_0\) 是允许初态，\(A(x)\) 是可执行动作，\(\Xi(x,a)\ne\varnothing\) 是必须考虑的响应集合，
\(\delta(x,a,\xi)\) 为后继状态，\(\mathcal O\) 是必须保留的观察族，\(G\) 是目标状态集合，\(U\) 为声明的收益或代价观察。
状态至少保留适用任务所需的主体、规则版本、资格、资源、到期权利和时间。
\end{definition}

响应包含登记风险边界内的外部扰动、实施误差及违规尝试；不得只保留执行者偏好的响应。
当动作产生随机后果时，可以在确定性鲁棒域中逐一登记后果，或者另行采用明确概率与损失界的随机任务。
本文的全称响应定理采用前一种语义。

\begin{definition}[允许执行与有限可达]
长度 \(n\) 的允许执行为
\[
x_0\xrightarrow{a_0,\xi_0}x_1\xrightarrow{a_1,\xi_1}\cdots
\xrightarrow{a_{n-1},\xi_{n-1}}x_n,
\]
其中 \(x_0\in X_0\)、\(a_j\in A(x_j)\)、\(\xi_j\in\Xi(x_j,a_j)\)，且
\(x_{j+1}=\delta(x_j,a_j,\xi_j)\)。可达集合由实际可执行路径的终态构成；鲁棒可达要求一个策略对全部登记响应达到目标。
\end{definition}

\begin{theorem}[有限执行不变量]
\label{base:induction}
若谓词 \(I:X\to\{0,1\}\) 满足
\[
\forall x_0\in X_0,\ I(x_0)=1,
\qquad
I(x)=1\land a\in A(x)\Longrightarrow
\forall\xi\in\Xi(x,a),\ I(\delta(x,a,\xi))=1,
\]
则任意允许有限执行的每个状态均满足 \(I\)。
\end{theorem}
\begin{certificate}
初态清单及初态谓词值，覆盖全部动作与响应的单步保持证书，以及实际执行与转移模型的对应记录。
\end{certificate}
\begin{proof}
\begin{semanticsteps}
\item 初态证书给出 \(x_0\in X_0\)；全称初态条件实例化为 \(I(x_0)=1\)。
\item 固定任意已执行前缀 \(x_0,\ldots,x_j\)，以 \(I(x_j)=1\) 为归纳假设。
\item 合法执行记录逐项给出 \(a_j\in A(x_j)\) 与 \(\xi_j\in\Xi(x_j,a_j)\)。
\item 将当前状态及动作代入单步证书，再对实际响应实例化：
\[
I(x_j)=1\land a_j\in A(x_j)
\Longrightarrow\forall\xi\in\Xi(x_j,a_j),\ I(\delta(x_j,a_j,\xi))=1
\Longrightarrow I(\delta(x_j,a_j,\xi_j))=1.
\]
\item 执行符合性给出 \(x_{j+1}=\delta(x_j,a_j,\xi_j)\)。等式替换得到 \(I(x_{j+1})=1\)。
\item 基例和归纳步相接，得到任意有限 \(n\) 的全部前缀均保持 \(I\)。
\end{semanticsteps}
\end{proof}

\begin{corollary}[资源、权限与既有权利的联合保持]
若 \(I=I_{\rm resource}\land I_{\rm scope}\land I_{\rm entitlement}\land I_{\rm time}\)，
且同一允许动作对全部响应同时保持这四个谓词，则其合取沿任意允许有限执行保持。
\end{corollary}
\begin{certificate}
共同动作的四项单步证书，以及资源与时点一致性记录。
\end{certificate}
\begin{proof}
\begin{semanticsteps}
\item 令 \(I_1,I_2,I_3,I_4\) 分别表示资源、用途、权利与时间谓词。初态证书给出 \(\bigwedge_{k=1}^4I_k(x_0)\)。
\item 对同一合法动作及任意登记响应，逐项使用
\(I_k(x)\Rightarrow I_k(\delta(x,a,\xi))\)。
\item 合取引入得到
\[
\bigwedge_{k=1}^4I_k(x)
\Longrightarrow\bigwedge_{k=1}^4I_k(\delta(x,a,\xi)).
\]
\item 该合取满足定理~\ref{base:induction} 的初态和单步条件，应用有限归纳得到结论。
\end{semanticsteps}
\end{proof}

\begin{remark}[安全保持与任务完成]
保持 \(I\) 允许系统持续处于安全但未完成的状态。按期完成需要目标可达性、下降秩、容量与时限的独立证明。
只要不能同时交付这些见证，就不能把不变量定理解释为生产任务已经完成。
\end{remark}

\section{规则版本与受保护任务的保持}

\begin{definition}[受保护基线]
声明的旧制度 \(\Sigma_0\) 中，记 \(\mathcal H_0\) 为须继续保留的合法人类任务、合同义务与市场选择集合。
变更制度 \(\Sigma_1\) 保持该基线，指存在保持观察、资源和响应的实现映射，使 \(\mathcal H_0\) 的允许执行可以在新制度中继续实现。
该要求只对 \(\mathcal H_0\) 生效，不要求新规则保留全部旧自动化部署权限。
\end{definition}

\begin{proposition}[任务保持与用途权限调整相容]
若 \(\mathcal H_0\subseteq\mathcal F(\Sigma_1)\)，并且用途变更只删除与 \(\mathcal H_0\) 不相交的部署事件集合 \(R\)，
则保留人类任务与收缩这些部署权限可以同时成立。
\end{proposition}
\begin{certificate}
受保护任务清单、删除事件清单、\(\mathcal H_0\cap R=\varnothing\) 的检查，以及新制度中的逐任务实现见证。
\end{certificate}
\begin{proof}
\begin{semanticsteps}
\item 由不相交得 \(\mathcal H_0\subseteq\mathcal F(\Sigma_0)\setminus R\)。
\item 逐任务实现见证又给出 \(\mathcal H_0\subseteq\mathcal F(\Sigma_1)\)。
\item 故权限删除不破坏被明确保护的子任务。
\end{semanticsteps}
\end{proof}

\section{局部成立与联合可行}

\begin{definition}[共同资源可行性]
设第 \(j\) 个制度模块需要资源向量 \(r_j\in\RR_+^d\)，共同可用资源为 \(R\)。
联合执行要求一组相容见证，使
\[
\sum_jr_j\le R,
\]
并使所有时间、主体、规则、动作接口和逐项权利相容。公共资金或设备只能按未重复承诺的份额计入。
\end{definition}

\begin{lemma}[局部存在性与共同见证的区别]
\label{base:joint}
对同一个共享变量 \(z\)，\(\bigwedge_j\exists z\,P_j(z)\) 一般不推出 \(\exists z\,\bigwedge_jP_j(z)\)。
若提交一个 \(z_\star\) 满足全部 \(P_j\)，则联合存在性成立。
\end{lemma}
\begin{certificate}
非蕴含的有限反例取 \(P_1(z):z=0\)、\(P_2(z):z=1\)。联合成立的证书为同一 \(z_\star\) 上逐项通过的检查。
\end{certificate}
\begin{proof}
\begin{semanticsteps}
\item 两个等式分别有解，但没有一个数同时等于零和一，故一般蕴含不成立。
\item 对共同见证，由所有 \(P_j(z_\star)\) 以合取规则得到 \(\bigwedge_jP_j(z_\star)\)，再作存在引入。
\end{semanticsteps}
\end{proof}

\begin{remark}[经济解释]
就业保护、参与支付、核验系统和高保障研究可能分别可行，却竞争同一预算、人员或期限。
联合定理必须提交经过合并核算的共同见证。本文中的同调结构、制度优化和现金流定理均服从这一要求。
\end{remark}

\chapter{可验证自动化的阶段转换}
\label{tra:chapter}

\section{构造、运行与市场扩张的不同经济角色}

技术的研发来源、交付后的运行条件与研发成本的回收渠道属于不同对象。
一个候选生成系统可以参与程序、模型和证明的构造；完成验证的成果可以由另一执行系统运行。
研发活动需要资金，也不决定资金必须来自某一种市场扩张方式。
本章分别登记这三类依赖，并给出阶段转换的充分条件。
相应的程序与证书分工、精确任务商及局部兼容要求见理论与经济学基础材料\cite{Necula1997,InternalTheory,InternalEconomics}。

\begin{definition}[生成器、执行器与验证器]
\label{tra:pwv}
固定任务 \(\mathcal M\)、规范 \(S\) 与实现版本 \(v\)。
候选生成器 \(P\) 输出候选程序和证明材料，记为
\[
P(S,d)\longmapsto (w,\pi),
\]
其中 \(d\) 是合法取得的构造资料，\(w\) 是可执行成果，\(\pi\) 是候选证书。
执行器 \(W_v\) 按被确认的程序、输入接口与运行环境实施动作。
独立验证器 \(V\) 检查
\[
V(S,w,\pi,v)\in\{0,1,\bot\},
\]
其中 \(1\) 表示全部登记证明义务通过，\(0\) 表示发现不满足项，\(\bot\) 表示未完成判定。
独立性指接受规则和检查过程能够依据规范及证据复算，不以生成器对自身答案的声明代替检查。
\end{definition}

带证明程序的经典构造表明，接收者可以检查不可信来源提供的程序是否满足指定安全策略\cite{Necula1997}。
本文采用这一有限含义：对指定对象和策略核验交付物。
程序安全、环境适用、任务按期完成与经济可持续分别承担证明义务。

\begin{definition}[来源依赖与运行依赖]
若 \(P\) 的输出实际参与 \(w\) 的构造，记 \(P\rightsquigarrow w\)。
若在声明任务域内，\(W_v\) 完成规定运行必须实时调用 \(P\) 或其不可替代服务，记 \(W_v\leadsto P\)。
运行依赖登记包括远程调用、授权服务、动态模型、必要数据和不可替代的维护接口。
已合法固定且可独立读取的程序、参数与证书属于交付物，不因其历史来源而自动构成实时依赖。
\end{definition}

\begin{proposition}[来源依赖不蕴含运行依赖]
\label{tra:provenance}
在允许保存、检查和独立执行交付物的系统类中，
\[
P\rightsquigarrow w\quad\not\Longrightarrow\quad W_v\leadsto P.
\]
\end{proposition}
\begin{certificate}
给出由 \(P\) 生成的有限程序、通过检查的证书、完整运行依赖清单，
以及关闭 \(P\) 的访问后仍可完成声明任务的执行语义与回放记录。
\end{certificate}
\begin{proof}
\begin{semanticsteps}
\item 取状态 \(x\in\{0,\ldots,m\}\)，动作 \(a\in\{0,1\}\)，响应 \(\xi\in\{0,1\}\)，并令
\[
\delta(x,a,\xi)=\max\{0,x-a(1+\xi)\},\qquad G=\{0\}.
\]
令 \(P\) 输出程序 \(w(x)=\mathbf1\{x>0\}\) 及该有限转移表的检查证书。
\item 保存 \(w\) 后，执行器仅需读取 \(x\)，无需调用 \(P\)。
\item 每个非目标状态对全部响应至少下降一，因而至多 \(m\) 步到达目标。
\item 这个实例同时满足来源依赖与运行独立，构成所述蕴含的反例。
\end{semanticsteps}
\end{proof}

\begin{remark}[独立性的经济含义]
命题~\ref{tra:provenance} 分离了研发服务与交付后生产服务。
它不把能源、设备、维护、许可证和数据成本设为零，也不保证所有任务都能采用同一种架构。
对于确实需要在线生成器的任务，应将该调用及其费用保留在运行合同中。
\end{remark}

\section{从有限检查到全部允许执行的不变量}

\begin{definition}[受控执行与安全回退]
\label{tra:guarded}
令 \(I\subseteq X\) 为安全状态集合。
候选动作 \(a\) 只有在当前状态、规范和版本相符且证书通过时进入执行。
记 \(A_V(x)\subseteq A(x)\) 为通过检查的动作集合。
当候选被拒绝或尚未判定时，系统执行预先认证的回退动作 \(a_{\rm f}(x)\)，使
\[
\forall x\in I,\quad a_{\rm f}(x)\in A(x),\qquad
\forall\xi\in\Xi(x,a_{\rm f}(x)),\quad
\delta(x,a_{\rm f}(x),\xi)\in I.
\]
实际候选生成、检查、等待和回退均占用时间与资源，并纳入状态。
\end{definition}

\begin{lemma}[有限枚举检查的健全性]
\label{tra:finite-check}
设声明模型的状态、动作及响应集合有限，模型转移可精确求值。
检查器接受动作策略 \(\sigma\)，当且仅当逐项验证
\[
\begin{aligned}
X_0&\subseteq I,\\
\forall x\in I,\quad &\sigma(x)\in A(x),\\
\forall x\in I\ \forall\xi\in\Xi(x,\sigma(x)),\quad
&\delta(x,\sigma(x),\xi)\in I.
\end{aligned}
\]
先检查动作合法性，再对该合法动作的非空响应集合检查全部后继。
则该检查对所声明的有限模型具有安全健全性。
\end{lemma}
\begin{certificate}
有限集合的完整清单、确定的转移求值规则、所有单步检查结果及无遗漏的枚举见证。
\end{certificate}
\begin{proof}
\begin{semanticsteps}
\item 若检查接受而存在一个不满足所列条件的初态或三元组，则该项属于完整清单。
\item 精确求值必在该项返回失败，与接受矛盾。
\item 所以接受意味着初态和所有单步条件同时成立。
\item 再由定理~\ref{base:induction}，任意允许有限执行保持 \(I\)。
\end{semanticsteps}
\end{proof}

\begin{definition}[实际实现与模型的符合关系]
\label{tra:realization}
以 \(Z\) 表示实际登记状态，\(\Delta\) 表示实际转移，\(q:Z\to X\) 表示任务解释。
实现符合性要求动作可实施、全部实际响应均有模型响应对应，并且
\[
q\Delta(z,a,\zeta)=\delta(qz,a,\eta(z,a,\zeta)).
\]
这里 \(\eta(z,a,\zeta)\in\Xi(qz,a)\)，且实际初态满足 \(qz_0\in X_0\)。
若物理环境只可得到保守集合包络，应对包络内全部后继证明保持，不能以单个估计值替代。
\end{definition}

\begin{theorem}[带拒绝与回退的有限执行安全保持]
\label{tra:safety}
若初态属于 \(I\)，每个 \(a\in A_V(x)\) 对全部登记响应保持 \(I\)，
回退动作也保持 \(I\)，且实际实现满足定义~\ref{tra:realization}，
则任意候选生成序列造成的允许有限执行均满足 \(qz_t\in I\)。
\end{theorem}
\begin{certificate}
初态证明、通过动作的单步证书、等待及回退证书、实际响应覆盖、实现交换等式和版本绑定记录。
\end{certificate}
\begin{proof}
\begin{semanticsteps}
\item 基例由 \(qz_0\in X_0\subseteq I\) 成立。
\item 假设 \(qz_t\in I\)。
\item 若候选通过，单步证书覆盖实际响应的像，故模型后继属于 \(I\)。
\item 若候选被拒绝或未完成判定，回退证书给出相同结论。
\item 实现交换等式将这两种结论变为 \(qz_{t+1}\in I\)。
\item 于是
\[
\operatorname{Initial}\land\operatorname{StepAllResponses}
\land\operatorname{SafeFallback}\land\operatorname{Conforms}
\Longrightarrow \forall t\le T,\ qz_t\in I.
\]
对执行长度归纳即完成证明。
\end{semanticsteps}
\end{proof}

\begin{remark}[检查与环境的边界]
有限检查证明的是所声明模型的性质。
实际环境超出响应集合、检查器实现错误或版本在检查后改变，都使相应符合性前件失效。
这些字段必须进入运行记录和变更控制，不能由候选生成器的性能分数代替。
\end{remark}

\section{安全健全性与生产活性}

\begin{definition}[完成秩与有界进展]
\label{tra:rank}
给定 \(\rho:I\to\mathbb N\)，要求 \(\rho(x)=0\) 当且仅当 \(x\in G\)。
一个进展动作对全部响应满足
\[
\rho(\delta(x,a,\xi))\le\rho(x)-1\qquad(x\notin G).
\]
非进展步骤须使 \(\rho\) 不增加。
有界进展证书另给出整数 \(K\ge1\)：只要尚未完成，
就在随后至多 \(K\) 个执行步骤内出现一次进展，并保证执行不会提前停滞。
\end{definition}

\begin{theorem}[有界进展下的按期完成]
\label{tra:liveness}
在定理~\ref{tra:safety} 的条件上，若定义~\ref{tra:rank} 成立，
则从 \(x_0\) 出发，至多经过 \(K\rho(x_0)\) 个步骤到达 \(G\)。
若每步总时延至多 \(\bar\tau\)、总资源成本至多 \(\bar c\)，则对应上界为
\[
\tau_{\rm finish}\le K\rho(x_0)\bar\tau,
\qquad C_{\rm finish}\le K\rho(x_0)\bar c.
\]
\end{theorem}
\begin{certificate}
自然数秩、全部响应下的下降和非增加证明、进展间隔上界、含检查及回退的时延与成本上界。
\end{certificate}
\begin{proof}
\begin{semanticsteps}
\item 设初始秩为 \(r_0=\rho(x_0)\)。若 \(r_0=0\)，由目标定义已经完成。
\item 对 \(r_0>0\)，实际调度证书保证每个未完成块至多含 \(K\) 步，且末尾包含被实际执行的进展动作。
\item 每块内非进展动作不增秩，进展动作对全部响应降秩至少一。
按块数 \(m\) 归纳，在尚未完成时有 \(\rho(x_{mK})\le r_0-m\)；
不足 \(K\) 步完成的块可用终止状态补齐观察。
\item 取 \(m=r_0\)，非负性与上界共同给出秩为零。由 \(\rho^{-1}(0)=G\) 得到目标到达。
\item 进展块数不超过 \(r_0\)，因此总步数不超过 \(Kr_0\)。
逐步时间和成本上界分别求和，得到
\(\tau_{\rm finish}\le Kr_0\bar\tau\)、\(C_{\rm finish}\le Kr_0\bar c\)。
\end{semanticsteps}
\end{proof}

\begin{proposition}[安全证书不足以签发完成证书]
\label{tra:safety-not-progress}
保持安全集合不蕴含任何有限完成期限。
\end{proposition}
\begin{certificate}
取 \(I=\{s,g\}\)、\(G=\{g\}\)，初态为 \(s\)，所有候选被拒绝，回退后继始终为 \(s\)。
\end{certificate}
\begin{proof}
\begin{semanticsteps}
\item 所有状态均在 \(I\) 中，故每个有限前缀安全。
\item 然而所有前缀终态都是 \(s\notin G\)，不存在有限完成时点。
\item 因此完成声明必须增加进展、容量及期限条件。
\end{semanticsteps}
\end{proof}

\section{版本更新与既有任务的保持}

\begin{definition}[版本切换证书]
\label{tra:update-certificate}
旧版本和新版本分别有状态空间 \(X_v,X_{v'}\)、安全集合 \(I_v,I_{v'}\) 与已验证执行器。
切换证书给出允许切换域 \(D_{vv'}\subseteq I_v\)、状态迁移 \(J_{vv'}\) 和受保护观察族 \(\mathcal O_H\)，满足
\[
J_{vv'}(D_{vv'})\subseteq I_{v'},
\qquad O_{v'}J_{vv'}=O_v\quad(O\in\mathcal O_H).
\]
观察包括已确认的人类资格、已产生的合同权利、到期义务与必要接入条件。
对须继续履行的旧任务，还须给出从切换后状态出发的合法实现策略，覆盖全部登记响应。
以 \(K_v\) 表示在当前剩余期限、资源和预算下，受保护未完任务具有共同鲁棒履行策略的状态集合。
切换证书还要求 \(J_{vv'}(D_{vv'}\cap K_v)\subseteq K_{v'}\)。
新版本未通过时保留已认证旧版本，或使用已认证退出与清偿流程。
\end{definition}

\begin{theorem}[运行与更新交替时的任务保持]
\label{tra:update-invariance}
设每个版本内部的运行步骤保持 \(I_v\cap K_v\)，并使既有权利按登记规则合法演化，
每次切换满足定义~\ref{tra:update-certificate}，初态属于对应的 \(I_v\cap K_v\)，
且旧任务的实现策略与其他任务共享的资源、时间和预算相容。
则任何有限次运行与切换的交替序列保持安全及已登记的既有权利，
并保留声明旧任务的可实现性。
\end{theorem}
\begin{certificate}
逐版本检查结果、每次迁移前后状态、受保护观察交换式、旧任务策略及共同资源分配见证。
每个运行和切换步骤同时保存未完任务的剩余期限及其可实现性证书。
\end{certificate}
\begin{proof}
\begin{semanticsteps}
\item 按交替序列长度归纳。
\item 运行步骤由所属版本的单步证书保持 \(I_v\cap K_v\) 和受保护观察规定的合法演化。
\item 切换步骤由 \(J_{vv'}(D_{vv'})\subseteq I_{v'}\) 保持安全，
由 \(O_{v'}J_{vv'}=O_v\) 保持切换瞬间的资格与权利数值。
\item 迁移对 \(K_v\) 的保持以及共同资源见证说明迁移后仍有可执行履约策略。
\item 每个步骤均满足同一合取义务，故有限交替序列成立。
\end{semanticsteps}
\end{proof}

\begin{remark}[权利保持与权利冻结]
支付可以依法消灭已清偿债权，主体也可以依合同自愿退出。
受保护观察保存的是相应权利及其合法演化规则，而非要求所有账目永久不变。
更新不能以改名、删字段或删除受损主体的方式，把未履行义务变成不存在。
\end{remark}

\begin{definition}[既有任务的精确表示]
在共同动作和响应标签下，任务表示 \(q:X\to\bar X\) 满足
\[
\begin{aligned}
A(x)&=\bar A(qx),&\Xi(x,a)&=\bar\Xi(qx,a),\\
q\delta(x,a,\xi)&=\bar\delta(qx,a,\xi),&G&=q^{-1}(\bar G),\\
O&=\bar Oq&& (O\in\mathcal O_H).
\end{aligned}
\]
它是对所列义务的精确任务商。
若动作或响应需要改名，必须另外给出双向覆盖的实现对应\cite{InternalTheory}。
\end{definition}

\begin{proposition}[更新中的任务策略复用]
\label{tra:reuse}
若更新前后任务均具有到同一核心 \(\bar{\mathcal M}\) 的精确任务商，
且状态迁移满足 \(q_{v'}J_{vv'}=q_v\)，则核心的有限步鲁棒策略可在两版本实现。
\end{proposition}
\begin{certificate}
两组动作与响应覆盖、后继交换、目标原像、迁移交换及核心策略的有限下降秩。
\end{certificate}
\begin{proof}
\begin{semanticsteps}
\item 对任一版本 \(w\in\{v,v'\}\)，精确任务商给出共同动作、全部响应和
\(q_w\delta_w=\bar\delta q_w\) 的逐步交换证书。
\item 令核心策略在 \(\bar x\) 选择 \(\bar a\)。双向动作覆盖给出实际选择 \(a_w\)，全部实际响应均映向核心响应。
\item 对核心下降秩 \(\bar\rho\)，策略证书和后继交换给出
\[
\bar\rho(q_w\delta_w(x,a_w,\xi))
\le\bar\rho(q_wx)-1
\]
于所有正秩状态和规定响应成立。
\item 迁移交换给出
\(q_{v'}J_{vv'}x=q_vx\)，因此迁移前后核心秩相同。
\item 对完成所需进展步数归纳，核心秩降至零；目标原像式
\(G_w=q_w^{-1}(\bar G)\) 把核心完成事实传回两个版本。
\end{semanticsteps}
\end{proof}

\section{阶段转换的四项成熟条件}

\begin{definition}[四项成熟证书]
\label{tra:maturity}
对固定任务域、版本、环境及期限，登记下列四类证书。
\begin{enumerate}
\item \textbf{构造成熟。}规范可明确表达，交付物可执行，关键性质有通过独立检查的证据。
\item \textbf{运行成熟。}实际环境与任务模型相符，成功率、时限、成本、可用性和安全回退满足预先登记的标准。
\item \textbf{更新成熟。}新版本重新验证，状态迁移保持安全，原人类任务与既有权利通过继续履行检查。
\item \textbf{财务独立。}全部构造、运行、更新、验证、治理及实验室辅助费用具有逐前缀可用资金，
且可行性不依赖受限普通市场的替代扩张收入。
\end{enumerate}
以 \(C_{\rm con},C_{\rm run},C_{\rm upd},C_{\rm fin}\) 表示四项通过谓词，令
\[
\operatorname{Mature}=C_{\rm con}\land C_{\rm run}\land C_{\rm upd}\land C_{\rm fin}.
\]
\end{definition}

运行指标必须保存分母、环境范围、测量窗口及误差处理规则。
在确定性鲁棒任务中，可以使用定理~\ref{tra:liveness} 的界；
在统计任务中，成功率和可用性证书只作其声明分布及置信范围内的结论。
财务独立指对指定收入渠道的独立，不表示没有外部资本、无需持续资金或必然已经盈利。

\begin{proposition}[成熟状态的受控转换]
\label{tra:phase-switch}
令制度只在四项成熟证书同时有效时把任务从受控构造阶段转入认证运行阶段，
并在证书失效时按已认证流程回退。
则所有按此规则完成的转换都有四项共同见证；单项性能改善不充分触发转换。
\end{proposition}
\begin{certificate}
共同任务域和版本、四项证书的有效期、转换记录、失效检测和回退记录。
\end{certificate}
\begin{proof}
\begin{semanticsteps}
\item 转换规则给出
\[
\operatorname{Switch}(x)=1
\Longrightarrow C_{\rm con}(x)\land C_{\rm run}(x)
\land C_{\rm upd}(x)\land C_{\rm fin}(x).
\]
\item 逐项合取消去得到同一 \(x\) 上四项有效证书，故具有共同见证。
\item 若某项 \(C_k(x)=0\)，则合取为假；即使另一项为真，也不能满足转换规则。
\item 失效后的动作采用独立有效的回退规则。
\(I(x)\land\operatorname{FallbackAllowed}(x)\Rightarrow I(x')\)
来自回退单步证书，再由定理~\ref{base:induction} 保持有限执行安全。
\end{semanticsteps}
\end{proof}

\section{全成本核算与逐前缀融资}

\begin{definition}[受限收入渠道之外的融资账本]
\label{tra:finance}
令 \(t=1,\ldots,T\)，初始可用资金为 \(F_0\ge0\)，时点 \(t\) 已到账且可支配的许可来源资金为 \(r_t\ge0\)。
在每个现金流发生事件后检查账本，使期内先后顺序也进入前缀。
全部实际支出写为
\[
c_t=c_t^{\rm build}+c_t^{\rm run}+c_t^{\rm update}
     +c_t^{\rm verify}+c_t^{\rm govern}+c_t^{\rm lab}.
\]
最后一项包含仍使用的候选生成、实验室搜索及相应计算费用。
资金余额和尚须保留的到期义务准备分别为
\[
B_t=F_0+\sum_{s=1}^t(r_s-c_s),\qquad L_t\ge0,
\]
并以 \(L_0\) 登记初始承诺。
同一资金不能重复抵押；未来融资意向、未经实现的估值和受限扩张收入不计入 \(r_t\)。
\end{definition}

\begin{theorem}[逐前缀融资的精确下界]
\label{tra:prefix-finance}
对给定的有限现金流与准备金路径，定义
\[
F_0^{\min}=\max_{0\le t\le T}
\left\{L_t+\sum_{s=1}^t(c_s-r_s)\right\},
\]
其中空和为零。
则全部前缀 \(B_t\ge L_t\) 成立，当且仅当 \(F_0\ge F_0^{\min}\)。
\end{theorem}
\begin{certificate}
全部费用及到账时点、资金使用限制、未清偿义务、现金流顺序、逐前缀余额和最大缺口。
\end{certificate}
\begin{proof}
\begin{semanticsteps}
\item 对任意 \(t\)，将余额式移项得
\[
B_t\ge L_t
\quad\Longleftrightarrow\quad
F_0\ge L_t+\sum_{s=1}^t(c_s-r_s).
\]
对有限个 \(t\) 同时要求右式，等价于大于等于这些数的最大值。
\item \(t=0\) 项为 \(L_0\ge0\)，所以该下界也覆盖初始偿付能力。
\end{semanticsteps}
\end{proof}

\begin{corollary}[不确定成本下的保守融资证书]
\label{tra:robust-finance}
若全部登记情景 \(\omega\) 中有 \(c_t(\omega)\le\bar c_t\)、\(r_t(\omega)\ge\underline r_t\)、
\(L_t(\omega)\le\bar L_t\)，并且
\[
F_0\ge\max_{0\le t\le T}
\left\{\bar L_t+\sum_{s=1}^t(\bar c_s-\underline r_s)\right\},
\]
则每个登记情景的全部前缀均有偿付准备。
\end{corollary}
\begin{certificate}
情景集合、费用上界、实际可用收入下界、义务上界及其共同可实现性。
\end{certificate}
\begin{proof}
\begin{semanticsteps}
\item 固定任意登记情景 \(\omega\) 和前缀 \(t\)。逐时点界给出
\(c_s(\omega)-r_s(\omega)\le\bar c_s-\underline r_s\)。
\item 同一货币单位内逐项相加，再加义务上界，得到
\[
L_t(\omega)+\sum_{s=1}^{t}(c_s(\omega)-r_s(\omega))
\le\bar L_t+\sum_{s=1}^{t}(\bar c_s-\underline r_s).
\]
\item 右侧不超过证书最大值，证书最大值又不超过 \(F_0\)。传递性给出每个真实前缀的缺口均被覆盖。
\item 使用定理~\ref{tra:prefix-finance} 的等价式得到
\(B_t(\omega)\ge L_t(\omega)\)。最后对任意 \(t,\omega\) 作全称引入。
\end{semanticsteps}
\end{proof}

\begin{remark}[融资与总期盈利]
最后一期累计收益为正，仍可能存在中途无钱支付工资、计算、验证或清偿义务的前缀。
因此总期盈利不能代替逐前缀条件。
同样，有限期限内可融资不推出无限期限永续，后续期限必须重新登记资金与义务。
\end{remark}

\section{普通市场替代扩张的非必要性}

\begin{definition}[不依赖受限扩张的转换路径]
记 \(R_t^{\rm ord}\) 为因新增受限普通市场替代部署而取得的收入。
一条阶段转换路径 \(p\) 不依赖该扩张，是指其真实可行账本在
\[
R_t^{\rm ord}=0\qquad(0\le t\le T)
\]
时仍通过构造、运行、更新与融资的共同检查。
研究资金、既有合法收入或许可用途合同可以作为候选来源，但必须经过到位与不重复承诺核验。
\end{definition}

\begin{theorem}[有共同见证的受限用途转换路径]
\label{tra:transition-existence}
设存在一个有限任务族及共同见证，使其具有已验证可执行成果、运行完成证书、
版本更新保持证书，以及在 \(R_t^{\rm ord}=0\) 时满足前缀下界的许可资金。
若资源、时限和用途条件相容，则存在不依赖受限普通市场替代扩张的成熟转换路径。
\end{theorem}
\begin{certificate}
同一任务族上的程序与证明、实际运行合同、版本迁移、保护义务、共同资源分配，
以及删除受限扩张收入后仍成立的逐前缀账本。
\end{certificate}
\begin{proof}
\begin{semanticsteps}
\item 按证书给出的构造计划产生并检查交付物，再按运行合同部署于许可任务。
\item 运行步骤使用定理~\ref{tra:safety} 与定理~\ref{tra:liveness}。
\item 发生更新时使用定理~\ref{tra:update-invariance}；发生现金流时使用定理~\ref{tra:prefix-finance}。
\item 共同资源与时限见证使这些步骤可以属于同一条路径，而非互相排斥的分别方案。
\item 四项成熟条件因此在所登记的转换时点同时成立，且全过程未使用 \(R_t^{\rm ord}\)。
\end{semanticsteps}
\end{proof}

\begin{proposition}[转换路径类非空的有限模型见证]
\label{tra:nonempty}
允许有限初始研究资金的模型类中，定理~\ref{tra:transition-existence} 的前件可以同时成立。
\end{proposition}
\begin{certificate}
采用命题~\ref{tra:provenance} 的有限任务，固定许可用途，
把既有人类权利作为不被程序写入的独立状态坐标。
以恒等迁移作为保持版本的合法更新，以覆盖全部有限成本和义务准备的初始资金融资。
\end{certificate}
\begin{proof}
\begin{semanticsteps}
\item 有限转移表可完全检查，程序在每个非目标状态降低自然数秩，
故构造成熟与所声明模型内的运行期限证书成立。
\item 恒等迁移保持全部权利观察；更新规则只接受重新验证的成果，因此更新成熟的规则可满足。
\item 令 \(r_t=R_t^{\rm ord}=0\)，并取
\[
F_0\ge\sum_{t=1}^T c_t+\max_{0\le t\le T}L_t.
\]
该数有限且覆盖全部前缀缺口，所以财务条件同时成立。
\item 所有步骤使用同一程序、资源及账本，构成一个共同见证。
\end{semanticsteps}
\end{proof}

\begin{corollary}[市场扩张不是逻辑必要条件]
\label{tra:not-necessary}
在包含上述见证的可行路径类中，
\[
\forall p\,[\operatorname{Mature}(p)\Rightarrow
\exists t\ R_t^{\rm ord}(p)>0]
\]
为假。
\end{corollary}
\begin{certificate}
一条 \(\operatorname{Mature}(p_\star)\) 成立而全部 \(R_t^{\rm ord}(p_\star)=0\) 的路径。
\end{certificate}
\begin{proof}
\begin{semanticsteps}
\item 把 \(p_\star\) 代入所列全称命题，其前件为真、后件为假，故全称命题不成立。
\end{semanticsteps}
\end{proof}

\begin{remark}[存在性与现实阶段]
有限模型见证排除了“必须先持续替代普通劳动才能取得可验证成果”的无条件必要性。
对某个实际任务能否转换，仍须交付其四项成熟证书。
简单任务的见证不替代复杂任务的构造证明；研究阶段和认证运行阶段可以在不同任务中并存。
\end{remark}

\chapter{用途准入与跨境规则保持}
\label{gov:chapter}

\section{作为登记情景的用途政策}

本章将用途边界作为明确选择的制度情景。
技术成果通过验证，使某些任务具备可实施条件；是否允许某次部署还取决于权利、资源与公共决策。
两种判断分别登记，使同一技术的构造研究、许可运行及普通市场部署可以采用不同规则。

\begin{definition}[部署事件与实际作用]
\label{gov:event}
一个部署事件记为
\[
e=(i,j,t,u,\mu,\chi,v,\ell,\pi),
\]
其中 \(i\) 是责任主体，\(j\) 是受影响对象或市场，\(t\) 是时点，\(u\) 是实际用途，
\(\mu\) 是软件、具身设备、远程服务或代理执行等实现方式，\(\chi\) 是作用与控制接口，
\(v\) 是技术及规则版本，\(\ell\) 是许可，\(\pi\) 是可核验材料。
实际用途由输出如何进入生产、交易或物理执行决定，不能只由服务名称或合同标题决定。
\end{definition}

\begin{definition}[受控范围、许可用途与传统工具基线]
\label{gov:scope}
给定版本 \(\nu\)，\(C_\nu(e)\in\{0,1\}\) 判断事件是否落入登记的高能力自主自动化范围。
\(\mathcal U_{\rm s}\) 是许可的高保障用途集合，\(\mathcal U_{\rm o}\) 是受限普通用途集合，二者不相交。
\(\mathcal B_0\) 是应继续保留的传统工具和已登记合法人类任务，且逐项确认 \(C_\nu(e)=0\)。
分类要求公开的功能、控制强度与实际作用标准；技术可验证性不自行决定其用途类别。
\end{definition}

高保障用途可以按明确任务登记，例如高后果设施的安全维护、基础研究或经许可的公共服务。
该枚举是情景变量，不由产业名称一次性决定。
同一企业可以同时具有许可研究事件、许可服务事件与被拒绝的替代部署事件。

\begin{definition}[严格用途准入情景]
\label{gov:strict-policy}
令 \(A_0(e)\) 表示范围外事件沿既有规则可以实施，\(Z_\nu(e)\) 表示责任、资金及必要程序条件通过。
定义
\[
\begin{aligned}
\operatorname{Admit}_\nu(e)\Longleftrightarrow{}&
[\neg C_\nu(e)\land A_0(e)]\\
&\lor[C_\nu(e)\land u(e)\in\mathcal U_{\rm s}
\land\operatorname{Licence}_\nu(e)
\land\operatorname{Verify}_\nu(e)\land Z_\nu(e)].
\end{aligned}
\]
在此情景下，属于 \(\mathcal U_{\rm o}\) 的受控部署不会因技术可验证而自动取得准入。
传统工具是否继续允许，按声明基线和既有合法条件判断。
\end{definition}

\begin{proposition}[技术验证与用途许可的逻辑分离]
\label{gov:verification-permission}
在严格用途情景中，存在 \(\operatorname{Verify}_\nu(e)=1\) 而 \(\operatorname{Admit}_\nu(e)=0\) 的事件。
\end{proposition}
\begin{certificate}
取已验证的受控执行程序及一个实际用途属于 \(\mathcal U_{\rm o}\) 的部署事件。
\end{certificate}
\begin{proof}
\begin{semanticsteps}
\item 受控事件使准入式第一分支为假，\(\mathcal U_{\rm o}\cap\mathcal U_{\rm s}=\varnothing\) 使第二分支为假。
\item 因此验证为真不改变准入为假的结果。
\end{semanticsteps}
\end{proof}

\section{实现方式变化与用途规则一致性}

\begin{definition}[用途事实标准化]
\label{gov:normalization}
令 \(n_\nu(e)\) 提取执行者权限、实际受影响任务、作用强度、责任、许可和证书等判定字段。
实现方式可以变化，但标准化过程必须覆盖这些变化造成的责任与风险差异。
规则可因子化为
\[
\operatorname{Admit}_\nu(e)=\bar A_\nu(n_\nu(e)).
\]
若两事件的风险或责任不同，标准化值应当不同，不能为取得同一决定而删除必要字段。
\end{definition}

\begin{theorem}[实现改名不改变同一用途的准入]
\label{gov:medium-invariance}
对任意事件变换 \(T\)，若 \(n_\nu(Te)=n_\nu(e)\)，则
\[
\operatorname{Admit}_\nu(Te)=\operatorname{Admit}_\nu(e).
\]
对任意有限串均满足该条件的变换，准入结论保持。
\end{theorem}
\begin{certificate}
软件服务、具身设备、远程调用或代理安排的实际控制链，及其全部判定字段的相等见证。
\end{certificate}
\begin{proof}
\begin{semanticsteps}
\item 单步由因子化直接得到
\(\bar A_\nu(n_\nu(Te))=\bar A_\nu(n_\nu(e))\)。
\item 长度一为基例；将长度 \(k\) 的保持式与第 \(k+1\) 次变换的保持式复合，
得到任意有限串的结论。
\end{semanticsteps}
\end{proof}

\begin{remark}[覆盖义务]
该定理要求识别实际事件并正确取得判定字段。
将机器人登记为租赁设备、将远程控制登记为咨询服务或插入代理人，都不能在字段相同的情形改变判定。
若事实识别或执行覆盖失败，应计入监管失效，不能由标准化公式直接宣布其已经解决。
\end{remark}

\begin{proposition}[全执行覆盖下的用途边界保持]
\label{gov:complete-gate}
若每一个实际实施的受控事件都经过同一版本的有效准入检查，
且无未登记执行路径，则严格用途情景中的受控实际事件均属于 \(\mathcal U_{\rm s}\)。
\end{proposition}
\begin{certificate}
全部执行入口的覆盖证明、许可与版本记录、实际作用回执及无旁路的实现证书。
\end{certificate}
\begin{proof}
\begin{semanticsteps}
\item 任取一个实际受控事件。
\item 覆盖条件给出 \(\operatorname{Admit}_\nu(e)=1\)。
\item 由于 \(C_\nu(e)=1\)，第一准入分支不成立，故第二分支成立，得到 \(u(e)\in\mathcal U_{\rm s}\)。
\item 该论证对每个实际事件成立；逐事件组成任何有限执行时，用途边界保持。
\end{semanticsteps}
\end{proof}

\section{基线保留与兼容变更支集}

\begin{definition}[用途变更的依赖支集]
\label{gov:support}
设现行制度坐标有限集为 \(\mathcal J\)，基线为 \(\Sigma_0\)，已确认的失灵记录为 \(\mathcal F\)。
直接修复坐标为 \(R(\mathcal F)\)，经预算、技术或法律关系证明需要联动的坐标为 \(\operatorname{Dep}(R)\)。
兼容的用途变更要求
\[
\Delta(\Sigma;\Sigma_0)\subseteq R(\mathcal F)\cup\operatorname{Dep}(R),
\]
并在共同资源条件下保持 \(\mathcal B_0\) 的合法执行。
\end{definition}

该支集约束把用途治理接入最小兼容制度选择\cite{InternalEconomics}。
候选族、权重、非空性及相对极小性的统一处理见第~\ref{opt:chapter}~章。
实际需要联动的传统工具应登记为必要依赖，并与明确保留的基线分别核验。

\section{跨境部署的共同任务与逐步交换}

跨境规则保持需要明确可携带的证书、承认范围和责任接口。
一个地方认可技术证明，不表示自动接受另一地方的全部用途权限。
以下结论只对各方明确登记的共同任务域成立。

\begin{definition}[共同任务域上的规则运输]
\label{gov:transport}
设辖区 \(j,k\) 的任务系统为 \(\mathcal M_j,\mathcal M_k\)。
运输记录包括状态映射 \(F_{jk}\)、动作对应 \(\alpha_{jk,x}\) 与响应对应 \(\beta_{jk,x,a}\)。
在共同任务域内，动作与全部响应对应为双射，并且
\[
\begin{aligned}
F_{jk}(X_{0,j})&\subseteq X_{0,k},\\
F_{jk}\delta_j(x,a,\xi)
&=\delta_k(F_{jk}x,\alpha_{jk,x}a,\beta_{jk,x,a}\xi),\\
O_kF_{jk}&=O_j,\qquad G_j=F_{jk}^{-1}(G_k),\\
\operatorname{Admit}_k(F_{jk}x,\alpha_{jk,x}a)
&=\operatorname{Admit}_j(x,a).
\end{aligned}
\]
观察 \(O\) 包括用途、责任、已有权利、计量单位以及声明的费用和时限。
上述各集合和等式均限制在登记共同域，且该域对所讨论的允许动作与响应保持。
本定义不要求两个辖区在共同域之外具有相同规则。
\end{definition}

\begin{theorem}[跨境规则的有限执行保持]
\label{gov:finite-transport}
若定义~\ref{gov:transport} 成立，任一共同域内的有限合法执行都可按所给映射运输，
保持准入、规定观察和完成判定。
\end{theorem}
\begin{certificate}
初态对应、逐状态动作及全部响应的双向覆盖、后继交换、权利观察、目标原像和许可等价。
\end{certificate}
\begin{proof}
\begin{semanticsteps}
\item 取源执行 \(x_{r+1}=\delta_j(x_r,a_r,\xi_r)\)，令 \(y_0=F_{jk}x_0\)。
\item 若 \(y_r=F_{jk}x_r\)，由许可等价与动作覆盖，\(\alpha_{jk,x_r}a_r\) 在目标系统合法。
\item 响应对应覆盖全部目标响应，后继交换给出
\[
y_{r+1}
=\delta_k(y_r,\alpha_{jk,x_r}a_r,\beta_{jk,x_r,a_r}\xi_r)
=F_{jk}x_{r+1}.
\]
归纳得到全部状态对应。
\item \(O_kF_{jk}=O_j\) 保持逐状态观察，目标原像式保持完成判定。
\end{semanticsteps}
\end{proof}

\begin{corollary}[有限辖区链的复合保持]
\label{gov:composition}
若每个相邻运输均满足定义~\ref{gov:transport}，且中间表示和版本一致，
则有限复合 \(F_{n-1,n}\cdots F_{01}\) 保持共同任务域的有限执行。
\end{corollary}
\begin{certificate}
相邻运输证书、中间状态类型和版本一致性，以及动作、响应对应的复合。
\end{certificate}
\begin{proof}
\begin{semanticsteps}
\item 长度一由定理~\ref{gov:finite-transport} 成立。
\item 假设前 \(n\) 段保持，末段的后继交换将已运输轨迹再次送入合法轨迹。
\item 观察等式、目标原像和许可等价逐项复合，双射的复合仍为双射。
\item 于是长度 \(n+1\) 成立，由归纳覆盖有限链。
\end{semanticsteps}
\end{proof}

\begin{remark}[单向认可与真实执行]
若接收方许可更严格，只能对通过接收方条件的子域给出运输，不能声明完整双向等价。
数学运输也不产生执法权、付款能力或条约义务。
跨境执行仍须提交有效责任主体、证书承认、实际资源和权利救济的制度记录。
\end{remark}

\section{核安全制度的可比较责任结构}

核安全公约提供一种已有的国际制度结构：国内许可与检查、监管职能分离、持证者责任、
安全优先、全生命周期评估，以及国家报告和缔约方审议\cite{IAEACNS}。
这些结构可以用于设计高后果技术的责任模型，相关法律效力仍按公约本身的适用范围判断。

\begin{center}
\begin{tabular}{p{0.20\textwidth}p{0.69\textwidth}}
\toprule
公约条款 & 可以抽取的制度关系 \\
\midrule
第 7 条 & 安全规则、许可、检查、评估及许可条件的执行。\\
第 8 条 & 监管机构具备必要能力和资源，监管职能与促进或使用核能的职能有效分离。\\
第 9 条 & 核设施安全的首要责任由相应持证者承担。\\
第 10 条 & 相关组织的政策给予核安全适当优先地位。\\
第 14 条 & 建造和运行前及全生命周期的安全评估、更新与核验。\\
第 5、20 条 & 提交履约报告，缔约方通过审议会议讨论并请求澄清。\\
\bottomrule
\end{tabular}
\end{center}

\begin{definition}[许可责任结构的比较映射]
\label{gov:nuclear-map}
对一个拟议自动化制度，分别登记规范、许可、运行主体、独立检查、持续更新和报告审议接口。
若映射保持“谁可许可、谁执行、谁核查、谁承担责任、何时更新”的关系，
则称其与上述核安全制度具有声明层次上的责任结构对应。
该对应不包含技术对象相同、全部风险相同或法律适用范围相同的断言。
\end{definition}

\Needspace{17\baselineskip}
\begin{proposition}[制度结构对应不推出法律同一]
\label{gov:no-legal-identity}
存在许可责任结构的比较映射，不蕴含核安全公约已对自动化用途作出限制。
\end{proposition}
\begin{certificate}
两套制度的对象类型、各自适用范围，以及映射实际保留的关系清单。
\end{certificate}
\begin{proof}
\begin{semanticsteps}
\item 记 \(S(F,A)\) 为两制度存在共同许可责任结构，\(L(A)\) 为某公约在其法律适用域中直接规制对象 \(A\)。
\item 结构证书只提供 \(S(F,A)\)；公约适用域由其条文给定，属于另一谓词。
\item 取一个具有许可、检查、责任安排但不属于该公约适用对象的制度模型。
该模型满足 \(S(F,A)=1\) 且 \(L(A)=0\)。
\item 此共同结构与异适用域见证使
\(S(F,A)\Rightarrow L(A)\) 不成立。实际用途规则须由其独立适用依据登记。
\end{semanticsteps}
\end{proof}

\begin{remark}[用途治理的经济任务]
用途准入、持续核验和跨境责任可把技术能力转化为边界明确的生产服务。
常规市场中的人类劳动量、收入份额及参与机会是否得到保护，还取决于需求、替代弹性、
执行覆盖和可兑现分配等经济条件。
这些结果不能单由许可结构或国际比较映射推出，必须与后续经济模块使用同一组资源和权利见证。
\end{remark}

\chapter{公共服务的签发与履约治理}
\label{integrated:public-services}
\label{svc:chapter}

公共服务的交付包含申请、核验、资源承诺、实际作用以及后续履约。
同一交付结果可以来自不同权限、成本和责任条件，
而尚未交付的申请也可能已经产生退款、说明理由或复核义务。
本章把这些对象登记为有类型的状态和原子语义转移，
将第~\ref{tra:chapter}~章的核验与运行条件
接入第~\ref{gov:chapter}~章的用途准入规则。
证明携带程序只为指定程序与策略之间的检查关系提供基础；
实际用途、履约和资源覆盖分别承担证书义务\cite{Necula1997}。
各证明的有类型分步推演以定理~\ref{atom:proof-induction}~为组合规则。

\section{申请状态与原子语义}

\begin{definition}[有界服务声明域]
\label{svc:domain}
固定有限主体集、请求标识集、任务类型集、输入编码集、
规范与实现版本集，以及有限核验窗口。
金额采用固定精度，预算、配额、计时和日志长度各有登记上界。
以 \(X\) 表示满足这些类型与上界的服务状态集合。
实际对象集合记为 \(Z\)，解释映射记为
\[
\theta:Z\longrightarrow X.
\]
有限性来自完整字段的有界编码；
仅有有限个服务阶段不构成 \(X\) 有限的证明。
若以有限抽象处理无界对象，
该抽象另须满足定义~\ref{tra:realization}~的实际响应覆盖与后继交换。
\end{definition}

\begin{definition}[完整服务状态]
\label{svc:state}
一个状态 \(x\in X\) 包含全体在册请求以及共享资源。
请求 \(r\) 的登记内容写为
\[
\begin{aligned}
x_r={}&
(\phi_r,i_r,\mathcal D_r,u_r,\sigma_r,v_r,S_r,\pi_r,\\
&d_r,b_r,k_r,n_r,t_r,\eta_r,L_r,A_r).
\end{aligned}
\]
其中 \(\phi_r\) 为阶段，\(i_r,\mathcal D_r\) 为实际主体及委托关系，
\(u_r,\sigma_r\) 为实际用途与作用范围，
\(v_r,S_r,\pi_r\) 为实现版本、规范及证明材料。
\(d_r\) 是输入与必要环境记录，
\(b_r,k_r\) 是预算及配额记录，
\(n_r,t_r\) 是操作标识与时间记录，
\(\eta_r\) 是实际作用及回执。
\(L_r\) 逐主体保存应付、实付、应退、实退与期限；
\(A_r\) 保存理由、复核、裁决和后续接入条件。
共享状态另保存可用资源、全部预留、已耗用资源、
当前规则版本和已经使用的操作标识。
\end{definition}

\begin{remark}[真实状态与可取得证据]
真实状态中的资格和实际作用，可能尚未被核验者完整观察。
证据读取写为 \(e:Z\to E\)，
不把 \(\theta z\) 中的所有事实默认成可即时读取的字段。
本章分别证明决定的健全性与实际执行的保持；
证据足以恢复类型的条件见第~\ref{std:chapter}~章。
\end{remark}

\begin{definition}[原子语义转移]
\label{svc:atom}
模型中的原子步骤为
\[
x\xrightarrow{a,\xi}x',
\qquad
a\in A(x),\quad
\xi\in\Xi(x,a),\quad
x'=\delta(x,a,\xi).
\]
采用定义~\ref{atom:record}~的前置、后置和字段保持记录。
原子性相对于声明的任务、作用接口和允许分解，
并不要求其等于单条机器指令。
必须一致提交的多个字段可以属于同一原子；
如声明不可再分，则按定义~\ref{atom:minimal}~核验允许细化。
每个合法动作的响应集非空，
并覆盖成功、拒绝、超时、崩溃及结果待确认等声明情形。
实际步骤及其响应解释满足
\[
\begin{aligned}
z'&=\Delta(z,a,\zeta),\\
\eta(z,a,\zeta)&\in\Xi(\theta z,a),\\
\theta z'&=\delta(\theta z,a,\eta(z,a,\zeta)).
\end{aligned}
\]
若外部作用不能与本地记录共同原子提交，
则分别登记发送意图、外部提交、回执取得和结果核对步骤。
本地回滚不解释为外部作用已经撤销。
\end{definition}

\begin{definition}[服务不变量]
\label{svc:invariant}
令 \(I:X\to\{0,1\}\) 为以下义务的合取：
字段有类型，实际余额与预留非负，
共享资源没有重复承诺，已发生作用均有对应的有效授权，
既有权利只按登记合同合法演化，
拒绝与救济记录没有因停止业务而被删除。
对每一项同单位资源，采用
\[
f+\sum_{r}h_r+s=C,\qquad
f\ge0,\quad h_r\ge0,\quad s\ge0
\]
表示自由资源、逐请求预留及已耗用量之间的核对关系。
不同资源分别使用该式，不把金额、时间和计算配额相加。
到期义务的实际履行另由完成目标与进展证书表达。
\end{definition}

\section{三值核验、授权证书与结果证书}

\begin{definition}[三值核验与申请决定]
\label{svc:three-values}
对请求 \(r\)，固定有限证明义务
\[
p_{rj}:X\longrightarrow\{0,1\},
\qquad j\in J_r.
\]
独立检查结果为
\[
v_{rj}\in\{1,0,?\}.
\]
健全性要求
\[
v_{rj}=1\Longrightarrow p_{rj}(x)=1,
\qquad
v_{rj}=0\Longrightarrow p_{rj}(x)=0.
\]
问号表示尚未完成该项判定。
定义决定 \(D_r(x)\in\{+,-,?\}\) 为
\[
D_r(x)=
\begin{cases}
- ,&\text{至少一项 }v_{rj}=0,\\
+ ,&\text{全部 }v_{rj}=1,\\
? ,&\text{其余情形}.
\end{cases}
\]
符号 \(+\)、\(-\)、\( ? \) 分别表示准许、拒绝与待核验。
缺少合同明确要求的材料可以构成一个已判假的义务；
未能判断某项事实则保留为问号。
\end{definition}

\begin{lemma}[三值决定的健全含义]
\label{svc:three-sound}
若各单项检查满足定义~\ref{svc:three-values}，
则
\[
D_r(x)=+
\Longrightarrow
\forall j\in J_r,\ p_{rj}(x)=1,
\]
\[
D_r(x)=-
\Longrightarrow
\exists j\in J_r,\ p_{rj}(x)=0.
\]
待核验本身不蕴含实际条件不成立。
\end{lemma}
\begin{certificate}
有限义务表、单项检查的解释与健全性证书、
决定合成表，以及一个实际条件成立而检查尚未完成的状态。
\end{certificate}
\begin{proof}
\begin{semanticsteps}
\item 固定 \(x\in X\) 及请求 \(r\)。
由决定表中准许分支的类型定义，
\[
D_r(x)=+\Longrightarrow
\forall j\in J_r,\ v_{rj}=1.
\]
\item 对上一项中的每个 \(j\)，应用单项健全性，
\[
v_{rj}=1\Longrightarrow p_{rj}(x)=1.
\]
对有限指标集取合取，得到第一式。
\item 若 \(D_r(x)=-\)，决定表给出某个
\(j_0\in J_r\) 满足 \(v_{rj_0}=0\)。
单项健全性给出
\[
p_{rj_0}(x)=0.
\]
把该指标作为存在见证，得到第二式。
\item 取 \(J_r=\{1\}\)、\(p_{r1}(x)=1\)、\(v_{r1}=?\)。
两个健全性蕴含的前件均不成立，因此没有违反检查条件。
决定表返回 \(D_r(x)=?\)，但实际条件仍为真。
故不能由待核验推出实际条件为假。
\end{semanticsteps}
\end{proof}

\begin{definition}[授权证书的绑定字段]
\label{svc:authorization-certificate}
授权证书 \(c_r\) 包含有类型的元组
\[
\begin{aligned}
c_r=(&r,i_r,\mathcal D_r,u_r,\sigma_r,v_r,S_r,\\
&h_{\rm in},h_{\rm env},\pi_r,
\bar b_r,\bar k_r,n_r,\nu,t_{\exp}).
\end{aligned}
\]
其中 \(h_{\rm in},h_{\rm env}\) 绑定实际输入及必要环境条件，
\(\bar b_r,\bar k_r\) 是作用上界，
\(\nu,t_{\exp}\) 是适用规则与有效期。
验证记录同时绑定检查器及规范版本。
认证机制证明所登记字段的来源与完整性；
独立核验另证明程序或动作满足声明规范。
实际输入、作用对象和当前版本须与这些字段匹配。
\end{definition}

\begin{definition}[结果与履约证书]
\label{svc:result-certificate}
结果证书 \(c_r^{\rm out}\) 绑定请求与操作标识、
实际执行对象、作用回执、实际资源使用及后继状态。
固定任务结果谓词 \(Q_r:X\to\{0,1\}\) 与履约谓词
\(H_r:X\to\{0,1\}\)。
结果检查器 \(V_{\rm out}\) 的健全性要求
\[
V_{\rm out}(c_r^{\rm out},x')=1
\Longrightarrow Q_r(x')=1.
\]
完整履约确认还要求
\[
Q_r(x')=1,\qquad H_r(x')=1.
\]
一份允许开始执行的证书，
不记录尚未发生的作用为已经发生，
也不把应付或应退金额记录为已实际支付。
\end{definition}

\begin{proposition}[授权与履约确认的分离]
\label{svc:authorization-not-result}
授权证书有效不蕴含结果证书已经成立，
也不蕴含付款、退款或申诉处理已经完成。
\end{proposition}
\begin{certificate}
给出已获准但尚未执行的状态 \(x\)，
其授权字段全部有效，
而 \(Q_r(x)=0\) 或 \(H_r(x)=0\)。
\end{certificate}
\begin{proof}
\begin{semanticsteps}
\item 取许可状态 \(x\in X\)，使全部授权义务成立。
根据定义~\ref{svc:three-values}，
完整通过的核验可以产生 \(D_r(x)=+\)。
\item 将服务结果字段设为尚未交付，
并按结果谓词定义登记
\[
Q_r(x)=0.
\]
该字段描述作用是否已经发生，
与开始作用前的资格字段具有不同类型。
\item 若从授权有效能够推出结果已成立，
则在此状态同时得到 \(Q_r(x)=1\) 与 \(Q_r(x)=0\)，矛盾。
\item 同理，取结果已产生而尚未实际结算的状态，
使 \(Q_r(x)=1\)、\(H_r(x)=0\)。
这排除从结果存在推断全部履约完成。
完成期限另须满足定理~\ref{tra:liveness}~的进展条件。
\end{semanticsteps}
\end{proof}

\section{作用提交、资源预留与重复请求}

\begin{definition}[实际作用的准入谓词]
\label{svc:gate}
对动作 \(a\)、证书 \(c\) 及当前状态 \(x\)，
分别以短谓词表示身份与委托、用途、
作用范围、对象绑定、证明、环境、
资源及操作标识检查：
\[
\begin{gathered}
\mathsf{Id},\ \mathsf{Use},\ \mathsf{Scp},\
\mathsf{Bind},\ \mathsf{Prf},\ \mathsf{Env},\
\mathsf{Res},\ \mathsf{Key}.
\end{gathered}
\]
另以 \(\mathsf{Now}(x,a,c)\) 表示实际提交时的规则、
撤销状态与有效期检查。
这些谓词均具有类型
\[
X\times A\times C\longrightarrow\{0,1\}.
\]
准入谓词为
\[
\begin{aligned}
G(x,a,c):={}&
\mathsf{Id}\land\mathsf{Use}\land\mathsf{Scp}
\land\mathsf{Bind}\land\mathsf{Prf}\\
&\land\mathsf{Env}\land\mathsf{Res}
\land\mathsf{Key}\land\mathsf{Now}.
\end{aligned}
\]
各短谓词均在同一个 \((x,a,c)\) 上求值。
作用提交还要求对应业务决定为准许。
退款、核对与申诉等处置动作使用各自的授权义务，
不继承已经失效的业务用途许可。
\end{definition}

\begin{definition}[原子预留与操作标识]
\label{svc:reservation}
操作键 \(\kappa\) 由请求标识、操作序号及绑定负载确定。
持久状态记录该键的负载、预留、作用提交状态与既有回执。
对一项资源，原子预留规则为
\[
\begin{gathered}
f\ge b,\qquad b\ge0,\qquad
\kappa\text{ 尚未登记},\\
(f,h_r,s)\longmapsto(f-b,h_r+b,s).
\end{gathered}
\]
余额比较、键的首次登记和预留更新属于同一已核验提交。
作用耗用 \(c\) 的规则为
\[
0\le c\le h_r,\qquad
(f,h_r,s)\longmapsto(f,h_r-c,s+c).
\]
释放未耗用预留 \(u\) 的规则为
\[
0\le u\le h_r,\qquad
(f,h_r,s)\longmapsto(f+u,h_r-u,s).
\]
同时到达的申请必须按实际隔离或等价串行化证书执行这些规则。
\end{definition}

\begin{lemma}[逐步预留不重复承诺资源]
\label{svc:resource-invariant}
若初态满足定义~\ref{svc:invariant}~的资源等式及非负条件，
且每次资源变化都是定义~\ref{svc:reservation}~中的合法提交，
则任意有限前缀保持该等式及全部余额非负。
\end{lemma}
\begin{certificate}
初始资源表、全部共享预留、
逐次原子提交的先决条件、更新记录及并发隔离见证。
\end{certificate}
\begin{proof}
\begin{semanticsteps}
\item 对预留步骤，由 \(f\ge b\ge0\) 得
\(f-b\ge0\) 与 \(h_r+b\ge0\)。
其他请求预留未变，故
\[
(f-b)+(h_r+b)+
\sum_{j\ne r}h_j+s
=f+\sum_j h_j+s.
\]
\item 对耗用步骤，由 \(0\le c\le h_r\) 得
\(h_r-c\ge0\) 与 \(s+c\ge0\)。
对应总量为
\[
f+(h_r-c)+
\sum_{j\ne r}h_j+(s+c)
=f+\sum_j h_j+s.
\]
\item 对释放步骤，由 \(0\le u\le h_r\) 得
\(f+u\ge0\) 与 \(h_r-u\ge0\)，并有
\[
(f+u)+(h_r-u)+
\sum_{j\ne r}h_j+s
=f+\sum_j h_j+s.
\]
\item 原子提交证书保证每个更新读取的是该提交的当前余额，
不是另一个并发更新之前已经失效的值。
因此以上三种单步推导覆盖实际更新。
\item 以初始资源表为归纳基例，
对提交次数逐次应用相应单步等式，
得到每个有限前缀的守恒与非负性。
\end{semanticsteps}
\end{proof}

\begin{definition}[重复操作的一致处理]
\label{svc:idempotency}
以 \(N_\kappa(x)\in\mathbb N\) 表示操作键 \(\kappa\)
已经造成的实际受控作用次数。
首次作用提交要求
\[
N_\kappa(x)=0,\qquad G(x,a,c)=1,
\]
并在同一实际作用提交语义中登记
\[
N_\kappa(x')=1.
\]
同键同负载的重复请求读取原回执或待确认状态，
不再触发新的实际作用。
同键异负载产生有理由的拒绝记录。
键的持久消费与实际作用之间须有原子性或外部幂等实现证书。
尚未收到回执不解释为该键尚未产生作用。
\end{definition}

\begin{theorem}[重复请求下的至多一次作用]
\label{svc:at-most-once}
若 \(N_\kappa(x_0)=0\)，
全部首次作用与重复请求满足定义~\ref{svc:idempotency}，
并且没有遗漏的作用入口，
则每个有限执行前缀满足
\[
N_\kappa(x_t)\in\{0,1\}.
\]
该结论不保证作用一定发生，也不保证回执按期到达。
\end{theorem}
\begin{certificate}
持久操作键表、负载绑定、全部作用入口、
实际作用与键消费的实现证书、
重试及崩溃恢复的转移表。
\end{certificate}
\begin{proof}
\begin{semanticsteps}
\item 在初态，已给出 \(N_\kappa(x_0)=0\)，
故结论对零步前缀成立。
\item 假设当前值为零。
若下一步为首次作用提交，
其合法后继按定义具有 \(N_\kappa(x_{t+1})=1\)。
若不是首次作用，该步骤不增加此计数，后继仍为零。
\item 假设当前值为一。
首次作用的守卫要求当前值为零，因此该分支不可合法进入。
同键重试、回执查询和恢复步骤均保留计数，
所以后继仍为一。
\item 全入口覆盖排除上述转移表之外另一次实际作用。
由第二、第三项，归纳步对全部允许响应成立。
\item 对 \(t\) 归纳得到所需范围。
全程没有使用某一步必然发生的前件，
故所得是次数上界，不是完成或可用性结论。
\end{semanticsteps}
\end{proof}

\begin{proposition}[提前核验不足以替代提交时核验]
\label{svc:current-binding}
检查时证书有效，不蕴含稍后实际作用时仍满足 \(G\)。
\end{proposition}
\begin{certificate}
取两个规则时期 \(0,1\)，
一份只适用于时期 \(0\) 的证书，
以及检查之后、作用之前发生的有效规则更新。
\end{certificate}
\begin{proof}
\begin{semanticsteps}
\item 在状态 \(x_0\) 中令当前时期为零，证书时期亦为零。
于是当前性谓词可以满足
\[
\mathsf{Now}(x_0,a,c)=1.
\]
\item 执行合法更新，得到时期为一的状态 \(x_1\)。
证书字段没有随该更新自动改变，仍登记时期零。
\item 按时期一致性的定义，
\[
\mathsf{Now}(x_1,a,c)=0.
\]
\item \(\mathsf{Now}\) 是 \(G\) 的一个合取项，
因此 \(G(x_1,a,c)=0\)。
这与先前检查通过相容，
故先前通过不能推出当前准入。
实际提交必须重新检查，或由覆盖整个作用期间的有效绑定协议保持该项。
\end{semanticsteps}
\end{proof}

\section{全路径覆盖与有限安全归纳}

\begin{definition}[全作用路径覆盖]
\label{svc:coverage}
令 \(\mathsf{Eff}(z,a,\zeta)\) 表示实际步骤产生受控作用。
覆盖证书要求
\[
\mathsf{Eff}(z,a,\zeta)=1
\Longrightarrow
\exists c,\ G(\theta z,a,c)=1
\]
且该 \(c\) 绑定该步骤实际使用的对象与当前条件。
直接入口、内部调用、后台任务、队列消费、回调、
委托代理及外部作用端点均列入覆盖。
在调用与作用图中，每条可实现的作用路径
都须经过适用于该实际作用的有效检查点。
对运行时增加的路径，先补充覆盖证书再允许产生作用。
\end{definition}

\begin{definition}[原子步骤的局部保持证书]
\label{svc:local-certificate}
把动作分成业务作用、核验与等待、资金处置、复核与退出四类。
对每个实际允许的原子动作保存
\[
I(x)=1\land a\in A(x)
\Longrightarrow
\forall\xi\in\Xi(x,a),\
I(\delta(x,a,\xi))=1.
\]
业务作用以准入谓词及程序证明支撑该式；
资源变化以引理~\ref{svc:resource-invariant}~支撑相应分量；
重复请求以定理~\ref{svc:at-most-once}~支撑计数分量。
等待、拒绝与救济仍须检查其实际费用、
资金变化及既有权利，不能作为无条件免费步骤。
\end{definition}

\begin{theorem}[公共服务有限执行的安全保持]
\label{svc:finite-safety}
若实际初态解释到 \(I\)，
全部实际原子步骤满足定义~\ref{svc:atom}~的符合关系，
业务作用满足定义~\ref{svc:coverage}，
且局部证书覆盖全部合法动作及响应，
则任意候选生成序列与允许响应造成的长度 \(m\) 的有限执行满足
\[
\forall t\in\{0,\ldots,m\},\qquad I(\theta z_t)=1.
\]
\end{theorem}
\begin{certificate}
初态、作用入口覆盖、局部保持证书、
实际响应解释、状态交换式及逐版本绑定。
候选生成可以提交动作和证明材料，
但接受规则、实际作用检查与结果回执分别独立登记。
\end{certificate}
\begin{proof}
\begin{semanticsteps}
\item 对实际初态 \(z_0\)，由初态证书取得
\[
x_0=\theta z_0\in X,\qquad I(x_0)=1.
\]
\item 假定 \(I(\theta z_t)=1\)，并取下一实际步骤
\[
z_{t+1}=\Delta(z_t,a_t,\zeta_t).
\]
全路径证书把实际业务作用送入已核验分支，
其余实际步骤送入各自的处置或等待分支。
\item 符合关系给出有类型响应
\[
\xi_t=\eta(z_t,a_t,\zeta_t)
\in\Xi(\theta z_t,a_t).
\]
局部证书中的全称响应量词因此覆盖 \(\xi_t\)。
\item 应用对应局部证书得到
\[
I(\delta(\theta z_t,a_t,\xi_t))=1.
\]
再以状态交换式替换括号内的后继，得到
\[
I(\theta z_{t+1})=1.
\]
\item 上述步骤核对了定理~\ref{tra:safety}~的全部服务实例前件。
按执行长度归纳得到任意有限前缀的结论；
候选从何产生不出现在该归纳的接受条件中。
\end{semanticsteps}
\end{proof}

\begin{remark}[核验语义与实际控制]
证书的密码真实性、核验程序的实现正确性、
输入及环境的真实性、作用端点的符合性分别登记。
本章的确定性结论以这些前件成立为条件。
若采用带失败概率的认证模型，
应另把有限窗口内的失败概率纳入声明的风险界。
应用定理~\ref{gov:complete-gate}~时，
实际用途与作用范围须进入同一准入谓词，
不能仅改变服务名称或插入代理关系便改变决定。
\end{remark}

\section{拒绝、退款、申诉与重新申请}

\begin{definition}[可观察的拒绝与待核验状态]
\label{svc:denied-state}
提交请求并收到拒绝，是模型内的业务转移：
\[
x\xrightarrow{\mathrm{submit},\,\mathrm{deny}(j)}x'.
\]
后继保存理由 \(j\)、原请求标识、既有资金义务与救济条件。
待核验转移则保存尚未完成的证明义务与后续检查动作。
二者均属于 \(X\)，不与定义~\ref{q:totalization}~中的
未定义状态 \(\bot\) 等同。
对可以实际发生的越权尝试，
其拒绝与记录动作也纳入允许响应表。
\end{definition}

\begin{definition}[拒绝后的权利登记]
\label{svc:rights}
对每个请求 \(r\) 与权利主体 \(i\)，分别保存
\[
L_{ri}=(d_{ri},p_{ri},u_{ri},f_{ri},\tau_{ri},a_{ri}),
\]
依次为应付、实付、应退、实退、期限与复核状态。
仅停止业务的拒绝动作可以改变阶段与理由，
但必须满足
\[
(d_{ri},p_{ri},u_{ri},f_{ri},\tau_{ri})'
=(d_{ri},p_{ri},u_{ri},f_{ri},\tau_{ri}).
\]
如果合同允许费用结算或义务调整，
则以具有独立依据、金额和回执的处置动作实现。
复核状态的变化不能无证据地充当实际退款。
\end{definition}

\begin{lemma}[停止业务不消灭既有退款义务]
\label{svc:denial-rights}
在定义~\ref{svc:rights}~下，
若拒绝前 \(u_{ri}-f_{ri}=a>0\)，
则仅停止业务的拒绝后仍有
\[
u'_{ri}-f'_{ri}=a.
\]
因此拒绝记录本身不能成为退款完成证书。
\end{lemma}
\begin{certificate}
拒绝前后的逐主体账本、
拒绝动作的字段更新清单，
以及实际退款与合同调整的独立动作类别。
\end{certificate}
\begin{proof}
\begin{semanticsteps}
\item 从权利登记读得拒绝前的两个有理金额
\[
u_{ri},f_{ri}\in\mathbb Q_{\ge0},
\qquad u_{ri}-f_{ri}=a>0.
\]
\item 字段保持式分别给出
\[
u'_{ri}=u_{ri},\qquad f'_{ri}=f_{ri}.
\]
\item 逐项代入差额，得到
\[
u'_{ri}-f'_{ri}
=u_{ri}-f_{ri}
=a>0.
\]
\item 退款完成谓词要求该项义务已依法清偿。
当前差额仍为正，因此单独的拒绝步骤不能满足该谓词。
\end{semanticsteps}
\end{proof}

\begin{proposition}[未交付结果不足以决定退款与救济]
\label{svc:failure-observation}
设结果表示 \(q_0\) 仅记录是否交付。
若存在两个未交付状态，
其应退金额或复核资格不同，
则这些权利不能全部通过 \(q_0\) 读回。
\end{proposition}
\begin{certificate}
取 \(q_0(x)=q_0(y)=0\)，
并登记 \(u_{ri}(x)=1\)、\(u_{ri}(y)=0\)；
两状态的金额来源分别有合法合同记录。
\end{certificate}
\begin{proof}
\begin{semanticsteps}
\item 假定存在读回函数
\[
g:q_0(X)\longrightarrow\mathbb Q_{\ge0},
\qquad u_{ri}=gq_0.
\]
\item 由 \(q_0(x)=q_0(y)\)，函数单值性给出
\[
g(q_0(x))=g(q_0(y)).
\]
\item 用因子化式分别替换两侧，得到
\[
u_{ri}(x)=u_{ri}(y).
\]
该式与证书中的 \(1\ne0\) 矛盾。
\item 因而引理~\ref{q:fiber}~的纤维条件不成立。
同样推导适用于同一结果纤维内不同的复核资格。
删除未交付请求也不能提供这些观察的总定义读回。
\end{semanticsteps}
\end{proof}

\begin{definition}[申诉与重新申请的接续]
\label{svc:reentry}
申诉可以产生维持拒绝、变更决定、
补正材料或允许重新申请等响应；
每种响应保留原请求的裁决与权利记录。
重新申请使用新标识 \(r'\ne r\)，
登记关联 \(\operatorname{parent}(r')=r\)，
并重新检查当期身份、用途、版本、预算和配额。
原请求未清偿的义务继续保留，
除非已由有依据的清偿或变更动作依法消灭。
允许申请的裁决不等于新申请已经取得实际作用许可。
\end{definition}

\begin{corollary}[救济路径与准入边界相容]
\label{svc:remedy-compatibility}
若复核及重新申请步骤均有局部保持证书，
新请求的实际作用仍经过谓词 \(G\)，
则加入这些救济步骤不破坏定理~\ref{svc:finite-safety}~的有限安全结论，
且原请求未清偿权利仍可按其标识读回。
\end{corollary}
\begin{certificate}
新增动作与完整响应表、逐主体权利的更新规则、
父请求关联、新请求核验及全部作用入口覆盖。
\end{certificate}
\begin{proof}
\begin{semanticsteps}
\item 将新增动作加入原合法动作集，
并为每个新增动作的全部响应给出保持 \(I\) 的局部证书。
\item 原业务动作的准入式不变，
新请求作用也要求 \(G=1\)，
所以原全路径覆盖条件在扩展动作表上继续成立。
\item 按定义~\ref{svc:reentry}，
父请求记录没有被新标识覆盖；
按定义~\ref{svc:rights}，
未清偿权利没有被仅改变阶段的动作删除。
\item 扩展系统因此满足定理~\ref{svc:finite-safety}~的同一组前件。
应用该定理取得有限安全保持，
再以原请求及主体索引读取其剩余权利。
\end{semanticsteps}
\end{proof}

\section{实际策略、调度与按期处置}

\begin{definition}[交付目标与处置目标]
\label{svc:targets}
服务交付目标 \(G_{\rm del}\) 要求规定结果与相应履约成立。
合法关闭目标 \(G_{\rm cls}\) 要求拒绝或终止已有依据，
到期资金义务完成，
且必须处理的复核已经裁决或依法结束。
登记的处置目标可取
\[
G=G_{\rm del}\cup G_{\rm cls}.
\]
证明到达 \(G\) 只表示完成合同规定的处置，
不能据此宣布每个申请均获得了所申请的业务结果。
如需保证实际交付，应独立选择 \(G=G_{\rm del}\)
并对该目标提交全部响应下的策略证书。
\end{definition}

\begin{definition}[被实际执行的进展策略]
\label{svc:implemented-policy}
令 \(\rho:\{x\in X:I(x)=1\}\to\mathbb N\) 满足
\[
\rho(x)=0\Longleftrightarrow x\in G.
\]
对正秩状态定义进展动作集合
\[
\begin{aligned}
P(x)=\{a\in A(x):\
\forall\xi\in\Xi(x,a),\ 
\rho(\delta(x,a,\xi))\le\rho(x)-1\}.
\end{aligned}
\]
控制器有明确的动作选择程序 \(\sigma\)。
其实际运行证书要求：
所有非进展步骤不增加秩；
未完成时至多经过 \(K\ge1\) 个实际原子步骤，
便执行一次由 \(\sigma\) 选取的 \(P(x)\) 中动作；
系统不会在目标外停滞。
检查、排队、重试、等待和复核均计入实际步数。
完成证书另登记每步总时延与总成本上界。
\end{definition}

\begin{proposition}[存在进展动作不保证实际完成]
\label{svc:selection-needed}
即使每个非目标状态都存在进展动作，
未约束实际选择和调度仍不能推出有限完成。
\end{proposition}
\begin{certificate}
取 \(X=\{s,g\}\)、\(G=\{g\}\)，
在 \(s\) 允许两个动作 \(a,w\)，响应均为单元素集。
规定
\[
\delta(s,a)=g,\qquad
\delta(s,w)=s,\qquad
\rho(s)=1,\quad\rho(g)=0.
\]
令实际控制器始终选择 \(w\)。
\end{certificate}
\begin{proof}
\begin{semanticsteps}
\item 对非目标状态 \(s\)，动作 \(a\) 的后继秩为零，
故 \(a\in P(s)\)，进展动作存在。
\item 实际策略满足 \(\sigma(s)=w\)；
由转移表，其实际后继仍为 \(s\)。
\item 以 \(x_0=s\) 为基例，
若 \(x_t=s\)，则
\[
x_{t+1}=\delta(s,w)=s.
\]
按步数归纳，所有有限前缀均留在 \(s\)。
\item 因为 \(s\notin G\)，不存在有限完成时点。
故完成结论需要实际选择及调度证书，
不能仅以 \(P(x)\ne\varnothing\) 代替。
\end{semanticsteps}
\end{proof}

\begin{theorem}[有界调度下的实际按期处置]
\label{svc:scheduled-completion}
在定理~\ref{svc:finite-safety}~的条件下，
若定义~\ref{svc:implemented-policy}~全部成立，
且每步时延至多 \(\bar t\)、成本至多 \(\bar c\)，
则从 \(x_0\) 出发至多
\[
N=K\rho(x_0)
\]
步进入声明目标 \(G\)，并满足
\[
t_{\rm finish}-t_0\le N\bar t,
\qquad
c_{\rm finish}\le N\bar c.
\]
若剩余期限和已预留的同单位资源覆盖这两个上界，
则此路径具有相应的期限与资源覆盖证书。
\end{theorem}
\begin{certificate}
实际策略程序及其绑定版本、
逐状态选择动作、全部响应下的秩变化、
调度间隔证书、等待及恢复时延、
共享资源与逐前缀预留记录。
\end{certificate}
\begin{proof}
\begin{semanticsteps}
\item 定理~\ref{svc:finite-safety}~保证实际状态留在 \(I\)，
所以每一步的秩 \(\rho(x_t)\) 均有定义。
\item 设尚未达到目标。
调度证书保证随后至多 \(K\) 个实际步骤中执行一个进展动作；
在它之前的非进展步骤均不增加秩。
\item 该被实际选取的动作属于当前 \(P(x)\)。
由 \(P\) 的全称响应定义，
实际响应造成的后继秩至少下降一。
该进展之后、下一次进展之前的步骤亦不增加秩。
\item 对完成进展块的个数 \(j\) 归纳，
在目标尚未到达的各块末端有
\[
\rho(x)\le\rho(x_0)-j.
\]
非负整数不能在 \(j=\rho(x_0)\) 后仍为正，
故至多经过此数目的块达到 \(G\)。
\item 每块至多 \(K\) 步，
得到 \(N=K\rho(x_0)\) 的总步数上界。
该前件核对同时满足定理~\ref{tra:liveness}~的完成条件。
\item 对全部实际步骤分别求和：
\[
\sum_{j<N}t_j\le N\bar t,\qquad
\sum_{j<N}c_j\le N\bar c.
\]
以剩余期限和资源预留覆盖这两个上界，
即得到对应可执行性条件。
\end{semanticsteps}
\end{proof}

\begin{corollary}[有限策略表的实施接口]
\label{svc:finite-policy}
若请求状态属于定理~\ref{q:reachability}~的有限鲁棒可达集 \(W_n\)，
转移表可精确计算，
则可以有限枚举选择一个逐步降秩策略表。
该表只有在非进展步骤保留此可达域且不增加秩、
实际执行绑定及有界调度条件全部成立后，
才产生定理~\ref{svc:scheduled-completion}~的实际完成结论。
\end{corollary}
\begin{certificate}
有限状态、动作与响应表，
可达递推的每轮集合，
各正秩状态的选定动作，
运行控制器与该表的一致性，以及调度见证。
\end{certificate}
\begin{proof}
\begin{semanticsteps}
\item 对 \(x\in W_n\)，取最小指标
\[
\rho(x)=\min\{j:x\in W_j\}.
\]
有限集合上的最小值可以由已列出的轮次逐项检查。
\item 若 \(\rho(x)>0\)，可达递推保证至少有一个动作
使全部响应后继属于 \(W_{\rho(x)-1}\)。
枚举有限动作表，逐项检查其非空响应表并保存首个通过者。
\item 对保存的动作 \(a\)，每个允许响应均满足
\[
\rho(\delta(x,a,\xi))\le\rho(x)-1.
\]
这就是推论~\ref{q:policy}~所需的选择见证。
\item 将实际控制器绑定到该表，
核验非进展步骤保留可达域且不增加秩，
再核验至多 \(K\) 步执行一次所选动作，
才取得定义~\ref{svc:implemented-policy}~的运行前件。
\item 应用定理~\ref{svc:scheduled-completion}~得到实际完成。
表的存在、表的计算和表的实际执行分别由上述不同步骤承担。
\end{semanticsteps}
\end{proof}

\begin{remark}[更新、共享容量与既有申请]
新版本、附加能力及新申请不得重复使用已承诺给旧申请的资源。
更新时应采用定义~\ref{tra:update-certificate}~中的迁移证书，
保持既有逐主体权利以及未完任务的共同可实现性。
再按定理~\ref{tra:update-invariance}~检查运行与更新的交替序列。
需要按期交付时，还须继续满足本章的实际调度条件；
仅保持未完任务仍有某条可行路径，并不限制其被推迟多久。
\end{remark}


\chapter{常规任务保护与劳动需求}
\label{lab:chapter}

本章把制度限制影响就业的命题拆成可分别检验的三个层次：
限制未经许可的直接替代，保留有真实需求的人类任务，
以及检验全部经济反馈后的劳动收入和实际购买力。
任务替代、新任务形成与生产率变化可以同时发生；
因此某一替代渠道的收缩不能直接代替净就业的比较
\cite{AcemogluRestrepo2019,Autor2015}。
以下制度构造沿用参与经济文献中关于任务对象、真实支付与最小兼容修复的区分
\cite{InternalEconomics}。

\section{比较范围与劳动投入的同口径计量}

\begin{definition}[常规任务比较域]
\label{lab:def-domain}
固定时期 \(t\)、地区、产品质量、规则版本与反事实基线。
令 \(\mathcal O_t\) 为事前登记的常规任务集合，
\(\mathcal S_t\) 为另行许可的专门任务集合。
任务的分类依据实际用途、控制权限与交付过程，
同一任务在本次比较中只归属于一个核算单元。
若同一生产过程兼含两类任务，先分解为可分别记录的任务单元，
再分别核算其劳动投入和最终产出。

用上标 \(0\) 表示比较基线，上标 \(G\) 表示实施待评制度后的状态。
本章的直接替代比较固定任务需求和单位劳动系数；
需求、价格、资本配置及新任务的变化在后面的总效应比较中另行纳入。
这个限定使直接替代具有确定的分母。
\end{definition}

\begin{definition}[劳动小时、全时当量与实际就业人数]
\label{lab:def-hours}
令 \(H_t\ge0\) 为同一核算范围内实际购买并完成的人类劳动小时。
给定固定的全时标准 \(\bar h^{F}>0\)，定义全时当量
\[
 J_t^{F}=\frac{H_t}{\bar h^{F}}.
\]
实际就业人数为
\[
 N_t^{\mathrm{emp}}
 =\left|\{i:h_{it}>0\}\right|,
 \qquad
 H_t=\sum_i h_{it}.
\]
其中 \(h_{it}\) 是主体 \(i\) 的实际有酬劳动小时。
除非另有共同工时约束，
\(J_t^{F}\) 的增减不等同于 \(N_t^{\mathrm{emp}}\) 的增减。
未执行的招聘计划、免费体验和重复登记不计入实际劳动小时。
\end{definition}

\section{覆盖率、执行率与直接替代的上界}

\begin{definition}[基线替代单元与剩余替代]
\label{lab:def-displacement}
令 \(\mathcal A_t\) 为有限的反事实直接替代单元集合。
每个单元 \(a\) 对应在指定基线、同一任务环境下将被自动化取代的劳动量
\(d_{at}\ge0\)，并且不同单元不重复计算同一劳动小时。
基线直接替代量为
\[
 D_t^0=\sum_{a\in\mathcal A_t}d_{at}.
\]
纯技术上能够替代的任务量可以另行登记，
但未进入指定基线的实际替代路径时，不计入 \(D_t^0\)。
令 \(\theta_{at}\in[0,1]\) 为其中进入制度有效覆盖范围的比例，
\(\varepsilon_{at}\in[0,1]\) 为已覆盖部分仍实际发生未经许可替代的比例。
在本项固定需求比较中，要求未覆盖部分的替代量不超过对应基线量。
制度诱发的新任务替代或任务规模变化进入后文总效应的相应更新步骤。
在这一比较边界内，剩余直接替代具有保守上界
\[
 \overline D_t^G
 =\sum_{a\in\mathcal A_t}
 d_{at}\bigl[(1-\theta_{at})+\theta_{at}\varepsilon_{at}\bigr],
 \qquad
 0\le D_t^G\le\overline D_t^G.
\]
这里的规避替代已经包含在
\(\theta_{at}\varepsilon_{at}d_{at}\) 中，
不在总效应核算时再次作为一项直接损失扣除。
覆盖率与剩余替代率均指实际任务范围内的数量，
规则文本的名义覆盖不代替这两个参数。
\end{definition}

\begin{theorem}[有效覆盖下的直接替代界]
\label{lab:thm-direct-bound}
若对全部 \(a\in\mathcal A_t\) 都有
\[
 \theta_{at}\ge\underline\theta_t,
 \qquad
 \varepsilon_{at}\le\overline\varepsilon_t,
 \qquad
 0\le\underline\theta_t,\overline\varepsilon_t\le1,
\]
则
\[
 D_t^G
 \le
 \bigl[1-\underline\theta_t(1-\overline\varepsilon_t)\bigr]D_t^0.
\]
故所避免的直接替代量
\(A_t^{\mathrm{dir}}:=D_t^0-D_t^G\) 满足
\[
 A_t^{\mathrm{dir}}
 \ge\underline\theta_t(1-\overline\varepsilon_t)D_t^0.
\]
当 \(D_t^0>0\)、\(\underline\theta_t>0\) 且
\(\overline\varepsilon_t<1\) 时，下界严格为正。
\end{theorem}
\begin{certificate}
记录比较基线、完整替代单元表、每个单元的劳动量、
覆盖范围、实际规避路径和去重标识。
参数界的证据须覆盖同一时期和相同用途；
跨境服务、外包和具身设备通过实际任务归属进入同一分母。
若只取得样本估计，则分别报告参数区间与其适用总体。
\end{certificate}
\begin{proof}
\begin{semanticsteps}
\item 逐项有
\[
 (1-\theta_{at})+\theta_{at}\varepsilon_{at}
 =1-\theta_{at}(1-\varepsilon_{at})
 \le1-\underline\theta_t(1-\overline\varepsilon_t).
\]
乘以非负 \(d_{at}\) 并求和，得到上界。
\item 用 \(D_t^0\) 减去该上界得到避免替代量的下界。
\item 三个严格条件使下界中的各因子均为正。
\item 该推演使用逐项界而不要求不同替代事件相互独立。
\end{semanticsteps}
\end{proof}

\begin{remark}
覆盖与规避按所涉劳动量加权。
少数部署可能承担大量任务，故违规设备数或违规部署数占比
不能直接替代 \(\varepsilon_{at}\)。
从部署记录估计劳动替代率时，
须逐一连接部署、被替代任务及其劳动量，再在相同分母上加总。
\end{remark}

\begin{corollary}[完全覆盖的直接保护特例]
\label{lab:cor-perfect-gate}
若 \(\underline\theta_t=1\) 且 \(\overline\varepsilon_t=0\)，
则登记任务域内 \(D_t^G=0\)。
若仅 \(\underline\theta_t=0\)，或仅能给出
\(\overline\varepsilon_t=1\)，则上述定理不给出正的避免替代下界。
\end{corollary}
\begin{certificate}
完整覆盖须逐条验证任务的实际入口和控制权限；
零规避须给出声明攻击域内的控制证据。
证书同时保留不在登记域内的任务边界。
\end{certificate}
\begin{proof}
\begin{semanticsteps}
\item 完全覆盖时上界为零，与 \(D_t^G\ge0\) 合并得等式。
\item 其余两种情形使
\(\underline\theta_t(1-\overline\varepsilon_t)=0\)，
所以只能得到非负下界。
\end{semanticsteps}
\end{proof}

\section{真实需求、人类份额与劳动需求}

\begin{definition}[购买人类任务的需求]
\label{lab:def-human-demand}
对 \(j\in\mathcal O_t\)，令 \(Q_{jt}\ge0\) 为真实购买且可履约的任务数量，
\(h_{jt}\in[0,1]\) 为其中要求人类完成的份额，
\(\ell_{jt}>0\) 为每单位合格人类任务所需的劳动小时。
由此产生的劳动需求为
\[
 H_t^{\mathcal O}=\sum_{j\in\mathcal O_t}\ell_{jt}h_{jt}Q_{jt}.
\]
数量 \(Q_{jt}\) 具有实际需求、付款来源与生产资源支持；
只发布人类份额要求而没有购买者时，允许 \(Q_{jt}=0\)。
工具辅助改变 \(\ell_{jt}\) 时，必须更新该系数，
不能把技术效率提升前的工时永久作为就业下界。
\end{definition}

\begin{theorem}[正需求与正人类份额的劳动下界]
\label{lab:thm-positive-demand}
设非空集合 \(\mathcal O_t^+\subseteq\mathcal O_t\) 满足
\[
 Q_{jt}\ge\underline Q_{jt}>0,\quad
 h_{jt}\ge\underline h_{jt}>0,\quad
 \ell_{jt}\ge\underline\ell_{jt}>0
 \qquad(j\in\mathcal O_t^+).
\]
若这些订单在同一状态中联合可履约，
则
\[
 H_t^{\mathcal O}
 \ge\sum_{j\in\mathcal O_t^+}
 \underline\ell_{jt}\underline h_{jt}\underline Q_{jt}>0.
\]
若它们进一步完成且得到约定报酬，
右端同时给出实际有酬劳动小时的下界。
\end{theorem}
\begin{certificate}
保存订单、质量标准、真实需求数量、付款预算、人类份额、
单位工时与联合资源配置。
用工需求和已完成用工分别列示；
同一劳动者在重叠时间承担的多个订单不得同时计为可履约。
\end{certificate}
\begin{proof}
\begin{semanticsteps}
\item 固定 \(j\in\mathcal O_t^+\)。实际人类交付量为 \(Q_{jt}h_{jt}\)。
\item 由 \(Q_{jt}\ge\underline Q_{jt}>0\) 与
\(h_{jt}\ge\underline h_{jt}>0\)，非负乘法保持次序：
\(Q_{jt}h_{jt}\ge\underline Q_{jt}\underline h_{jt}\)。
\item 再使用单位劳动系数下界
\(\ell_{jt}\ge\underline\ell_{jt}>0\)，得到
\[
Q_{jt}h_{jt}\ell_{jt}
\ge\underline Q_{jt}\underline h_{jt}\underline\ell_{jt}>0.
\]
\item 对非空任务集合求和，其余任务的非负劳动量不会降低该下界。
\item 若另有同一合同下的实际完成回执与最终支付回执，定义中的完成与支付谓词为真，
则所需工时进入已完成且已支付工时口径。没有这两项回执时，仅保留劳动需求下界。
\end{semanticsteps}
\end{proof}

\begin{corollary}[有工时上限时的就业人数下界]
\label{lab:cor-headcount}
若实际完成上述劳动的每名主体满足
\(0<h_{it}\le\bar h_t<\infty\)，则
\[
 N_t^{\mathrm{emp}}
 \ge
 \left\lceil H_t^{\mathcal O}/\bar h_t\right\rceil.
\]
\end{corollary}
\begin{certificate}
核对实际工时、主体去重、劳动报酬与共同工时上限。
全时当量的标准值不自动充当实际工时上限。
\end{certificate}
\begin{proof}
\begin{semanticsteps}
\item 由
\(H_t^{\mathcal O}\le N_t^{\mathrm{emp}}\bar h_t\)
得人数不小于 \(H_t^{\mathcal O}/\bar h_t\)。
\item 人数为整数，向上取整即得。
\end{semanticsteps}
\end{proof}

\begin{remark}
直接保护和正需求共同承担经济结论。
前者限制特定任务被无资格自动化替代，
后者保证人类任务仍有人购买。
二者不要求人类在相同技术指标上超过自动化系统；
但人类资格条款也不单独产生订单、工资或消费能力。
\end{remark}

\section{净就业效应的无重复分解}

\begin{definition}[固定顺序的比较路径]
\label{lab:def-decomposition}
令 \(\mathcal H(x)\) 为声明经济模型在状态 \(x\) 下的总劳动小时。
选择并事前固定有限比较路径
\[
 x^{(0)}=x^0,\ x^{(1)},\ldots,x^{(m)}=x^G.
\]
第一步只改变实际直接替代，包括已经计入的规避替代，
其余步骤依次改变真实产出需求、新任务及单位劳动系数、
组织与资源配置和其他已登记的外部状态。
要求每个中间状态属于同一模型的定义域，
每个经济坐标在规定步骤中更新一次。
记
\[
 \Delta H_k=\mathcal H(x^{(k)})-\mathcal H(x^{(k-1)}).
\]
在第一步符合直接替代比较定义时，
\(\Delta H_1=A_t^{\mathrm{dir}}\)。
不同排序可以改变各项归属，故排序与模型均属于证书内容。
\end{definition}

\begin{theorem}[直接保护与净劳动效应的精确接续]
\label{lab:thm-net-effects}
令
\[
 G_t=\sum_{k=2}^{m}[\Delta H_k]_+,\qquad
 C_t=\sum_{k=2}^{m}[-\Delta H_k]_+,\qquad
 [z]_+=\max\{z,0\}.
\]
则
\[
 H_t^G-H_t^0=A_t^{\mathrm{dir}}+G_t-C_t.
\]
因而
\[
 C_t\le
 \underline\theta_t(1-\overline\varepsilon_t)D_t^0+G_t
 \quad\Longrightarrow\quad H_t^G\ge H_t^0.
\]
严格小于时，净劳动小时严格增加。
\end{theorem}
\begin{certificate}
提交全部中间反事实、同口径劳动量、更新次序与每项差分。
价格变化经需求影响劳动时，只在相应状态更新中计入；
同一成本引起的产出收缩不能再以独立损失重复扣除。
本证书区分模型恒等式、参数估计和实际因果识别。
\end{certificate}
\begin{proof}
\begin{semanticsteps}
\item 相邻差分望远镜相消，给出
\[
 \mathcal H(x^G)-\mathcal H(x^0)
 =\sum_{k=1}^{m}\Delta H_k
 =A_t^{\mathrm{dir}}+G_t-C_t.
\]
代入定理~\ref{lab:thm-direct-bound} 的下界得到充分条件。
\item 若充分条件具有严格裕量，右端也严格为正。
\item 固定 \(\bar h^{F}>0\) 后，同除该数即得全时当量的同向比较。
\end{semanticsteps}
\end{proof}

\begin{remark}
这里的 \(G_t\) 与 \(C_t\) 是所声明比较路径的正、负差分之和。
它们不是不经估计即可赋值的社会收益与社会成本，
也不声称各机制之间没有交互。
交互作用由预先固定的更新次序归属。
只证明 \(A_t^{\mathrm{dir}}>0\) 时，
净劳动小时的符号仍由未确定的 \(G_t-C_t\) 决定。
\end{remark}

\section{实际收入、价格与需求回流}

\begin{definition}[可支配现金与固定生活束]
\label{lab:def-real-income}
对家庭 \(i\)，以最终到账的工资、参与收入与其他收入，
扣除同一期间的税费和必要现金支出，定义可支配现金 \(Y_i\ge0\)。
已在净到账中扣除的费用不再次扣除。
固定物品束 \(\mathbf b_i\ge0\)，令其在状态 \(a\in\{0,G\}\) 的价格为
\[
 P_i^a=\sum_k p_k^a b_{ik}>0,\qquad
 R_i^a=Y_i^a/P_i^a.
\]
\(R_i^a\) 衡量该固定物品束的可购买数量。
它是公开声明的实际收入指标，不把家庭偏好不变作为经验结论。
\end{definition}

\begin{proposition}[收入保护转化为购买力保护的精确条件]
\label{lab:prop-real-income}
令 \(\Delta Y_i=Y_i^G-Y_i^0\)、
\(\Delta P_i=P_i^G-P_i^0\)。
则
\[
 R_i^G-R_i^0
 =\frac{P_i^0\Delta Y_i-Y_i^0\Delta P_i}{P_i^0P_i^G}.
\]
因此实际收入不下降当且仅当
\[
 \Delta Y_i\ge (Y_i^0/P_i^0)\Delta P_i.
\]
\end{proposition}
\begin{certificate}
记录逐项最终到账、税费和现金成本，以及同一品质生活束的价格。
工资承诺、未到期应收款和封闭平台积分不替代当期可支配现金。
涉及技术降低物价的主张时，另提交实际价格传导证据。
\end{certificate}
\begin{proof}
\begin{semanticsteps}
\item 对两个比值通分，再代入
\(Y_i^G=Y_i^0+\Delta Y_i\)、
\(P_i^G=P_i^0+\Delta P_i\)。
\item 分母严格为正，故差值的符号等于分子的符号。
\end{semanticsteps}
\end{proof}

\Needspace{21\baselineskip}
\begin{proposition}[给定价格下收入重分配的需求条件]
\label{lab:prop-demand-transfer}
在有限家庭集合及一个声明的局部比较区间内，
若消费支出函数具有形式 \(C_i(Y_i)=a_i+m_iY_i\)，
\(0\le m_i\le1\)，且全部资金提供者和接受者的收入变化
\(\delta_i\) 均已纳入，则
\[
 \Delta C=\sum_i m_i\delta_i.
\]
若 \(\sum_i\delta_i=0\)，净需求增加的条件仍是
\(\sum_i m_i\delta_i>0\)，而非仅有某一接受者取得正收入。
\end{proposition}
\begin{certificate}
保存消费函数的适用区间、真实收入变化及全部提供资金者。
组织转移在合并收入中一正一负对应；
不得只保留接受者支出而删除提供者被挤出的支出。
\end{certificate}
\begin{proof}
\begin{semanticsteps}
\item 逐人相减得
\(C_i(Y_i+\delta_i)-C_i(Y_i)=m_i\delta_i\)，求和即得。
\item 总收入变化为零并不使加权和必为零；
其符号取决于收入流向与各消费系数。
\item 这个局部模型没有把转移支付当作新增真实产出。
\end{semanticsteps}
\end{proof}

在这些就业、购买力与资金约束成立的候选制度之间，
最小兼容修复及净优势的比较另见第~\ref{opt:chapter}~章。

\chapter{有效参与、资金来源与履约连续性}
\label{part:chapter}

参与制度的经济对象是具有真实交付内容的合同。
本章依次建立有效参与、应付金额、实际资金、最终到账和下一期入口之间的条件链。
资格核验维护请求权归属，资金覆盖维护支付能力，
执行时序把支付能力转化为到期履行。
这些环节依据参与经济文献中的基本份额、差异分配与开放支付结构展开
\cite{InternalEconomics}。

\section{真实参与事件与最终确认集合}

\begin{definition}[规则版本下的有效参与]
\label{part:def-valid}
固定结算期 \(t\) 和事前公开的规则版本 \(R\)。
参与事件为 \(e=(i,t,\pi,\ell,R)\)，
其中 \(i\) 是规范主体，\(\pi\) 是实际行为过程，\(\ell\) 是可核验记录。
定义
\[
\begin{aligned}
 \operatorname{Valid}_R(e)\ \Longleftrightarrow{}&
 \operatorname{Admitted}_R(e)\land\operatorname{HumanAct}_R(e)\\
 &{}\land\operatorname{Compliant}_R(e)
 \land\operatorname{Completed}_R(e)
 \land\operatorname{Unique}_R(e).
\end{aligned}
\]
各谓词分别表示已接纳、约定的人类实际参与、工具及过程合规、
约定任务完成，以及主体额度和请求权标识唯一。
允许何种辅助工具由 \(R\) 明定；
符合已允许辅助方式的行为不因使用工具而自动失去资格。
账户登记为自然人和一次活动确由其按要求控制，分别核验。

本章采用每主体每期一份基本请求权。
同一主体的多个账户归并到同一个主体标识；
同一合同可以包含多次合法动作。
请求权唯一性不因已经付款而消灭真实完成事实，
重复付款由未付义务账本另行禁止。
\end{definition}

\begin{definition}[真实完成集合与确认集合]
\label{part:def-eligible-set}
组织在订约前声明最大规模 \(N_t^{\max}<\infty\)，
接纳集合为 \(\mathcal U_t\)，满足
\(|\mathcal U_t|\le N_t^{\max}\)。
真实有效参与集合与最终确认集合分别记为
\[
 \mathcal E_t^\ast
 =\{i\in\mathcal U_t:\operatorname{Valid}_R(e_{it})\},
 \qquad
 \widehat{\mathcal E}_t.
\]
\(\mathcal E_t^\ast\) 由真实事实定义，
\(\widehat{\mathcal E}_t\) 由核验与复核程序输出。
两者相等是需要证书支持的性质，不是重新命名得到的事实。
错误排除、尚未付款和提出申诉不把真实完成者从保护分母中删除。
\end{definition}

\section{正基本份额、差异份额与空集合}

\begin{definition}[事前覆盖的分配合同]
\label{part:def-allocation}
固定基本份额 \(\beta_t>0\)、差异池 \(D_t\ge0\)，
以及可在承诺期限内使用的支付专款 \(B_t^P\)，要求
\[
 B_t^P\ge N_t^{\max}\beta_t+D_t.
\]
运营、核验、复核和风险赔付资金分别列示；
同一资金不得同时覆盖这些用途与参与支付。
本章的应付额指合同约定的到账口径，
已列入其他成本项的费用不在结算时再次冲减该应付额。

令 \(N_t=|\mathcal E_t^\ast|\)。
当 \(N_t>0\) 时，对可比较的真实贡献 \(s_{it}\ge0\)，
取严格递增函数 \(g_t\)，满足 \(g_t(0)=0\)，并定义
\[
 \lambda_{it}=
 \begin{cases}
 \displaystyle
 \frac{g_t(s_{it})}{\sum_{j\in\mathcal E_t^\ast}g_t(s_{jt})},
 &\displaystyle\sum_{j\in\mathcal E_t^\ast}g_t(s_{jt})>0,\\[1.2ex]
 1/N_t,&\displaystyle\sum_{j\in\mathcal E_t^\ast}g_t(s_{jt})=0,
 \end{cases}
 \qquad
 b_{it}=\beta_t+D_t\lambda_{it}.
\]
当 \(N_t=0\) 时，不定义归一权重，不形成当期参与分配；
未使用专款按预先约定留存或返还。
空集合状态不被记为已经实现一次人类参与循环。
\end{definition}

\begin{theorem}[基本份额、贡献次序与支付覆盖]
\label{part:thm-allocation}
当 \(N_t>0\) 时，
\[
 b_{it}\ge\beta_t>0,\qquad
 \sum_{i\in\mathcal E_t^\ast}b_{it}
 =N_t\beta_t+D_t\le B_t^P.
\]
同一期同一比较域内，等贡献给出等应付额；
当 \(D_t>0\) 时，严格较高的贡献给出严格较高的应付额。
若最终确认集合正确、贡献记录正确，且逐人最终到账
\(y_{it}=b_{it}\)，则上述结论同时成立于实际到账。
\end{theorem}
\begin{certificate}
保存接纳名单、人数上限、真实过程证据、完整应得者分母、
贡献计量、资金来源及用途限制、到期时间和逐人最终性回执。
应付额证书与实际到账证书分别核对。
争议记录保留原始事件和相应资金预留。
\end{certificate}
\begin{proof}
\begin{semanticsteps}
\item 两个非空分支均满足
\(\lambda_{it}\ge0\) 和 \(\sum_i\lambda_{it}=1\)；
逐人下界及总额恒等式随之成立。
\item 因 \(N_t\le N_t^{\max}\)，总额不超过预留款。
\item 若 \(s_{it}>s_{jt}\)，分母严格为正，且
\[
 b_{it}-b_{jt}
 =D_t\frac{g_t(s_{it})-g_t(s_{jt})}
 {\sum_{\ell\in\mathcal E_t^\ast}g_t(s_{\ell t})}.
\]
当且仅当此处的差异池严格为正时，该表达式严格为正。
\item 等贡献时分子为零；全零贡献时按等份分配。
\item 最后，正确主体对应与 \(y_{it}=b_{it}\)
把应付命题接续为实际支付命题。
\end{semanticsteps}
\end{proof}

\begin{corollary}[基本份额与无资格自动体的表现分离]
\label{part:cor-machine-performance}
在规则、真实需求、资金、接纳规模与核验条件保持的比较中，
提高无资格自动体的任务表现不改变有效参与者的基本下界
\(b_{it}\ge\beta_t\)；
实际履约条件成立时，还有 \(y_{it}\ge\beta_t\)。
\end{corollary}
\begin{certificate}
分别登记自动体的表现指标与其合同资格，
并记录比较中保持不变的需求、资金和人数条件。
\end{certificate}
\begin{proof}
\begin{semanticsteps}
\item 无资格自动体不属于 \(\mathcal E_t^\ast\)，
其表现不进入真实参与集合的分母。
\item 基本下界由有效完成与固定的 \(\beta_t\) 决定，
直接应用定理~\ref{part:thm-allocation}。
\end{semanticsteps}
\end{proof}

\section{真实资金来源与合并现金恒等式}

\begin{definition}[来源分账与现金边界]
\label{part:def-funds}
在同一结算期内，令 \(R_t^{\mathrm{op}}\) 为核算参与者集合之外
的真实服务购买、赞助或其他经营合同所形成的已到账收入；
\(K_t^{\mathrm{in}}\) 为已到账新增融资，
\(K_t^{\mathrm{out}}\) 为融资退出、归还本金及其他已单列的融资流出；
\(U_t\) 为已到账外部转移。
参与者向组织支付的相关费用记为 \(f_{it}\ge0\)，
不再计入上述经营收入。
令 \(C_t\ge0\) 为组织的真实运营现金支出，
不包含参与者分配和已经单列的融资流出。
期初、期末可用现金分别为 \(B_t^{\mathrm{beg}},B_t^{\mathrm{end}}\)。
所有流量具有唯一交易标识，并使用相同币种、时期和合并边界。
融资所得不因进入支付池而成为经营收入；
融资新增的未来义务继续进入对应负债记录。
\end{definition}

\begin{theorem}[参与净流入的现金约束]
\label{part:thm-source-balance}
在来源和用途完整、没有重复计数时，
\[
\begin{aligned}
 \sum_i y_{it}+C_t+B_t^{\mathrm{end}}
 ={}&B_t^{\mathrm{beg}}+R_t^{\mathrm{op}}
 +K_t^{\mathrm{in}}-K_t^{\mathrm{out}}+U_t+\sum_i f_{it}.
\end{aligned}
\]
令 \(R_t=R_t^{\mathrm{op}}+K_t^{\mathrm{in}}-K_t^{\mathrm{out}}+U_t\)，
\(\Delta B_t=B_t^{\mathrm{end}}-B_t^{\mathrm{beg}}\)，则
\[
 \sum_i(y_{it}-f_{it})=R_t-C_t-\Delta B_t.
\]
若 \(R_t=0\)、\(C_t>0\) 且 \(\Delta B_t\ge0\)，
参与者整体现金净流入严格为负。
\end{theorem}
\begin{certificate}
逐笔保存来源账户、交易标识、到账时间、用途和融资性质。
组织内部转账在合并时抵销；未到账应收款不计入当期现金。
参与者支出与组织收入的同一交易只保留一组对应项。
\end{certificate}
\begin{proof}
\begin{semanticsteps}
\item 将各账户的期初余额加真实流入、减真实流出，
分别等于期末余额的恒等式相加，
内部转账一正一负抵销，得到第一式。
\item 移项得到净流入式，再代入三个符号条件得严格负值。
\item 储备消耗对应 \(\Delta B_t<0\)，可以支持有限期支付，
但不改变期初储备有限这一事实。
\end{semanticsteps}
\end{proof}

\section{已成立义务与偿付覆盖的数学归纳}

\begin{definition}[逐主体逐事件义务账本]
\label{part:def-ledger}
以最终确认的结算事件为原子操作。
状态保存可及时履约资产 \(B_n\ge0\)，
主体 \(i\) 对合同事件 \(k\) 的真实未付义务 \(L_{ik,n}\ge0\)，
累计最终到账 \(Y_{ik,n}\ge0\)，
以及累计真实成立权利 \(\Gamma_{ik,n}\)。
令 \(L_n=\sum_{i,k}L_{ik,n}\)，初态满足
\[
 B_0\ge L_0,\qquad
 L_{ik,0}+Y_{ik,0}=\Gamma_{ik,0}.
\]
支付只允许针对已成立且主体、事件、版本相匹配的义务：
\[
 0\le u\le L_{ik,n},\qquad
 B'=B-u,\quad L'_{ik}=L_{ik}-u,\quad Y'_{ik}=Y_{ik}+u.
\]
其他权利人的义务不变。
付款发起但尚未最终结算时，不先行增加 \(Y_{ik}\)。

新增真实义务 \(v_{ik}\ge0\)，总额 \(v=\sum_{i,k}v_{ik}\)；
尚未计入资产的新入金为 \(w\ge0\)。
接受该操作的条件及更新为
\[
 B-L+w\ge v,\qquad
 B'=B+w,\quad L'_{ik}=L_{ik}+v_{ik},\quad
 \Gamma'_{ik}=\Gamma_{ik}+v_{ik}.
\]
已经计入 \(B\) 的预付款不得再次计作 \(w\)。
从同一账户支付运营支出 \(c\) 时，要求 \(0\le c\le B-L\)；
已从支付专款口径扣除的成本不得再次作该项支出。
无资格请求、虚假临时登记与真实义务分账，
其拒绝或撤销不减少真实 \(L_{ik}\)。
\end{definition}

\begin{theorem}[偿付覆盖与权利连续性的生成保持]
\label{part:thm-ledger-invariant}
由上述原子操作构成的任意有限合法执行满足
\[
 B_n\ge L_n\ge0,\qquad
 L_{ik,n}+Y_{ik,n}=\Gamma_{ik,n}.
\]
真实支付保持 \(B-L\)，新增义务与允许支出保持其非负性。
仅改变未来规则不取消已经成立的真实权利。
\end{theorem}
\begin{certificate}
保存初态、每笔匹配义务、实际金额、资金变化、最终性回执及新入金唯一标识。
版本迁移保留旧合同的权利人与欠付金额。
资金须在承诺期限内可用，尚未实现的追偿不作为当前可用资产。
\end{certificate}
\begin{proof}
\begin{semanticsteps}
\item 支付后 \((B-u)-(L-u)=B-L\)，
且 \(u\le L_{ik}\) 保证未付义务非负；
同一义务的 \(L_{ik}+Y_{ik}\) 保持。
\item 新增义务后
\[
 B'-L'=B-L+w-v\ge0,
\]
而 \(L_{ik}\) 与 \(\Gamma_{ik}\) 同增 \(v_{ik}\)。
\item 允许支出后 \(B'-L'=B-L-c\ge0\)，权利明细不变。
\item 拒绝无资格请求和未来版本更新不改变真实义务。
\item 初态满足两项不变量，每个允许操作保持它们，
故按操作次数进行数学归纳，任意有限前缀都满足结论
\cite{InternalTheory,InternalEconomics}。
\end{semanticsteps}
\end{proof}

\begin{theorem}[指定支付次序的资金前缀条件]
\label{part:thm-prefix-funding}
固定有限顺序 \(n=1,\ldots,M\)。
令 \(\rho_n\ge0\) 为该步开始前实际到达的新资金，
\(u_n\ge0\) 为该步到期且匹配已成立义务的支付，
\(c_n\ge0\) 为未在其他项扣除的必要现金支出。
在权限、非现金资源和支付设施均满足时，
该次序无新增信用地完成全部支付与支出，
当且仅当对每个 \(m\le M\) 都有
\[
 B_0+\sum_{n=1}^{m}\rho_n
 \ge\sum_{n=1}^{m}(u_n+c_n).
\]
\end{theorem}
\begin{certificate}
按真实到账先后排列事件；
同一步中的资金只有在支出前可用时才能作为 \(\rho_n\)。
保存全部到期义务，不能通过遗漏未付款人使条件成立。
该证书与未到期义务的资产预留证书分别核验。
\end{certificate}
\begin{proof}
\begin{semanticsteps}
\item 完成前 \(m\) 步后的余额恰为左端减右端。
\item 可执行性要求每个余额非负，给出必要性。
\item 反之，若全部前缀满足不等式，
则每步可用现金覆盖本步应支出额；
结合非现金前件逐步执行，归纳得到充分性。
\item 只检查 \(m=M\) 不足以控制此前已经出现的资金缺口。
\end{semanticsteps}
\end{proof}

\begin{proposition}[有限队列中的及时履约]
\label{part:prop-timely-payment}
从时刻 \(\tau_0\) 起，一个完整批次至多含 \(M\) 个处理项，
包括全部真实请求、无效请求及必要复核。
若批次未完成时调度不空转，批外请求不能插队占用其窗口，
每项完整处理至最终结算的时间不超过 \(d\)，
且所需资金与设施持续可用，则批次至迟在 \(\tau_0+Md\) 完成。
\end{proposition}
\begin{certificate}
记录完整队列、消息与重试时间、服务上界、流量隔离和执行责任。
若承诺截止时间为 \(\tau^\ast\)，另核对 \(\tau_0+Md\le\tau^\ast\)。
\end{certificate}
\begin{proof}
\begin{semanticsteps}
\item 第一项至多耗时 \(d\)。
\item 前 \(k\) 项至多耗时 \(kd\) 时，
下一项连同全部必要通信至多再用 \(d\)，
故前 \(k+1\) 项至多耗时 \((k+1)d\)。
\item 无空转及无批外插队使归纳没有遗漏的等待时间。
\end{semanticsteps}
\end{proof}

\section{再接入、联合容量与有限期连续性}

\begin{definition}[同一可达状态中的接续计划]
\label{part:def-reentry}
令 \(x_t\) 为本期真实履约完成后的实际状态。
\(\kappa_{it}\ge0\) 是尚需由本期收入承担、
且尚未在到账额 \(y_{it}\) 中扣除的现金成本。
令 \(\pi_{i,t+1}\) 为从同一个 \(x_t\) 出发的下一期合法接入计划，
其任何前缀的累计必要现金支出不超过 \(k_{i,t+1}\ge0\)。
计划另包含真实可用的接纳容量、时间、必要能力与允许工具。
时间和努力的效用成本单列；
只有对应实际现金义务时才计入现金预算。
参与者可以选择退出，退出不消灭旧合同债权。
\end{definition}

\begin{theorem}[资源与容量共同支持的再接入]
\label{part:thm-reentry}
若真实到账满足 \(y_{it}\ge\beta_t\)，且
\[
 y_{it}-\kappa_{it}\ge k_{i,t+1},
\]
到账早于接续支出的期限，
并且 \(\pi_{i,t+1}\) 的非现金前件在 \(x_t\) 中成立，
则参与者具有至少一条可负担的下一期合法接入路径。
若只能使用基本份额的保证，则 \(\beta_t\ge\kappa_{it}+k_{i,t+1}\) 是上述现金条件的一个充分条件。
若要使全部本期参与者能够同时接入，
还须这些路径在同一个状态中具有不冲突的联合容量与资源配置。
\end{theorem}
\begin{certificate}
记录同一末态、逐人到账时间、未扣成本、路径的各前缀支出，
以及下一期资格、席位、时间和设备配置。
联合结论另给出完整席位分配；
分别存在于不同反事实状态中的单人入口不合并为联合入口。
\end{certificate}
\begin{proof}
\begin{semanticsteps}
\item 可用于接续的现金由前件满足 \(y_{it}-\kappa_{it}\ge k_{i,t+1}\)。
\item 在仅使用基本份额的分支中，
\[
 y_{it}-\kappa_{it}
 \ge\beta_t-\kappa_{it}\ge k_{i,t+1}.
\]
每个路径前缀的累计必要现金支出均不超过此数，
并且非现金前件与时序成立，故逐步执行该有限计划可行。
\item 联合容量配置使多个路径的资源使用不会相互否定。
\item 是否实际执行仍是参与者的选择。
\end{semanticsteps}
\end{proof}

\Needspace{16\baselineskip}
\begin{corollary}[有限期自愿参与的连续性]
\label{part:cor-finite-continuity}
固定有限期限 \(T\)。
若从初态开始，每期实际到达状态均满足有效参与、
足额资金、及时履约及下一期联合接入条件，
则选择继续参与的主体能够沿这些证书完成至第 \(T\) 期的接续。
\end{corollary}
\begin{certificate}
逐期保存实际末态、当期外部资金、未付义务和下一期容量配置。
后续资金不以尚未发生的参与费或未经承诺的融资补写为已到账。
\end{certificate}
\begin{proof}
\begin{semanticsteps}
\item 初期入口由初态证书给出。
\item 假设已经合法到达第 \(t\) 期，履约与再接入定理给出到达第 \(t+1\) 期的可行路径。
\item 对 \(t<T\) 归纳即得。
\item 有限证书不包含期限以后的资金与容量，因而不推出无限期连续运行。
\end{semanticsteps}
\end{proof}

\chapter{劳动合同权利与按需探索服务}
\label{integrated:labour-access}
\label{rights:chapter}

本章将已经成立的劳动请求权与申请公共探索服务的资格分别建模。
前者由真实履约、合同规则和实际付款连接；
后者由明确的服务对象、真实需求、用途条件及有限资源连接。
两种权利可以属于同一个人，但不因此采用同一分配分母或同一资格条件。
本章沿用参与支付与资源相容的分析接口\cite{InternalEconomics}，
并把自愿选择、既得债权、复核可达和探索资格独立性落实为动作及证书。
这些性质是声明制度中的可检验要求；正收入本身不替代全部要求。

\section{两类权利、真实来源与证据状态}

\begin{definition}[劳动请求权与探索资格]
\label{rights:two-objects}
固定规则版本 \(\nu\)、有限主体集合 \(\mathcal I\)、结算期和共同资源边界。
同一自然人的账户归并到唯一主体标识 \(i\)。
劳动事件 \(e=(i,k,\nu,\pi,\ell)\) 包含合同标识 \(k\)、真实过程 \(\pi\) 和证据 \(\ell\)。
其真实有效性采用定义~\ref{part:def-valid} 的合取：
\[
\begin{aligned}
T^L_\nu(e)\Longleftrightarrow{}&
\operatorname{Admitted}_\nu(e)\land\operatorname{HumanAct}_\nu(e)\\
&{}\land\operatorname{Compliant}_\nu(e)
\land\operatorname{Completed}_\nu(e)\land\operatorname{Unique}_\nu(e).
\end{aligned}
\]
劳动账本保存真实成立额 \(\Gamma_{ik}\)、未付额 \(L_{ik}\) 与最终到账 \(Y_{ik}\)。
允许的辅助方式由合同明定，不能仅因使用已允许工具否认真实劳动。

探索资格另写为
\[
E^X_\nu(i)=
\operatorname{Person}(i)\land\operatorname{UniqueId}_\nu(i)
\land\operatorname{Scope}_\nu(i)\land\operatorname{SafeUse}_\nu(i).
\]
其中各项采用事前公开的主体、服务范围与用途条件。
真实需求 \(n_i\ge0\) 在相同服务类别和期限内登记。
本章选择的基本探索资格不要求完成某份劳动合同、
取得某个劳动评分、拥有某个财富水平或接受另一项服务部署。
资格不等于无限数量的交付承诺；具体配额还受需求、资金与容量约束。
\end{definition}

\begin{definition}[事实判断与程序判断]
\label{rights:source-status}
对任一真实谓词 \(T(e)\in\{0,1\}\)，程序状态记为
\(\widehat T(e)\in\{1,0,?\}\)，问号表示尚待核验。
真实来源包括实际行为者、控制过程、用途、事件与规则版本；
证据状态包括收集到的材料、检查输出和复核阶段。
检查器健全性是独立条件
\[
\operatorname{Sound}_\nu
\Longrightarrow
\bigl(\operatorname{Check}_\nu(\ell,e)=1\Longrightarrow T_\nu(e)=1\bigr).
\]
及时识别全部真实事件还需要另外的完整性、证据可得性和处理时限条件。
尚未核验与真实不成立不作同义处理。
\end{definition}

\begin{proposition}[结果、来源与检查结果分别承担证明义务]
\label{rights:source-evidence}
相同交付结果不一般蕴含相同劳动资格。
在检查器健全、证据真实且解释映射有效时，通过检查可以支持真实资格结论；
单独持有通过标记不能代替这些前件。
\end{proposition}
\begin{certificate}
保存规范主体、合同、规则版本、控制过程、结果及其对应证据。
非蕴含证书取两个结果相同、其余资格条件相同，
但 \(\operatorname{HumanAct}_\nu\) 分别为一和零的事件。
检查结论的适用证书另核验来源、解释和实现。
\end{certificate}
\begin{proof}
\begin{semanticsteps}
\item 设两个事件 \(e,e'\) 的结果观察满足 \(U(e)=U(e')\)。
结果相等只给出关于 \(U\) 的等式。
\item 在所给反例中，其余四个资格谓词均为真，
而 \(\operatorname{HumanAct}_\nu(e)=1\)、\(\operatorname{HumanAct}_\nu(e')=0\)。
逐项代入合取，得到 \(T^L_\nu(e)=1\) 和 \(T^L_\nu(e')=0\)。
\item 健全性前件与 \(\operatorname{Check}_\nu(\ell,e)=1\) 合取，
由定义~\ref{rights:source-status} 的蕴含得到 \(T_\nu(e)=1\)。
\item 来源、模型和实现条件按命题~\ref{base:certificate} 逐项成立时，
上述模型结论才适用于实际事件。删除任一所需前件不能保留该推演。
\end{semanticsteps}
\end{proof}

\begin{remark}[待复核权利]
真实债权由已履行合同及适用规则确定，程序误判不能使该事实消失。
账面识别不完整时，应为尚待处理的已接纳合同保留其最大可能合法义务所需资金。
这种保留是风险准备，不是把每个未经核实的请求直接登记为真实债权。
复核改正识别错误时，不再次生成同一事件已经成立的 \(\Gamma_{ik}\)。
\end{remark}

\section{自愿选择、债权保持与可操作制度尊严}

\begin{definition}[权利状态与可操作判据]
\label{rights:state}
状态保存每份合同的已接受条款 \(\nu_{ik}\)、本人同意记录、
新劳动是否仍可请求的标志 \(a_{ik}\)，以及最近退出记录。
另保存 \((\Gamma_{ik},L_{ik},Y_{ik})\)、复核请求及已预留的资金和处理时段、
探索资格的允许判定事实 \(z_i^X\)。
对本章声明的制度，定义
\[
\begin{aligned}
\mathsf D_i:={}&
\mathsf{ChoiceAvailable}_i
\land\mathsf{TermsPreserved}_i
\land\mathsf{EarnedRightsPreserved}_i\\
&{}\land\mathsf{AffordableReview}_i
\land\mathsf{BasicAccessIndependent}_i .
\end{aligned}
\]
\(\mathsf{ChoiceAvailable}\) 要求拒绝新合同和按约退出均有真实可执行路径；
不能仅以界面中出现一个选项代替该路径。
\(\mathsf{TermsPreserved}\) 要求已接受条款能够读回，且变更有明示的适用边界。
\(\mathsf{EarnedRightsPreserved}\) 要求真实债权被保留并按到期义务处理。
其余两项分别要求可负担的复核路径和不以劳动评分、财富为条件的基本探索资格。
该合取是本章对制度尊严的操作性判据，不声称穷尽人的全部价值。
\end{definition}

\begin{definition}[允许的权利动作]
\label{rights:actions}
允许动作按以下守卫与更新执行，所有响应均须保持相应后置条件。
\begin{enumerate}
\item \textbf{提出与接受。}
提出合同只登记完整条款，不产生已完成劳动的请求权。
接受新合同要求可核验的本人选择、条款可得及相应人数和资金预留。
曾经退出的主体重新进入时，必须提供退出之后针对新合同的同意记录。
\item \textbf{拒绝新合同。}
本人拒绝时，不启动该合同的新劳动，不产生未约定费用，
并保持其他已成立的 \(\Gamma_{ik},L_{ik},Y_{ik}\) 及 \(z_i^X\)。
\item \textbf{退出。}
退出按照已接受的有限通知、交接或结算程序关闭未来劳动请求标志 \(a_{ik}\)。
已完成劳动的债权不因退出减少；既有清偿义务也不因退出被删除。
必要交接的时限、资源和费用必须事前列明并具有实际路径证书。
\item \textbf{履约与支付。}
\(T^L_\nu(e)=1\) 的真实完成生成一次相匹配的新义务；
资金守卫和支付更新采用定义~\ref{part:def-ledger}。
支付额 \(u\) 必须满足 \(0\le u\le L_{ik}\)，并经过主体、事件与版本匹配。
未最终到账时不增加 \(Y_{ik}\)。
\item \textbf{复核。}
有复核资格的主体可以提交请求，系统按预留资源登记、处理并交付理由和决定。
提出复核不取消旧债权，不降低基础探索资格，也不自动扣减已经约定的到账额。
正确复核改变的是错误识别；真实权利仍按原事件计数一次。
\item \textbf{未来规则更新。}
新版本不得覆盖旧合同的已接受条款、抹除真实债权或撤销已成立的复核义务。
当允许判定事实 \(z_i^X\) 没有变化时，
仅改变劳动选择、劳动评分或财富不得改变基本探索资格。
\end{enumerate}
没有写入这些动作的处罚、费用和资格联动，不能在执行时临时加入。
实际用途或安全事实改变时的独立审查另作登记，不能伪装为对拒绝劳动的处罚。
\end{definition}

\begin{theorem}[选择记录、真实债权与资格分离的有限保持]
\label{rights:invariant}
若初态满足已接受条款可读回、真实义务有资金覆盖，
且每个实际入口仅执行定义~\ref{rights:actions} 的动作，
则任意允许有限执行保持：
新劳动有相匹配的本人同意；退出后没有未经新同意启动的劳动；
\[
L_{ik}+Y_{ik}=\Gamma_{ik},\qquad
B^L\ge\sum_{i,k}L_{ik}\ge0;
\]
拒绝、退出、复核及仅涉及劳动的规则变更不自行降低基本探索资格。
\end{theorem}
\begin{certificate}
初态权利明细、条款快照、同意和退出的时间顺序、
全部动作入口、资金守卫、响应覆盖及最终到账回执。
对资格保持另给出动作前后 \(z_i^X\) 相等的检查。
对新义务保存唯一事件标识和未重复计入的新入金。
\end{certificate}
\begin{proof}
\begin{semanticsteps}
\item 将合同同意、条款保持、债权恒等式、资金覆盖和资格分离写为合取 \(I(x)\)。
初态证书逐项给出 \(I(x_0)\)。
\item 接受动作的守卫给出本人同意与已接受条款；
退出后重新接受的时间守卫排除使用退出前的旧同意。
提出或拒绝合同不启动新劳动，因而保持同意条件。
\item 退出动作将未来请求标志关闭，且不修改旧 \(\Gamma_{ik},L_{ik},Y_{ik}\)。
所以债权恒等式保持；后续新劳动仍受新同意守卫约束。
\item 对相匹配的实际支付，更新给出
\[
(L_{ik}-u)+(Y_{ik}+u)=L_{ik}+Y_{ik},
\quad
(B^L-u)-\bigl(\textstyle\sum L-u\bigr)=B^L-\sum L.
\]
守卫 \(u\le L_{ik}\) 保证未付额非负。
\item 新真实义务和允许运营支出分别满足原账本的资金守卫。
由定理~\ref{part:thm-ledger-invariant}，它们保持资金覆盖和权利恒等式。
修正程序记录不能被当作同一真实事件的第二次新增。
\item 拒绝、退出与复核动作保持 \(z_i^X\)；
资格规则对相同 \(z_i^X\) 给出相同 \(E^X_\nu(i)\)。
未来规则更新按守卫保留既有合同和债权。
\item 上述逐项保持覆盖每个允许动作及全部登记响应。
应用定理~\ref{base:induction}，按实际动作次数归纳得到任意有限前缀的 \(I\)。
\end{semanticsteps}
\end{proof}

\begin{lemma}[已预留权利程序的有限完成]
\label{rights:procedure}
设退出或复核程序有整数剩余秩 \(r\in\{0,\ldots,m\}\)；
\(r=0\) 当且仅当规定的终结回执已经交付。
若尚未完成时，每次处理对全部允许响应都使 \(r\) 至少下降一，
下一次处理连同必要等待至多耗时 \(d\)，且资金、人员和权限持续预留，
则程序至多经过 \(m\) 次处理、在 \(md\) 时间内完成。
\end{lemma}
\begin{certificate}
实际程序、初始秩上界、每项响应的秩下降、处理与等待上界、
不可被其他任务重复使用的人员和资金，以及终结回执的定义。
\end{certificate}
\begin{proof}
\begin{semanticsteps}
\item 初始证书给出 \(0\le r_0\le m\)。
\item 若第 \(j\) 次处理前 \(r_j>0\)，下降证书给出
\(0\le r_{j+1}\le r_j-1\)。
\item 归纳得到尚未完成时 \(r_j\le m-j\)；
故第 \(m\) 次之后不可能仍有 \(r_m>0\)。
\item \(r_m=0\) 按定义给出终结回执。
每次处理和等待至多耗时 \(d\)，故总时长不超过 \(md\)。
\item 持续预留排除推演中缺钱、缺少权限或队列被其他任务无限占用的分支。
只保存安全状态而没有下降证书时，不能作出本项完成推论。
\end{semanticsteps}
\end{proof}

\section{复核的现金可负担性与实际入口}

\begin{definition}[事前覆盖的复核批次]
\label{rights:review-budget}
固定一个完整批次，包含全部应当处理的请求及必要的身份确认和补充材料程序，
请求数至多为 \(M_A<\infty\)。
对请求 \(j\)，预留的现金上界 \(\bar c_j^A\) 覆盖依法必须支付且未在其他项列支的全部程序费用。
另以 \(C_A\) 记录不属于这些逐请求费用的共同处理支出。
可用专款满足
\[
B_A\ge\sum_{j=1}^{M_A}\bar c_j^A+C_A.
\]
付款安排不得要求申请者先垫付其无力承担的必要费用：
可以直接支付程序提供者，或者在费用发生前提供相应专款。
个人时间、无障碍条件和证据取得能力另以真实可执行路径核验，
不能由现金补助自动推定。
同一费用若已经由程序预算承担，不再从劳动到账额中重复扣除。
\end{definition}

\begin{theorem}[有资金和容量的可负担复核路径]
\label{rights:affordable-review}
若请求属于定义~\ref{rights:review-budget} 的完整批次，
每个费用发生前都有相应专款可用，
且同一实际状态中具有相容的证据取得、时间和处理容量，
则每位请求者存在不依赖新增个人垫资的复核路径。
若各程序满足引理~\ref{rights:procedure}，且调度上界覆盖完整队列，
则每项请求在其登记期限内取得终结回执。
\end{theorem}
\begin{certificate}
完整请求清单、逐项费用和支付时点、专款来源、队列与席位分配、
申请者实际可用的提交方式、必要协助及每项程序的秩证书。
若要求决定与真实事实一致，另提交核验健全性和纠错完整性证书。
\end{certificate}
\begin{proof}
\begin{semanticsteps}
\item 对固定请求 \(j\)，每个必要费用前缀均不超过其预留上界 \(\bar c_j^A\)。
付款时点证书使该前缀所需资金在支出前可用。
\item 直接支付或预先供款使请求者在该前缀的必要自有垫资为零。
对所有前缀重复这一检查，得到整条路径的现金可负担性。
\item 共同容量和时间见证在同一个状态中成立，
故不同请求的可执行路径不会占用同一不可共享席位。
这一步使用引理~\ref{base:joint} 的共同见证要求。
\item 各路径的非现金守卫成立，因而可以实际进入处理程序。
完整调度与秩下降满足引理~\ref{rights:procedure}，推出有限期限完成。
\item 终结回执证明程序已经完成。
若还要确认裁决正确，必须再合取独立的事实识别证书；
完成程序本身不蕴含申请者必然胜诉。
\end{semanticsteps}
\end{proof}

\begin{remark}[可操作判据的边界]
拒绝与退出的权利不是强制当事人退出；复核权也不是强制申诉。
制度只需保留真实可达的选项并履行已触发的义务。
现金费用得到覆盖不表示时间成本为零，也不表示任何证据都能获得。
因此，\(\mathsf D_i\) 的完成证书同时读取动作、债权、资源和回执，
不由同意按钮、正付款或一项程序名称单独签发。
\end{remark}

\section{可比较劳动、分配完备性与净收益}

\begin{definition}[有声明单位的劳动贡献]
\label{rights:comparable-labour}
在同一结算期和可比较合同域内，令 \(\mathcal E^L\) 为真实有效参与者集合，
\(M_L=|\mathcal E^L|\le M_L^{\max}<\infty\)。
对事先定义的质量等级 \(g\in\mathcal G_L\)，记录实际完成且不重复计数的劳动小时 \(H_{ig}\ge0\)。
若合同提供了适用的等级折算系数 \(\omega_g>0\)，定义
\[
s_i^L=\sum_{g\in\mathcal G_L}\omega_g H_{ig}.
\]
同一劳动小时不得因为出现在多个记录中而重复进入该和。
若跨等级折算没有依据，只在同一等级或其他已声明的可比较域内排序。
贡献分数是合同计量，不是对人的价值或探索资格的评价。
\end{definition}

\begin{definition}[与既有支付规则相容的完备分支]
\label{rights:allocation}
固定 \(\beta>0\)、\(D\ge0\)，并预留
\[
B^L_P\ge M_L^{\max}\beta+D.
\]
取严格递增函数 \(g_L\)，满足 \(g_L(0)=0\)。
当 \(M_L>0\) 时，令 \(S_L=\sum_{j\in\mathcal E^L}g_L(s_j^L)\)，并取
\[
\lambda_i=
\begin{cases}
g_L(s_i^L)/S_L,&S_L>0,\\
1/M_L,&S_L=0,
\end{cases}
\qquad b_i=\beta+D\lambda_i.
\]
当 \(M_L=0\) 时，不定义归一权重，不产生本批次的劳动分配义务；
专款按已约定的留存或返还规则处理。
贡献和为零的非空分支不要求 \(D=0\)，而是仍按等份分配差异池。
这与定义~\ref{part:def-allocation} 使用同一合同规则。
\end{definition}

\begin{theorem}[正基本份额、贡献次序及实际到账]
\label{rights:pay-coverage}
在定义~\ref{rights:allocation} 的非空分支中，
\[
b_i\ge\beta>0,\qquad
\sum_i b_i=M_L\beta+D\le B^L_P.
\]
相同贡献对应相同应付额；当 \(D>0\) 时，
严格较高的可比较贡献对应严格较高的应付额。
若真实主体集合和贡献被正确确认，且匹配的义务已经最终足额支付，
则 \(y_i=b_i\)，上述比较同时成立于实际到账。
\end{theorem}
\begin{certificate}
真实接纳及履约名单、唯一主体、小时与等级记录、折算规则、
空集合或贡献和分支、预留款、逐人应付及最终性回执。
未付款、待复核或提出申诉不成为删除真实应得者的理由。
\end{certificate}
\begin{proof}
\begin{semanticsteps}
\item \(H_{ig}\ge0\) 与 \(\omega_g>0\) 给出 \(s_i^L\ge0\)。
函数条件给出 \(g_L(s_i^L)\ge0\)。
\item 若 \(S_L>0\)，归一式给出 \(\lambda_i\ge0\) 和 \(\sum_i\lambda_i=1\)。
若 \(S_L=0\)，非空性给出 \(M_L\ge1\)，故等份分支也满足这两式。
\item 代入 \(b_i=\beta+D\lambda_i\)，得到逐人下界与总额恒等式。
再以 \(M_L\le M_L^{\max}\) 和预留条件得到支付覆盖。
\item 若 \(s_i^L>s_j^L\)，则 \(S_L>0\)，且
\[
b_i-b_j
=D\,\frac{g_L(s_i^L)-g_L(s_j^L)}{S_L}.
\]
严格递增性与 \(D>0\) 使右端严格为正；\(D=0\) 时该差为零。
\item 正确确认与最终付款给出 \(y_i=b_i\)。
这正是定理~\ref{part:thm-allocation} 从合同额接续到到账额的实际实例。
空集合分支不使用上述归一式，也不声称已经产生一轮劳动收入。
\end{semanticsteps}
\end{proof}

\begin{proposition}[净现金、机会成本与实际购买力分别核验]
\label{rights:net-income}
令 \(\kappa_i\ge0\) 为尚未从 \(y_i\) 扣除的必要参与现金成本，
\(o_i\ge0\) 为在声明比较尺度下另行计量、未包含于 \(\kappa_i\) 的机会成本。
定义合同净现金 \(C_i^L=y_i-\kappa_i\)，相对于该机会成本的私人净收益
\(V_i^L=y_i-\kappa_i-o_i\)。
则
\[
C_i^L\ge\beta-\kappa_i,\qquad
V_i^L\ge\beta-\kappa_i-o_i.
\]
所以 \(\beta>\kappa_i+o_i\) 是严格正私人净收益的充分条件。
对同一家庭的完整可支配现金 \(Y_i^a\) 和固定生活束价格 \(P_i^a>0\)，
\(a\in\{0,1\}\)，实际购买力不下降的充要条件仍为
\[
\Delta Y_i\ge (Y_i^0/P_i^0)\Delta P_i.
\]
\end{proposition}
\begin{certificate}
实际到账、未扣现金成本、机会成本的反事实与计量尺度、
家庭全部收入和支出的无重复清单，以及同一品质生活束的价格。
若其他收入的损失已经进入反事实差额，不在同一净收益指标中再次扣除。
\end{certificate}
\begin{proof}
\begin{semanticsteps}
\item 实际付款条件给出 \(y_i\ge\beta\)。
两边减去相同 \(\kappa_i\)，得到净现金下界。
\item 再减去未重复计入的 \(o_i\)，得到私人净收益下界。
若 \(\beta-\kappa_i-o_i>0\)，则 \(V_i^L>0\)。
\item 对完整家庭指标直接通分：
\[
\frac{Y_i^1}{P_i^1}-\frac{Y_i^0}{P_i^0}
=\frac{P_i^0\Delta Y_i-Y_i^0\Delta P_i}{P_i^0P_i^1}.
\]
\item 分母为正，故购买力差非负等价于分子非负。
除以 \(P_i^0>0\)，得到所列条件，
也即命题~\ref{lab:prop-real-income} 在本章支付情形中的应用。
\end{semanticsteps}
\end{proof}

\begin{remark}[正报酬的经济边界]
正基本份额不单独保证扣除全部机会成本后仍有正收益，
也不单独保证生活物价变化后的购买力提高。
自愿参加的选项与参加后的实际结果分别登记。
劳动收入、服务补助、提供者收入和总福利不能相互代替；
其中货币转移的合并关系继续服从定理~\ref{part:thm-source-balance}。
\end{remark}

\section{有限同类可分探索服务的按需配额}

\begin{definition}[需求、公共购买预算与可执行容量]
\label{access:divisible}
固定一种可分的同类探索服务，其质量、单位、用途许可和完成期限相同。
令 \(\mathcal I_X\) 为在完整程序中已正确确认资格、并自愿登记本批次需求的有限集合。
每人需求 \(n_i\in[0,\infty)\) 须满足公开规则的真实条件
\(\operatorname{NeedValid}_\nu(i,n_i)\)，并通过相应正确核验；
申报数量不自动等于真实需求。总需求与人数分别记为
\[
N=\sum_{i\in\mathcal I_X}n_i,\qquad M_X=|\mathcal I_X|.
\]
\(N\) 是服务需求总量，不能在需求不同的情形用人数 \(M_X\) 代替。
令 \(K\ge0\) 为实际可用容量，\(p>0\) 为合同单位价格，
\(B\ge0\) 为已经到位、可在交付期限使用的公共购买预算。
身份核验、申请协助、复核及共同管理费用另有覆盖，
或者已经从本项预算口径扣除；同一费用不再扣第二次。
本批次允许个人无需付款取得所分配的服务，不以先垫资作为准入条件。

当 \(N>0\) 时，定义
\[
\alpha=\min\left\{1,\frac K N,\frac B{pN}\right\},
\qquad a_i=\alpha n_i.
\]
当 \(N=0\) 时，定义 \(\alpha=0\) 和全部 \(a_i=0\)，不计算含 \(N\) 的比值。
配额 \(a_i\) 是允许交付的服务量；实际交付另由记录 \(d_i\) 确认。
\end{definition}

\begin{theorem}[按需配额的完备边界与资源覆盖]
\label{access:proportional}
上述规则在全部登记的非负有限参数上有定义，并满足
\[
0\le a_i\le n_i,\qquad
\sum_i a_i\le K,\qquad
\sum_i pa_i\le B.
\]
若 \(N>0\) 且 \(K=0\) 或 \(B=0\)，则全部配额为零。
若 \(N>0\)、\(K>0\)、\(B>0\)，则每个 \(n_i>0\) 的主体都有 \(a_i>0\)。
该严格正结论首先是配额结论，实际使用须另接续交付证书。
\end{theorem}
\begin{certificate}
正确资格和真实需求、同类服务及可分性、单位价格、
购买资金到账、实际容量、辅助费用的独立覆盖，
以及按照总需求选择的计算分支。
\end{certificate}
\begin{proof}
\begin{semanticsteps}
\item 若 \(N=0\)，有限个非负需求之和为零，故每个 \(n_i=0\)。
定义直接给出 \(a_i=0\)，三项弱不等式均成立。
\item 若 \(N>0\)，则 \(pN>0\)，所以全部比值有定义。
非负性给出 \(0\le\alpha\le1\)，继而 \(0\le a_i\le n_i\)。
\item 有限求和给出 \(\sum_i a_i=\alpha N\)。
由 \(\alpha\le K/N\) 得 \(\alpha N\le K\)；
由 \(\alpha\le B/(pN)\) 得 \(p\alpha N\le B\)。
\item 若 \(K=0\) 或 \(B=0\)，最小值中出现零，而其余项非负，
所以 \(\alpha=0\)。
\item 若 \(K,B,N,p\) 均严格为正，
则参与取最小值的三个数都严格为正，故 \(\alpha>0\)。
再由 \(n_i>0\) 得 \(a_i=\alpha n_i>0\)。
\end{semanticsteps}
\end{proof}

\begin{proposition}[同需同量与主体置换对称]
\label{access:symmetry}
在同一批次中，\(n_i=n_j\) 蕴含 \(a_i=a_j\)。
对主体集合的任意双射 \(\sigma\)，若身份对应、资格、
需求及其他适用事实一并置换，且 \(B,K,p\) 不变，则
\[
a'_{\sigma(i)}=a_i.
\]
主体置换只改变标识的对应关系，不删除逐人的权利。
\end{proposition}
\begin{certificate}
置换前后的主体对应、完整资格与需求记录、
共同价格、资金和容量，以及没有使用姓名顺序或财富排序的配额规则。
\end{certificate}
\begin{proof}
\begin{semanticsteps}
\item 同一批次中的 \(\alpha\) 对全部主体相同，
因此 \(n_i=n_j\Rightarrow\alpha n_i=\alpha n_j\)。
\item 双射给出
\(N'=\sum_i n'_{\sigma(i)}=\sum_i n_i=N\)。
\item 若 \(N=0\)，两份计算都进入零分支。
若 \(N>0\)，共同的 \(B,K,p,N\) 给出 \(\alpha'=\alpha\)。
\item 对每个 \(i\)，代入 \(n'_{\sigma(i)}=n_i\)，
得 \(a'_{\sigma(i)}=\alpha'n'_{\sigma(i)}=\alpha n_i=a_i\)。
\end{semanticsteps}
\end{proof}

\begin{remark}[可分性与价格边界]
严格正的任意小配额只有在服务确实可分且该数量可交付时才有实际含义。
存在最小启动规模、整段设备时段或不可拆分名额时，使用后文的不可分模型。
本节限定 \(p>0\)；若合同价格为零，应删除货币比值并独立核验提供者资金，
不能把除以零的表达式当作无限供给证明。
\end{remark}

\section{从服务配额到实际交付和匹配支付}

\begin{definition}[交付准备与服务账单账本]
\label{access:delivery-ledger}
固定已经签发的配额 \(a_i\)，以 \(d_i\) 记录实际交付，
\(O_i\) 记录尚未交付的配额，\(Z_i\) 记录本人自愿撤回的未交付配额。
初态为 \(d_i=Z_i=0\)、\(O_i=a_i\)。
令 \(K^X\) 为剩余服务容量，\(B^X\) 为本批次可用购买资金。
对具有唯一事件标识 \(k\) 的提供者账单，
保存真实成立额 \(\Gamma_k^X\)、未付额 \(L_k^X\) 和最终到账 \(Y_k^X\)。
新批次初态取 \(L_k^X=Y_k^X=\Gamma_k^X=0\)，
\(B_0^X=B\)、\(K_0^X=K\)。

实际交付动作须有正确的服务回执、相匹配的资格与用途，
且 \(0<u\le O_i\)。
该交付事件的唯一标识 \(k\) 此前不得用于成立另一项相同账单；
新账单子项的 \(\Gamma_k^X,L_k^X,Y_k^X\) 初值均为零。
同一实际交付不能凭复制回执再记一次。
它执行
\[
\begin{aligned}
d_i'&=d_i+u,&O_i'&=O_i-u,&K^{X\prime}&=K^X-u,\\
\Gamma_k^{X\prime}&=\Gamma_k^X+pu,&
L_k^{X\prime}&=L_k^X+pu,&B^{X\prime}&=B^X.
\end{aligned}
\]
未交付的计划不能作为该动作的回执。
账单支付以 \(0\le v\le L_k^X\) 为守卫，
实际最终结算后同时令
\[
B^{X\prime}=B^X-v,\qquad
L_k^{X\prime}=L_k^X-v,\qquad
Y_k^{X\prime}=Y_k^X+v.
\]
本人可以撤回尚未交付且尚未产生不可撤销费用的 \(0\le z\le O_i\)，
此时 \(O_i'=O_i-z\)、\(Z_i'=Z_i+z\)，其余上述账本不变。
已经成立的提供者费用及劳动债权继续保留。
提供者在收款前所需的履约资金另有前缀证书；
购买方准备充分不能代替提供者自身的中途流动性。
\end{definition}

\begin{theorem}[服务权利、交付容量与购买准备的归纳保持]
\label{access:delivery-invariant}
若初态采用定义~\ref{access:delivery-ledger}，
则其允许动作组成的任意有限执行保持
\[
\begin{gathered}
d_i+O_i+Z_i=a_i,\qquad d_i,O_i,Z_i\ge0,\\
K^X\ge\sum_iO_i,\qquad
B^X-\sum_kL_k^X\ge p\sum_iO_i,\\
L_k^X+Y_k^X=\Gamma_k^X,\qquad L_k^X,Y_k^X\ge0.
\end{gathered}
\]
因此，真实已交付账单和尚未交付的配额不占用同一份购买准备两次。
\end{theorem}
\begin{certificate}
初态配额及预算、逐服务实际回执、唯一账单、相匹配的最终支付、
自愿撤回记录、可用容量和完整动作覆盖。
资金与容量不得被外部任务占用；
任何允许的其他支出须另外证明不侵占上述准备。
\end{certificate}
\begin{proof}
\begin{semanticsteps}
\item 定理~\ref{access:proportional} 给出
\(\sum_i a_i\le K\) 和 \(p\sum_i a_i\le B\)。
与初态定义合取，得到全部初态不变量。
\item 对交付动作，真实且未使用的事件回执给出一次实际交付；
\(u\le O_i\) 给出 \(O_i'\ge0\)，且
\((d_i+u)+(O_i-u)+Z_i=a_i\)。
容量与待交付总量同减 \(u\)，故其差保持。
\item 令 \(L^X=\sum_kL_k^X\)、\(O=\sum_iO_i\)。
交付后有
\[
B^{X\prime}-L^{X\prime}-pO'
=B^X-(L^X+pu)-p(O-u)=B^X-L^X-pO.
\]
账单 \(\Gamma_k^X,L_k^X\) 同增 \(pu\)，所以账单权利恒等式保持。
\item 实际支付使 \(B^X,L^X\) 同减 \(v\)，
故 \(B^X-L^X-pO\) 保持。
同一账单的 \(L_k^X+Y_k^X\) 保持，且 \(v\le L_k^X\) 排除超付。
\item 自愿撤回使 \(O_i\) 减少 \(z\)、\(Z_i\) 增加 \(z\)，
故逐人配额恒等式保持。
容量余量增加 \(z\)，购买准备余量增加 \(pz\)，两者仍非负。
\item 每个允许动作保持合取，且旧劳动账本未被修改。
按定理~\ref{base:induction} 对动作数归纳，得到全部有限前缀的结论。
\end{semanticsteps}
\end{proof}

\begin{corollary}[有限批次的实际取得与足额结算]
\label{access:actual-delivery}
若本批次的提供者具有真实履约资源和资金，
对每个仍愿接受服务的主体预留一个完成其全部 \(O_i\) 的处理时段，
全部时段与支付窗口共同可行且具有有限时限，
则存在实际执行，使最终
\[
O_i=0,\qquad d_i=a_i-Z_i,\qquad L_k^X=0.
\]
对未撤回服务的主体，有 \(d_i=a_i\)。
若另外 \(n_i>0\)、\(K>0\)、\(B>0\)，则该主体实际取得正服务量。
\end{corollary}
\begin{certificate}
共同的提供者资金前缀、服务时段、资格和用途许可、
交付上界、全部未付账单及支付最终性。
时段覆盖完整批次，不允许其他请求无限插队。
\end{certificate}
\begin{proof}
\begin{semanticsteps}
\item 有限主体集合给出有限个待处理配额。
按预留次序逐个执行完整交付，非现金条件使每一步实际可执行。
\item 定理~\ref{access:delivery-invariant} 保证此前步骤不会耗尽后续已预留容量和购买资金。
每完成一个主体，其 \(O_i\) 变为零；有限次后全部 \(O_i=0\)。
\item 对全部已成立账单按匹配标识付款。
资金准备和实际支付设施满足定理~\ref{part:thm-prefix-funding}，
有限处理条件接续命题~\ref{part:prop-timely-payment}，故全部到期账单完成结算。
\item 配额恒等式给出 \(d_i=a_i-Z_i\)。
本人未撤回时 \(Z_i=0\)，得到 \(d_i=a_i\)。
再应用定理~\ref{access:proportional} 的严格正分支，得到实际正交付。
\end{semanticsteps}
\end{proof}

\section{公共采购、个人补助与提供者费用的一次核算}

\begin{definition}[同一服务账单的两种付款来源]
\label{access:invoice}
对已实际交付的 \(d_i\)，唯一账单金额为 \(pd_i\)。
个人实际支付 \(h_i\ge0\)、公共补助实际支付 \(s_i\ge0\)，满足
\[
h_i+s_i=pd_i.
\]
若已足额交付 \(d_i=a_i\)，则得到 \(h_i+s_i=pa_i\)。
本章的财富独立服务路径允许 \(h_i=0\)、\(s_i=pd_i\)。
个人自愿承担部分金额时，仍不改变本批次事先冻结的资格、需求和配额。
若存在其他公共直接购买或固定维护合同，须具有不同的实际交付或维护义务，
分别登记合同和预算；已经包含在 \(p\) 中的费用不再以另一个名称收费。
不得把同一账单重新登记为第二笔采购。
\end{definition}

\begin{proposition}[低个人付款与完整服务收入的同账本相容]
\label{access:payment-identity}
若定义~\ref{access:invoice} 的两笔付款均已最终到账，
则提供者取得的该笔收入恰为 \(pd_i\)，而不是 \(pd_i+s_i\)。
在同时包含个人、公共付款方与提供者的现金边界内，
这两笔转移在支付发生时的合计现金变化为零。
在本章配额约束下，公共补助合计不超过 \(B\)。
\end{proposition}
\begin{certificate}
同一账单的个人和公共付款标识、来源账户、到账时间、实际交付、
未付余额和退款或撤回处理。
提供者的其他收入与成本另列，不混入该账单金额。
\end{certificate}
\begin{proof}
\begin{semanticsteps}
\item 最终到账使该账单收入等于 \(h_i+s_i\)。
代入账单恒等式，得到收入 \(pd_i\)。
\item 三方余额变化分别为 \(-h_i,-s_i,h_i+s_i\)；
相加得到零。这是定理~\ref{part:thm-source-balance} 的逐笔合并实例。
\item 非负个人付款给出 \(s_i\le pd_i\)。
由交付恒等式和非负撤回量，得到 \(d_i\le a_i\)。
\item 有限求和并应用定理~\ref{access:proportional}，
得到 \(\sum_i s_i\le p\sum_i d_i\le p\sum_i a_i\le B\)。
\item 当 \(h_i=0\) 时，提供者仍收到 \(s_i=pd_i\)。
这个分支减少个人支付，不减少实际使用的设备、能源和人员。
\end{semanticsteps}
\end{proof}

\begin{remark}[提供者成本与真实资源]
提供者核算收入 \(R^V\) 时，将 \(h_i+s_i\) 作为同一服务收入。
若其可变支出为 \(cQ\)、另列固定支出为 \(F\)，则
\(\Pi^V=R^V-cQ-F\)；已经包含在 \(cQ\) 或 \(F\) 的工资、认证或维护费用不再重复扣除。
当 \(\Pi^V\) 解释为经济利润时，\(cQ+F\) 按定义~\ref{market:full-cost}
纳入正常资本报酬等完整机会成本；实际现金结余仍按付款时点另行计算。
购买者支付 \(pQ\) 与提供者收到 \(pQ\) 是同一转移的两端；
不能在合并福利计算中把它与提供者全部投入机械相加为两份相同资源损失。
真实资源按唯一作业标识记录，每一机器时、能源量、人员时段只计入一次。
公共补助能够改变付款来源，但不能由账本恒等式创造这些资源。
若提供者用采购收入支付劳动债权，应记录真实的先收款、后付款路径，
而不能把同一采购款再次称为新增外部资金。
\end{remark}

\section{覆盖身份核验和资金渠道的全流程非干扰}

\begin{definition}[两份合法共同执行见证]
\label{access:paired-witness}
比较两个各自可达的制度状态 \(x,x'\)：
主体财富向量 \(w,w'\) 与劳动评分向量 \(s^L,s^{L\prime}\) 可以不同。
共同执行见证 \(\mathcal C,\mathcal C'\) 分别覆盖身份核验、申请、
需求确认、资源准备、配额计算、实际交付与付款，并满足：
\begin{enumerate}
\item 规范主体对应、真实资格事实、服务请求、需求 \(n_i\)、
规则版本、服务质量和期限相同；财富与劳动评分不被替换为隐含的资格代理字段。
\item 两份见证中均存在实际可用的身份和材料取得路径。
必要费用均有到期前资金，所需时间、设备和协助均可获得；
没有一份执行依赖仅另一份执行中存在的个人垫资。
\item 真实接纳程序、用途许可和正确确认条件分别通过；
初核错误及复核所需资源已包括在完整程序中。
\item 本批次 \(B,K,p\)、申请处理容量与分配优先规则相同。
筹资、信用、资格证明与设备配置不因财富或劳动评分变化而失去有效性。
私人自愿分担不成为准入、排序或本批次配额的条件。
\item 每份执行中的劳动支付、复核和探索服务共同可行，
各自通过合并资金前缀及非现金资源检查。
相同公共输入下的实际程序步骤采用相同规则；涉及外部响应时保留相应的配对见证。
\end{enumerate}
上述条件是对整个执行的核验要求，不只是要求末端配额函数不读取 \(w\) 或 \(s^L\)。
它们不声称财富变化后税收、总预算或实际需求必然不变；
这些变化若发生，就超出本次保持比较的前件。
\end{definition}

\begin{theorem}[财富与劳动评分变化下的服务权利非干扰]
\label{access:noninterference}
若存在定义~\ref{access:paired-witness} 的两份合法共同执行见证，
则两份执行确认的基本探索资格相同，签发的配额满足
\[
E_\nu^{X\prime}(i)=E_\nu^X(i),\qquad a_i'=a_i.
\]
若两份执行均足额交付、主体作出相同的接受或撤回选择，
则实际交付满足 \(d_i'=d_i\)。
该结论不要求两份执行中的劳动报酬相等。
\end{theorem}
\begin{certificate}
两份可达状态、完整共同执行、逐阶段主体与事实对应、
必要证据路径、申请费用和资金前缀、准入与用途许可、
共同预算容量、配额及实际交付回执。
至少逐一检查身份、申请、筹资、资源和交付五个渠道没有产生新的排除。
\end{certificate}
\begin{proof}
\begin{semanticsteps}
\item 初态共同事实给出相同的规范主体与 \(z_i^X\)。
资格定义的每个谓词分别相同，因此合取 \(E_\nu^X(i)\) 相同。
\item 身份和证据路径在两份状态中都实际可执行。
各自的正确确认条件把相同真实事实送入相同的最终确认结果；
不能以一份执行中的证明材料替另一份执行补齐缺口。
\item 对公共程序阶段编号归纳。
假设前一阶段的主体、资格、需求和适用规则对应，
共同规则和配对响应给出下一阶段相同的相关字段。
两份各自成立的费用、权限与容量守卫保证这一步在两边均可执行。
\item 有限阶段归纳得到相同的 \(\mathcal I_X,n_i,N\)。
共同的 \(B,K,p\) 使两份配额计算进入同一分支；
零分支均为零，正分支均有 \(\alpha'=\alpha\)。
所以 \(a_i'=\alpha'n_i'=a_i\)。
\item 实际交付证书给出
\(d_i=a_i-Z_i\) 和 \(d_i'=a_i'-Z_i'\)。
相同选择给出 \(Z_i'=Z_i\)，故 \(d_i'=d_i\)。
\item 劳动评分可以改变劳动账本内的差异份额。
公共服务结论只使用各自已经通过的共同资金检查，
并未把劳动报酬相等作为隐藏前件。
\end{semanticsteps}
\end{proof}

\begin{remark}[末端函数独立不足以保证完整入口]
设身份核验要求申请者自付 \(f>0\)，且没有其他可用付款路径。
若 \(w_i<f\le w_i'\)，则一份状态无法完成核验，另一份状态可以完成。
即使最后的公式不读取财富，两份最终接纳集合仍可能不同。
这正是必须核验两份合法共同执行，而不能仅检查函数参数列表的原因。
同样，某一劳动评分若决定申领补助、信用额度或设备席位，
也必须作为潜在联动渠道审查。
\end{remark}

\section{不可分同需服务的可实现随机机会}

\begin{definition}[等需求类中的整数名额]
\label{access:indivisible}
固定一个同资格、同需求、同质量的有限申请类 \(J\)，人数为 \(m\)。
每名申请者需要一个不可分服务名额，申请在抽签前完成并冻结。
整数容量为 \(K_{\mathrm{int}}\in\mathbb N\)，单名额合同价格为 \(p>0\)，
已经到位的本项预算为 \(B_{\mathrm{int}}\ge0\)。
该预算与容量不与另一服务批次重复承诺。
当 \(m>0\) 时，取
\[
k=\min\left\{m,K_{\mathrm{int}},
\left\lfloor B_{\mathrm{int}}/p\right\rfloor\right\}.
\]
从全部 \(k\) 元子集中均匀抽取 \(S\subseteq J\)，令 \(X_i=\1_{\{i\in S\}}\)。
当 \(m=0\) 时，取 \(k=0\)、\(S=\varnothing\)，不定义逐人的归一概率。
\end{definition}

\begin{theorem}[每次可行、同需同概率及期望服务量]
\label{access:lottery}
若冻结名单和随机源均符合所声明的均匀抽签机制，
则每次抽签都满足
\[
|S|=k\le K_{\mathrm{int}},\qquad pk\le B_{\mathrm{int}}.
\]
当 \(m>0\) 时，每人被提供名额的概率和期望配额均为
\[
\Pr(X_i=1)=\EE X_i=\frac{k}{m}.
\]
若 \(K_{\mathrm{int}}\ge1\) 且 \(B_{\mathrm{int}}\ge p\)，则该概率严格为正。
这不保证每人在同一次实现中都取得名额。
\end{theorem}
\begin{certificate}
等需求类、完整冻结名单、规范主体、整数容量、可用预算、
随机机制、实际抽签记录和无法事后选择性重抽的执行记录。
一种有限实现是取得均匀随机排列并选前 \(k\) 人；
排列可以用有限次均匀整数抽取与交换形成。
随机源的分布正确性是独立前件，单次抽签日志不能自行证明该分布。
\end{certificate}
\begin{proof}
\begin{semanticsteps}
\item 若 \(m=0\)，定义给出空集合与零付款，容量和资金约束成立。
以下仅对 \(m>0\) 计算概率。
\item \(k\) 为三个非负整数的最小值，
故 \(0\le k\le m\)、\(k\le K_{\mathrm{int}}\)。
由 \(k\le\lfloor B_{\mathrm{int}}/p\rfloor\)，得到 \(pk\le B_{\mathrm{int}}\)。
\item 均匀排列的每个前 \(k\) 人集合出现次数相同，
因此该有限抽签实现对全部 \(k\) 元子集给出相同概率。
\item 当 \(1\le k\le m\) 时，包含固定主体 \(i\) 的子集数为
\(\binom{m-1}{k-1}\)，全部子集数为 \(\binom m k\)。
因此
\[
\Pr(X_i=1)=
\frac{\binom{m-1}{k-1}}{\binom m k}=\frac{k}{m}.
\]
\item 当 \(k=0\) 时，\(S=\varnothing\)，故概率为零，也等于 \(k/m\)。
对指示变量，期望定义给出
\(\EE X_i=1\Pr(X_i=1)+0\Pr(X_i=0)=k/m\)。
\item \(m>0\)、\(K_{\mathrm{int}}\ge1\)、\(B_{\mathrm{int}}\ge p\)
使三个整数均至少为一，故 \(k\ge1\)，继而 \(k/m>0\)。
\end{semanticsteps}
\end{proof}

\begin{corollary}[抽签机会与自愿实际使用的区别]
\label{access:lottery-delivery}
若所有中签者均愿接受，并且选中集合的服务及支付资源在每种抽签结果下都可执行，
则实际服务指示量 \(D_i=X_i\)，从而 \(\EE D_i=k/m\)。
允许本人拒绝时，一般只能保证名额机会的等概率；
实际使用量为 \(D_i=X_i A_i\)，其中 \(A_i\in\{0,1\}\) 表示接受且实际完成。
\end{corollary}
\begin{certificate}
全部可能选中集合的可执行服务策略、实际接受记录和交付回执。
抽签后退出不删除既有劳动债权或未来基本探索资格。
\end{certificate}
\begin{proof}
\begin{semanticsteps}
\item 每种结果中，容量与资金约束由定理~\ref{access:lottery} 成立；
非现金执行证书使选中名额能够交付。
\item 全部中签者接受时，\(X_i=1\Rightarrow D_i=1\)，
未中签时 \(X_i=0\Rightarrow D_i=0\)，故 \(D_i=X_i\)。
取期望得到第一项结论。
\item 允许拒绝时，实际服务同时要求选中和接受完成，
所以 \(D_i=X_iA_i\)。
没有关于 \(A_i\) 的额外条件，不能把 \(\EE D_i\) 替换为 \(k/m\)。
\end{semanticsteps}
\end{proof}

\begin{remark}[不可分服务的零预算边界]
当 \(0<B_{\mathrm{int}}<p\) 时，即使容量为正，仍有 \(k=0\)。
因此，可分服务中正资金与正容量推出正配额的结论不能直接搬到不可分名额。
等概率是同一资格和同需求类中的事前性质；
不同需求类的名额安排还需要公开的类间规则及共同资源证书。
抽签不得绕过高保障操作所需的用途许可。
\end{remark}

\section{两套分配规则共同成立的有限数值见证}

\begin{proposition}[不同劳动报酬与相同探索服务的相容实例]
\label{access:finite-witness}
存在一个有限构造，同时满足劳动差异支付、同需同量探索、
独立复核准备及提供者收款前的现金约束。
以下数字仅为构造参数，不是经验测量。
取两名真实有效劳动参与者，其可比较贡献为 \(1,3\)，
\[
\beta=4,\quad D=4,\quad g_L(s)=s,\quad
M_L^{\max}=2,\quad B_P^L=12.
\]
两人的真实探索需求均为 \(2\)，取
\[
p=2,\quad K=3,\quad B=6.
\]
另外为两项申请及必要复核分别预留现金 \(1\)，总额 \(B_A=2\)；
提供者持有履约前现金 \(B_0^V=3\)，覆盖本例交付所需的全部现金投入 \(3\)。
上述四份初始资金由不同专款覆盖，合计 \(23\)，不重复承诺。
\end{proposition}
\begin{certificate}
两项真实劳动事件、有效且无重复的贡献记录、两项相同探索需求；
四个实际可用资金账户；
相互不冲突的劳动、申请复核和服务时段；
提供者投入、实际服务和各笔最终付款的唯一标识。
本构造声明提供者全部交付现金投入为 \(3\)，
若另有固定或维护支出，就必须修改该完整成本证书。
\end{certificate}
\begin{proof}
\begin{semanticsteps}
\item 贡献和为 \(4\)，故
\(\lambda_1=1/4\)、\(\lambda_2=3/4\)。
代入劳动合同，得到 \(b_1=5\)、\(b_2=7\)，合计恰为 \(12\)。
真实履约及最终支付证书使 \(y_1=5,y_2=7\)。
\item 探索总需求为 \(N=2+2=4\)。
因此
\[
\alpha=\min\{1,3/4,6/(2\cdot4)\}=3/4,
\qquad a_1=a_2=3/2.
\]
总配额为 \(3\)，账单合计为 \(6\)，分别等于容量与购买预算。
\item 申请及复核费用的每个前缀分别由预留的 \(1\) 覆盖，
不要求申请者从工资中先行垫付，也不将该费用再从 \(5,7\) 中扣除。
共同时间和证据路径满足定理~\ref{rights:affordable-review} 的非现金前件。
\item 提供者先用其现金 \(3\) 支付全部交付投入，余额为零；
实际交付 \(3\) 单位服务后，按唯一账单取得购买方付款 \(6\)，余额为 \(6\)。
该过程每个现金前缀非负，且投入只支付一次。
\item 取个人付款 \(h_1=h_2=0\)，公共付款为 \(s_1=s_2=3\)。
于是每笔 \(h_i+s_i=pa_i\)，提供者本批次收入为 \(6\)，不是 \(12\)。
\item 同一计划中的专款和时段相容，故上述步骤可以实际接续。
不同劳动贡献给出不同劳动报酬，而相同探索需求给出相同探索服务；
两项结论没有使用同一个资格条件或分配分母。
\end{semanticsteps}
\end{proof}

\section{两类权利的联合执行与结论范围}

\begin{theorem}[劳动合同权利与按需探索的共同实现]
\label{access:joint-rights}
设一个有限期限内的同一实际制度具有以下共同见证：
权利动作满足定理~\ref{rights:invariant}；
退出与复核的真实路径和时限已获覆盖；
劳动分配满足定理~\ref{rights:pay-coverage} 并最终付款；
探索服务采用可分规则或相应等需求类的不可分规则，
且全部身份、准入、资金和实际交付条件成立。
对所声明的财富与劳动评分比较，
另具有定义~\ref{access:paired-witness} 的两份合法共同执行见证。
若各模块在同一现金前缀、资源和时间表中相容，
则可共同实现本章的 \(\mathsf D_i\)、有效劳动的正到账，
以及所声明范围内的按需服务量或同需同概率机会。
财富与劳动评分的非干扰结论仅在定理~\ref{access:noninterference} 的比较域内适用。
\end{theorem}
\begin{certificate}
同一主体与期限的共同计划；每笔资金的来源、用途和实际先后；
不同任务对同一资源的唯一作业标识；
劳动、复核、探索和提供者履约的完整义务；逐人实际回执。
各账户的专款在每个前缀均不得超过可用资金扣除既有预留后的余额。
若提供者用采购收入支付劳动费用，合并计划必须保存这两笔不同交易的接续，
不将内部转账重复计作外部融资。
\end{certificate}
\begin{proof}
\begin{semanticsteps}
\item 固定证书中的同一共同状态与计划 \(z_\star\)。
各模块所用资源、现金、主体和时点均从 \(z_\star\) 读出。
\item 权利动作的归纳保持给出真实债权、同意和资格分离；
退出与复核的有限路径给出真实可用的选择及复核入口。
两份共同执行再由定理~\ref{access:noninterference} 给出所声明的基本探索独立性。
这些结论逐项满足 \(\mathsf D_i\) 所列操作性要求。
\item 有效劳动的正确确认与最终付款给出 \(y_i=b_i\ge\beta>0\)。
该结论的分母是实际有效劳动集合，不是探索申请集合。
\item 可分分支由定理~\ref{access:proportional} 和推论~\ref{access:actual-delivery}
给出所承诺的实际服务；不可分分支由定理~\ref{access:lottery}
给出每次可行与同需同概率，并在接受和交付条件成立时得到实际服务。
\item 共同现金前缀与资源计划使上述动作可以相容接续，
而非分别依赖互相冲突的预算或时段。
因此满足引理~\ref{base:joint} 的共同见证条件，得到全部结论的合取。
\item 有限期限之外的需求、资金、容量和核验性能没有包含于证书，
故本定理不推出无限期连续服务，也不推出完整社会福利已经提高。
\end{semanticsteps}
\end{proof}

\begin{remark}[需要分别保持的经济含义]
按劳分配依据已约定的真实贡献，公共探索资格依据另一组制度条件。
提高劳动贡献可以改变劳动报酬，但不因此提高人的基本探索资格；
公共服务的实际使用又不能被改记为一笔未发生的劳动收入。
两套制度能够共同运行，依靠的是可执行的规则、真实资金和共同资源，
而不是对任何一种权利作无限资源解释。
\end{remark}

\chapter{真实性核验与可执行激励}
\label{inc:chapter}

真实性约束既是交付对象的组成部分，也是支付制度的实施条件。
本章不以工具名称判断资格，而以事前规则、实际控制和事件证据判断。
机会主义分析采用风险中性的条件模型，
将发现概率、实际执行能力与私人偏离收益分别列出
\cite{InternalEconomics,Becker1968}。

\section{同结果与不同来源的经济含义}

\begin{proposition}[结果记录对真人资格的不充分性]
\label{inc:prop-origin}
若存在两个事件 \(e_H,e_A\)，其结果记录满足
\(y(e_H)=y(e_A)\)，而
\(\operatorname{HumanAct}_R(e_H)=1\)、
\(\operatorname{HumanAct}_R(e_A)=0\)，
则不存在仅依赖结果的函数 \(v\)，
能在这两个事件上同时正确判断真人资格。
\end{proposition}
\begin{certificate}
记录相同结果、不同实际来源以及合同对来源的要求。
后续表示必须保留足以区分该资格的证据；
没有取得的来源信息不以输出质量代替。
\end{certificate}
\begin{proof}
\begin{semanticsteps}
\item 若存在这样的 \(v\)，相同输入必给出相同输出；
正确性却同时要求该输出为一和零，矛盾。
\item 因此需要额外来源信息，或承认当前观察不能完成该判定。
\end{semanticsteps}
\end{proof}

\begin{definition}[规则版本下的违规偏离]
\label{inc:def-deviation}
固定 \(R\)、他人行动与主体 \(i\) 的合法备选集合。
在此集合中选取效用最高的基准行动 \(a_i^0\)；
若最高值不取得，则本章只比较事前指定且可执行的合法基准。
登记偏离集合 \(\mathcal V_i(R)\) 包含违反真实性或工具义务的行为，
例如未经许可的自动替代、重复申领、来源伪造和版本篡改。
在真实符合条件的多个项目之间选择，不属于本定义的违规。
检测输出 \(\widehat V\in\{1,0,?\}\)
与真实 \(\operatorname{Valid}_R\) 分别记录；
问号表示尚待核验，不能解释为既有真实权利已经消失。
\end{definition}

\section{追偿、处罚与私人偏离收益}

\begin{definition}[可执行责任下的偏离效用]
\label{inc:def-utility}
对 \(v\in\mathcal V_i(R)\)，
令 \(g_i(v)\) 为相对合法基准、在执行责任之前的毛增益，
\(c_i(v)\ge0\) 为相对基准增加的私人行为成本。
令 \(d_i(v)\) 为发现概率，
\(s_i(v)\) 为发现后责任实际执行的条件概率，并令
\[
 p_i(v)=d_i(v)s_i(v).
\]
令 \(r_i(v)\ge0\) 为执行时实际追回的利益，
\(f_i(v)\ge0\) 为另行成立且实际可执行的处罚。
追回与处罚不能包含同一责任金额两次。
在风险中性、概率与所用信念一致的模型中，
\[
 \Delta U_i(v)
 =g_i(v)-c_i(v)-p_i(v)\bigl(r_i(v)+f_i(v)\bigr).
\]
若回收金额随机，\(r_i,f_i\) 取执行事件条件下的相应期望值。
名义罚额、无法执行的财产请求和未成立的责任不计入这些金额。
\end{definition}

\begin{theorem}[登记偏离域内的严格激励约束]
\label{inc:thm-deterrence}
若对全部 \(i,v\) 都有
\[
 p_i(v)\bigl(r_i(v)+f_i(v)\bigr)
 >g_i(v)-c_i(v),
\]
则每个登记偏离严格劣于其合法基准。
若有限偏离集合中的每项差值至少有共同裕量 \(\eta>0\)，
任何只在这些偏离之间随机化的策略也严格劣于该共同基准。
\end{theorem}
\begin{certificate}
保存完整登记偏离集合、合法基准、毛增益与额外成本、
发现和执行的条件概率、可执行追偿及处罚、
独立裁决与复核路径。
非货币动机若进入比较，须先以同一效用尺度计入 \(g_i-c_i\)。
对合谋作结论时，合谋行动与联合收益必须显式进入偏离域。
\end{certificate}
\begin{proof}
\begin{semanticsteps}
\item 将条件代入定义，逐项得到 \(\Delta U_i(v)<0\)。
\item 若各项不超过 \(-\eta\)，对任意概率权重 \(\mu(v)\)
有
\[
 \sum_v\mu(v)\Delta U_i(v)
 \le-\eta\sum_v\mu(v)=-\eta<0.
\]
这里 \(p=ds\) 使用条件概率乘法，
不要求发现和执行相互独立。
\end{semanticsteps}
\end{proof}

\begin{proposition}[有限执行能力的可行边界]
\label{inc:prop-enforcement-bound}
若某偏离满足 \(g-c>0\)，且全部可行方案均有
\[
 p\le\bar p,\quad r\le\bar r,\quad f\le\bar f,\qquad
 \bar p(\bar r+\bar f)\le g-c,
\]
则该域内不存在严格威慑方案。
特别地，\(p=0\) 或 \(r+f=0\) 时，
正私人净增益不能被该责任机制消除。
\end{proposition}
\begin{certificate}
上界来自实际权限、可执行资产、处理能力及允许的责任范围。
审计预算与执行资源共同决定可行集合；
分别达到的参数最大值不自动构成可联合达到的方案。
\end{certificate}
\begin{proof}
\begin{semanticsteps}
\item 对每个可行方案，
\(p(r+f)\le\bar p(\bar r+\bar f)\le g-c\)，
故 \(\Delta U\ge0\)，严格威慑条件不成立。
\item 两个特例均使预期可执行责任为零。
\end{semanticsteps}
\end{proof}

\section{误判复核与保护概率}

\begin{definition}[初核与复核的误差]
\label{inc:def-errors}
对声明事件域，令 \(\alpha_0\) 为合格事件被初核拒绝的概率，
\(\varepsilon_0\) 为不合格事件被初核接受的概率。
令 \(q_H\) 为初核误拒后经复核改正的条件概率，
\(s_H\) 为初核正确接受后被错误推翻的条件概率；
\(q_A\) 为初核误受后经复核改正的条件概率，
\(s_A\) 为初核正确拒绝后被错误接受的条件概率。
零概率条件事件对应的数值可任取于 \([0,1]\)，
因为其在下式中的系数为零。
\end{definition}

\begin{theorem}[复核的净纠错条件]
\label{inc:thm-review}
最终误拒和误受率分别为
\[
 \alpha_f=\alpha_0(1-q_H)+(1-\alpha_0)s_H,\qquad
 \varepsilon_f=\varepsilon_0(1-q_A)+(1-\varepsilon_0)s_A.
\]
因此 \(\alpha_f\le\alpha_0\) 当且仅当
\((1-\alpha_0)s_H\le\alpha_0q_H\)；
误受率下降具有对应条件
\((1-\varepsilon_0)s_A\le\varepsilon_0q_A\)。
\end{theorem}
\begin{certificate}
保留完整初核总体、复核资格、被纠正和被错误推翻的记录。
只分析主动申诉者时另报样本选择范围，
不能据此自动替代全部有效参与者的误差率。
\end{certificate}
\begin{proof}
\begin{semanticsteps}
\item 在真实合格总体内，令 \(R_0\) 表示初核拒绝，\(R_f\) 表示最终拒绝。
互斥分解为 \(R_f=(R_f\cap R_0)\sqcup(R_f\cap R_0^c)\)。
\item 条件概率定义给出
\[
\Pr(R_f\cap R_0)=\alpha_0(1-q_H),\qquad
\Pr(R_f\cap R_0^c)=(1-\alpha_0)s_H.
\]
\item 互斥相加得到 \(\alpha_f\) 公式；对真实不合格总体及接受事件作相同分解得到 \(\varepsilon_f\) 公式。
\item 逐项移项：
\[
\alpha_f-\alpha_0=(1-\alpha_0)s_H-\alpha_0q_H,
\quad
\varepsilon_f-\varepsilon_0=(1-\varepsilon_0)s_A-\varepsilon_0q_A.
\]
\item 两个差额分别不大于零的充要条件即为所列条件；严格下降要求对应不等式严格。
\end{semanticsteps}
\end{proof}

\begin{theorem}[有限执行的联合可靠性下界]
\label{inc:thm-union-bound}
一次完整执行至多含 \(M\) 个必须正确的判断。
若对第 \(j\) 次判断的全部允许历史，
其复核后条件失败概率均不超过 \(\epsilon_j\)，
且所有判断成功时资格、过程和支付证书成立，
则完整保护的概率至少为
\[
 \left[1-\sum_{j=1}^{M}\epsilon_j\right]_+.
\]
\end{theorem}
\begin{certificate}
记录每次调用的任务、允许历史、误拒和误受定义及概率界来源。
抽样置信度与系统执行失败概率分别报告；
数据不能支持历史一致界时，不签发本定理的数值结论。
\end{certificate}
\begin{proof}
\begin{semanticsteps}
\item 将不足 \(M\) 次的执行补为失败概率零的空判断。
\item 由全期望公式，第 \(j\) 次失败事件 \(F_j\) 的无条件概率不超过 \(\epsilon_j\)。
\item 再由
\(\mathbf 1_{\cup_jF_j}\le\sum_j\mathbf 1_{F_j}\)
取期望，得到失败并集概率不超过其和。
\item 全部判断成功蕴含完整保护，故得到所列下界。
\item 该证明不使用各判断相互独立。
\end{semanticsteps}
\end{proof}

\section{纠错资源与资金来源的保持}

\begin{proposition}[不可即时追回错误支付的备用覆盖]
\label{inc:prop-error-reserve}
在一个声明的结算窗口内，
真实参与义务为 \(L\)，原支付专款满足 \(B^P\ge L\)。
若错误支付造成且尚未实际追回的现金损失为 \(e\ge0\)，
另有及时可用的独立备用资金 \(R^E\)，
则继续足额支付真实义务的现金覆盖条件为
\[
 B^P-e+R^E\ge L.
\]
若已知 \(e\le\bar e\)，取 \(R^E\ge\bar e\) 是充分条件。
复核运营成本由其他已覆盖账户承担；
若由同一备用账户承担，应把该成本另加到所需覆盖额。
\end{proposition}
\begin{certificate}
保存实际错误流出、已追回金额、备用资产的可用时间与用途。
已经使用来覆盖真实义务的资金不再计为独立备用金。
追偿承诺只在真正到账后抵减 \(e\)。
\end{certificate}
\begin{proof}
\begin{semanticsteps}
\item 支付损失后的可用现金为 \(B^P-e+R^E\)，
足额覆盖既有真实义务等价于其不小于 \(L\)。
\item 当 \(R^E\ge\bar e\ge e\) 时，
\(B^P-e+R^E\ge B^P\ge L\)，得到充分性。
\item 同账户另付成本会进一步减少现金，故必须另计。
\end{semanticsteps}
\end{proof}

\begin{proposition}[真实性核验不产生付款资源]
\label{inc:prop-no-funds-from-verification}
若期初可用现金、当期全部外部流入和参与费用均为零，
并且不能新增信用，则无论资格核验的正确率多高，
实际非负支付总额都只能为零。
正确排除虚假请求可以保护既有份额，
但不能单独创造正的可支付资金。
\end{proposition}
\begin{certificate}
将核验输出、请求权金额与现金来源分别记录。
治理支出、实际追偿以及获赔转移均有独立现金凭证，
不把一条通过记录记作一笔入金。
\end{certificate}
\begin{proof}
\begin{semanticsteps}
\item 将全部初始资金和到账流入设为零。没有新增信用时，任一合法前缀均须满足
\[
0\ge\sum_{n\le m}(u_n+c_n),\qquad u_n,c_n\ge0.
\]
\item 非负和不大于零，推出每个已经执行的 \(u_n=c_n=0\)。
\item 核验输出改变资格观察，未产生具有真实到账凭证的现金增量。
故 \(\partial_{\rm verify}B=0\)，不改变上述前缀约束。
\item 对所有前缀作全称引入，得到实际非负支付总额为零。
\end{semanticsteps}
\end{proof}

\begin{remark}
本章的经济接续关系为
\[
\begin{aligned}
 \text{真实有效参与}
 &\Longrightarrow\text{合同应付份额}\\
 &\Longrightarrow\text{匹配义务和实际资金}\\
 &\Longrightarrow\text{及时最终到账}\\
 &\Longrightarrow\text{有容量的可负担接入}.
\end{aligned}
\]
每个箭头分别使用前述证书；
第一项不单独推出后面所有项。
这一结构使真实性控制与收入制度相容，
同时保留真实需求、资金更新、资源容量与自主退出的独立条件。
\end{remark}

\chapter{任务类型与最粗动态充分商}
\label{q:chapter}

经济制度需要保存的不仅是产出数额，还包括取得产出的过程、
参与资格、收益归属、付款义务和后续进入条件。
当两个事件产生相同结果而具有不同权利状态时，
只保存结果的表示不足以承担资格核验和收益结算。
本章先确定必须保留的经济观察，再构造相对于这些观察的最粗表示。
所用的观察因子化、分型细化和有限执行保持接口，
分别接续任务表示与参与分配的形式化研究
\cite{InternalTheory,InternalEconomics}。
\section{经济任务、真实类型与证据状态}
\begin{definition}[有限经济任务域]
\label{q:task}
固定规则版本 \(R\)、核算期间和有限状态集合 \(X\)。
经济任务写为
\[
\mathcal M=(X,A,\Xi,\delta,G,U,\mathcal O).
\]
其中 \(A(x)\) 是状态 \(x\) 下的合法动作集合；
\(\Xi(x,a)\) 是执行动作 \(a\) 后允许出现的外部响应集合；
\(\delta(x,a,\xi)\) 是后继状态；
\(G\subseteq X\) 是任务接受集合；
\(U(x,a,\xi)\) 保存规定的收益、费用或其他逐步评价；
\(\mathcal O\) 是必须保持的状态观察族。
存在有限的共同标签集合 \(\mathsf A,\mathsf Z\)，使
\[
A(x)\subseteq\mathsf A,\qquad
\varnothing\ne\Xi(x,a)\subseteq\mathsf Z
\quad(a\in A(x)).
\]
每个合法动作均有非空响应域。
状态包含本任务需要的身份、规则、余额、应付、到账、
时间和历史信息；某项信息只有在不影响任何规定观察及其后继时，
才允许从表示中省略。
\end{definition}
\begin{remark}
有限性可以来自已声明的账期、数量范围和精度。
若实际模型使用连续金额、无界历史或无限任务族，
须另证有限抽象的充分性或登记近似误差。
本章的有限稳定次数不直接适用于未经处理的无限状态空间。
\end{remark}
\begin{definition}[事件类型与参与资格]
\label{q:typing}
令 \(q_0:X\to Q_0\) 为已有的满射表示，
保存原任务规定的结果、资源或其他观察。
令 \(\tau:X\to\mathcal T\) 表示真实经济类型，
包括当前任务需要区分的用途、参与过程、授权及权利类别。
检测结果 \(\widehat\tau\) 可以另取待复核值 \(?\)；
待复核属于证据状态，不等同于真实类型已经确定。
对于参与事件 \(e\)，定义
\[
\begin{aligned}
\operatorname{Valid}_R(e)
={}&\operatorname{Admitted}_R(e)
\land\operatorname{HumanAct}_R(e)\\
&\land\operatorname{Compliant}_R(e)
\land\operatorname{Completed}_R(e)
\land\operatorname{Unique}_R(e).
\end{aligned}
\]
其中人类实际参与与账户身份分别核验，
唯一性按照已经声明的事件及申领单位判断。
其他合法经济合同可以采用相应的过程类型；
不同合同的类型条件分别登记。
\end{definition}
\begin{lemma}[观察因子化的充要条件]
\label{q:fiber}
设 \(q:X\to Q\) 满射，\(f:X\to Y\) 为任意观察。
存在唯一 \(\bar f:Q\to Y\) 使 \(f=\bar f q\)，
当且仅当
\[
q(x)=q(y)\Longrightarrow f(x)=f(y).
\]
\end{lemma}
\begin{certificate}
通过证书是每个商纤维内观察值恒定的记录。
否决证书是一对满足 \(q(x)=q(y)\)、
但 \(f(x)\ne f(y)\) 的状态。
\end{certificate}
\begin{proof}
\begin{semanticsteps}
\item 若因子化成立，则同一商值经过同一函数取得同一观察值。
\item 反向定义 \(\bar f(q(x))=f(x)\)；
纤维恒定性保证定义不依赖代表元。
\item 满射性保证每个商值都有代表，并使 \(\bar f\) 唯一。
\end{semanticsteps}
\end{proof}
\begin{corollary}[结果观察与资格观察分别保持]
\label{q:outcome-type}
若 \(U_0=\bar U_0q_0\)，但存在
\[
q_0(x)=q_0(y),\qquad \tau(x)\ne\tau(y),
\]
则原表示完整保存 \(U_0\)，却不能完整读回真实类型 \(\tau\)。
因此，相同结果不足以替代合同规定的过程资格。
\end{corollary}
\begin{certificate}
记录共同结果、不同真实类型、
产生差异的合同条款及支持真实类型判断的独立证据。
\end{certificate}
\begin{proof}
\begin{semanticsteps}
\item 对 \(U_0\) 使用已给定的因子化。
\item 对 \(\tau\) 使用引理~\ref{q:fiber} 的必要性，
同像异型状态立即否定其因子化。
\end{semanticsteps}
\end{proof}
\section{联合分型的最粗静态细化}
\begin{theorem}[同时保留原任务与真实类型的最粗表示]
\label{q:joint}
定义
\[
q^+(x)=(q_0(x),\tau(x)),\qquad
Q^+=\operatorname{im}(q_0,\tau).
\]
则 \(q^+:X\to Q^+\) 满射，
同时保留 \(q_0\)、\(\tau\) 及所有通过 \(q_0\) 因子化的观察。
对任意满射 \(p:X\to P\)，若存在映射 \(a,c\) 使
\[
q_0=ap,\qquad \tau=cp,
\]
则存在唯一 \(h:P\to Q^+\) 使 \(q^+=hp\)。
因此，\(q^+\) 是同时满足这两项保持义务的最粗静态细化。
\end{theorem}
\begin{certificate}
保存两坐标投影及公式
\[
h(p(x))=(q_0(x),\tau(x)).
\]
值域使用实际像 \(Q^+\)，不加入从未出现的类型组合。
对每个旧商纤维列出实际类型及其证据来源。
\end{certificate}
\begin{proof}
\begin{semanticsteps}
\item 坐标投影给出
\[
q_0=\operatorname{pr}_1q^+,\qquad
\tau=\operatorname{pr}_2q^+.
\]
若 \(f=\bar f q_0\)，则
\(f=\bar f\operatorname{pr}_1q^+\)。
\item 当 \(p(x)=p(y)\) 时，
\(q_0(x)=a(p(x))=a(p(y))=q_0(y)\)，
并且 \(\tau(x)=c(p(x))=c(p(y))=\tau(y)\)。
\item 故 \(h\) 良定义。
\item 每个 \(P\) 中元素均由某个 \(x\) 表示，
所以 \(h\) 唯一且满射。
\item 任何满足两项保持义务的表示都能再映向 \(Q^+\)，
这正是所声明细化次序中的最粗性。
\end{semanticsteps}
\end{proof}
\begin{corollary}[新增区分的精确计数]
\label{q:count}
当 \(X\) 有限时，令
\[
K_y=\left|\tau(q_0^{-1}(y))\right|.
\]
则
\[
|Q^+|=\sum_{y\in Q_0}K_y.
\]
若全部 \(K_y=1\)，旧表示已经保存真实类型；
只有 \(K_y>1\) 的纤维需要增加区分。
\end{corollary}
\begin{certificate}
逐旧商值保存其非空类型清单，
并核对清单与原始状态集合的覆盖关系。
\end{certificate}
\begin{proof}
\begin{semanticsteps}
\item 不同 \(y\) 对应互不相交的表示块；
第 \(y\) 块恰含 \(K_y\) 个实际类型。
\item 求和得到计数式。
\item 全部纤维类型恒定时，应用引理~\ref{q:fiber}。
\end{semanticsteps}
\end{proof}
\begin{proposition}[独立证据恢复真实类型的条件]
\label{q:evidence}
令 \(e:X\to E\) 为允许取得的附加证据。
存在解码器 \(g\) 满足
\[
\tau(x)=g(q_0(x),e(x)),
\]
当且仅当 \(\tau\) 在 \((q_0,e)\) 的每个纤维上恒定。
若存在同一联合证据下的异型状态，
则任何只读取该联合证据的确定性程序都不能在两状态上同时正确。
\end{proposition}
\begin{certificate}
保存证据产生过程、访问范围、规则版本及解码表；
否决记录保存两状态的全部允许输入与不同真实类型。
证据的正确取得与类型码的存储长度分别核验。
\end{certificate}
\begin{proof}
\begin{semanticsteps}
\item 解码器存在的充要条件是引理~\ref{q:fiber}
在联合映射 \((q_0,e)\) 上的应用。
\item 若两状态的允许输入相同，
则对程序执行步数归纳，其内部状态及下一步操作相同，
终止时输出亦相同。
\item 两种不同真实类型需要不同正确输出，故不能同时恢复。
\end{semanticsteps}
\end{proof}
\section{全部有限后继上的观察饱和}
\begin{definition}[动作与响应的总化]
\label{q:totalization}
令 \(\Lambda=\mathsf A\times\mathsf Z\)，
并加入未定义状态
\[
X_\bot=X\sqcup\{\bot\}.
\]
对 \(\ell=(a,\xi)\) 定义
\[
\widehat\Phi_\ell(x)=
\begin{cases}
\delta(x,a,\xi),&
a\in A(x),\ \xi\in\Xi(x,a),\\
\bot,&\text{其他情形},
\end{cases}
\qquad
\widehat\Phi_\ell(\bot)=\bot.
\]
未定义状态用于保留合法域信息。
它不表示某项尚未完成的核验已经给出否定结论。
\end{definition}
\begin{definition}[初始测试与演化饱和]
\label{q:saturation}
取有限初始测试族 \(\mathcal P_0\)，
包括 \(q_0,\tau\)、目标指示、全部规定经济观察、
未定义状态指示，以及每个动作和响应的可用性。
原本只定义在 \(X\) 上的测试，在 \(\bot\) 处
赋予各自区别于合法值的未定义标记。
对逐步效用 \(U(x,a,\xi)\)，
固定标签后将其扩为状态测试，
在不合法处使用区别于合法值的标记。
递归定义
\[
\mathcal P_{n+1}
=\mathcal P_n\cup
\{f\circ\widehat\Phi_\ell:
f\in\mathcal P_n,\ \ell\in\Lambda\},
\]
\[
x\sim_n y
\Longleftrightarrow
\forall f\in\mathcal P_n,\quad f(x)=f(y).
\]
词 \(w=\ell_1\cdots\ell_k\) 的后继约定为
\[
\widehat\Phi_w
=\widehat\Phi_{\ell_k}\cdots\widehat\Phi_{\ell_1},
\qquad
\widehat\Phi_{\varnothing}=\operatorname{id}.
\]
\end{definition}
\begin{lemma}[有限词与递归测试等价]
\label{q:word-tests}
\(\mathcal P_n\) 所产生的等价关系，
恰由全部长度不超过 \(n\) 的测试
\[
f\circ\widehat\Phi_w,\qquad f\in\mathcal P_0,
\]
共同决定。
\end{lemma}
\begin{certificate}
为每个新测试保存其初始测试及标签词；
为每个标签词保存从递归规则得到该测试的有限推导。
\end{certificate}
\begin{proof}
\begin{semanticsteps}
\item \(n=0\) 时只含空词，结论成立。
\item 由长度不超过 \(n\) 的词再复合一个标签，
得到长度不超过 \(n+1\) 的词；
反之，每个非空词都能去掉一次复合，
还原为前次迭代已有测试。
\item 对 \(n\) 归纳即得结论。
\end{semanticsteps}
\end{proof}
\begin{theorem}[最粗动态充分商与有限稳定]
\label{q:dynamic}
设初始分区含 \(b_0\) 个非空块。
上述分区至多发生
\[
|X_\bot|-b_0
\]
次严格细化后稳定。
稳定关系可写为
\[
x\sim_*y
\Longleftrightarrow
\forall w\in\Lambda^*,\ \forall f\in\mathcal P_0,\quad
f(\widehat\Phi_wx)=f(\widehat\Phi_wy).
\]
稳定商 \(q_*:X_\bot\to Q_*=X_\bot/{\sim_*}\)
保持所有有限后继上的初始测试，
并存在良定义转移 \(\bar\Phi_\ell\) 满足
\[
q_*\widehat\Phi_\ell=\bar\Phi_\ell q_*.
\]
它是保持这些观察与标签后继的最粗商。
\end{theorem}
\begin{certificate}
每次分裂保存同一旧块中的异值状态对、
区分它们的初始测试及有限标签词。
终止时保存每块的初始测试恒定性、
每个标签的后继同块性及全部状态的覆盖表。
\end{certificate}
\begin{proof}
\begin{semanticsteps}
\item 每次严格细化至少增加一个非空块，
而分区块数不超过 \(|X_\bot|\)，故严格细化次数有界。
\item 若某一轮不再细化，则
\[
x\sim_ny\Longrightarrow
\widehat\Phi_\ell x\sim_n\widehat\Phi_\ell y,
\]
否则某个 \(f\circ\widehat\Phi_\ell\) 会在下一轮分开二者。
\item 因此该轮关系对每个标签稳定。
\item 对词长归纳，任意有限复合均保持同块性，
所以此后不会再产生新的分裂。
\item 初始测试在每块恒定，全部有限后继观察遂可读回。
\item 商转移定义为
\(\bar\Phi_\ell([x])=[\widehat\Phi_\ell x]\)，
稳定性保证代表元无关。
\item 最后，任意保持初始测试及后继的表示 \(p\)，
都由相同词长归纳保持全部
\(f\circ\widehat\Phi_w\)。
\item 故 \(p(x)=p(y)\) 蕴含 \(x\sim_*y\)，
存在唯一 \(h\) 使 \(q_*=hp\)。
\item 由此得到最粗性。
\end{semanticsteps}
\end{proof}
\begin{corollary}[新增经济义务后的最粗细化]
\label{q:new-observation}
若新增观察 \(g\) 已通过 \(q_*\) 因子化，
原稳定商同时保持 \(g\) 的全部有限后继。
若存在同商像而异 \(g\) 值状态，
则任何保留新观察的商必须区分该状态对。
以 \(\mathcal P_0\cup\{g\}\) 重新饱和，
得到加入新义务后的最粗动态充分商。
\end{corollary}
\begin{certificate}
保存 \(g=\bar gq_*\) 的读回式，
或保存否决旧表示的同像异值对；
需要细化时另保存新的稳定分区。
\end{certificate}
\begin{proof}
\begin{semanticsteps}
\item 因子化情形下，
\[
g\widehat\Phi_w=\bar g\bar\Phi_wq_*,
\]
故全部后继可读回。
\item 异值情形由引理~\ref{q:fiber} 排除旧商；
新商的存在和最粗性由定理~\ref{q:dynamic} 给出。
\end{semanticsteps}
\end{proof}
\section{精确任务商与鲁棒有限可达性}
\begin{definition}[精确经济任务商]
\label{q:exact}
在共同动作与响应标签下，满射 \(q:X\to\bar X\)
称为精确经济任务商，若
\[
\begin{aligned}
A(x)&=\bar A(qx),&
\Xi(x,a)&=\bar\Xi(qx,a),\\
q\delta(x,a,\xi)&=\bar\delta(qx,a,\xi),&
G&=q^{-1}(\bar G),\\
U(x,a,\xi)&=\bar U(qx,a,\xi),
\end{aligned}
\]
并且每个规定观察都通过相应的状态或执行上下文投影因子化。
\end{definition}
\begin{proposition}[稳定商给出完整任务接口]
\label{q:stable-exact}
按定义~\ref{q:saturation} 纳入全部动作、
响应、目标、效用与经济观察时，
稳定商在 \(X\) 上的限制满足定义~\ref{q:exact}。
\end{proposition}
\begin{certificate}
逐块核对动作集合、每个动作的完整响应集合、
目标真值、各标签效用及规定观察；
逐标签保存后继交换式。
\end{certificate}
\begin{proof}
\begin{semanticsteps}
\item 未定义状态指示使 \(\bot\) 单独成类。
\item 动作与响应的可用性测试使它们在每个合法状态块内恒定，
因此可在商上定义同一量词域。
\item 目标、效用和经济观察分别由初始测试恒定性下降。
\item 后继交换来自定理~\ref{q:dynamic}。
\item 这些恰为精确任务商的全部条件。
\end{semanticsteps}
\end{proof}
\begin{theorem}[有限执行保持的归纳证书]
\label{q:execution}
精确任务商保持任意有限合法执行。
将 \(q\) 按 \(q(\bot)=\bar\bot\) 扩展到未定义状态，
以下仍以 \(q\) 表示该扩展。
对共同标签词 \(w=\ell_1\cdots\ell_n\)，有
\[
q\widehat\Phi_w=\bar\Phi_wq.
\]
执行中每一步的合法性、规定效用与终态接受结果均对应。
\end{theorem}
\begin{certificate}
保存初态投影及逐步交换方块；
支付、费用、期限和资格按各自的观察逐步读回。
\end{certificate}
\begin{proof}
\begin{semanticsteps}
\item 空词的等式为 \(q=q\)。
\item 若长度 \(n\) 成立，则追加标签 \(\ell\) 后
\[
q\widehat\Phi_\ell\widehat\Phi_w
=\bar\Phi_\ell q\widehat\Phi_w
=\bar\Phi_\ell\bar\Phi_wq.
\]
这完成词长归纳。
\item 共同合法域排除非法标签；
目标原像与效用因子化给出接受和评价保持。
\end{semanticsteps}
\end{proof}
\begin{theorem}[鲁棒可达集的双向保持]
\label{q:reachability}
令 \(W_0=G\)，并递归定义
\[
W_{n+1}
=W_n\cup
\{x:\exists a\in A(x),\
\forall\xi\in\Xi(x,a),\
\delta(x,a,\xi)\in W_n\}.
\]
商系统采用相同递推得到 \(\bar W_n\)。
对每个 \(n\ge0\)，
\[
W_n=q^{-1}(\bar W_n).
\]
\end{theorem}
\begin{certificate}
保存目标原像、动作和响应的完整对应、
每轮可达集合及归纳中的量词替换。
合法动作的响应域非空，避免以空量词代替实际执行。
\end{certificate}
\begin{proof}
\begin{semanticsteps}
\item \(n=0\) 由目标原像式成立。
\item 假设 \(W_n=q^{-1}(\bar W_n)\)。
\item 共同动作、共同响应和后继交换给出
\[
\begin{aligned}
&\exists a\in A(x)\ \forall\xi\in\Xi(x,a):
\delta(x,a,\xi)\in W_n\\
\Longleftrightarrow\quad&
\exists a\in\bar A(qx)\ \forall\xi\in\bar\Xi(qx,a):
\bar\delta(qx,a,\xi)\in\bar W_n.
\end{aligned}
\]
两侧鲁棒前驱互为原像；
原像与并集交换，故结论在 \(n+1\) 成立。
\item 数学归纳法完成证明。
\end{semanticsteps}
\end{proof}
\begin{corollary}[有限步策略的合法回落]
\label{q:policy}
在 \(\bar W_n\) 上令
\[
\rho(\bar x)=\min\{k:\bar x\in\bar W_k\}.
\]
为每个正秩状态选择一个使全部后继秩严格下降的动作，
并令原系统采用相同动作。
该策略对全部允许响应在至多 \(n\) 步内进入 \(G\)。
\end{corollary}
\begin{certificate}
逐状态保存选择动作、完整响应清单与后继秩；
原系统同时核对动作的实际权限和输入可用性。
\end{certificate}
\begin{proof}
\begin{semanticsteps}
\item 可达递推保证每个正秩状态存在所需动作。
\item 精确任务商保证该动作可以原层执行，
且全部响应的投影秩至少下降一。
\item 自然数秩至多下降 \(n\) 次到零，
此时商状态属于目标，原状态由目标原像式属于 \(G\)。
\end{semanticsteps}
\end{proof}
\section{不同经济制度之间的共同任务}
\begin{definition}[共同的许可、履约与权利任务]
\label{q:three-institutions}
记 \(\mathsf F\) 为高后果设施的许可与履约治理，\(\mathsf A\) 为可验证自动化的用途治理，
\(\mathsf H\) 为真人参与合同治理。
选择共同任务时，保留资格、实际过程、责任、资源、证据、具体权利、修复和延续条件。
核安全制度只为 \(\mathsf F\) 的许可责任结构提供公开参照；行业实际物理状态及不可转让权利仍保留在各自完整任务中。
\end{definition}

\begingroup
\small
\renewcommand{\arraystretch}{1.16}
\begin{longtable}{p{0.15\textwidth}p{0.24\textwidth}p{0.24\textwidth}p{0.24\textwidth}}
\toprule
共同观察 & 设施许可与履约 & 自动化用途治理 & 真人参与合同\\
\midrule
\endfirsthead
\toprule
共同观察 & 设施许可与履约 & 自动化用途治理 & 真人参与合同\\
\midrule
\endhead
资格与用途 & 责任主体、持证活动与许可范围 & 系统、用途及实际控制权限 & 真人身份、接纳及允许工具\\
过程与证据 & 运行条件、核查及履约记录 & 规范、证明、执行和实际响应 & 真实参与、完成及复核记录\\
资源与权利 & 运行资源、到期合同义务与具体债权 & 执行预算、责任与交付义务 & 人数上限、逐人应付及实际到账\\
纠正与履约 & 对声明违约和资源缺口的合法纠正 & 越权、错误执行和履约缺陷修复 & 冒领纠正、追回和逐人补付\\
继续接入 & 具备资源、许可和安全条件后续行 & 版本与任务条件保持后继续部署 & 具备资格、现金和容量后再参与\\
\bottomrule
\end{longtable}
\endgroup

表中的对应只登记需要证明的共同任务字段，不直接签发同构。
只有共同字段的全部动作、响应与接受判定实现精确对应时，才能应用以下最粗商定理。
资源账户之间的支付映射也仅适用于可依法计量和转移的项目，不能把物理安全、不可逆损害或人的身份权利改写为任意货币转移。
\begin{definition}[双向覆盖的动作与响应实现]
\label{q:adapters}
当原系统与共同系统使用不同动作名称时，
对每个原状态 \(x\) 给定满射
\[
\alpha_x:A(x)\twoheadrightarrow\bar A(qx),
\]
并对每个合法原动作给定满射
\[
\beta_{x,a}:\Xi(x,a)
\twoheadrightarrow\bar\Xi(qx,\alpha_x(a)).
\]
要求
\[
q\delta(x,a,\xi)
=\bar\delta(qx,\alpha_x(a),\beta_{x,a}(\xi)),
\]
并保持目标及所有规定观察。
凡属任务的实际费用、时间和权限，也须由该实现保存。
\end{definition}
\begin{proposition}[双向覆盖保持鲁棒量词]
\label{q:adapter-reachability}
定义~\ref{q:adapters} 的条件成立时，
定理~\ref{q:reachability} 的可达集原像结论仍成立。
\end{proposition}
\begin{certificate}
保存每个共同动作的至少一个原实现、
全部原动作的去向及全部响应的双向覆盖。
如一次共同动作由多个微观步骤实现，
另保存有限终止、累计费用和中间状态合法性。
\end{certificate}
\begin{proof}
\begin{semanticsteps}
\item 固定原动作 \(a\) 且假定它对全部原响应成功。
对任意共同响应 \(\bar\xi\)，满射 \(\beta_{x,a}\) 给出原响应 \(\xi\) 使 \(\beta_{x,a}(\xi)=\bar\xi\)。
\item 后继交换与归纳前件给出
\[
\delta(x,a,\xi)\in W_n
\Longrightarrow
\bar\delta(qx,\alpha_x(a),\bar\xi)\in\bar W_n.
\]
因 \(\bar\xi\) 任意，该共同动作对全部共同响应成功。
\item 反向，若共同动作 \(\bar a\) 成功，动作满射选择 \(a\) 满足 \(\alpha_x(a)=\bar a\)。
任意原响应映向一个共同响应，因此交换式与归纳前件给出所有原后继均属于 \(W_n\)。
\item 两方向合并得到鲁棒前驱的原像等式；目标为基例，按可达轮次归纳得
\(W_n=q^{-1}(\bar W_n)\)。
\end{semanticsteps}
\end{proof}
\begin{theorem}[共同观察下的最粗商同构]
\label{q:common-image}
设有限个制度实现 \(X_j\) 已采用共同的测试语义、
动作和完整响应接口。
在不交并 \(\bigsqcup_jX_j\) 上加入同一个未定义状态，
执行定理~\ref{q:dynamic} 的饱和，得到 \(q_*\)。
记 \(Y_j=q_*(X_j\cup\{\bot\})\)。
若
\[
Y_1=\cdots=Y_s=Y,
\]
则各制度相对于这组共同经济义务的最粗动态任务商，
均与 \(Y\) 同构。
\end{theorem}
\begin{certificate}
保存共同测试编码、
全部标签后继表、稳定分区及每个制度对 \(Y\) 的覆盖。
动作改名时附定义~\ref{q:adapters} 的双向覆盖证书。
共同像必须包含声明的非失败经济任务状态。
\end{certificate}
\begin{proof}
\begin{semanticsteps}
\item 两状态在全体不交并中的等价，
恰由相同的全部有限后继测试决定。
\item 限制到任一 \(X_j\) 后，
它仍是该制度同一观察族的最粗稳定关系。
\item 所以各自商与其像 \(Y_j\) 存在典范双射，
且该双射保持测试与转移。
\item 共同像相同使这些双射的值域相同，
复合得到所需任务商同构。
\end{semanticsteps}
\end{proof}
\begin{remark}
若各 \(Y_j\) 不同，可另研究其共同且对后继封闭的子任务。
不同制度各自存在最粗商，并不自动推出这些商相同。
用途名称可以通过共同语义编码比较；
凡身份、权利或规则差异仍影响任务观察的，
都应保留在商中。
仅有失败状态的共同像不构成经济循环的有效共同实例。
\end{remark}
\begin{proposition}[规则变化与合法子任务保持]
\label{q:rule-change}
设新规则排除某些旧规则允许的行为。
完整旧任务与新任务的合法动作集合因而不同。
若已经指定的旧合法参与子任务具有到新制度的
目标、观察及动作响应保持映射，
则该子任务的有限执行与鲁棒可达性可以继续保持。
\end{proposition}
\begin{certificate}
保存两版规则、所选子任务、被保留的合法动作、
全部响应及目标原像。
新增的资格核验和申诉步骤按实际费用、期限与接续关系登记。
\end{certificate}
\begin{proof}
\begin{semanticsteps}
\item 记受保护旧子任务为 \(\mathcal H\)，新实现映射为 \(j\)。证书给出
\(j\delta_{\mathcal H}(x,a,\xi)=\delta_{\rm new}(jx,\alpha a,\beta\xi)\)。
\item 对一个旧合法原子，动作和全部响应覆盖保证新侧可以实现同一任务步。
目标、权利与费用观察由证书保持。
\item 空路径保持初态；在长度 \(n\) 的路径后追加一个原子，用第一步交换式得到长度 \(n+1\) 的端点及观察对应。
\item 因而任意有限受保护路径保持。目标原像和完整响应覆盖又使鲁棒可达递推逐轮相同。
\item 被新规则删除的域外动作不属于 \(\mathcal H\)，不进入以上归纳量词域。
\end{semanticsteps}
\end{proof}
\chapter{保持经济观察的链商}
\label{ch:chapter}

动态任务商保存经济决策与资格判断，
链复形则记录守恒约束、允许修复及其组合关系。
二者通过明确的实现接口接续。
本章先给出保持线性经济观察的最粗链商，
再说明它与完整状态任务的关系。
链复形及其商采用通常的代数拓扑定义
\cite{Hatcher2002}；
观察保持与执行保持分别使用任务接口
\cite{InternalTheory,InternalEconomics}。
\section{链复形、经济读出与观察核}
\begin{definition}[有限经济链实现]
\label{ch:complex}
令 \(C_\bullet\) 为有限维有理链复形，
仅有有限多个非零次数：
\[
\cdots\xrightarrow{d_{n+1}}C_n
\xrightarrow{d_n}C_{n-1}\xrightarrow{}\cdots,
\qquad d_{n-1}d_n=0.
\]
在每个次数给定线性观察
\[
O_n:C_n\to V_n.
\]
不同金额单位、权利类别及核算窗口分别编码；
不能将量纲不同的数值直接相加。
约定 \(d_0=0\)、\(V_{-1}=0\)。
链元表示形式效果及其组合；
实际余额非负、权限和期限属于额外执行条件。
\end{definition}
\begin{lemma}[线性观察下降的核判据]
\label{ch:factor}
对满射线性映射 \(p:C_n\to D_n\)，
存在唯一线性 \(\bar O_n\) 使 \(O_n=\bar O_np\)，
当且仅当
\[
\ker p\subseteq\ker O_n.
\]
\end{lemma}
\begin{certificate}
保存 \(\ker p\) 的基及逐基观察值，
并按代表元构造诱导读出 \(\bar O_n\)。
\end{certificate}
\begin{proof}
\begin{semanticsteps}
\item 若因子化成立，核元必有零观察。
\item 反向定义 \(\bar O_n(p(c))=O_n(c)\)；
代表之差属于 \(\ker p\)，故读出不变。
\item 线性性由定义得到，满射性保证唯一。
\end{semanticsteps}
\end{proof}
\begin{theorem}[最大不可观察子复形]
\label{ch:maximal-kernel}
定义
\[
\boxed{
N_n=\ker O_n\cap d_n^{-1}(\ker O_{n-1}).
}
\]
则 \(N_\bullet\) 是包含于逐次观察核中的最大子复形。
因此，典范链商
\[
q_C:C_\bullet\twoheadrightarrow C_\bullet/N_\bullet
\]
是保持所声明线性观察的最粗链商。
\end{theorem}
\begin{certificate}
逐次数保存 \(O_n,d_n\)、
\(N_n\) 的基及两项零式
\[
O_nN_n=0,\qquad O_{n-1}d_nN_n=0.
\]
另保存 \(d^2=0\) 及与任一候选链商的核包含证明。
\end{certificate}
\begin{proof}
\begin{semanticsteps}
\item 若 \(c\in N_n\)，则 \(O_{n-1}(d_nc)=0\)，
且
\[
O_{n-2}d_{n-1}d_nc=0.
\]
故 \(d_nc\in N_{n-1}\)，所以 \(N\) 为子复形。
\item 若另一子复形 \(K\) 满足 \(K_n\subseteq\ker O_n\)，
则对 \(c\in K_n\)，有
\[
d_nc\in K_{n-1}\subseteq\ker O_{n-1},
\]
从而 \(c\in N_n\)。
\item 故 \(K\subseteq N\)，最大性成立。
\item 任一保观察的满射链映射 \(p:C\to D\)
都有子复形核 \(\ker p\subseteq\ker O\)，
因而 \(\ker p\subseteq N\)。
\item 定义 \(h(p(c))=c+N\)，得到唯一满射链映射
\[
h:D\to C/N,\qquad q_C=hp.
\]
这给出最粗链商的普遍性质。
\end{semanticsteps}
\end{proof}
\begin{theorem}[观察与边界观察的显式商表示]
\label{ch:explicit-image}
令
\[
S_n(c)=(O_nc,\ O_{n-1}d_nc).
\]
则
\[
\ker S_n=N_n,\qquad
(C/N)_n\cong\operatorname{im}S_n.
\]
在该像上，商边界为
\[
\bar d_n(a,b)=(b,0).
\]
\end{theorem}
\begin{certificate}
保存矩阵
\[
S_n=
\begin{pmatrix}
O_n\\ O_{n-1}d_n
\end{pmatrix},
\]
其核基、像基以及
\(S_{n-1}d_n=\bar d_nS_n\)。
\end{certificate}
\begin{proof}
\begin{semanticsteps}
\item 核等式直接来自两坐标同时为零。
\item 第一同构定理给出像表示。
\item 对任意 \(c\in C_n\)，
\[
S_{n-1}d_nc
=(O_{n-1}d_nc,\ O_{n-2}d_{n-1}d_nc)
=(O_{n-1}d_nc,0).
\]
故所列边界良定义且与原边界交换。
\item 再次应用该式得到 \(\bar d^2=0\)。
\end{semanticsteps}
\end{proof}
\section{任意模板商与最粗性的区别}
\begin{proposition}[模板商向最粗观察商的因子化]
\label{ch:template}
若 \(F:C\twoheadrightarrow D\) 为满射链映射，
且 \(O=\bar OF\)，则
\[
\ker F\subseteq N,\qquad
C/\ker F\twoheadrightarrow C/N.
\]
两者相等的充要条件是 \(\ker F=N\)。
若 \(D_n\) 上的映射
\[
v\longmapsto
(\bar O_nv,\ \bar O_{n-1}d_nv)
\]
在每个次数单射，则该相等条件成立。
\end{proposition}
\begin{certificate}
保存 \(dF=Fd\)、观察因子化、
模板上两组观察组成的矩阵及其满列秩证书。
\end{certificate}
\begin{proof}
\begin{semanticsteps}
\item 核包含由定理~\ref{ch:maximal-kernel} 得到。
\item 若模板上的联合观察单射，而 \(c\in N_n\)，
则 \(F_nc\) 的观察及边界观察均为零，
所以 \(F_nc=0\)。
\item 于是 \(N_n\subseteq\ker F_n\)，与反向包含合并。
\item 一般情形下，第一同构定理给出两商之间的满射；
它为同构恰当且仅当两个核相等。
\end{semanticsteps}
\end{proof}
\begin{theorem}[共同观察像给出的链商同构]
\label{ch:common-image}
设若干经济链实现 \(C^{(j)}\)
使用相同语义与计量单位的观察空间 \(V_n\)。
按定理~\ref{ch:explicit-image} 定义 \(S_n^{(j)}\)。
若逐次数满足
\[
\operatorname{im}S_n^{(1)}
=\cdots=\operatorname{im}S_n^{(s)},
\]
则它们的最粗观察链商存在保持观察的典范同构。
\end{theorem}
\begin{certificate}
在共同坐标中保存各像空间的基，
并以双向线性表示验证像空间相等。
同时核对单位、主体索引和权利字段的对应。
\end{certificate}
\begin{proof}
\begin{semanticsteps}
\item 各商均同构于自己的观察像，
而像上的边界均为 \((a,b)\mapsto(b,0)\)。
\item 相同像空间上的恒等映射保持边界和第一坐标观察。
\item 与各自像表示同构复合，即得结论。
\end{semanticsteps}
\end{proof}
\begin{remark}
逐制度构造最大不可观察核，
并不能单独证明这些制度具有相同观察像。
同调群维数相同也不提供上述逐坐标同构。
当某个制度保留更多逐人权利、付款期限或合法动作时，
应将这些差异纳入共同任务或明确限定比较的子任务。
\end{remark}
\section{动态链观察与完整状态任务}
\begin{proposition}[链作用下的动态观察饱和]
\label{ch:dynamic-observations}
设给定有限个链自映射 \(T_\ell:C\to C\)，
满足 \(dT_\ell=T_\ell d\)。
将观察扩为全部
\[
O_nT_{w,n},\qquad w\in\Lambda^*.
\]
以该观察族构造的最大不可观察子复形 \(N^{\rm dyn}\)
对每个 \(T_\ell\) 稳定，
故这些作用可下降到 \(C/N^{\rm dyn}\)。
在有限维、有限次数支撑情形，
每个次数的全部观察可用有限个线性无关行向量表示。
\end{proposition}
\begin{certificate}
保存链作用矩阵、与边界的交换式、
饱和观察的行空间基及其对右复合 \(T_\ell\) 的稳定性。
\end{certificate}
\begin{proof}
\begin{semanticsteps}
\item 若 \(c\in N_n^{\rm dyn}\)，则每个
\(O_nT_{w,n}c\) 和
\(O_{n-1}T_{w,n-1}d_nc\) 均为零。
\item 对 \(T_\ell c\) 的相同观察，
只是再增加一个有限复合标签；
边界项由 \(dT_\ell=T_\ell d\) 作相同处理。
\item 故 \(T_\ell c\in N_n^{\rm dyn}\)。
\item 全部行向量位于有限维对偶空间 \(C_n^*\)，
因此存在有限基。
\item 逐次加入新行时，严格增加的秩有有限上界。
\end{semanticsteps}
\end{proof}
\begin{remark}
本命题要求已经给出与边界相容的线性作用。
一般状态转移可以是非线性的，
不能仅由金额矩阵推定其存在此种线性实现。
此时仍须独立使用第~\ref{q:chapter}章的动态状态商，
并验证状态表示与链表示之间的接口。
\end{remark}
\begin{proposition}[线性关系与状态等价关系的不同类型]
\label{ch:linear-state}
令自由向量空间以三个状态 \(x_0,x_1,x_2\) 为基，
观察矩阵为
\[
O=
\begin{pmatrix}
1&1&1\\
0&1&2
\end{pmatrix},
\qquad d=0.
\]
则
\[
\ker O=\mathbb Q(e_0-2e_1+e_2).
\]
该核不是由一个非平凡状态分区的同类基向量之差生成的核。
保留这些线性观察也不保证保留混合状态上的任意非线性谓词。
\end{proposition}
\begin{certificate}
保存三个互异观察列及核向量。
取两个非负、总质量为一的混合状态
\[
v=\tfrac12e_0+\tfrac12e_2,\qquad w=e_1.
\]
它们有相同线性观察，
但“状态 \(x_0\) 的质量为正”在二者上的真值不同。
\end{certificate}
\begin{proof}
\begin{semanticsteps}
\item 直接相乘得到核向量，
矩阵秩为二，故核恰由该向量张成。
\item 三个观察列互异，所以任何两个不同基向量之差均不属于核。
\item 非平凡状态分区必识别至少两个基状态，
其自由线性延拓的核必含对应基向量之差，故不可能等于本核。
\item 对所列 \(v,w\) 直接计算观察和非线性谓词，得到最后结论。
\end{semanticsteps}
\end{proof}
\begin{theorem}[状态商的自由图链实现]
\label{ch:graph-realization}
设 \(q:X\to Q\) 为已经通过全部任务观察核验的精确状态商。
记 \(\Phi_\ell\) 为合法域上的标签后继。
以状态为顶点，以每个合法标签转移为有向边，构造图链复形：
\[
C_0=\mathbb Q[X],\qquad
C_1=\mathbb Q[(x,\ell)\text{ 为合法转移}],\qquad
C_{n\ge2}=0,
\]
\[
d_1e_{(x,\ell)}
=e_{\Phi_\ell x}-e_x.
\]
商图以 \(Q\) 及其合法标签边构造。
定义
\[
F_0e_x=e_{q(x)},\qquad
F_1e_{(x,\ell)}=e_{(q(x),\ell)}.
\]
则 \(F\) 是逐次数满射链映射。
以商顶点及商边的全部坐标为观察时，
该图链商是相应观察族的最粗链商，
且原任务谓词可以从商顶点读回。
\end{theorem}
\begin{certificate}
保存顶点对应、合法标签对应及后继交换。
商顶点和商边坐标采用指示函数编码；
任何新增高阶关系都另给边界相容证书。
\end{certificate}
\begin{proof}
\begin{semanticsteps}
\item 后继交换给出
\[
\begin{aligned}
F_0d_1e_{(x,\ell)}
&=e_{q(\Phi_\ell x)}-e_{q(x)}\\
&=e_{\bar\Phi_\ell q(x)}-e_{q(x)}
=\bar d_1F_1e_{(x,\ell)}.
\end{aligned}
\]
状态满射及合法标签保持给出两个次数的满射。
\item 当 \(O_n=F_n\) 时，
\(F_n c=0\) 已蕴含 \(F_{n-1}d_nc=0\)，
所以最大不可观察核恰为 \(\ker F\)。
\item 每个任务谓词先通过状态商因子化，
商顶点的指示编码又能恢复该商状态，
因此谓词仍可读回。
\end{semanticsteps}
\end{proof}
\begin{definition}[状态与账本链的相容接口]
\label{ch:interface}
对于一般账本链实现，给定状态编码
\(\iota:X\to C_k\)、状态商 \(q_X:X\to Q_X\)
及链商 \(q_C:C_k\to\bar C_k\)。
若存在 \(\bar\iota\) 和读回映射 \(\sigma\)，使
\[
q_C\iota=\bar\iota q_X,\qquad
\sigma\bar\iota=\operatorname{id}_{Q_X},
\]
并且实际动作的实现与这些映射交换，
则称两类表示对该任务相容。
只用于账本子任务时，\(Q_X\) 也限于已声明的账本观察；
完整资格任务需要保留完整状态信息。
\end{definition}
\begin{proposition}[相容接口下的任务读回]
\label{ch:readback}
定义~\ref{ch:interface} 成立时，
每个 \(f=\bar f q_X\) 均满足
\[
f=\bar f\,\sigma\,q_C\iota.
\]
若动作实现亦交换，则该等式沿任意有限合法执行保持。
\end{proposition}
\begin{certificate}
保存两个接口等式、每个规定谓词的状态读回式、
以及逐步动作实现的交换方块。
\end{certificate}
\begin{proof}
\begin{semanticsteps}
\item 由两个接口等式，
\[
\sigma q_C\iota=\sigma\bar\iota q_X=q_X,
\]
再复合 \(\bar f\) 即得读回。
\item 有限执行部分对步骤数归纳，
每增加一步使用一次动作交换式，
与定理~\ref{q:execution} 的归纳相同。
\end{semanticsteps}
\end{proof}
\Needspace{24\baselineskip}
\section{逐人权利与不可删除的观察}
\begin{theorem}[群体总额不能替代逐人权利]
\label{ch:individual}
设账本包含至少两名权利主体。
聚合映射
\[
L(h_1,\ldots,h_m,o,x)
=\left(\sum_{i=1}^m h_i,o,x\right)
\]
保存群体总额，却不能保存每个人的金额读出。
若第 \(i\) 人与第 \(j\) 人的权利均属于任务，
则非零方向 \(e_{H_i}-e_{H_j}\)
不得属于保持这些权利的观察核。
\end{theorem}
\begin{certificate}
对 \(i\ne j\)，取
\[
\nu=a(e_{H_i}-e_{H_j}),\qquad a\ne0.
\]
保存 \(L\nu=0\) 及第 \(i\) 人读出
\(\ell_i(\nu)=a\)。
应付、到账、期限和申诉状态分别登记；
余额读出不自动替代其他权利字段。
\end{certificate}
\begin{proof}
\begin{semanticsteps}
\item 两个人类坐标相加为零，故 \(L\nu=0\)；
第 \(i\) 人的金额变化为 \(a\)。
\item 由引理~\ref{ch:factor}，
\(\ell_i\) 不能通过 \(L\) 因子化。
\item 若所选观察包含 \(\ell_i\)，则
\(\nu\notin\ker O_1\)，从而 \(\nu\notin N_1\)。
\end{semanticsteps}
\end{proof}
\begin{remark}
最粗性始终相对于已经明确的观察族。
保留逐人权利的最粗商，
可以比只保留群体总额的最粗商更细。
非线性资格、到账最终性和时间条件，
分别由状态纤维判据或相容接口证明保持。
代数压缩本身不产生新的支付权、兑换权或执行权限。
\end{remark}

\chapter{标准化服务接口与共同观察链}
\label{integrated:standard-service}
\label{std:chapter}

标准化接口规定不同实现如何提交相同类型的请求、
给出可检查的决定，并履行相同的权利与资源义务。
共同语法只确定字段可以传递；
共同任务还要求这些字段、全部动作及后继响应具有相同语义。
本章依次核对治理决定的可表达性、证据恢复与有限计算、
合同级动作实现、最粗动态商及其自由图链。
所用任务充分性与最粗商次序承接第~\ref{q:chapter}~章，
链商的普遍性质承接第~\ref{ch:chapter}~章\cite{InternalTheory}。

\section{共同合同与版本语义}

\begin{definition}[共同服务合同]
\label{std:contract}
固定合同任务
\[
\mathcal T=(\mathsf I,\mathsf U,\mathsf V,
\mathsf A,\mathsf Z,\mathcal O,G).
\]
各分量分别登记主体、用途、合同版本、
动作与响应标签、必须保留的观察以及接受目标。
这些观察与目标接口在各 \(X_i\) 上具有带共同值域的函数解释；
对应目标集合记为 \(G_i\subseteq X_i\)。
观察包括身份与委托、实际用途、作用范围、
证书绑定、预算配额、实际费用与时间、
执行回执和逐请求逐主体的权利。
拒绝、待核验、退款、申诉和重新接入均属于合同状态。
对两种实现 \(i\in\{1,2\}\)，
以 \(X_i\) 表示完成该合同编码后的有限状态空间。
底层实际系统另通过符合关系解释到 \(X_i\)。
\end{definition}

\begin{definition}[合同级版本与主体对应]
\label{std:normalization}
第 \(i\) 种实现给出主体对应 \(\chi_i\)、
版本对应 \(\mu_i\) 及状态标准化 \(n_i\)：
\[
\chi_i:I_i\to\mathsf I,\qquad
\mu_i:V_i\to\mathsf V,\qquad
n_i:X_i^{\rm loc}\to X_i.
\]
对应中的主体必须具有共同合同所指的同一责任与权利身份。
本地程序、模型或证明对象的标识可以不同，
但实际对象与批准对象的绑定关系必须保持。
对每个共同观察 \(f\)，要求存在统一读出 \(\bar f\)，使
\[
f_i=\bar f n_i.
\]
若原始标识、证明材料或撤销记录本身属于必保留观察，
则另保存其读回映射；
不能仅留下合同版本名称而删除这些任务差异。
版本更新的全部动作与响应也进入后续共同任务核验。
\end{definition}

\begin{remark}[版本名称与后续责任]
两个实现当前使用相同合同版本名称，
并不证明其下一次撤销、收费或复核决定相同。
标准化必须保存实际版本绑定及每项已登记后继义务。
同一用途在不同技术载体中的准入保持，
采用定理~\ref{gov:medium-invariance}~的字段因子化条件；
该定理不替代这些字段取得正确及完整的证据。
\end{remark}

\section{治理决定的严格识别改进与有限算法}

\begin{definition}[表示能够表达的决定]
\label{std:decision-family}
对有限状态集合 \(X\) 上的满射表示 \(p:X\to P\)，定义
\[
\mathcal D(p)=
\{d:X\to\{0,1\}:
\exists\bar d:P\to\{0,1\},\ d=\bar d p\}.
\]
这表示决定由该表示的值唯一确定。
另给定合同实际要求的决定族 \(\mathcal F\)，并记
\[
\mathcal D_{\mathcal F}(p)=\mathcal F\cap\mathcal D(p).
\]
上述定义是函数因子化条件，
没有预先提供实际状态的读取程序、
决定函数的算法或其时间与资金上界。
\end{definition}

\begin{theorem}[结果与治理类型联合表示的严格改进]
\label{std:strict-decisions}
令 \(q_0:X\to Q_0\) 保存结果，
\(\tau:X\to T\) 保存规定治理类型，
并使用实际像定义
\[
q^+=(q_0,\tau),\qquad
Q^+=\operatorname{im}(q_0,\tau).
\]
则
\[
\mathcal D(q_0)\subseteq\mathcal D(q^+).
\]
若同一结果纤维内存在不同类型，该包含严格。
对合同决定族 \(\mathcal F\)，
只有当新增区分能区分其中某项实际要求的决定时，
才得到
\(\mathcal D_{\mathcal F}(q_0)
\subsetneq\mathcal D_{\mathcal F}(q^+)\)。
\end{theorem}
\begin{certificate}
两个坐标投影、原决定的因子化、
同结果异类型状态对，
以及用于证明合同任务严格改进的具体决定谓词。
\end{certificate}
\begin{proof}
\begin{semanticsteps}
\item 取 \(d\in\mathcal D(q_0)\)。
按定义存在 \(\bar d\) 使 \(d=\bar d q_0\)。
定理~\ref{q:joint}~的第一坐标投影给出
\[
q_0=\operatorname{pr}_1q^+.
\]
\item 以等式替换取得
\[
d=(\bar d\operatorname{pr}_1)q^+.
\]
右侧函数的类型为
\(X\to Q^+\to\{0,1\}\)，
因此 \(d\in\mathcal D(q^+)\)。
\item 设 \(q_0(x)=q_0(y)\) 而 \(\tau(x)\ne\tau(y)\)。
定义
\[
d_\star(z)=\mathbf1\{\tau(z)=\tau(x)\}.
\]
该函数由 \(q^+\) 的第二坐标读出，
所以 \(d_\star\in\mathcal D(q^+)\)。
\item 对见证状态，有
\[
d_\star(x)=1,\qquad d_\star(y)=0.
\]
引理~\ref{q:fiber}~排除
\(d_\star\in\mathcal D(q_0)\)，
故全体可表达决定的包含严格。
\item 与同一 \(\mathcal F\) 取交保留包含方向。
若证书另外给出
\[
d\in\mathcal F\cap
\bigl(\mathcal D(q^+)\setminus\mathcal D(q_0)\bigr),
\]
该 \(d\) 便是合同决定族严格包含的见证。
没有这一见证时，不把一般表达能力的增加宣布为合同效果已经改善。
\end{semanticsteps}
\end{proof}

\begin{corollary}[退款义务提供具体治理见证]
\label{std:refund-witness}
在命题~\ref{svc:failure-observation}~的两个请求状态中，
若治理类型记录尚有退款义务，
则决定
\[
d(x)=\mathbf1\{u_{ri}(x)-f_{ri}(x)>0\}
\]
可以由相应联合表示读回，
并且在相同未交付结果下可能无法由结果表示读回。
\end{corollary}
\begin{certificate}
取两个未交付请求的实际回执与应退依据，
分别登记未清偿金额为 \(1\) 和 \(0\)。
治理类型明确含有 \(d\) 的真值。
\end{certificate}
\begin{proof}
\begin{semanticsteps}
\item 两个未清偿金额属于同一货币、主体和合同口径，
故与零比较得到的 \(d\) 具有布尔类型。
\item 治理类型中的登记坐标使
\(d=\bar d q^+\)，从而
\(d\in\mathcal D(q^+)\)。
\item 两状态的结果表示相同，
而该决定分别为 \(1\) 和 \(0\)。
按引理~\ref{q:fiber}，
\(d\notin\mathcal D(q_0)\)。
\item 若该退款判断属于合同要求的 \(\mathcal F\)，
则它满足定理~\ref{std:strict-decisions}~最后一项的严格性条件。
\end{semanticsteps}
\end{proof}

\begin{lemma}[证据恢复治理类型的纤维条件]
\label{std:evidence-recovery}
设实际可取得的记录为
\[
h=(q_0,e):X\longrightarrow Q_0\times E.
\]
存在类型解码器 \(\gamma:h(X)\to T\)
使 \(\tau=\gamma h\) 的充要条件是
\[
h(x)=h(y)\Longrightarrow\tau(x)=\tau(y).
\]
当条件不成立时，
相同输入记录的确定性程序不能同时正确恢复这两个不同类型。
\end{lemma}
\begin{certificate}
记录映射的类型、合法证据来源、
每个非空记录纤维的类型表；
失败时保存一个同记录异类型的状态对。
\end{certificate}
\begin{proof}
\begin{semanticsteps}
\item 将命题~\ref{q:evidence}~中的联合记录
具体代入为 \(h=(q_0,e)\)。
其解码器存在条件正是所列纤维恒定式。
\item 若解码器存在且 \(h(x)=h(y)\)，
函数单值性依次给出
\[
\gamma(h(x))=\gamma(h(y)),
\qquad \tau(x)=\tau(y).
\]
\item 反向假设纤维恒定。
对实际出现的记录 \(w=h(x)\)，赋值
\[
\gamma(w)=\tau(x).
\]
另一个代表 \(y\) 满足 \(h(y)=w\) 时，
纤维条件保证所赋值不变。
\item 若同一记录具有不同真实类型，
确定性程序在相同输入及同一初始内部状态下，
每步计算与终止输出相同。
该输出不能同时等于两个不同类型。
由程序步骤的有限归纳得到所述限制。
\end{semanticsteps}
\end{proof}

\begin{theorem}[有限证据表上的可实施决定算法]
\label{std:finite-decision}
进一步假设 \(X=\{x_1,\ldots,x_N\}\) 有完整有限表，
记录编码相等性可判定，
每行 \((h(x_j),\tau(x_j))\) 已有独立有效证据，
且满足引理~\ref{std:evidence-recovery}~的纤维条件。
若所需决定的联合表示读出表也已给定，
则可构造一个在实际记录上恢复类型并计算该决定的有限算法。
实际执行还须取得正确的当前记录，
并覆盖该算法的时间、存储与费用。
\end{theorem}
\begin{certificate}
完整行表、类型证据、编码比较程序、
按记录分组后的字典、决定表、
记录取得程序及各操作资源上界。
\end{certificate}
\begin{proof}
\begin{semanticsteps}
\item 顺序读取 \(N\) 行，
按可判定的记录相等关系分组。
对每组比较其所有类型标签；
纤维条件使每组只有一个类型值。
\item 把每个非空组写为一个字典项
\[
(w,\gamma(w)).
\]
完整行表保证所有实际可出现的 \(w\in h(X)\)
都具有一个字典项。
\item 对实际状态 \(x\)，记录程序取得准确的 \(w=h(x)\)。
逐项扫描字典在至多 \(N\) 次比较后找到相同键，
并返回
\[
\gamma(w)=\tau(x).
\]
\item 原记录的第一坐标已给出 \(q_0(x)\)。
将其与返回的类型配对，构造
\[
(q_0(x),\gamma(h(x)))=q^+(x).
\]
\item 按已给定的决定表求值 \(\bar d(q^+(x))\)。
因 \(d=\bar d q^+\)，程序输出为 \(d(x)\)。
\item 若取得记录耗时至多 \(t_e\)，
每次比较耗时至多 \(t_c\)，
决定表读取耗时至多 \(t_d\)，
则运行阶段总时延至多
\[
t_e+Nt_c+t_d.
\]
费用和存储分别按同类有限操作计算。
预先建立字典的成本另计入准备阶段，
并不由表示的存在自动免除。
\end{semanticsteps}
\end{proof}

\begin{remark}[表示、计算与正确取证]
定理~\ref{std:strict-decisions}~说明新增观察能够表达哪些决定；
引理~\ref{std:evidence-recovery}~说明证据能否恢复所需类型；
定理~\ref{std:finite-decision}~才给出带有限资源步骤的算法。
若实际记录不在已经核验的字典范围，
系统进入补充核验或范围外处理状态，
不以任意默认类型代替未知事实。
版本变化使表内条件失效时，须先更新证据与决定表。
\end{remark}

\section{全部动作与响应的合同级对应}

\begin{definition}[双向覆盖的服务实现]
\label{std:adapters}
设候选共同状态空间为 \(Y_0\)，
实现映射为 \(p_i:X_i\to Y_0\)。
对每个 \(x\in X_i\) 给出满射
\[
\alpha_{i,x}:A_i(x)\twoheadrightarrow A_0(p_ix),
\]
对每个合法动作给出满射
\[
\beta_{i,x,a}:\Xi_i(x,a)
\twoheadrightarrow\Xi_0(p_ix,\alpha_{i,x}a).
\]
要求
\[
p_i\delta_i(x,a,\xi)
=\delta_0(p_ix,\alpha_{i,x}a,\beta_{i,x,a}\xi).
\]
全部规定观察、目标以及初态范围具有相同对应，
实际累计费用、时间和权限也被保存。
动作若由多个内部步骤实现，
须附有限终止、完整响应、
中间权利保持及累计资源证书。
影响这些观察的内部选择不得作为无关细节删除。
\end{definition}

\begin{proposition}[双向覆盖保存鲁棒前驱]
\label{std:robust-predecessor}
在定义~\ref{std:adapters}~下，
对任意 \(W\subseteq Y_0\) 与 \(x\in X_i\)，
原实现存在一个动作使全部响应后继进入 \(p_i^{-1}(W)\)，
当且仅当共同系统存在一个动作使全部响应后继进入 \(W\)。
\end{proposition}
\begin{certificate}
每个共同动作的实际实现、
每个实际动作的共同标签、
响应表的双向覆盖以及后继交换式。
\end{certificate}
\begin{proof}
\begin{semanticsteps}
\item 假设实际动作 \(a\in A_i(x)\) 对全部实际响应成功。
令
\[
\bar a=\alpha_{i,x}(a)\in A_0(p_ix).
\]
\item 任取 \(\bar\xi\in\Xi_0(p_ix,\bar a)\)。
响应满射给出 \(\xi\in\Xi_i(x,a)\)
使 \(\beta_{i,x,a}(\xi)=\bar\xi\)。
实际成功假设给出
\[
p_i\delta_i(x,a,\xi)\in W.
\]
\item 用后继交换式替换左侧，
得到 \(\delta_0(p_ix,\bar a,\bar\xi)\in W\)。
因为 \(\bar\xi\) 任意，共同动作对全部响应成功。
\item 反向设共同动作 \(\bar a\) 对全部共同响应成功。
动作满射给出实际动作 \(a\) 满足
\(\alpha_{i,x}(a)=\bar a\)。
\item 任取实际响应 \(\xi\in\Xi_i(x,a)\)。
其像属于共同响应集，故
\[
\delta_0(p_ix,\bar a,\beta_{i,x,a}\xi)\in W.
\]
后继交换把该式改写为
\(\delta_i(x,a,\xi)\in p_i^{-1}(W)\)。
\item 两方向分别保留存在动作与全称响应。
这核对了命题~\ref{q:adapter-reachability}~的服务接口前件，
允许逐轮应用鲁棒可达递推。
\end{semanticsteps}
\end{proof}

\begin{remark}[共同标签系统的构造范围]
以下动态商在已经核验的合同级状态系统上构造。
该系统具有共同动作与完整响应标签，
每个合法标签在其状态上的后继是确定的。
原实现的不同方法名、内部协议和外部资源调用，
须先通过上述双向覆盖及实现证书接入。
没有这一接入证书时，不能将不同软件的原始调用图
视为已经拥有同一转移语义。
\end{remark}

\section{有限演化饱和与精确共同任务商}

\begin{definition}[服务任务的初始观察]
\label{std:initial-tests}
在不交并 \(X_1\sqcup X_2\) 上使用共同观察族 \(\mathcal P_0\)。
该族包括共同合同全部规定观察、
初态和目标指示、各动作与响应的可用性、
各标签的成本、转移金额与时延，以及未定义状态指示。
加入唯一的 \(\bot\)，
对未定义标签采用定义~\ref{q:totalization}~的总化。
真实拒绝、待核验与结果待确认状态仍位于各自 \(X_i\) 中。
令
\[
\Lambda=\mathsf A\times\mathsf Z,\qquad
\widehat X=X_1\sqcup X_2\sqcup\{\bot\}.
\]
\end{definition}

\begin{lemma}[有限接口表的饱和算法]
\label{std:finite-saturation}
若全部初始测试及总化后继表均可精确计算，
则可以有限计算保持所有有限后继观察的稳定分区。
其严格细化次数至多为
\[
|\widehat X|-b_0,
\]
其中 \(b_0\) 为初始观察分区的非空块数。
\end{lemma}
\begin{certificate}
状态清单、标签清单、全部测试及后继表、
逐轮分区、分裂见证与最终稳定检查。
记录相等性比较使用已登记的共同类型编码。
\end{certificate}
\begin{proof}
\begin{semanticsteps}
\item 为每个状态计算初始测试值元组 \(\chi_0(x)\)。
以元组相等定义初始分区，得到有限块数 \(b_0\)。
\item 已计算第 \(n\) 轮时，
按有限标签表为每个状态计算
\[
\chi_{n+1}(x)=
\bigl(\chi_n(x),
(\chi_n(\widehat\Phi_\ell x))_{\ell\in\Lambda}\bigr).
\]
每个坐标都由已给定有限表取得，
故这一轮是有限算法。
\item 原元组作为新元组的首坐标，
所以新分区只可能保持或细化旧分区。
一次严格细化至少增加一个非空块；
总块数不超过 \(|\widehat X|\)。
\item 当块内后继的块编号对每个标签相同时停止。
该条件正是定理~\ref{q:dynamic}~中的稳定性条件。
该定理给出全部有限后继观察保持及最粗性，
而本构造提供其有限计算与逐轮检查证书。
\end{semanticsteps}
\end{proof}

\begin{theorem}[服务实现的最粗动态共同商]
\label{std:common-task}
在引理~\ref{std:finite-saturation}~的条件下，定义
\[
x\sim y
\Longleftrightarrow
\forall w\in\Lambda^*\ \forall f\in\mathcal P_0,\
f(\widehat\Phi_wx)=f(\widehat\Phi_wy).
\]
令 \(q_*:\widehat X\to Q_*\) 为稳定商，
并记
\[
Y_i=q_*(X_i),\qquad
Y_{0i}=q_*(X_{0i}).
\]
若
\[
Y_1=Y_2=Y,\qquad Y_{01}=Y_{02},
\]
且共同像包含声明的非失败服务任务，
则两种合同级实现的最粗动态任务商同构于 \(Y\)。
该同构保存初态、全部合法动作与响应、
费用、期限、逐主体权利和接受目标。
\end{theorem}
\begin{certificate}
共同观察编码、动作与响应接口、
稳定分区、每个实现对 \(Y\) 的覆盖表、
初态像相等见证及非失败任务见证。
若仍有本地动作名称，附定义~\ref{std:adapters}~的实现证书。
\end{certificate}
\begin{proof}
\begin{semanticsteps}
\item 初始测试含全部合法域、目标、效用及经济观察，
故命题~\ref{q:stable-exact}~适用。
对 \(q_i=q_*|_{X_i}\) 得到精确合同任务接口。
\item 不交并上的等价关系由同一组有限词测试决定。
将这些测试限制到 \(X_i\)，
仍得到该实现同一观察语义下的最粗稳定关系。
\item 定理~\ref{q:common-image}~的像条件在此由
\(Y_1=Y_2\) 提供。
分别定义
\[
\psi_i:X_i/{\sim}\longrightarrow Y,\qquad
\psi_i([x])=q_i(x).
\]
同块性保证良定义，实际像定义保证满射，
而类的定义保证单射。
\item 由稳定后继交换式，
\[
\psi_i([\widehat\Phi_\ell x])
=\bar\Phi_\ell\psi_i([x]).
\]
初始观察在类上恒定，故各观察与合法域亦被保存。
\item 取
\[
\Psi=\psi_2^{-1}\psi_1.
\]
该双射保持上述观察与转移。
另用 \(Y_{01}=Y_{02}\) 核对允许初态，
得到包含初态范围的共同任务同构。
\item 最粗性来自同一观察族的稳定商，
不是由两种实现使用相同接口名称推得。
不同共同像或不同初态像均不能省略该项核验。
\end{semanticsteps}
\end{proof}

\begin{corollary}[有限履约路径与鲁棒选择的保持]
\label{std:task-consequences}
在定理~\ref{std:common-task}~下，
任意有限合同执行及其逐步权利观察在两种商系统间对应。
对同一目标及同一合同步数上界，
其鲁棒可达集也在共同商上对应。
实际完成仍须使用第~\ref{svc:chapter}~章的实施与调度条件。
\end{corollary}
\begin{certificate}
有限标签词、初态对应、逐步交换方块、
鲁棒可达轮次及实际策略选择。
宏观合同动作另附内部原子步数与成本上界。
\end{certificate}
\begin{proof}
\begin{semanticsteps}
\item 定理~\ref{std:common-task}~给出精确任务接口，
核对了定理~\ref{q:execution}~的前件。
因此对任意有限词 \(w\) 有
\[
q_i\widehat\Phi_w=\bar\Phi_wq_i.
\]
\item 每个逐主体权利观察通过 \(q_i\) 因子化，
故在该式两侧施加同一读出函数，
得到每步义务与到账记录的对应。
\item 动作与非空响应域相同、目标为原像，
核对了定理~\ref{q:reachability}~的前件。
按该定理取得
\[
W_n^{(i)}=q_i^{-1}(\bar W_n).
\]
\item 由共同像及初态像相同，商上的履约能力一致。
实际执行还需选择并实施这些动作，
其时间与资源约束分别由定理~\ref{svc:scheduled-completion}~检查。
\end{semanticsteps}
\end{proof}

\section{共同任务的自由图链与最粗观察核}

\begin{definition}[合同状态与事件的自由图链]
\label{std:graph}
令 \(E_i\) 为 \(X_i\) 的全部合法标签边，
\(E_Y\) 为共同商 \(Y\) 的全部合法标签边。
定义
\[
C^{(i)}_0=\mathbb Q[X_i],\quad
C^{(i)}_1=\mathbb Q[E_i],\quad
C^{(i)}_{n\ge2}=0,
\]
\[
D_0=\mathbb Q[Y],\qquad
D_1=\mathbb Q[E_Y],\qquad D_{n\ge2}=0.
\]
边界矩阵按逐列事件定义：
\[
B_ie_{(x,\ell)}
=e_{\Phi_\ell x}-e_x,\qquad
B_Ye_{(y,\ell)}
=e_{\bar\Phi_\ell y}-e_y.
\]
零次链表示形式状态组合，一次链表示形式事件组合。
实际申请状态仅通过对应基向量编码；
任意带负系数的链元不解释为可执行请求。
\end{definition}

\begin{theorem}[服务状态商的逐列链证书]
\label{std:graph-map}
在定理~\ref{std:common-task}~下，定义
\[
F_{i,0}e_x=e_{q_i(x)},\qquad
F_{i,1}e_{(x,\ell)}=e_{(q_i(x),\ell)}.
\]
则 \(F_i:C^{(i)}\to D\) 是逐次数满射链映射，并满足
\[
B_YF_{i,1}=F_{i,0}B_i.
\]
\end{theorem}
\begin{certificate}
顶点映射表、带完整响应的边映射表、
共同合法域以及两侧边界矩阵。
每列均保存起点、终点和共同后继交换式。
\end{certificate}
\begin{proof}
\begin{semanticsteps}
\item 取一次链基向量 \(e_{(x,\ell)}\)。
先求其原边界再映射，得到
\[
F_{i,0}B_ie_{(x,\ell)}
=e_{q_i(\Phi_\ell x)}-e_{q_i(x)}.
\]
\item 精确任务商的后继交换式给出
\[
q_i(\Phi_\ell x)=\bar\Phi_\ell q_i(x).
\]
把该等式替入上一式，得到
\[
F_{i,0}B_ie_{(x,\ell)}
=e_{\bar\Phi_\ell q_i(x)}-e_{q_i(x)}.
\]
\item 按商图边界的定义，
右侧等于
\[
B_Ye_{(q_i(x),\ell)}
=B_YF_{i,1}e_{(x,\ell)}.
\]
基上的等式经线性延拓得到矩阵交换式。
\item \(q_i\) 对 \(Y\) 满射，
所以每个 \(e_y\) 有某个 \(e_x\) 作为原像，
得到零次满射。
\item 任取共同合法边 \((y,\ell)\)，
选择 \(q_i(x)=y\)。
精确合法域保证 \((x,\ell)\) 也是合法边，
其基向量映到所选共同边。
因此一次映射也满射。
\item 高次数均为零，零次以下边界亦为零。
结合矩阵交换式，核对了
定理~\ref{ch:graph-realization}~的全部图链条件。
\end{semanticsteps}
\end{proof}

\begin{theorem}[保持完整服务观察的共同链商]
\label{std:chain-isomorphism}
以共同顶点及带标签边的全部指标坐标为线性观察，
即取 \(O_{i,n}=F_{i,n}\)。
则最大不可观察子复形为
\[
N_{i,n}
=\ker O_{i,n}\cap
d_n^{-1}(\ker O_{i,n-1})
=\ker F_{i,n}.
\]
两种实现的最粗观察链商因此同构，
其显式同构为
\[
\Phi=\widetilde F_2^{-1}\widetilde F_1,\qquad
\widetilde F_i:C^{(i)}/N_i\xrightarrow{\cong}D.
\]
因此逐次数有
\[
H_n(C^{(1)}/N_1)\cong H_n(D)\cong H_n(C^{(2)}/N_2).
\]
\end{theorem}
\begin{certificate}
全部 \(F_i,B_i,B_Y\) 矩阵、
观察核基、逐次数满射证明、
商类到共同坐标的映射及逆映射。
\end{certificate}
\begin{proof}
\begin{semanticsteps}
\item 取 \(c\in\ker F_{i,n}\)。
链交换式给出
\[
F_{i,n-1}d_nc=d_nF_{i,n}c=0.
\]
所以该 \(c\) 同时属于所列两个核条件。
\item 反向取 \(c\in N_{i,n}\)。
其第一项核条件已经给出
\(F_{i,n}c=0\)。
两方向合并得到 \(N_{i,n}=\ker F_{i,n}\)。
\item 定理~\ref{ch:maximal-kernel}~
说明这正是该完整观察族的最大不可观察子复形。
因此相应商具有最粗观察链商的普遍性质。
\item 定义
\[
\widetilde F_i(c+N_i)=F_i(c).
\]
若更换代表，差属于 \(\ker F_i\)，像不变；
所以此映射良定义。
核等式给出单射，逐次数满射给出满射。
\item 原 \(F_i\) 与边界交换，
故诱导映射及其逆均与商边界交换。
组合 \(\widetilde F_2^{-1}\widetilde F_1\)
即为所声明的链同构。
\item 链同构及其逆均与边界交换，故把每次的循环空间与边界空间分别双射到另一侧对应空间。
在 \(\ker d_n/\operatorname{im}d_{n+1}\) 上取商，得到所列逐次数同调同构。
\item 同构的两侧均保留共同指标观察。
若改用更少的金额或阶段观察，
须重新计算最大不可观察核；
不能继续无条件沿用本定理的完整任务读回结论。
\end{semanticsteps}
\end{proof}

\begin{proposition}[实际状态的任务谓词读回]
\label{std:readback}
令 \(\iota_i(x)=e_x\)、
\(\bar\iota(y)=e_y\)，
并在实际商状态编码上定义 \(\sigma(e_y)=y\)。
每个通过 \(q_i\) 因子化的任务谓词
\(f_i=\bar f q_i\) 均可通过共同图链的零次编码读回：
\[
f_i=\bar f\,\sigma\,F_{i,0}\iota_i.
\]
该结论只在实际状态编码的像上解释为资格真值。
\end{proposition}
\begin{certificate}
状态基编码、两个接口等式、
逐谓词的共同读出与动作交换证书。
\end{certificate}
\begin{proof}
\begin{semanticsteps}
\item 对任意 \(x\in X_i\)，基上的映射定义给出
\[
F_{i,0}\iota_i(x)
=F_{i,0}e_x
=e_{q_i(x)}
=\bar\iota(q_i(x)).
\]
\item 读回函数在该像上满足
\[
\sigma\bar\iota(q_i(x))=q_i(x).
\]
因此 \(\sigma F_{i,0}\iota_i=q_i\)。
\item 将 \(f_i=\bar f q_i\) 中的 \(q_i\) 替换为上一式，
得到所列读回公式。
这些正是定义~\ref{ch:interface}~
与命题~\ref{ch:readback}~的接口条件。
\item 任意有理线性组合未必属于
\(\bar\iota(Y)\)。
在其上计算某个线性观察，
不构成某个实际请求的布尔资格判断。
故公式的资格解释保留在已经声明的状态编码域内。
\end{semanticsteps}
\end{proof}

\begin{remark}[逐主体权利与链压缩]
若主体 \(i\) 与主体 \(j\) 的退款均属任务，
则必须分别保留其观察。
仅保存总退款金额会消去转移方向
\(e_i-e_j\)，但该方向改变两人的独立请求权。
定理~\ref{ch:individual}~的核判据因此排除这种合并。
同调群维数相同也不能替代
定理~\ref{std:common-task}~的全部后继条件。
\end{remark}

\section{生命周期阶段矩阵的解释范围}

\begin{definition}[七阶段流程投影]
\label{std:phase}
以 \(R,A,E,S,D,F,U\) 依次表示登记、
获准与预留、执行回执、结算、
拒绝或失败、退款核对以及复核裁决。
取九条阶段边
\[
\begin{array}{lll}
e_1:R\to A,&e_2:A\to E,&e_3:E\to S,\\
e_4:R\to D,&e_5:A\to D,&e_6:D\to F,\\
e_7:F\to U,&e_8:S\to U,&e_9:U\to R.
\end{array}
\]
最后一条仅在裁决允许并建立新申请时具有实际实现。
没有资金预留的拒绝可以采用金额为零的退款核对；
理由、金额和裁决结果仍留在完整状态中。
该阶段链取 \(P_0=\mathbb Q^7\)、\(P_1=\mathbb Q^9\)、
\(P_{n\ge2}=0\)，边界为
\[
B_P=
\begin{pmatrix}
-1&0&0&-1&0&0&0&0&1\\
1&-1&0&0&-1&0&0&0&0\\
0&1&-1&0&0&0&0&0&0\\
0&0&1&0&0&0&0&-1&0\\
0&0&0&1&1&-1&0&0&0\\
0&0&0&0&0&1&-1&0&0\\
0&0&0&0&0&0&1&1&-1
\end{pmatrix}.
\]
\end{definition}

\begin{proposition}[阶段链的可计算同调证书]
\label{std:phase-homology}
上述阶段链满足
\[
H_0(P)\cong\mathbb Q,\qquad H_1(P)\cong\mathbb Q^3.
\]
一次循环的一组基为
\[
\begin{aligned}
z_1&=e_1+e_5-e_4,\\
z_2&=e_1+e_2+e_3+e_8-e_4-e_6-e_7,\\
z_3&=e_4+e_6+e_7+e_9.
\end{aligned}
\]
\end{proposition}
\begin{certificate}
边界矩阵、一个秩为六的生成树子矩阵、
三列循环向量及其零边界和线性独立证书。
\end{certificate}
\begin{proof}
\begin{semanticsteps}
\item 每列包含一个 \(1\) 和一个 \(-1\)，
故七行之和为零，
得到 \(\operatorname{rank}B_P\le6\)。
\item 取列 \(e_1,e_2,e_3,e_4,e_6,e_7\)，
并删除 \(R\) 行。
所得六阶矩阵的行列式为 \(1\)，
因此 \(\operatorname{rank}B_P\ge6\)。
\item 合并两项秩界得到秩为六。
零次边界为零，故
\[
\dim H_0(P)=7-6=1,\qquad
\dim\ker B_P=9-6=3.
\]
\item 对三个向量逐项应用边界。
每条路线的内部顶点相消，
两端亦相同，得到
\[
B_Pz_1=B_Pz_2=B_Pz_3=0.
\]
\item 三个向量在坐标 \(e_5,e_8,e_9\) 上组成单位矩阵，
所以线性独立。
它们是三维核的一组基。
因 \(P_2=0\)，没有来自二次链的边界，
故 \(H_1(P)=\ker B_P\)。
\end{semanticsteps}
\end{proof}

\begin{proposition}[阶段闭合不等于完整经济状态重复]
\label{std:phase-boundary}
若两个同阶段状态具有不同的必保留观察，
则阶段投影不满足完整服务任务的纤维条件。
上述非零一次同调也不蕴含实际付款、
退款或可无限接续的服务循环。
\end{proposition}
\begin{certificate}
保存两个处于同一阶段但具有不同主体权利或新旧请求标识的状态；
对 \(z_1,z_2\) 保存其负边系数，
对 \(z_3\) 保存重新申请时的新标识与推进时间。
\end{certificate}
\begin{proof}
\begin{semanticsteps}
\item 取相同阶段的两个状态 \(x,y\)，
使某个必保留观察 \(f\) 满足 \(f(x)\ne f(y)\)。
阶段投影 \(\pi\) 满足 \(\pi(x)=\pi(y)\)。
\item 引理~\ref{q:fiber}~因此排除
\(f=\bar f\pi\)。
故阶段投影不能作为该完整观察任务的充分商。
\item \(z_1,z_2\) 中的负系数表示形式路线之差。
实际有向执行不能以负次数实施某条边，
所以这两个核向量本身不是执行计划。
\item 对非负阶段路线 \(z_3\)，
定义~\ref{svc:reentry}~要求新请求 \(r'\ne r\)。
因此即使终点阶段重新成为 \(R\)，
完整请求标识观察也没有回到原值。
\item 为同一阶段图配置可用资金为零、应退金额为一的账本。
矩阵及其同调不变，但退款的资金守卫不成立。
这给出非零同调而不能实际退款的见证。
实际接续还须核对权限、资金、回执与时间，
因而阶段核等式不足以保证无限延续。
\end{semanticsteps}
\end{proof}

\section{资金效果矩阵与实际非负执行}

\begin{definition}[三账户的服务资金效果]
\label{std:money}
在固定货币和合同口径下，
令 \(b=(u,h,s)^\top\in\mathbb Q_{\ge0}^3\)
分别表示客户可用余额、托管余额及服务方到账。
预留、结算、退款的单位效果为矩阵三列：
\[
T=
\begin{pmatrix}
-1&0&1\\
1&-1&-1\\
0&1&0
\end{pmatrix}.
\]
对实际事件 \(e:x\to x'\)，
登记非负操作金额 \(a_e\in\mathbb Q_{\ge0}^3\)，并核对
\[
b(x')-b(x)=Ta_e.
\]
操作类型与金额进入事件观察；
若它们在同一候选商边的原像上不恒定，须先细化该边。
对不改变这三个余额的事件，取 \(a_e=0\)。
真实费用、退款依据及最终到账按合同另行核验；
总额守恒不替代逐主体义务。
\end{definition}

\begin{lemma}[资金效果的守恒与形式关系]
\label{std:money-kernel}
矩阵 \(T\) 满足
\[
(1,1,1)T=0,\qquad
\operatorname{rank}T=2,\qquad
\ker T=\mathbb Q(1,0,1)^\top.
\]
其核向量表示等额预留与退款的净余额效果相抵，
不表示相关事件已经实际发生。
\end{lemma}
\begin{certificate}
三列的逐坐标金额效果、
前两列的独立性及第三列等于第一列的相反数。
\end{certificate}
\begin{proof}
\begin{semanticsteps}
\item 分别相加每列的三个坐标，得到
\(-1+1+0=0\)、\(0-1+1=0\)、\(1-1+0=0\)。
所以 \((1,1,1)T=0\)。
\item 设前两列的线性组合为零。
第一行给出第一系数为零，
第三行给出第二系数为零，
故这两列独立。
\item 第三列等于第一列的相反数，
因此秩既至少为二又至多为二。
\item 对 \(a=(a_1,a_2,a_3)^\top\)，方程 \(Ta=0\)
的第一行给出 \(a_3=a_1\)，
第三行给出 \(a_2=0\)。
第二行随之满足，
故核恰为所列一维空间。
\item 上述方程没有动作发生、权限或回执谓词，
因而它证明净效果关系，
不单独签发实际履约事实。
\end{semanticsteps}
\end{proof}

\begin{corollary}[资金效果复形的同调]
\label{std:money-homology}
以操作金额、账户余额和资金总额分别作为二次、一次和零次链，构造
\[
M_2=\mathbb Q^3\xrightarrow{\ T\ }M_1=\mathbb Q^3
\xrightarrow{\ b\ }M_0=\mathbb Q,
\qquad b=(1\ 1\ 1),
\]
其余次数取零。则
\[
H_0(M)=0,\qquad H_1(M)=0,\qquad
H_2(M)=\mathbb Q[(1,0,1)^\top].
\]
\end{corollary}
\begin{certificate}
矩阵 \(T,b\)、前两列独立性、\(b\) 的满射见证、
核基 \((1,0,1)^\top\) 及全部链群的次数声明。
\end{certificate}
\begin{proof}
\begin{semanticsteps}
\item 引理~\ref{std:money-kernel} 给出 \(bT=0\)，
故相邻边界的复合为零，所列对象确为链复形。
\item 因 \(b(1,0,0)^\top=1\)，映射 \(b\) 满射。
所以 \(H_0(M)=M_0/\operatorname{im}b=0\)。
\item 秩零化度定理给出 \(\dim\ker b=3-1=2\)。
\(bT=0\) 使 \(\operatorname{im}T\subseteq\ker b\)，
而 \(\operatorname{rank}T=2\)，因此两空间相等。
\item 以同调定义取商，得到
\(H_1(M)=\ker b/\operatorname{im}T=0\)。
\item 因 \(M_3=0\)，二次同调等于 \(\ker T\)。
代入已计算的核基，得到所列一维空间。
\item 该结论陈述有理线性效果及其形式关系。
非负余额、实际权利和执行先后仍由下述执行定理逐项核验，
不从同调为零直接推出可履约或有利润。
\end{semanticsteps}
\end{proof}

\begin{theorem}[预留、结算与退款的非负执行见证]
\label{std:funded-settlement}
设初始余额为 \((u,h,s)\)，
新申请预留 \(p\)，合同合法费用为 \(c\)，并满足
\[
u\ge p\ge c\ge0,\qquad
h\ge L_{\rm other}\ge0,\qquad s\ge0.
\]
这里 \(L_{\rm other}\) 是原托管资金须覆盖的其他未付义务。
假设相应操作均获授权，
实际支付具有最终性回执，费用只结算一次，
且退款义务为 \(p-c\)。
则依次预留 \(p\)、结算 \(c\)、退款 \(p-c\)
构成余额逐前缀非负的执行，
终态为
\[
(u-c,h,s+c).
\]
原托管余额 \(h\) 在全部这些步骤中得到保留。
\end{theorem}
\begin{certificate}
客户可用余额、逐请求预留、
费用依据、退款债权、一次性结算标识、
三个支付原子的权限、时限及最终到账回执。
其他主体已有义务仍使用其原托管覆盖。
\end{certificate}
\begin{proof}
\begin{semanticsteps}
\item 预留原子的金额向量为
\[
a_1=(p,0,0)^\top.
\]
将 \(Ta_1\) 加到初态，取得
\[
b_1=(u-p,h+p,s)^\top.
\]
由 \(u\ge p\ge0\)，三个坐标均非负。
\item 结算原子的金额向量为
\(a_2=(0,c,0)^\top\)。
状态变化为
\[
b_2=b_1+Ta_2
=(u-p,h+p-c,s+c)^\top.
\]
由 \(p-c\ge0\)，有 \(h+p-c\ge h\ge0\)，
其余两个坐标也非负。
\item 退款原子的金额向量为
\(a_3=(0,0,p-c)^\top\)。
再次更新得到
\[
b_3=b_2+Ta_3
=(u-c,h,s+c)^\top.
\]
其中 \(u-c\ge u-p\ge0\)，
所以终态亦非负。
\item 每次动作的金额均非负，
且托管坐标在各前缀分别为
\(h+p,h+p-c,h\)，均不低于原有 \(h\)。
故本次请求不动用其他主体原有的这部分覆盖。
\item 预留后，托管资产为 \(h+p\)，
登记未付义务为 \(L_{\rm other}+c+(p-c)\)，
两者之差为 \(h-L_{\rm other}\ge0\)。
把结算与退款分别绑定其真实债权与最终到账，
逐次应用命题~\ref{atom:payment}~的一致支付字段检查。
合法费用支付 \(c\)，退款支付 \(p-c\)，
由相应回执确认两项已登记义务得到清偿。
\item 按定理~\ref{atom:telescope}~
累计三个相接的实际增量，得到
\[
b_3-b_0=T(a_1+a_2+a_3)=(-c,0,c)^\top.
\]
这把矩阵效果、非负前缀和实际履约证书连接起来。
\end{semanticsteps}
\end{proof}

\begin{proposition}[事件图与资金观察的交换接口]
\label{std:ledger-interface}
将共同图链顶点上的资金观察线性延拓为
\[
\ell_0:D_0\longrightarrow\mathbb Q^3,
\qquad \ell_0e_y=b(y).
\]
若每个实际合同边均满足定义~\ref{std:money}~的更新，
则逐边有
\[
\ell_0B_Ye=Ta_e.
\]
对任意实际相接的有限事件序列，
资金变化等于这些逐边效果的和。
\end{proposition}
\begin{certificate}
完整状态中的余额读出、
每条事件的起终点、非负金额与实际回执，
以及相接序列中的一致主体、合同及计量单位。
\end{certificate}
\begin{proof}
\begin{semanticsteps}
\item 取实际合同边 \(e:y\to y'\)。
其图边界具有类型
\[
B_Ye=e_{y'}-e_y\in D_0.
\]
\item 对边界施加线性观察，得到
\[
\ell_0B_Ye
=\ell_0e_{y'}-\ell_0e_y
=b(y')-b(y).
\]
\item 该实际事件的资金证书给出
\(b(y')-b(y)=Ta_e\)。
以等式替换取得逐边交换式。
\item 对相接序列 \(y_0\to\cdots\to y_n\)，
各步使用同一类型余额观察。
定理~\ref{atom:telescope}~因此适用，得到
\[
b(y_n)-b(y_0)=\sum_{j=0}^{n-1}Ta_{e_j}.
\]
\item 此处 \(D_0\) 的基是状态，
\(D_1\) 的基是事件；
余额是对零次链的观察。
第~\ref{r:chapter}~章则以账户和修复操作承担不同次数。
两种次数与对象不能互换解释，
其同调也不能不经接口便认定为同一种经济障碍。
\end{semanticsteps}
\end{proof}

\section{公开平台规范的接口例证}

作为标准化接口的公开例证，
Jakarta EE Platform 11 定义应用的标准承载平台\cite{JakartaPlatform11}。
其中身份、安全与授权规范可以为主体及权限记录提供实现接口，
服务接口规范可以承载请求与响应
\cite{JakartaSecurity4,JakartaAuthorization3,JakartaREST4}。
事务、并发及连接规范可以为资源操作与外部调用提供实现安排
\cite{JakartaTransactions2,JakartaConcurrency31,JakartaConnectors21}。
以下对应是本章合同证书的实现分工，
不是对任何具体产品已经满足全部数学前件的确认。

\begingroup
\small
\renewcommand{\arraystretch}{1.14}
\begin{longtable}{p{0.20\textwidth}p{0.29\textwidth}p{0.41\textwidth}}
\toprule
合同义务 & 规范接口例证 & 仍须提交的任务证书\\
\midrule
\endfirsthead
\toprule
合同义务 & 规范接口例证 & 仍须提交的任务证书\\
\midrule
\endhead
主体与授权 &
Security、Authorization &
真实主体及委托对应、用途事实、作用范围、当前规则与全路径检查。\\
申请与响应 &
RESTful Web Services &
完整动作与响应编码；拒绝、待核验和救济状态不得遗漏。\\
资源一致提交 &
Transactions &
真实参与资源、隔离条件、逐前缀覆盖、操作键与外部作用的符合关系。\\
异步工作 &
Concurrency &
任务实际主体、作用权限、重试一致性，以及有界调度与期限证明。\\
外部作用 &
Connectors &
外部资源的实际响应、重复作用控制、回执最终性及无遗漏路径。\\
任务证明与结果 &
独立核验接口 &
规范适用性、程序与输入绑定、核验器健全性及实际结果证书。\\
\bottomrule
\end{longtable}
\endgroup

\begin{remark}[规范对应的结论范围]
使用某项规范接口，不单独蕴含任务证明通过，
也不单独建立实际用途真实性、最终付款或按期完成。
共同任务与共同链商的结论，
以定理~\ref{std:common-task}~
及定理~\ref{std:chain-isomorphism}~的证书条件为准。
当两种实现满足这些条件时，
其声明合同具有相同可观察行为；
不同实现的全部内部状态、所有用途或所有版本并未因此被宣告等价。
\end{remark}

\chapter{收益侵占的最小可执行修复}
\label{r:chapter}

资格核验确定谁可以取得收益，资金覆盖确定承诺能否履行，
修复机制则处理已经发生的错误归属。
本章在固定逐人债权和同一核算单位下，
将合法支付与有权追回写为守恒操作。
以此证明原制度的具体修复障碍、所需新增方向及其执行次序。
这种构造接续参与分配与有限修复的形式接口
\cite{InternalEconomics}，
同调及相对同调采用标准定义\cite{Hatcher2002}。
\section{逐人账本与资源相容的操作}
\begin{definition}[逐人收益修复账本]
\label{r:ledger}
固定 \(m\ge1\) 名已经核定权利的主体。
一次链群及零次链群为
\[
C_1=\mathbb Q^{m+2},\qquad C_0=\mathbb Q,
\]
账户按
\[
(H_1,\ldots,H_m,O,X)
\]
排列，分别表示逐人权利账户、组织托管账户和无资格所得账户。
边界 \(B:C_1\to C_0\) 为求和映射：
\[
B=(1,\ldots,1).
\]
实际余额取非负坐标；链元允许有符号形式效果。
所有金额使用同一计价单位和结算窗口，
费用已计入另行预留的费用资金，
下列转账数量均指这些账户的净收支。
\end{definition}
\begin{definition}[支付与追回方向]
\label{r:directions}
以 \(e_{H_i},e_O,e_X\) 表示账户标准基，定义
\[
p_i=e_{H_i}-e_O,\qquad r=e_O-e_X.
\]
原支付操作矩阵与扩展操作矩阵为
\[
A_-=(p_1,\ldots,p_m)
=\begin{pmatrix}
I_m\\ -\mathbf{1}_m^\top\\ 0_{1\times m}
\end{pmatrix},
\]
\[
A_+=(p_1,\ldots,p_m,r)
=\begin{pmatrix}
I_m&0_{m\times1}\\
-\mathbf{1}_m^\top&1\\
0_{1\times m}&-1
\end{pmatrix}.
\]
其中 \(\mathbf{1}_m\) 是 \(m\) 维全一列向量。
令
\[
C^-:\quad\mathbb Q^m\xrightarrow{A_-}
\mathbb Q^{m+2}\xrightarrow{B}\mathbb Q,
\]
\[
C^+:\quad\mathbb Q^{m+1}\xrightarrow{A_+}
\mathbb Q^{m+2}\xrightarrow{B}\mathbb Q.
\]
两复形在其他次数均为零。
新增 \(r\) 表示已经具备合法权限和现实资产条件的追回操作。
\end{definition}
\begin{lemma}[守恒复形与有限残差保持]
\label{r:balance}
有 \(BA_-=BA_+=0\)。
对任一矩阵 \(A\) 满足 \(BA=0\) 的情形，
若 \(B\omega_0=0\)，并按有限操作更新
\[
\omega_{n+1}=\omega_n-Au_n,
\]
则所有有限 \(n\) 均满足 \(B\omega_n=0\)。
\end{lemma}
\begin{certificate}
逐列保存 \(Bp_i=0\)、\(Br=0\)，
并保存初始缺陷与每次实际效果记录。
\end{certificate}
\begin{proof}
\begin{semanticsteps}
\item 每个支付或追回向量恰有一个正一和一个负一，
所以列和为零，矩阵乘积为零。
\item 基例已由 \(B\omega_0=0\) 给定；
若第 \(n\) 步成立，则
\[
B\omega_{n+1}=B\omega_n-BAu_n=0.
\]
对执行次数归纳即得结论。
\end{semanticsteps}
\end{proof}
\begin{definition}[形式修复与实际修复]
\label{r:execution-domain}
对同口径的参照配置和实际配置，
令 \(\omega=x^{\rm ref}-x^{\rm act}\)，并要求 \(B\omega=0\)。
形式修复是有理向量 \(u\) 满足 \(Au=\omega\)。
实际修复还须有有限操作序列，
使每步权限有效、余额非负、费用已覆盖、期限满足，
并取得具有最终性的执行回执。
可实际执行的系数与顺序组成集合 \(\mathcal U(x^{\rm act})\)。
\end{definition}
\begin{theorem}[有限线性修复的原始与对偶证书]
\label{r:dual}
对有限有理矩阵 \(A\) 与向量 \(\omega\)，
以下两类证书恰有一类存在：
\[
Au=\omega;
\qquad\text{或}\qquad
\lambda A=0,\quad\lambda\omega\ne0.
\]
在 \(B\omega=0\)、\(BA=0\) 的复形中，
第一类等价于 \([\omega]=0\in\ker B/\operatorname{im}A\)；
第二类证明该同调类非零。
\end{theorem}
\begin{certificate}
保存增广矩阵 \([A\mid\omega]\) 的精确有理消元记录。
可解时提交回代向量；
不可解时提交矛盾行对应的线性组合 \(\lambda\)。
实际修复另提交 \(u\) 的合法顺序实现。
\end{certificate}
\begin{proof}
\begin{semanticsteps}
\item 若消元后没有矛盾行，回代给出解。
\item 若出现矛盾行，其原行线性组合对 \(A\) 为零，
对 \(\omega\) 非零，即为所需 \(\lambda\)。
\item 两类同时成立会导致
\(\lambda\omega=\lambda Au=0\)，矛盾。
\item 有限维消元覆盖全部可能，
而同调结论直接来自闭链模去边界的定义。
\end{semanticsteps}
\end{proof}
\begin{corollary}[对偶读出给出的不可消残差]
\label{r:residual}
若 \(\lambda A=0\)、\(\lambda\omega\ne0\)，
则对每个形式或实际候选 \(u\)，
\[
\|\omega-Au\|_2
\ge\frac{|\lambda\omega|}{\|\lambda\|_2}>0.
\]
\end{corollary}
\begin{certificate}
保存非零对偶读出及固定金额单位下的范数。
\end{certificate}
\begin{proof}
\begin{semanticsteps}
\item \(\lambda(\omega-Au)=\lambda\omega\)。
\item 应用 Cauchy--Schwarz 不等式并除以
\(\|\lambda\|_2>0\)，即得下界。
\end{semanticsteps}
\end{proof}
\section{逐人侵占缺陷的非零类与实际归零}
\begin{theorem}[逐人账本的修复维数]
\label{r:ranks}
有
\[
\operatorname{rank}B=1,\qquad
\operatorname{rank}A_-=m,\qquad
\operatorname{rank}A_+=m+1.
\]
因此
\[
H_1(C^-)\cong\mathbb Q,\qquad H_1(C^+)=0,
\]
并且两复形的 \(H_0,H_2\) 均为零。
\end{theorem}
\begin{certificate}
保存 \(A_-\) 前 \(m\) 行的单位子矩阵，
以及 \(A_+\) 的前 \(m\) 行与最后一行组成的可逆子矩阵。
\end{certificate}
\begin{proof}
\begin{semanticsteps}
\item \(B\) 非零，故秩为一，\(\dim\ker B=m+1\)。
\item \(A_-\) 的前 \(m\) 行是单位矩阵，故其列独立。
\item 若 \(\sum_i u_ip_i+vr=0\)，
逐个 \(H_i\) 坐标给出 \(u_i=0\)，
再由 \(X\) 坐标给出 \(v=0\)，所以 \(A_+\) 也列独立。
\item 两者的像均在 \(\ker B\) 内，
故第一同调维数分别为一和零。
\item 两矩阵单射给出 \(H_2=0\)；
\(B\) 满射给出 \(H_0=0\)。
\item 这里 \(H_0=0\) 描述形式约束复形的正合性，
并不表示实际账户总资源为零。
\end{semanticsteps}
\end{proof}
\begin{definition}[固定逐人债权的侵占向量]
\label{r:claim}
令 \(d_i\in\mathbb Q_{\ge0}\) 为第 \(i\) 人尚未取得的已核定收益，
并令 \(\Delta=\sum_i d_i>0\)。
这些收益已被错误记入无资格所得账户，定义
\[
\omega=(d_1,\ldots,d_m,0,-\Delta)^\top.
\]
逐人债权及其来源在修复期间固定；
任何债权变更都另外经过规则规定的复核。
\end{definition}
\begin{theorem}[侵占类的非零性与新增边界]
\label{r:obstruction}
对定义~\ref{r:claim} 的缺陷，
\[
[\omega]_{C^-}=\Delta[r]\ne0,\qquad
[\omega]_{C^+}=0.
\]
具体证书为
\[
\lambda_X=(0,\ldots,0,1),\quad
\lambda_XA_-=0,\quad
\lambda_X\omega=-\Delta,
\]
以及
\[
A_+(d_1,\ldots,d_m,\Delta)^\top=\omega.
\]
\end{theorem}
\begin{certificate}
逐人保存 \(d_i\)、合计 \(\Delta\)、债权核验结果，
并计算上述对偶与修复等式。
\end{certificate}
\begin{proof}
\begin{semanticsteps}
\item 由于
\[
\omega=\sum_{i=1}^m d_ip_i+\Delta r,
\]
而每个 \(p_i\) 已是 \(C^-\) 的边界，
故 \([\omega]_{C^-}=\Delta[r]\)。
\item 对偶读出非零，由定理~\ref{r:dual} 得到原类非零。
\item 扩展后 \(r\) 也成为边界，
所列系数向量直接证明扩展类为零。
\end{semanticsteps}
\end{proof}
\begin{theorem}[先追回、再逐人付款的可执行修复]
\label{r:constructive}
设实际余额为
\[
x_0=(h_1,\ldots,h_m,o,x),\qquad
h_i,o,x\ge0,\quad x\ge\Delta.
\]
假定追回与各笔付款的权限、净额、费用和期限均已通过核验，
且执行回执确认这些净额实际发生。
则先执行 \(\Delta r\)，再依次执行 \(d_ip_i\)，
形成余额始终非负的有限路径，
其最终净效果为 \(\omega\)，每个人恰取得 \(d_i\)。
\end{theorem}
\begin{certificate}
保存起始可用余额、追回授权、逐人债权、
独立费用储备、实际执行顺序、时间戳及最终性回执。
同一资金和同一债权均以唯一标识防止重复使用。
\end{certificate}
\begin{proof}
\begin{semanticsteps}
\item 追回后余额为
\[
(h_1,\ldots,h_m,o+\Delta,x-\Delta).
\]
前 \(k\) 人付款完毕后，余额为
\[
\left(
h_1+d_1,\ldots,h_k+d_k,
h_{k+1},\ldots,h_m,\
o+\Delta-\sum_{i=1}^k d_i,\
x-\Delta
\right).
\]
基例 \(k=0\) 为追回后的余额。
\item 从 \(k\) 到 \(k+1\) 的支付，
可由托管余额覆盖，因为
\[
o+\Delta-\sum_{i=1}^k d_i
\ge d_{k+1}.
\]
执行后公式继续成立，故对 \(k\) 归纳，
全部中间余额非负。
\item \(k=m\) 时托管余额恢复为 \(o\)，
逐人增加 \(d_i\)，无资格所得减少 \(\Delta\)。
\item 最终净变化为 \(\omega\)。
\item 权限、期限和最终性由所声明执行证书保证。
\end{semanticsteps}
\end{proof}
\begin{remark}
追回、由独立储备补偿、以及发行新增资源分别记账。
如果只能补偿而不能追回，
仍须记录无资格所得尚未消除的事实。
外部注资、手续费或资源发行进入本窗口时，
应增加来源及费用账户，重新核验守恒矩阵。
本章没有把有符号线性系数直接视为实际可执行转账。
\end{remark}
\section{相对同调与最少新增修复方向}
\begin{theorem}[新增治理方向与原侵占类的连接同构]
\label{r:relative}
以恒等映射包含 \(C_1^-,C_0^-\)，
并将旧支付生成元包含到扩展操作群，
得到短正合列
\[
0\longrightarrow C^-\xrightarrow{i}C^+
\xrightarrow{\pi}\mathcal R\longrightarrow0.
\]
其中 \(\mathcal R=C^+/C^-\) 仅在二次非零：
\[
\mathcal R_2=\mathbb Q\bar e_r,\qquad
H_2(\mathcal R)\cong\mathbb Q.
\]
连接同态满足
\[
\boxed{
\delta:H_2(\mathcal R)\xrightarrow{\cong}H_1(C^-),
\qquad \delta[\bar e_r]=[r].
}
\]
特别地，\([\omega]=\Delta\delta[\bar e_r]\)。
\end{theorem}
\begin{certificate}
保存包含与商映射的逐次数矩阵，
将 \(\bar e_r\) 提升为 \(e_r\)，
并核对其边界 \(A_+e_r=r\in C_1^-\)。
\end{certificate}
\begin{proof}
\begin{semanticsteps}
\item 商复形在一次、零次为零，
二次只剩新增生成元，且无三次边界。
\item 短正合列的同调长正合列
\cite{Hatcher2002} 含片段
\[
H_2(C^+)\longrightarrow H_2(\mathcal R)
\xrightarrow{\delta}H_1(C^-)
\longrightarrow H_1(C^+).
\]
定理~\ref{r:ranks} 使两端均为零，故 \(\delta\) 同构。
\item 按连接映射的提升后取边界定义，
\(\delta[\bar e_r]=[r]\)。
\item 最后使用定理~\ref{r:obstruction} 的类等式。
\end{semanticsteps}
\end{proof}
\begin{theorem}[新增独立修复方向的维数下界]
\label{r:minimal-dimension}
令 \(U=\operatorname{im}A\subseteq\ker B\)，
目标缺陷空间为 \(\mathcal Z\subseteq\ker B\)。
若新增 \(k\) 个守恒方向组成矩阵 \(R\)，并要求
\[
\mathcal Z\subseteq U+\operatorname{im}R,
\]
则
\[
k\ge
\dim\frac{U+\mathcal Z}{U}.
\]
在纯线性层面，该下界可以达到。
对于本章的侵占缺陷族，下界为一，
新增 \(r\) 即达到下界；
定理~\ref{r:constructive} 另外给出其实际执行见证。
\end{theorem}
\begin{certificate}
保存原操作像空间、目标空间及商空间的基。
达到下界时，保存各商基的守恒提升。
若声称实际最小修复，再提交提升方向的权限、资金与顺序证书。
\end{certificate}
\begin{proof}
\begin{semanticsteps}
\item 将全部方向投影到 \(\ker B/U\)。
\item \(k\) 个新增列至多张成 \(k\) 维空间，
而目标商空间必须包含于其中，所以得到下界。
\item 选取 \((U+\mathcal Z)/U\) 的一组基，
再在 \(\mathcal Z\) 中选取代表作为新增列，
即可在线性层面达到该维数。
\item 本章所有侵占向量模去 \(U\) 后均为 \(\Delta[r]\)，
且 \([r]\ne0\)，所以目标商空间一维。
\item 新增 \(r\) 足够，并已有非负顺序执行证书。
\end{semanticsteps}
\end{proof}
\begin{remark}
本定理的最小性计量独立修复方向。
制度改动数量、执行成本、时间和福利影响使用各自的比较口径。
一个方向可以对应多项实际控制步骤；
相同线性方向的不同实现也可能具有不同费用与权限条件。
\end{remark}
\section{具体账本的最粗逐人观察模板}
\begin{theorem}[逐人角色映射与链交换]
\label{r:ledger-map}
设具体账本复形 \(\mathcal D\) 以账户为一次基，
以允许的单位转账为二次基：
\[
d_2[e:u\to v]=e_v-e_u,\qquad
d_1z=\sum_v z_v.
\]
令 \(\mathcal D_0=\mathbb Q\)，
并令 \(\mathcal D_n=0\) 对所有 \(n\notin\{0,1,2\}\) 成立。
角色映射 \(\rho(v)\in\{H_1,\ldots,H_m,O,X\}\)
只合并同一人的内部账户，不合并不同权利主体。
令
\[
F_1e_v=e_{\rho(v)},\qquad F_0=\operatorname{id}.
\]
对一笔转账，记其各 \(H_i\) 角色净变化为 \(s_i\)，
\(X\) 角色净变化为 \(s_X\)，定义
\[
F_2[e]=\sum_i s_i e_{p_i}-s_Xe_r.
\]
则 \(F:\mathcal D\to C^+\) 为链映射。
若各次数满射，则 \(\mathcal D/\ker F\cong C^+\)。
\end{theorem}
\begin{certificate}
保存账户所有者、角色、金额单位及每笔转账的实际净效果。
满射性通过逐个共同生成元的原账本实现验证。
实际模拟另核验非负操作、先后次序和完整响应覆盖。
\end{certificate}
\begin{proof}
\begin{semanticsteps}
\item 资源守恒使托管角色的净变化为
\(-\sum_i s_i-s_X\)，故
\[
A_+F_2[e]
=(s_1,\ldots,s_m,-\sum_i s_i-s_X,s_X)^\top
=F_1d_2[e].
\]
账户求和与同一角色合并交换，
故 \(BF_1=F_0d_1\)。
\item 线性延拓得到全部链交换。
\item 核是子复形，逐次数第一同构定理得到所述链同构。
\end{semanticsteps}
\end{proof}
\begin{theorem}[固定逐人债权下净额账本的最粗链商]
\label{r:coarsest-ledger}
在定理~\ref{r:ledger-map} 的满射条件下，
固定逐人债权及其核验结果，
取净额观察
\[
O_1=F_1,\qquad O_0=\operatorname{id},\qquad O_2=0.
\]
则定理~\ref{ch:maximal-kernel} 的最大不可观察子复形满足
\[
N=\ker F.
\]
因此 \(C^+\) 是保留全部逐人净额、托管净额及无资格净额的最粗链商。
\end{theorem}
\begin{certificate}
保存 \(A_+\) 单射证书、逐人观察的完整坐标、
链交换式及各次数核等式。
如果任务还观察费用、期限、操作类型或债权变化，
将其加入观察族后重新核验最粗性。
\end{certificate}
\begin{proof}
\begin{semanticsteps}
\item 在二次，
\[
N_2=\ker(F_1d_2)
=\ker(A_+F_2)=\ker F_2,
\]
最后一步使用 \(A_+\) 单射。
\item 在一次，\(F_1c=0\) 已蕴含
\(d_1c=BF_1c=0\)，故 \(N_1=\ker F_1\)。
\item 零次观察为恒等，故 \(N_0=0=\ker F_0\)。
\item 所有其他次数均为零，得到 \(N=\ker F\)。
\item 结合定理~\ref{ch:maximal-kernel} 的普遍性质即得最粗性。
\item 每个 \(H_i\) 坐标单独保留，
因而定理~\ref{ch:individual} 所述逐人差异不能被商去。
\end{semanticsteps}
\end{proof}
\Needspace{16\baselineskip}
\begin{corollary}[共同逐人账本的链商对应]
\label{r:common-template}
若不同经济实现采用同一 \(m\)、相同计量语义和逐人权利对应，
并分别满足定理~\ref{r:coarsest-ledger}，则
\[
\mathcal D_1/N_1\cong\cdots\cong\mathcal D_s/N_s\cong C^+.
\]
完整经济任务的等价，另由定理~\ref{q:common-image} 的
共同像与动作响应条件保证。
\end{corollary}
\begin{certificate}
保存各实现向共同模板的满射链映射，
以及主体、债权和单位的对应表；
完整任务比较同时提交动态商证书。
\end{certificate}
\begin{proof}
\begin{semanticsteps}
\item 各观察链商均与同一 \(C^+\) 同构，
复合同构即可。
\item 金额链交换保存已声明净额观察；
完整动作、资格和期限的保持由相应动态任务证书分别给出。
\end{semanticsteps}
\end{proof}
\chapter{合法参与循环的同伦保持}
\label{cyc:chapter}

分配缺陷被修复以后，还需证明合法参与过程能够继续。
本章把阶段重复、经济资金条件与链同伦保持分别形式化，
然后在共同证书下合成。
阶段环表达一轮过程的结构，
逐人账本保存收益和修复，
完整状态任务则保留身份、历史、期限及实际进入条件
\cite{InternalTheory,InternalEconomics}。
\section{阶段循环与实际路径的接续}
\begin{definition}[四阶段过程复形]
\label{cyc:phase}
取四个阶段顶点 \(v_0,v_1,v_2,v_3\)，
依次表示准入准备、实际履约、核验结算和再进入准备。
边 \(e_i:v_i\to v_{i+1}\) 的下标模四。
令
\[
P_1=\mathbb Q^4,\qquad P_0=\mathbb Q^4,\qquad P_{n\ge2}=0,
\]
\[
d_1^P=
\begin{pmatrix}
-1&0&0&1\\
1&-1&0&0\\
0&1&-1&0\\
0&0&1&-1
\end{pmatrix}.
\]
实际执行的阶段投影可以形成闭环；
完整状态仍保存期次、余额、身份和义务，
不要求一轮之后返回同一个完整经济状态。
\end{definition}
\begin{theorem}[合法阶段环的一维第一同调]
\label{cyc:homology}
令 \(z=e_0+e_1+e_2+e_3\)，则
\[
d_Pz=0,\qquad H_1(P)=\mathbb Q[z].
\]
\end{theorem}
\begin{certificate}
保存有向边界矩阵和四条边的阶段解释。
实际循环另提交同一完整执行的合法路径提升。
失败、退出或资源不足的分支不作为成功循环的边。
\end{certificate}
\begin{proof}
\begin{semanticsteps}
\item 四列之和为零，所以 \(z\) 闭合。
\item 前三列线性无关，且全部列的和为零，
故矩阵秩为三、核为一维。
\item 由于 \(P_2=0\)，不存在二次边界，
第一同调恰由 \([z]\) 张成。
\end{semanticsteps}
\end{proof}
\begin{definition}[可连续提升的阶段协议]
\label{cyc:lifting}
对每个阶段 \(v_i\)，给定可作为提升的完整状态集合 \(\mathcal L_i\)。
要求从每个声明可达的 \(s\in\mathcal L_i\) 出发，
存在有限合法协议到达 \(\mathcal L_{i+1}\)，
其阶段投影为 \(e_i\)，并保持主体、规则和账本接续。
协议的资金、时间及中间状态条件均被核验。
若要求对外部扰动鲁棒，则协议须对全部允许响应
在给定有限步界内到达下一提升集合。
\end{definition}
\begin{theorem}[有限轮合法循环的接续归纳]
\label{cyc:finite-lifting}
设初态属于 \(\mathcal L_0\)，
且定义~\ref{cyc:lifting} 在所考察的全部执行前缀上成立。
则每个所考察的有限轮数 \(n\) 都存在可接续的合法执行，
其阶段投影经过 \(n\) 轮四阶段循环。
局部协议具有鲁棒条件时，结论对全部允许响应成立。
\end{theorem}
\begin{certificate}
保存每条阶段协议的起点、实际终点、主体和账本连续性，
以及有限步界和费用界。
后一协议必须从前一协议实际产生的状态开始。
\end{certificate}
\begin{proof}
\begin{semanticsteps}
\item 空前缀由初态给出。
\item 若前 \(k\) 条阶段边已经实现，
其实际终点按条件属于下一阶段的提升集合。
\item 从这个终点应用对应局部协议，
即可追加第 \(k+1\) 条边，且仍可继续。
\item 对 \(k=0,\ldots,4n-1\) 归纳得到 \(n\) 轮。
\item 鲁棒版本在每一个实际响应终点应用相同接续条件；
有限步界逐段相加得到整体有限步界。
\end{semanticsteps}
\end{proof}
\begin{remark}
单独为每条阶段边找到一条可行路径，
并不保证这些路径的终点和起点能够连接。
阶段复形是循环见证，完整经济任务的最粗充分商
仍由资格、金额、期限及其全部有限后继观察构造。
\end{remark}
\section{资金覆盖、实际支付与下一期接入}
\begin{corollary}[已融资履约与阶段循环的接续]
\label{cyc:funded-continuity}
\label{cyc:solvency}
\label{cyc:base-payment}
固定同一实际执行与有限期限。
若定理~\ref{part:thm-ledger-invariant}、
定理~\ref{part:thm-allocation}、
定理~\ref{part:thm-reentry}
及推论~\ref{part:cor-finite-continuity} 的全部前件
在该执行中共同成立，则已成立的逐主体逐事件权利保持，
非空有效参与集合中的合格者实际取得正基本份额，
选择继续参与的主体具有期限内可接续的合法路径：
\[
B_n\ge L_n\ge0,\qquad
L_{ik,n}+Y_{ik,n}=\Gamma_{ik,n},\qquad
y_{it}\ge\beta_t>0.
\]
再接入可以使用实际到账条件
\(y_{it}-\kappa_{it}\ge k_{i,t+1}\)，
并按上述定理同时核验时序和联合容量。
若阶段协议另满足定义~\ref{cyc:lifting}，
这些真实接续路径投影为所考察有限轮数的合法阶段循环。
\end{corollary}
\begin{certificate}
共同保存逐主体逐事件义务、实际资金前缀、
正确且完整的有效参与集合、最终性回执与下一期联合接入计划。
接续协议使用前一协议的真实末态，
资金、费用和权利不得跨模块重复计入。
\end{certificate}
\begin{proof}
\begin{semanticsteps}
\item 初始账本证书给出 \(B_0\ge L_0\ge0\) 和
\(L_{ik,0}+Y_{ik,0}=\Gamma_{ik,0}\)。
\item 每个结算原子同时减少现金与同一债权，且增加该债权的实际到账量；
新增义务原子满足资金覆盖。因此单步保持
\[
B_n-L_n\ge0,\qquad L_{ik,n}+Y_{ik,n}=\Gamma_{ik,n}.
\]
按执行次数归纳，上式在每个有限前缀成立。
\item 对同一有效参与集合，分配证书给出 \(b_{it}\ge\beta_t>0\)，
按期履约证书给出 \(y_{it}=b_{it}\)。等式替换得到 \(y_{it}\ge\beta_t\)。
\item 下一期资金证书给出 \(y_{it}-\kappa_{it}\ge k_{i,t+1}\)。
将其与资格、容量、时间及同一后继状态上的可达路径证书合取，得到实际再接入。
\item 阶段投影仅在实际路径提升证书成立时使用；
\(q(x_{n+1})=\bar\delta(q(x_n),\bar a_n,\bar\xi_n)\)
逐步保持，故完整实际执行映向声明的参与阶段。
\end{semanticsteps}
\end{proof}
\section{扩展修复复形的显式收缩}
\begin{lemma}[逐人修复复形的收缩算子]
\label{cyc:contraction}
对 \(C^+\) 定义升一次数的线性映射
\[
t_0(a)=ae_O,\qquad
t_1(h_1,\ldots,h_m,o,x)=(h_1,\ldots,h_m,-x),
\qquad t_2=0.
\]
则
\[
d_Ct+td_C=\operatorname{id}_{C^+},
\qquad t^2=0.
\]
\end{lemma}
\begin{certificate}
逐次数核验
\[
Bt_0=1,\qquad A_+t_1+t_0B=I_{m+2},
\qquad t_1A_+=I_{m+1}.
\]
\end{certificate}
\begin{proof}
\begin{semanticsteps}
\item 零次上 \(Bt_0(a)=a\)。
\item 一次向量上，
\[
A_+t_1(h_1,\ldots,h_m,o,x)
=(h_1,\ldots,h_m,-\sum_i h_i-x,x),
\]
而 \(t_0B\) 只在托管坐标增加 \(\sum_i h_i+o+x\)。
\item 相加恢复原向量。
\item 对二次向量 \((u_1,\ldots,u_m,v)\)，
先用 \(A_+\) 得到
\[
(u_1,\ldots,u_m,-\sum_i u_i+v,-v),
\]
再用 \(t_1\) 恢复原向量。
\item 三个次数的等式证明收缩恒等式。
\item 此外 \(t_1t_0=0\)，其余复合亦为零，故 \(t^2=0\)。
\end{semanticsteps}
\end{proof}
\begin{remark}
收缩是形式链关系的见证。
实际追回和支付仍使用定理~\ref{r:constructive} 的
非负余额路径及法律权限。
收缩算子不承担现实资金调度或直接执行的职责。
\end{remark}
\section{非零耦合下的合法循环保持}
\begin{definition}[阶段与修复的相容耦合]
\label{cyc:coupling}
令 \(E_n=P_n\oplus C_n^+\)。
给定次数为负一的映射
\[
\kappa_n:C_n^+\to P_{n-1},
\qquad d_P\kappa+\kappa d_C=0,
\]
定义总边界
\[
D(p,c)=(d_Pp+\kappa c,\ d_Cc).
\]
耦合记录修复步骤对阶段链表示的影响；
实际动作和经济观察由相容的状态实现另外保留。
\end{definition}
\begin{theorem}[带耦合的链同伦回缩]
\label{cyc:retraction}
定义
\[
j(p)=(p,0),\qquad
r_E(p,c)=p-\kappa tc,\qquad
T(p,c)=(0,tc).
\]
则 \(D^2=0\)，\(j,r_E\) 为链映射，并满足
\[
r_Ej=\operatorname{id}_P,\qquad
\operatorname{id}_E-jr_E=DT+TD.
\]
因此
\[
H_n(E)\cong H_n(P)
\quad\text{对全部次数 }n,
\qquad H_1(E)\cong\mathbb Q[z].
\]
\end{theorem}
\begin{certificate}
保存 \(d_P\kappa+\kappa d_C=0\)、收缩 \(t\)
及 \(j,r_E,T\) 的逐次数矩阵。
实际循环另保存定义~\ref{cyc:lifting} 的路径接续证书。
\end{certificate}
\begin{proof}
\begin{semanticsteps}
\item 直接平方给出
\[
D^2(p,c)=
\bigl(d_P^2p+(d_P\kappa+\kappa d_C)c,\ d_C^2c\bigr)=0.
\]
\(Dj=jd_P\) 立即成立。
\item 利用 \(d_Ct+td_C=\operatorname{id}\)，有
\[
\begin{aligned}
r_ED(p,c)
&=d_Pp+\kappa c-\kappa td_Cc\\
&=d_Pp+\kappa d_Ctc
=d_P(p-\kappa tc)
=d_Pr_E(p,c).
\end{aligned}
\]
故 \(r_E\) 也是链映射，且 \(r_Ej=\operatorname{id}_P\)。
\item 再计算
\[
(DT+TD)(p,c)
=(\kappa tc,\ (d_Ct+td_C)c)
=(\kappa tc,c)
=(\operatorname{id}_E-jr_E)(p,c).
\]
对闭链，上式右侧与恒等映射的差是边界；
故两链映射诱导的同调映射互逆。
\item 最后使用定理~\ref{cyc:homology} 得到第一同调结论。
\end{semanticsteps}
\end{proof}
\begin{proposition}[非零耦合实例]
\label{cyc:nonzero}
对二次向量 \((u_1,\ldots,u_m,v)\)，取
\[
\kappa_2(u_1,\ldots,u_m,v)=u_1z,\qquad
\kappa_1=\kappa_0=0.
\]
则耦合非零且满足定义~\ref{cyc:coupling} 的相容式，
因而定理~\ref{cyc:retraction} 适用。
\end{proposition}
\begin{certificate}
保存 \(m\ge1\)、\(d_Pz=0\)
及其他耦合分量为零的逐次记录。
\end{certificate}
\begin{proof}
\begin{semanticsteps}
\item 唯一可能非零的相容项是 \(d_P\kappa_2\)，
它由 \(d_Pz=0\) 消失。
\item 取 \(u_1=1\)、其余坐标为零即可见耦合非零。
\end{semanticsteps}
\end{proof}
\begin{proposition}[合法循环保持与旧障碍消除]
\label{cyc:separate}
扩展后的 \(E\) 与阶段复形 \(P\) 链同伦等价。
原修复复形 \(C^-\) 与扩展复形 \(C^+\) 不链同伦等价。
在零耦合实例中，
\[
H_1(P\oplus C^-)\cong\mathbb Q\oplus\mathbb Q,\qquad
H_1(P\oplus C^+)\cong\mathbb Q.
\]
被保留的是合法参与循环类，
被消除的是已经明确指定的收益侵占类。
\end{proposition}
\begin{certificate}
保存链同伦回缩恒等式与前后同调维数。
零耦合时另保存总边界的分块对角结构。
\end{certificate}
\begin{proof}
\begin{semanticsteps}
\item 第一项由定理~\ref{cyc:retraction} 给出。
\item 链同伦等价必诱导同调同构，
而 \(H_1(C^-)\cong\mathbb Q\)、\(H_1(C^+)=0\)，
故二者不链同伦等价。
\item 零耦合下核与像逐块分解，
同调与有限直和交换，得到两式。
\end{semanticsteps}
\end{proof}
\section{经济观察、修复与再参与的联合结论}
\begin{theorem}[逐人权益修复与有限参与循环的联合保持]
\label{cyc:joint}
固定完整经济任务及有限执行范围。
假定：真实资格有有效证据；
动态任务商保持所有规定观察和动作响应；
状态与账本链接口相容；
逐人债权、实际资金、追回权限和支付条件满足
定理~\ref{r:constructive}；
阶段耦合满足定义~\ref{cyc:coupling}；
所考察的每一期满足资金覆盖、接入和路径接续条件。
则每个所考察的有限执行前缀同时具有：
逐人资格与权利可读回、指定侵占缺陷实际修复、
合法阶段循环类保持，
以及资金和入口允许范围内的再次参与路径。
\end{theorem}
\begin{certificate}
联合证书包含动态分区及双向动作响应覆盖、
逐人债权与独立资格证据、
最大不可观察核及状态链接口、
原始与对偶矩阵记录、非负执行回执、
同伦回缩矩阵、资金不变量及连续路径提升。
各证书使用同一规则版本、主体标识、金额单位和执行期间。
\end{certificate}
\begin{proof}
\begin{semanticsteps}
\item 精确任务表示对每个必要权利观察给出 \(f=\bar f q\)，
状态与账本表示的交换证书把同一主体的债权及到账事实读回。
\item 合法修复证书给出 \(A_+u=\omega\)，且 \(u\) 属于有权限、资源和非负前缀的执行集合。
真实回执确认同一缺陷的实际纠正，因此 \([\omega]_{C^+}=0\) 与实际修复同时成立。
\item 收缩与耦合相容证书给出
\(\operatorname{id}_E-jr=DT+TD\) 及 \(rj=\operatorname{id}_P\)。
在同调上链同伦诱导相同映射，故 \(H_*(E)\cong H_*(P)\)。
\item 对每期账本和有限接续分别应用保持归纳，得到
\(B_n\ge L_n\ge0\)、\(L_{ik,n}+Y_{ik,n}=\Gamma_{ik,n}\)
和实际到账下界。下一期成本、容量和合法路径证书给出实际接续。
\item 全部步骤使用共同主体、期次、资源及后继状态，因此可以合取；
任意有限执行前缀同时满足权利读回、指定修复、阶段同调保持与实际再接入。
\end{semanticsteps}
\end{proof}
因此，完整经济推导采用下列有条件的谓词链：
\[
\begin{aligned}
&\text{真实资格与履约}
\ \Longrightarrow\ \text{已成立逐人收益权}\\
&\text{收益权与资金覆盖、合法支付}
\ \Longrightarrow\ \text{实际取得正份额}\\
&\text{实际份额与必要费用条件}
\ \Longrightarrow\ \text{金融接入条件}\\
&\text{金融接入、下一期资格与接纳容量}
\ \Longrightarrow\ \text{可接续的合法参与路径}.
\end{aligned}
\]
最粗任务表示使必要区分得到保留，
可执行修复使错误归属得到纠正，
同伦回缩使合法阶段结构得到保持。
每个实际期间的资金、权限、证据及入口条件，
构成上述数学结论进入经济运行的共同前提。

\chapter{高保障任务的可达能力与生产可能集}
\label{front:chapter}

\section{能力扩张的任务含义}

限制某些常规用途的部署机会与扩张高保障任务能力，可以分别作用于不同的可行集合。
经济学分析需要同时登记被保留的技术任务、被收缩的商业机会以及实现两者所需的资源。
本章的能力扩张指在固定目标观察下出现真实的新可达任务；它不由程序规模、系统名称或账本操作数量定义。

\begin{definition}[旧能力、扩展能力与效果表示]
固定初态 \(x_0\)、目标观察 \(g:X\to Y\) 与共同约束域。
旧、新允许路径集合分别为 \(\mathcal P_-\)、\(\mathcal P_+\)，其可达状态集合为
\(\mathcal R_-(x_0)\)、\(\mathcal R_+(x_0)\)。
给定效果表示 \(\Delta(x_0,x)\in\QQ^n\)，守恒或相容矩阵 \(B\)，以及旧操作效果矩阵 \(A\)，满足
\[
BA=0,
\qquad
\forall\gamma\in\mathcal P_-,\quad
\Delta(x_0,\operatorname{end}\gamma)\in\im A.
\]
后一条件称为旧效果覆盖证书：它覆盖全部旧允许路径的累计效果。
\end{definition}

\begin{lemma}[旧能力的对偶排除证书]
\label{front:separation}
对指定状态 \(x_\star\)，令 \(\omega_\star=\Delta(x_0,x_\star)\)。若存在行向量 \(\lambda\) 使
\[
\lambda A=0,\qquad\lambda\omega_\star\ne0,
\]
则 \(x_\star\notin\mathcal R_-(x_0)\)。
\end{lemma}
\begin{certificate}
旧效果覆盖证书、目标效果 \(\omega_\star\)、行向量 \(\lambda\) 和两项精确矩阵乘积。
\end{certificate}
\begin{proof}
\begin{semanticsteps}
\item 若 \(x_\star\) 旧可达，则旧效果覆盖给出 \(\omega_\star=Au\)。
\item 于是
\[
x_\star\in\mathcal R_-(x_0)
\Longrightarrow\omega_\star\in\im A
\Longrightarrow\lambda\omega_\star=\lambda Au=0,
\]
与证书矛盾，故目标旧不可达。
\end{semanticsteps}
\end{proof}

\begin{remark}[全局覆盖的证明责任]
覆盖条件只要求旧路径的效果落在 \(\im A\)，不要求 \(\im A\) 中每个有符号组合均可执行。
若 \(A\) 仅来自某个状态附近的线性化，覆盖范围就是该局部域；不能据此排除域外或未登记的非线性路径。
有限线性系统的解与分离函数构成可核验双证书\cite{InternalEconomics}。
\end{remark}

\section{新增操作与严格可达扩张}

\begin{theorem}[共同约束下的严格可达扩张]
\label{front:strict}
在旧效果覆盖成立时，设新增能力矩阵 \(R\) 满足 \(BR=0\)，并有
\[
[A\ R]u_\star=\omega_\star,
\qquad
\lambda A=0,\quad\lambda\omega_\star\ne0.
\]
若 \(u_\star\) 有一条从 \(x_0\) 到 \(x_\star\) 的真实允许执行，且每条旧允许路径都有在新制度中保持相同初态、相同终态与全部约束的实现，
则在共同状态表示下
\[
\boxed{\mathcal R_-(x_0)\subsetneq\mathcal R_+(x_0).}
\]
\end{theorem}
\begin{certificate}
矩阵 \(A,B,R\)，旧覆盖与分离函数，新增解、非负资源、实际权限、执行顺序、时限、安全条件和终态回执，
以及旧路径的端点保持映射。新旧系统使用同一状态表示、资源边界和目标解释。
\end{certificate}
\begin{proof}
\begin{semanticsteps}
\item 旧路径保持给出 \(\mathcal R_-\subseteq\mathcal R_+\)。
\item 引理~\ref{front:separation} 给出
\(x_\star\notin\mathcal R_-\)；新增允许执行给出 \(x_\star\in\mathcal R_+\)。
\item 故差集非空，包含严格成立。
\item 谓词链为
\[
\begin{aligned}
&\operatorname{OldCovered}\land\operatorname{DualWitness}
\Longrightarrow\operatorname{OldUnreachable}(x_\star),\\
&\operatorname{PrimalWitness}\land\operatorname{Executable}
\Longrightarrow\operatorname{NewReachable}(x_\star),\\
&\operatorname{OldPreserved}\land\operatorname{OldUnreachable}(x_\star)
\land\operatorname{NewReachable}(x_\star)
\Longrightarrow\mathcal R_-\subsetneq\mathcal R_+.
\end{aligned}
\]
\end{semanticsteps}
\end{proof}

\begin{corollary}[严格任务能力扩张]
\label{front:taskimage}
若另有
\[
g(x_\star)\notin g(\mathcal R_-(x_0)),
\]
则可实现任务结果集合满足
\[
g(\mathcal R_-(x_0))\subsetneq g(\mathcal R_+(x_0)).
\]
\end{corollary}
\begin{certificate}
结果观察 \(g\)、旧结果的完整覆盖或排除证书，以及新结果的实际验收记录。
\end{certificate}
\begin{proof}
\begin{semanticsteps}
\item 旧路径保持给出 \(\mathcal R_-(x_0)\subseteq\mathcal R_+(x_0)\)。
\item 对任意 \(y\in g(\mathcal R_-(x_0))\)，存在旧可达 \(x\) 使 \(g(x)=y\)。同一个 \(x\) 也新可达，故 \(y\in g(\mathcal R_+(x_0))\)。
\item 新执行证书给出 \(x_\star\in\mathcal R_+(x_0)\)，故 \(g(x_\star)\) 属于新结果像。
\item 额外排除证书给出 \(g(x_\star)\notin g(\mathcal R_-(x_0))\)。包含与差集见证合取，得到严格包含。
\end{semanticsteps}
\end{proof}

\begin{remark}[生产能力与实施记录]
新的内部状态可以只改变记账方式或执行记录，未必提供新的生产结果。
因此，只有在 \(g\) 保留指定的科学、能源、基础设施或其他高保障任务验收量时，上一推论才解释为该任务能力边界扩张。
资金修复的同调计算本身不证明新的物理产出已经实现。
\end{remark}

\section{新增修复方向的数量与实际成本}

\begin{proposition}[线性修复方向的维数下界]
\label{front:dimension}
令旧效果空间为 \(U\subseteq\ker B\)，待修复效果张成空间为 \(Z\subseteq\ker B\)。
若新增 \(k\) 个独立候选方向的张成空间为 \(V\)，并使 \(Z\subseteq U+V\)，则
\[
k\ge\dim\bigl((U+Z)/U\bigr).
\]
允许选取任意相容有符号方向时，该下界可以达到。
\end{proposition}
\begin{certificate}
\(U\)、\(Z\) 的有理基，商空间基与对应提升，以及新增方向的实际候选清单。
\end{certificate}
\begin{proof}
\begin{semanticsteps}
\item 在 \(\ker B/U\) 中取像，\((U+Z)/U\) 包含于 \((U+V)/U\)。
\item 后者由至多 \(k\) 个像生成，故维数不超过 \(k\)。
\item 反向，取 \((U+Z)/U\) 的一个基并在 \(Z\) 中逐一提升；这些提升与 \(U\) 一起覆盖 \(Z\)，新增方向数恰为该维数。
\end{semanticsteps}
\end{proof}

\begin{remark}[不同最小性]
维数下界描述缺失的独立线性能力数量。某个方向可能难以合法取得，可能需要高额设备成本或不可满足的执行时间。
因此，方向最小性、规则变更最小性和总成本最小性分别核验，不能用一个秩计算代替全部制度比较。
\end{remark}

\section{常规机会收缩与高保障任务扩张的相容性}

\begin{proposition}[固定收益函数下的机会集合比较]
\label{front:profit}
设常规部署机会 \(\Omega_C^\pi\subseteq\Omega_C^0\)，并使用同一收益函数 \(\Pi\)，则在扩展实数意义下
\[
\sup_{a\in\Omega_C^\pi}\Pi(a)
\le\sup_{a\in\Omega_C^0}\Pi(a).
\]
同时，若定理~\ref{front:strict} 成立，高保障任务的可达集合可以严格扩张。
\end{proposition}
\begin{certificate}
共同口径下的机会清单和收益函数，以及高保障任务的严格扩张证书。空机会集合按 \(\sup\varnothing=-\infty\) 约定。
\end{certificate}
\begin{proof}
\begin{semanticsteps}
\item 取任意 \(a\in\Omega_C^\pi\)。集合包含给出 \(a\in\Omega_C^0\)。
\item 上确界定义给出 \(\Pi(a)\le\sup_{b\in\Omega_C^0}\Pi(b)\)。
\item 右侧是全部受限候选收益的共同上界，因此
\(\sup_{a\in\Omega_C^\pi}\Pi(a)\le\sup_{a\in\Omega_C^0}\Pi(a)\)。
\item 空集合采用扩展实数约定 \(\sup\varnothing=-\infty\)，结论仍成立。
该比较固定收益函数和机会域，不推导其他新增任务的市场收入。
\end{semanticsteps}
\end{proof}

\begin{remark}[经济解释的边界]
该命题描述固定收益函数下的受限机会价值，不直接预测资产价格、资本流向或长期创新速度。
后者还取决于需求、竞争、融资、预期和制度执行。对新增高保障用途的研究收益与维护成本，应结合第~\ref{integrated:market}~章的完整成本与第~\ref{opt:chapter}~章的共同口径比较。
\end{remark}

\chapter{标准复用、服务需求与完整成本回收}
\label{integrated:market}
\label{market:chapter}

\section{标准接口与可复用成果的经济对象}

命题~\ref{front:profit} 比较固定收益函数下的常规部署机会。
本章进一步研究另一个问题：在保留用途、责任和真实参与约束时，
可验证成果能否通过服务合同形成一个有实际支付的非空市场。
服务可以包含经过核验的程序执行、证明检查、接口适配、
版本维护和约定期限内的故障处理。
单次技术展示、名义订单和未来融资预期均不直接计为已实现的销售收入。

标准兼容性可以减少重复开发，也可能产生认证、迁移和共同维护成本。
关于兼容性、网络效应与系统竞争的经济分析，
需要同时考虑这些收益与供应商的竞争选择\cite{KatzShapiro1994}。
以下结论不把兼容性直接等同于成本下降、正价格或正利润。
任务正确性与证明组合采用已经登记的任务观察、
程序语义和可核验接口\cite{InternalTheory}；
完整成本、真实履约和制度兼容性采用共同经济核算边界
\cite{InternalEconomics}。

\begin{definition}[同等任务覆盖下的复用比较]
\label{market:reuse-data}
固定 \(n\in\NN\)、\(n\ge1\) 个部署项目，
以及有限模块集合 \(L\)。
对每个 \(\ell\in L\)，
令 \(f_\ell\ge0\) 为一次完成指定模块构造与证明的真实成本。
在独立重复构造的对照中，每个项目都承担全部 \(f_\ell\)。
在复用方案中，共同成果只构造一次，
项目 \(i\) 承担适配成本 \(a_i\ge0\)，
共同标准和维护安排另承担 \(F_0\ge0\)。
两方案比较同一期间内相同的任务覆盖、
可靠性、时限、责任、更新范围和资源机会成本。
下式登记发生差异的可比成本；
两方案逐项相同的其余成本已经同额抵销，
进行完整利润核算时仍须加回。
\[
 C_{\rm repeat}=n\sum_{\ell\in L}f_\ell,
 \qquad
 C_{\rm reuse}=\sum_{\ell\in L}f_\ell+\sum_{i=1}^n a_i+F_0.
\]
适配成本包括环境前提核验、语义转换、版本绑定、
接口组合证明及必要运行验收。
未满足这些条件的接口连接，不计为相同任务的成功复用。
\end{definition}

\begin{definition}[有限顺序组合的证明接口]
\label{market:proof-interface}
对项目 \(i\)，登记实际执行轨迹
\(s_0^i,s_1^i,\ldots,s_m^i\) 和环境条件 \(E_i\)。
第 \(\ell\) 个模块具有前提 \(A_\ell\)、后件 \(G_\ell\)
及经过版本绑定的实际转移关系 \(T_\ell\)，并提交
\[
 \forall x,y,\quad
 A_\ell(x)\land T_\ell(x,y)\Longrightarrow G_\ell(y).
\]
项目接口另外提交
\[
 E_i\land\bigwedge_{r<\ell}G_r(s_r^i)
 \Longrightarrow A_\ell(s_{\ell-1}^i),
 \qquad 1\le\ell\le m,
\]
其中 \(\ell=1\) 时的空合取为真。
还须有实际模块执行
\(T_\ell(s_{\ell-1}^i,s_\ell^i)\)，以及任务读回
\[
 E_i\land\bigwedge_{\ell=1}^mG_\ell(s_\ell^i)
 \Longrightarrow \operatorname{Task}_i(s_0^i,\ldots,s_m^i).
\]
环境条件覆盖相应轨迹的实际输入、扰动和资源边界，
不以未经核验的测试环境代替部署环境。
\end{definition}

\begin{lemma}[满足前提的有限证明复用]
\label{market:composition}
若定义~\ref{market:proof-interface} 的环境、
模块执行和证明接口全部成立，
则实际组合轨迹满足 \(\operatorname{Task}_i\)。
已经证明的模块蕴涵可以复用，
但项目接口、实际执行和任务读回仍须单独核验。
\end{lemma}
\begin{certificate}
模块版本与证明的绑定表、验证器健全性、
实际环境证据、每个转移的实现回执、
接口蕴涵及最终任务读回。
涉及并发、循环或异常恢复时，
另提供覆盖相应行为的不变量和完成证书；
本引理的有限顺序模型不自动覆盖这些行为。
\end{certificate}
\begin{proof}
\begin{semanticsteps}
\item 由 \(E_i\) 和第一项接口蕴涵，
得到 \(A_1(s_0^i)\)。
实际执行给出 \(T_1(s_0^i,s_1^i)\)。
模块证明因此推出 \(G_1(s_1^i)\)。
\item 假设已经得到
\(\bigwedge_{r<\ell}G_r(s_r^i)\)。
环境条件与第 \(\ell\) 项接口蕴涵给出
\(A_\ell(s_{\ell-1}^i)\)。
\item 将这个前提与实际转移
\(T_\ell(s_{\ell-1}^i,s_\ell^i)\) 合取，
模块证明推出 \(G_\ell(s_\ell^i)\)。
这一步保留此前已经成立的后件。
\item 对有限序号 \(\ell=1,\ldots,m\) 归纳，
得到全部后件。
最后应用任务读回蕴涵，得到
\(\operatorname{Task}_i(s_0^i,\ldots,s_m^i)\)。
\end{semanticsteps}
证明分别使用模块语义、实际转移与环境接口。
缺少任一项时，接口名称相同不足以完成上述谓词链。
\end{proof}

\begin{proposition}[复用节约的充要成本条件]
\label{market:reuse-saving}
在定义~\ref{market:reuse-data} 的共同口径下，
复用方案严格节约成本，当且仅当
\[
 \sum_{i=1}^n a_i+F_0
 <(n-1)\sum_{\ell\in L}f_\ell.
\]
等号对应成本相同，反向严格不等式对应复用成本更高。
\end{proposition}
\begin{certificate}
相同任务覆盖的两套实际实施方案、
通过引理~\ref{market:composition} 的复用接口、
完整成本清单、共同折算时点和无重复计量表。
\end{certificate}
\begin{proof}
\begin{semanticsteps}
\item 将两方案的成本定义相减：
\[
 C_{\rm repeat}-C_{\rm reuse}
 =(n-1)\sum_{\ell\in L}f_\ell-\sum_{i=1}^n a_i-F_0.
\]
\item 严格节约的定义为左边大于零。
将右边的适配成本和共同成本移至不等式另一侧，
恰得到命题中的严格不等式。
\item 移项可逆，故该条件也充分。
把大于零分别替换为等于零或小于零，
得到其余两种比较。
\end{semanticsteps}
\end{proof}

\begin{corollary}[同额适配成本下的规模阈值]
\label{market:reuse-scale}
令 \(f=\sum_{\ell\in L}f_\ell\)，并设每项适配成本均为 \(a\)。
当 \(f>a\) 时，严格节约等价于
\[
 n>\frac{f+F_0}{f-a}.
\]
当 \(f\le a\) 时，任意有限 \(n\ge1\) 均不能在该模型中严格节约。
\end{corollary}
\begin{certificate}
各项目确有相同适配成本的证据，
以及非负的 \(f,a,F_0\)。
\end{certificate}
\begin{proof}
\begin{semanticsteps}
\item 上一命题给出节约条件
\((n-1)f-na-F_0>0\)，
即 \(n(f-a)>f+F_0\)。
\item 若 \(f-a>0\)，除以这个正数保持不等号方向，
得到所述阈值。
\item 若 \(f-a\le0\)，则左边非正、右边非负，
严格大于不可能成立。
\end{semanticsteps}
\end{proof}

\begin{remark}[标准平台的有限证明范围]
平台的身份、授权、事务、组件和接口契约，
可以承担服务部署中的一部分执行义务。
平台规范要求执行已经配置的访问控制规则，
并不替部署者选择全部用途政策
\cite{JakartaPlatform11}。
模块版本、前提、读回和实际任务覆盖，
仍属于引理~\ref{market:composition} 的独立证书。
同一平台的兼容实现不自动取得相同的人工智能输出正确性，
也不自动消除适配与维护成本。
\end{remark}

\section{可兑现需求与完整经济成本}

\begin{definition}[与价格选择分离的合同量]
\label{market:contract-demand}
固定服务任务、验收规则、期限和有限合同组合 \(J\)。
合同 \(j\) 指定数量 \(q_j>0\)，总量为
\[
 Q=\sum_{j\in J}q_j>0.
\]
在选择本期实际单位价格之前，登记共同价格上界
\(\bar p>0\)。
每份合同明确：只要所选价格
\(p\in[0,\bar p]\)、
服务按约交付且验收条件成立，
其对应数量仍为 \(q_j\)，付款额为 \(p q_j\)。
合同方具有未重复承诺的支付能力，
并有覆盖最高付款额 \(\bar p q_j\) 的可用预算、
授权及到期支付安排。
数量和上述条件在价格选择后不被任意改写。
\end{definition}

该定义登记的是一个可以兑现的有条件合同组合。
它不是把未知需求函数 \(Q(p)\) 代入随后选择的价格，
也不是把调查中的支付意愿直接作为销售事实。
如果实际需求会随价格变化而撤回，
则应另提交需求函数或数量下界证书；
本章有关固定 \(Q\) 的结论只适用于已经覆盖的合同量。
付款能力、付款义务和实际付款是不同谓词。

\begin{definition}[完整经济成本与利润]
\label{market:full-cost}
在同一核算期间和折算时点，
令 \(C^{\rm econ}(Q)\) 为交付合同组合的完整经济成本。
其中包括实际生产与运行、验证、证明适配、更新、
维护、责任履行、必要研究辅助和治理所消耗的资源，
以及自有资源和资本的正常机会成本。
采用可核验上界
\[
 C^{\rm econ}(Q)\le F+cQ,
 \qquad F\ge0,\quad c\ge0.
\]
令实际服务收入为 \(R\)，经济利润定义为
\[
 \Pi^{\rm econ}=R-C^{\rm econ}(Q).
\]
正常资本报酬已经计入成本；
另设的 \(\pi_0\ge0\) 是超过完整成本后的经济利润目标。
资本购置、折旧、残值与正常资本报酬使用一致的跨期口径，
不把同一资源支出重复计入成本。
\end{definition}

供应商的成本核算与全社会的资源核算具有不同边界。
外购投入的付款可以计入供应商成本；
合并付款人与收款人时，还须分清资源耗用与内部转移。
实际资金收支另按到期时点登记，
不能用经济成本的期间分摊替代现金支付计划。

\begin{proposition}[营业现金与经济利润互不替代]
\label{market:cash-not-profit}
在上述定义下，正营业现金结余不蕴涵正经济利润；
正的期间经济利润测算也不蕴涵每个支付前缀均有足够现金。
\end{proposition}
\begin{certificate}
两组有限金额与时间安排，
分别登记已经发生的营业收付、非现金机会成本和到期义务。
\end{certificate}
\begin{proof}
\begin{semanticsteps}
\item 第一组安排的收入为 \(10\)，
营业现金支出为 \(4\)，
另有未重复计量的自有资源和正常资本机会成本 \(8\)。
营业现金结余为 \(10-4=6>0\)，
经济利润为 \(10-(4+8)=-2<0\)。
\item 第二组安排的合同收入为 \(12\)，
完整成本为 \(10\)，故按全部履约结算计算的利润为 \(2\)。
其中成本 \(6\) 必须先支付，
收入 \(12\) 与其余成本 \(4\) 在后一时点发生。
\item 若第二组安排的初始可用现金和先期融资都为零，
首个支付前缀有 \(0-6<0\)。
该安排虽有正的期间利润，
仍不能据此完成先期支付。
\item 两个见证分别否定两个全称蕴涵。
因此成本回收与逐前缀资金覆盖必须分别证明。
\end{semanticsteps}
\end{proof}

\section{正价格区间与逐家成本覆盖}

\begin{theorem}[可兑现合同量上的非空价格区间]
\label{market:positive-price}
设定义~\ref{market:contract-demand} 与
定义~\ref{market:full-cost} 的条件成立，并有
\[
 Q(\bar p-c)>F+\pi_0.
\]
则区间
\[
 \mathcal P
 =\left(c+\frac{F+\pi_0}{Q},\,\bar p\right]
\]
非空，且其中每个价格严格为正。
对任意 \(p\in\mathcal P\)，
若合同全部实际交付、验收并足额结算，
则交易量为 \(Q>0\)，且
\(\Pi^{\rm econ}>\pi_0\)。
\end{theorem}
\begin{certificate}
价格选择前已经固定的合同数量、
支付上界与支付能力，
同等任务口径下的完整成本上界，
容量、交付与验收证据，
以及实际结算金额和期限的回执。
\end{certificate}
\begin{proof}
\begin{semanticsteps}
\item 因 \(Q>0\)，将条件除以 \(Q\) 得
\(\bar p>c+(F+\pi_0)/Q\)。
故下端点严格小于上端点，区间非空。
\item 下端点非负。
区间不包含下端点，因此任一
\(p\in\mathcal P\) 都满足 \(p>0\)。
\item 定义中的合同数量不随本区间内的价格改变，
所以足额交付的总数量仍为 \(Q\)。
足额实际结算给出 \(R=pQ\)。
\item 完整成本上界推出
\[
 \Pi^{\rm econ}
 \ge pQ-(F+cQ)
 =Q(p-c)-F.
\]
\item 由 \(p>c+(F+\pi_0)/Q\)，
得到 \(Q(p-c)-F>\pi_0\)。
将这一步与上一步合成即得结论。
\end{semanticsteps}
区间的非空性是代数结论。
实际正收入还使用了独立的交付、验收和支付前件。
\end{proof}

\begin{proposition}[选定服务商集合的逐家覆盖]
\label{market:each-provider}
固定有限非空的活跃服务商集合 \(I\)。
对每个 \(i\in I\)，有未重复分配的合同量 \(Q_i>0\)、
成本上界 \(F_i+c_iQ_i\)、目标 \(\pi_i^0\ge0\)
及价格上界 \(\bar p_i\)，并满足
\[
 a_i^\star:=c_i+\frac{F_i+\pi_i^0}{Q_i}<\bar p_i.
\]
允许分别定价时，可同时选择
\(p_i\in(a_i^\star,\bar p_i]\)。
若要求统一价格，则同时落入上述全部保证区间、
从而保证全部逐家严格利润目标的价格存在，当且仅当
\[
 \max_{i\in I}a_i^\star<\min_{i\in I}\bar p_i.
\]
所有结算和成本结论均以共同资源可行为前提。
\end{proposition}
\begin{certificate}
合同的唯一归属、逐家成本与支付能力，
共同设施成本的无重复分摊、
每家交付容量、价格规则、
共同支付时间表及逐家实际收款记录。
\end{certificate}
\begin{proof}
\begin{semanticsteps}
\item 分别定价时，每个区间均非空。
有限次选择各区间的中点，
得到全部 \(p_i\) 的共同价格向量。
\item 对每家分别应用定理~\ref{market:positive-price}，
在其合同实际完成时取得
\(\Pi_i^{\rm econ}>\pi_i^0\)。
共同资源证书保证这些合同能够同时完成。
\item 若统一价格 \(p\) 属于全部保证区间，
则对全部 \(i\) 有 \(p>a_i^\star\)、
\(p\le\bar p_i\)。
有限性给出
\(\max_i a_i^\star<p\le\min_i\bar p_i\)，
故必要条件成立。
\item 反之，若最大下端点严格低于最小上端点，
取两者的中点即同时属于每家价格区间。
这证明充分性。
\end{semanticsteps}
\end{proof}

\begin{remark}[总利润与市场范围]
总利润为正不保证每家成本得到覆盖：
两家利润分别为 \(10\) 和 \(-5\) 时，总利润仍为 \(5\)。
分别非空的价格区间也未必允许统一价格。
上述命题只证明所选服务商集合可共同履约，
不保证任意新增服务商均取得订单或永久留在市场。
\end{remark}

\section{两部收费、维护义务与进入边界}

\begin{definition}[固定维护费与使用费]
\label{market:two-part-data}
本节单独固定整数客户数 \(M\ge1\)。
每个客户的可行使用量属于非空紧集
\(K\subseteq[0,\bar q]\)，且 \(0\in K\)。
其服务价值 \(u:K\to\RR\) 连续并满足 \(u(0)=0\)。
完整增量成本为 \(c\ge0\)，
固定成本 \(F\ge0\) 包括合同期间已经承诺的维护、
更新、验证和正常资本机会成本。
两部收费由固定准入费 \(a\ge0\) 与单位使用费 \(c\) 组成。
订阅后的支付为 \(a+cq\)，不订阅的价值与支付均归零。
客户在预算可行的选择中最大化净价值，
并按预先登记的规则选择并列最优使用量。
总容量至少为 \(M\bar q\)。
每个客户的未重复承诺可用预算为 \(B_j\)。
\end{definition}

\begin{lemma}[使用量最优值的存在与预算约束]
\label{market:usage-maximum}
在定义~\ref{market:two-part-data} 下，
令
\[
 S(q)=u(q)-cq,\qquad S^\star=\max_{q\in K}S(q).
\]
最大值存在，且 \(S^\star\ge0\)。
按登记规则选定一个最大化使用量 \(q^\star\)，并定义
\[
 \widehat a
 =\min\left\{S^\star,\ \min_{1\le j\le M}(B_j-cq^\star)\right\}.
\]
若 \(0\le a<\widehat a\)，
则 \(q^\star\) 对每位订阅者可负担且仍为最优使用量，
订阅的净价值严格为正。
\end{lemma}
\begin{certificate}
紧可行域、连续价值函数、
固定增量成本、实际可用预算、
最大值见证与数量选择规则。
\end{certificate}
\begin{proof}
\begin{semanticsteps}
\item \(S\) 连续，\(K\) 非空紧致，
故极值定理给出最大值和至少一个最大化点。
又因 \(0\in K\)，有
\(S^\star\ge S(0)=0\)。
\item 由 \(a<\widehat a\le B_j-cq^\star\)，
得到 \(a+cq^\star<B_j\)。
故 \(q^\star\) 位于客户 \(j\) 的实际预算可行域。
\item 对任意其他预算可行使用量 \(q\in K\)，
最大化性质给出 \(S(q^\star)\ge S(q)\)。
两边减去同一固定费 \(a\)，最优性不变。
\item 由 \(a<\widehat a\le S^\star\)，
得到 \(S^\star-a>0\)。
它严格高于不订阅的零值。
\end{semanticsteps}
\end{proof}

\begin{theorem}[两部收费的完整成本回收]
\label{market:two-part}
在上一引理的条件下，若
\[
 M\widehat a>F+\pi_0,
 \qquad \pi_0\ge0,
\]
则可以选择
\[
 \frac{F+\pi_0}{M}<a<\widehat a.
\]
对遵守所登记数量选择与支付合同的客户，
订阅量为 \(M\)，总使用量严格为正。
若全部服务交付和付款实际完成，
完整经济利润满足 \(\Pi^{\rm econ}>\pi_0\)。
\end{theorem}
\begin{certificate}
整数客户清单、非重复预算、
上一引理的最大值与可负担性证书、
客户选择和服务接纳规则、
固定维护义务、完整成本与实际支付记录。
\end{certificate}
\begin{proof}
\begin{semanticsteps}
\item 因 \(M>0\)，条件等价于
\((F+\pi_0)/M<\widehat a\)。
故所列开区间非空，并且其中任意 \(a>0\)。
\item 上一引理使 \(M\) 位客户都能负担最优使用量，
并严格偏好履行订阅合同。
合同接纳与选择规则给出实际订阅数量 \(M\)。
\item 条件还给出
\(S^\star\ge\widehat a>(F+\pi_0)/M\ge0\)，
故 \(S^\star>0\)。
由于 \(S(0)=0\)，任何最大化点都不等于零。
因此 \(Mq^\star>0\)。
\item 足额实际支付给出
\(R=Ma+cMq^\star\)。
完整成本上界为 \(F+cMq^\star\)，
所以
\[
 \Pi^{\rm econ}\ge Ma-F>\pi_0.
\]
\item 使用费恰用于覆盖对应的增量成本；
固定费承担固定维护与利润目标。
收入与成本中的同一使用量项相减一次，
不把这部分收入再次视作固定成本回收。
\end{semanticsteps}
\end{proof}

\begin{remark}[实际付款与固定责任]
严格订阅价值和支付能力使合同具有可行的支付基础，
但不构成实际付款回执。
维护责任在已承诺期间内不因客户当期少用或不用而自动消失。
若收费安排不足以承担这些固定义务，
不能通过删除必要维护项目来宣称完整成本已经覆盖。
固定费提前收取可以改变现金时序，
但收费前的构造支出仍须有先期资金。
\end{remark}

\begin{proposition}[有限客户下的可维持服务商数量]
\label{market:entry-bound}
设共有 \(M\in\NN\) 个客户，
每个客户在本期至多向一家服务商支付固定费 \(a>0\)。
每家服务商的固定完整成本均为 \(F>0\)，
单位使用费仅覆盖单位使用成本。
若有 \(N\) 家服务商分别覆盖完整成本，
则
\[
 N\left\lceil\frac{F}{a}\right\rceil\le M,
 \qquad\text{从而}\qquad
 N\le\frac{Ma}{F}.
\]
\end{proposition}
\begin{certificate}
唯一订阅关系、整数客户数、
逐家固定费实际收入、
固定完整成本和使用费收支抵销表。
\end{certificate}
\begin{proof}
\begin{semanticsteps}
\item 令服务商 \(i\) 的付费客户数为整数 \(M_i\)。
成本覆盖要求 \(aM_i-F\ge0\)，
所以 \(M_i\ge F/a\)。
\item 整数性进一步给出
\(M_i\ge\lceil F/a\rceil\)。
对 \(N\) 家求和，左边至少为
\(N\lceil F/a\rceil\)。
\item 唯一订阅关系保证
\(\sum_iM_i\le M\)，得到第一式。
也可直接将 \(aM_i\ge F\) 求和，
得到 \(aM\ge NF\)，即第二式。
\end{semanticsteps}
\end{proof}

自由进入还需要登记新增服务商实际能够争取的客户、
进入成本、现有服务商的响应和退出规则。
若一个进入者在这些响应之后仍可复制严格正的净经济利润，
并且没有其他进入约束，则“无人愿意进入”的候选状态不能成立。
反之，已有企业的正会计结余不自动提供这样的进入见证。
固定服务商数量下的正经济利润结论，
与自由进入后的长期利润结论应当分别证明。

\section{逐前缀资金、采购与实际结算}

\begin{definition}[按真实先后排列的支付前缀]
\label{market:cash-data}
将合同期间全部现金事件按实际先后排列为
\(t=1,\ldots,T\)。
令 \(B_0\ge0\) 为初始未重复承诺余额，
\(I_t\ge0\) 为当步已经到账且可用的收入，
\(A_t\ge0\) 为当步已经可用的其他资金，
\(O_t\ge0\) 为当步必须履行的支付。
同一序号内的收入只有先于付款可用时才可抵付该笔付款；
否则必须拆成不同序号。
全部付款后的余额定义为
\[
 B_t=B_0+\sum_{s=1}^t(I_s+A_s-O_s).
\]
借款本金可以属于 \(A_t\)，但不因此成为经济利润。
相应偿还、利息和其他到期义务必须进入后续 \(O_t\)。
\end{definition}

\begin{theorem}[有限支付序列的资金覆盖归纳]
\label{market:prefix-cash}
在定义~\ref{market:cash-data} 的固定事件顺序下，
全部付款可在不出现负余额的条件下获得资金支持，
当且仅当每个前缀满足
\[
 B_0+\sum_{s=1}^t(I_s+A_s)
 \ge\sum_{s=1}^tO_s,
 \qquad 1\le t\le T.
\]
若另有合法支付通道、按序执行和到期到账证据，
则这些资金支持的付款实际完成。
\end{theorem}
\begin{certificate}
真实事件顺序、已经可用的资金、
无重复抵押与无重复承诺记录、
全部到期义务、逐前缀不等式、
支付权限、结算通道、执行指令和收款回执。
\end{certificate}
\begin{proof}
\begin{semanticsteps}
\item 必要性：若全部付款后余额均非负，
将余额定义中的付款项移至另一侧，
逐项得到所列前缀不等式。
\item 充分性的初态为 \(B_0\ge0\)。
假设前 \(t-1\) 步已经按记录完成，
则付款前可用余额为
\(B_{t-1}+I_t+A_t\)。
\item 将余额定义代入第 \(t\) 个前缀不等式，
得到
\(B_{t-1}+I_t+A_t\ge O_t\)。
故当步付款具有足够资金，
付款后余额 \(B_t\ge0\)。
\item 对 \(t=1,\ldots,T\) 归纳，
得到全部付款的资金可行性。
资金可行性没有推出支付操作已经发生。
\item 实际完成还须将每步已授权的付款操作执行到
规定的收款账户，并核对到账回执。
这些附加证据逐步建立
\(\operatorname{Funded}\land\operatorname{Executed}
\land\operatorname{Received}\)，
从而把可支付金额转化为实际收入。
\end{semanticsteps}
\end{proof}

\begin{proposition}[采购与补助的唯一收入归属]
\label{market:transfer-once}
设所有实际收款具有唯一事件标识，
服务购买、公共采购与补助只是这些事件的分类。
同一收款事件至多计入一次总收入。
在同时包含付款人与收款人的核算边界内，
该笔货币转移本身不增加合并后的净货币资源。
\end{proposition}
\begin{certificate}
实际收款事件集合、金额、
付款人与收款人、合同和补助归属，
以及发票总额与分担金额之间的核对表。
\end{certificate}
\begin{proof}
\begin{semanticsteps}
\item 对唯一事件集合 \(E\)，收入定义为
\(R=\sum_{e\in E}r_e\)。
把 \(E\) 划分为互不相交的类别后，
分组求和仍等于该总和。
\item 若一张发票由客户支付 \(b\)、公共资金支付 \(g\)，
总收入为 \(b+g\)。
发票总额已经包含 \(g\)，
再加一次 \(g\) 会把同一事件重复计入。
\item 每笔付款在付款人账上为负、
在收款人账上为等额正数。
合并抵销由引理~\ref{opt:transfer} 得到。
采购形成的真实服务价值、征集与管理成本，
以及资金机会成本另外登记。
\end{semanticsteps}
\end{proof}

\begin{corollary}[受用途约束的非空服务市场]
\label{market:nonempty}
若一个有限非空服务商集合具有相容的合法用途、
已经核验的任务交付、定义~\ref{market:contract-demand}
的合同量、定理~\ref{market:positive-price}
的价格余量、逐家完整成本覆盖和
定理~\ref{market:prefix-cash} 的资金及实际结算证书，
则该共同见证构成一个正价格、正交易量、
逐家超过所定完整成本目标的合规服务市场。
\end{corollary}
\begin{certificate}
上述全部证书在同一主体、合同、任务、规则、
期间、货币单位和真实资源上的共同实现；
受保护常规用途的限制及既有参与权利保持不变。
\end{certificate}
\begin{proof}
\begin{semanticsteps}
\item 令活跃服务商集合为有限非空集 \(I\)。对每份实际合同 \(r\)，
用途与交付证书分别给出
\(\operatorname{Permit}(r)=1\) 与 \(\operatorname{Delivered}(r)=1\)。
二者合取使该交易属于所声明的合法交付集合。
\item 对每个 \(i\in I\)，价格区间和真实合同量给出
\(p_i>0\)、\(Q_i>0\)。逐家证书分别保留这些量，
不以总利润替代任一主体的收入条件。
\item 资金前缀及匹配支付回执给出 \(R_i=p_iQ_i\)。
代入完整成本界，得到
\[
\Pi_i^{\rm econ}=R_i-C_i^{\rm econ}(Q_i)
\ge(p_i-c_i)Q_i-F_i>\pi_i^0\ge0.
\]
\item 所有前件由同一个资源与执行见证满足，因此上述逐家结论能够合取。
从 \(I\ne\varnothing\) 取 \(i_0\in I\)，则
\(\sum_{i\in I}Q_i\ge Q_{i_0}>0\)，故合规市场非空。
\end{semanticsteps}
\end{proof}


\chapter{服务差异、有效竞争与无效投入}
\label{integrated:competition}
\label{comp:chapter}

\section{兼容标准、真实服务差异与收入归属}

服务使用同一接口，并不意味着客户对所有服务商的评价相同。
响应时间、经过验收的任务范围、
维护责任、实际适配距离和资源可用性，
都可以构成真实服务差异。
这些差异是否值得其成本，应由客户价值和完整资源核算确定。
标准化与兼容性可以改变竞争条件，
却不独立决定价格或企业利润\cite{KatzShapiro1994}。

本章分别研究真实差异形成的价格竞争，
以及不增加真实产出但争夺收入归属的投入。
低价格、低经济利润和高生产率竞争不被定义为无效投入。
无效投入必须由明确的价值不增、
真实成本和可删除性证书识别。
后文采用固定位置的空间竞争模型说明有差异的服务定价；
其分析传统见 \cite{Hotelling1929}。
这里全部均衡结论由本章给定假设直接推导。

\begin{remark}[存在结论的经济边界]
推论~\ref{market:nonempty} 证明给定证书下可以保留商业服务空间，
并不证明需求证书已在任何现实行业成立。
它也不保证每家进入者盈利、
企业价值不受限制影响，或市场规模持续扩张。
命题~\ref{front:profit} 的受限机会比较仍然成立；
新市场能否补偿放弃的收益，
须进入第~\ref{opt:chapter}章的完整制度比较。
\end{remark}

\begin{definition}[真实服务选择与无效投入]
\label{comp:effort-data}
设服务商 \(i\) 的报价为 \(p_i\)，
真实且可核验的服务属性为 \(z_i\)，
其他投入为 \(e_i\in E_i\)，其中 \(0\in E_i\)。
将联合向量分别记为 \(p,z,e\)。
固定真实交付和 \(z\) 时，
投入 \(e_i\) 不增加所登记的服务总价值，
却消耗真实资源 \(k_i(e_i)\)，其中
\[
 k_i(0)=0,\qquad k_i(e_i)>0\quad(e_i\ne0).
\]
投入前后的收入归属仍可能不同：
\[
 \Pi_i^0(p,z,e)
 =R_i^0(p,z,e)-C_i(p,z)-k_i(e_i).
\]
\(C_i\) 已包括其余完整经济成本。
仅有服务价值不增，尚未要求 \(R_i^0\) 与 \(e\) 无关。
\end{definition}

\begin{proposition}[价值不增并不排除私人投入激励]
\label{comp:rent-example}
存在一个两服务商有限模型，
其中投入不增加服务价值，
但双方的唯一均衡都采用正投入。
\end{proposition}
\begin{certificate}
固定已支付收入池 \(B>0\)，
取 \(0<\kappa<\delta\le B/2\)，
令 \(e_i\in\{0,1\}\)，并规定
\[
 R_1^0=\frac B2+\delta(e_1-e_2),\qquad
 R_2^0=\frac B2+\delta(e_2-e_1).
\]
两家除投入外的成本固定，
各自投入成本为 \(\kappa e_i\)，
真实交付与客户价值在四种投入组合下相同。
\end{certificate}
\begin{proof}
\begin{semanticsteps}
\item 由 \(\delta\le B/2\)，
两家在所有组合中的收入均非负。
相加得到 \(R_1^0+R_2^0=B\)，
故投入只改变同一收入池的归属。
\item 固定对方任意选择，
本方将 \(e_i\) 从 \(0\) 改为 \(1\)，
收入增加 \(\delta\)，真实成本增加 \(\kappa\)。
净利润增加 \(\delta-\kappa>0\)。
\item 因而对方的两种选择下，
本方均严格选择 \(1\)。
唯一相互最优的组合是 \((1,1)\)。
\item 在该组合与 \((0,0)\) 中，
两家的收入均为 \(B/2\)，真实服务价值相同；
前者却多耗费 \(2\kappa\)。
这给出所需有限见证。
\end{semanticsteps}
\end{proof}

\begin{definition}[删除投入的可行性与非干扰证书]
\label{comp:deletion}
在治理后的共同可行集合 \(\mathcal A\) 中，
对任意联合选择 \((p,z,e)\) 和任意 \(i\)，
登记一个允许的单方删除操作
\[
 D_i(p,z,e)=(p,z,e_{-i},0)\in\mathcal A.
\]
它保持生产、交付、安全、用途准入、
固定维护责任、他人权利和真实服务属性 \(z\)。
同时，全部收入相关输入及其验证结果保持不变，
包括真实需求归属、资格、队列、验收记录、
结算规则和被纳入期间内的后继状态。
治理后收入因此具有表示
\[
 R_i^G(p,z,e)=\overline R_i(p,z).
\]
这一表示要求实际非干扰证书；
只在函数的书写中省略符号 \(e\)，不满足定义。
\end{definition}

若投入可以改变被称作 \(z\) 的报告值，
则须核验报告值与真实属性的读回，
并证明删除操作不会通过该通道改变收入。
若投入影响未来报价、准入或客户转移，
这些后继也必须纳入收入相关观察。
未被覆盖的间接通道不属于本定义的已证明范围。

\begin{theorem}[可行删除与收入非干扰下的零无效投入]
\label{comp:zero-effort}
若定义~\ref{comp:effort-data}
与定义~\ref{comp:deletion} 的条件成立，
且治理后的利润为
\[
 \Pi_i^G(p,z,e)
 =\overline R_i(p,z)-C_i(p,z)-k_i(e_i),
\]
则任何含 \(e_i\ne0\) 的可行选择，
都不是服务商 \(i\) 相对于固定其他行为的利润最大化选择。
若以允许的可行单方偏离定义一个均衡，
则任何这样的均衡都满足 \(e_i=0\) 对全部 \(i\) 成立。
\end{theorem}
\begin{certificate}
真实价值与成本分解、
每个非零投入的正成本、
不破坏生产及准入条件的删除操作，
覆盖直接与间接通道的收入非干扰证明，
以及所使用的单方偏离和期间边界。
\end{certificate}
\begin{proof}
\begin{semanticsteps}
\item 固定任一可行选择，设 \(e_i\ne0\)。
删除证书给出另一个可行选择 \(D_i(p,z,e)\)。
该选择保持其他参与者行为及本方真实属性 \(z_i\)。
\item 非干扰证书使删除前后收入相同；
其余真实成本 \(C_i(p,z)\) 也相同。
删除后的投入成本为 \(k_i(0)=0\)。
\item 因而利润差为
\[
 \Pi_i^G(D_i(p,z,e))-\Pi_i^G(p,z,e)
 =k_i(e_i)>0.
\]
存在这个可行的严格改进，
故原选择不可能是利润最大化选择。
\item 若一个均衡含有非零 \(e_i\)，
上一步提供允许的获利偏离，与均衡定义矛盾。
所以任何存在的均衡均有全部 \(e_i=0\)。
\end{semanticsteps}
本定理不以排除无效投入来证明均衡存在；
均衡存在及真实服务选择的最优性仍需各自条件。
\end{proof}

\begin{corollary}[真实创新选择的保留]
\label{comp:innovation}
若某项改进改变了真实服务属性、可实现任务、
交付成本或客户价值，
则不能仅凭定理~\ref{comp:zero-effort}
将其归入必须删除的 \(e_i\)。
在该定理的实际适用范围内，
删除 \(e_i\) 保留已经选择的真实属性 \(z\)。
\end{corollary}
\begin{certificate}
改进对真实任务和资源的作用记录，
以及删除操作前后的属性、权利与履约对照。
\end{certificate}
\begin{proof}
\begin{semanticsteps}
\item 定理的删除操作要求保持真实属性与生产可行性。
若某项改进是这些属性成立的必要条件，
删除它就不满足该前提。
\item 若改进增加真实服务价值，
它也不满足无效投入的价值不增条件。
\item 对确实满足全部条件的 \(e_i\)，
删除定义直接给出 \(z\) 不变。
因此结论只排除已定义的无效投入，
不排除通过 \(z\) 表示的真实创新。
\end{semanticsteps}
\end{proof}

\section{固定两端服务差异的价格模型}

\begin{definition}[固定位置的双服务商模型]
\label{comp:hotelling-data}
本节令 \(M>0\) 为连续客户总体的规模，
客户 \(x\) 在 \([0,1]\) 上均匀分布，密度为 \(M\)。
它与上一章两部收费模型中的整数客户数分别定义。
两服务商的位置固定为 \(0\) 与 \(1\)，
每家容量均为 \(M\)，价格允许取任意 \(p_i\ge0\)。
单位完整增量成本为 \(c\ge0\)，
活跃服务商的完整固定成本为 \(F\ge0\)。
令真实匹配耗费参数 \(\theta>0\)，服务价值满足
\[
 v\ge c+2\theta.
\]
每位客户至多购买一单位，其净价值为
\[
 U_1(x)=v-p_1-\theta x,\qquad
 U_2(x)=v-p_2-\theta(1-x),\qquad U_0(x)=0.
\]
客户选择净价值最高的可用服务；
非负价值与不购买并列时购买。
两服务并列的单个位置按任意固定规则分配，
其客户测度为零。
活跃服务商的利润为
\[
 \Pi_i=(p_i-c)D_i-F.
\]
\end{definition}

固定位置表示已经存在的真实服务匹配差异，
本模型不包含位置选择或进入选址的均衡。
匹配耗费计量客户的真实时间、转换或适配损失，
并与已经进入固定成本的证明适配项目区分。
参数 \(F\) 包含正常资本报酬，
故此处利润是超过完整成本后的经济利润。

\begin{lemma}[面对指定对手价格的分段需求]
\label{comp:demand}
在定义~\ref{comp:hotelling-data} 下，
若对方价格为 \(c+\theta\)，则本方价格 \(p\ge0\) 的需求为
\[
 D(p)=
 \begin{cases}
 M,&0\le p\le c,\\[2pt]
 \displaystyle\frac{M(c+2\theta-p)}{2\theta},
       &c\le p\le c+2\theta,\\[6pt]
 0,&p\ge c+2\theta.
 \end{cases}
\]
边界处两种表达相同。
每种价格下的需求均不超过本方容量。
\end{lemma}
\begin{certificate}
连续客户分布、完整容量、
固定两端的位置和价值函数、
对方价格与外部选项，
以及实际采取的并列处理规则。
\end{certificate}
\begin{proof}
\begin{semanticsteps}
\item 先考虑左端服务商。
对方价格为 \(c+\theta\) 时，
\[
 U_2(x)=v-c-2\theta+\theta x\ge0.
\]
故所有客户至少有一个非负价值的服务选项。
这一点对本方全部非负价格都成立。
\item 两服务的净价值差为
\[
 U_1(x)-U_2(x)=c+2\theta-p-2\theta x.
\]
因为 \(\theta>0\)，左端服务被选择的条件为
\[
 x\le x^\star(p):=\frac{c+2\theta-p}{2\theta}.
\]
\item 若 \(p\le c\)，则 \(x^\star(p)\ge1\)，
本方取得全部测度 \(M\)。
若 \(c\le p\le c+2\theta\)，
则 \(x^\star(p)\in[0,1]\)，
需求为 \(Mx^\star(p)\)。
\item 若 \(p\ge c+2\theta\)，
则 \(x^\star(p)\le0\)，
本方至多取得零测度的端点客户，故需求为零。
\item 对右端服务商作变量替换 \(x\mapsto1-x\)，
得到相同公式。
三段公式均属于 \([0,M]\)，容量条件因此满足。
\end{semanticsteps}
\end{proof}

\begin{lemma}[价格全域上的平方配方上界]
\label{comp:global-best-response}
面对对方价格 \(c+\theta\)，
活跃服务商在全部 \(p\ge0\) 中的唯一最优价格为
\[
 p^\star=c+\theta.
\]
最大变动利润为 \(M\theta/2\)，
扣除固定成本后的利润为 \(M\theta/2-F\)。
\end{lemma}
\begin{certificate}
引理~\ref{comp:demand} 的全部三段需求，
非负成本和正的 \(M,\theta\)，
以及在价格比较中保持不变的固定成本。
\end{certificate}
\begin{proof}
\begin{semanticsteps}
\item 若 \(0\le p\le c\)，
变动利润为 \(M(p-c)\le0\)。
若 \(p\ge c+2\theta\)，需求为零，变动利润也是零。
\item 在中间区间令 \(y=p-c\in[0,2\theta]\)。
将需求代入得到
\[
 (p-c)D(p)
 =\frac{M}{2\theta}y(2\theta-y)
 =\frac{M}{2\theta}\bigl(\theta^2-(y-\theta)^2\bigr).
\]
\item 平方非负，所以中间区间的变动利润不超过
\(M\theta/2\)。
等号成立当且仅当 \(y=\theta\)，即 \(p=c+\theta\)。
\item 因 \(M\theta/2>0\)，
两侧区间的非正值或零值不能达到这一上界。
故该点是全体非负价格中的唯一最大化点。
\item 减去相同固定成本 \(F\) 不改变最大化价格，
最大完整经济利润为 \(M\theta/2-F\)。
\end{semanticsteps}
这证明的是全部可行价格上的最优反应，
而不只是某个内点一阶条件。
\end{proof}

\begin{theorem}[价格均衡与可避免固定成本的退出条件]
\label{comp:hotelling-equilibrium}
在固定两家活跃服务商的价格子博弈中，
\[
 p_1=p_2=c+\theta,\qquad D_1=D_2=M/2
\]
构成纳什均衡，两家各得经济利润
\[
 \Pi_1=\Pi_2=\frac{M\theta}{2}-F.
\]
若服务商还可以在承担固定成本前退出，
退出收益为零，且 \(F\) 完全可避免，
则该双活跃候选构成包含退出选择的均衡，
当且仅当
\[
 F\le\frac{M\theta}{2}.
\]
\end{theorem}
\begin{certificate}
模型参数、两家完整容量、
全域价格最优反应、
固定成本的发生时点与可避免性、
零退出收益及消费者选择规则。
商业执行另外提交服务验收和实际结算证书。
\end{certificate}
\begin{proof}
\begin{semanticsteps}
\item 对手报价 \(c+\theta\) 时，
引理~\ref{comp:global-best-response}
给出本方唯一最佳活跃报价也是 \(c+\theta\)。
该事实对两家对称成立。
\item 将该价格代入需求公式，
得到每家 \(D_i=M/2\)；
代入利润式，得到 \(M\theta/2-F\)。
因而任何单方活跃价格偏离都不能提高利润。
\item 若 \(F\le M\theta/2\)，
均衡活跃收益非负，退出收益为零，
故退出也不能构成严格获利偏离。
全部单方选择均已覆盖，候选成为均衡。
\item 若 \(F>M\theta/2\)，
活跃收益严格为负，而退出取得零，
退出是严格获利偏离。
因此该候选不能是包含退出的均衡。
\end{semanticsteps}
边界相等时，经济利润为零但完整成本已经覆盖。
命题不声称任意数量进入者下仍保持相同价格或利润。
\end{proof}

\begin{corollary}[有限数值见证]
\label{comp:numeric}
取
\[
 (M,c,\theta,F,v)=(100,2,2,60,6).
\]
则上述均衡价格为 \(4\)，
每家成交量为 \(50\)，
每家合同收入为 \(200\)，
变动成本为 \(100\)，
完整固定成本为 \(60\)，经济利润为 \(40\)。
\end{corollary}
\begin{certificate}
声明的五个参数、每家容量 \(100\)、
相同的收入与成本单位、
固定成本包含正常资本报酬的核算口径。
\end{certificate}
\begin{proof}
\begin{semanticsteps}
\item 检查覆盖条件：
\(v=6=c+2\theta\)。
检查退出条件：
\(F=60<100\cdot2/2=100\)。
\item 价格为 \(c+\theta=4\)，
每家数量为 \(M/2=50\)。
因此合同收入为 \(4\cdot50=200\)，
变动成本为 \(2\cdot50=100\)。
\item 完整经济利润为
\(200-100-60=40>0\)。
实际收到这笔收入仍以按约交付和足额结算为条件。
\end{semanticsteps}
\end{proof}

\section{同质定价与下一期间固定成本}

\begin{definition}[完全同质的单位价格子博弈]
\label{comp:homogeneous-data}
另设两家服务在全部客户重视的属性上完全相同，
客户总量为 \(M>0\)，每家容量为 \(M\)，
单位完整增量成本为 \(c\ge0\)，
客户共同价值 \(v>c\)。
各家价格 \(p_i\ge0\)。
较低且不超过价值的价格取得全部需求，
同价时均分需求。
本子博弈中每家固定成本 \(F\ge0\) 已经发生，
不能通过当期停业收回。
这些同质性条件独立于标准接口是否兼容。
\end{definition}

\begin{proposition}[同质定价的变动成本边界]
\label{comp:homogeneous-price}
在定义~\ref{comp:homogeneous-data} 下，
\(p_1=p_2=c\) 构成价格子博弈的纳什均衡。
每家变动利润为零；
若将已经发生的固定成本计入完整期间损益，
每家经济利润为 \(-F\)。
\end{proposition}
\begin{certificate}
全部属性同质、无容量短缺、
价格与需求分配规则、
增量成本、已发生且不可收回的固定成本。
\end{certificate}
\begin{proof}
\begin{semanticsteps}
\item 对方价格为 \(c\) 时，
本方取 \(p=c\) 得到需求 \(M/2\)，
变动利润为 \((c-c)M/2=0\)。
\item 取任意可行的 \(p<c\)，
虽可取得全部需求，但变动利润为
\((p-c)M<0\)。
若 \(c=0\)，这一低价区间为空。
\item 取任意 \(p>c\)，
对方提供相同属性且容量充足，
本方需求为零，变动利润也是零。
\item 因而本方无任何价格偏离取得正变动利润；
两家对称，构成纳什均衡。
扣除已发生的 \(F\)，完整期间利润为 \(-F\)。
\end{semanticsteps}
\end{proof}

\begin{corollary}[可避免固定成本的下一期缺口]
\label{comp:next-period-gap}
若下一期间重新承担 \(F>0\) 是进入或继续经营的必要条件，
而退出可取得零，且仍仅依靠上一命题的双服务商单位价格安排，
则双方同时承担固定成本不能通过无亏损进入条件。
这一结论不排除固定费、真实采购、
成本下降或有价值服务差异所支持的其他市场安排。
\end{corollary}
\begin{certificate}
下一期完整固定成本、其可避免性、
预期价格子博弈、全部其他合法收入来源的清单，
以及零退出收益。
\end{certificate}
\begin{proof}
\begin{semanticsteps}
\item 对预期单位价格 \(p=c\)，
单位价格收入只覆盖对应增量成本，
无法贡献固定成本的回收额。
\item 加入下一期必须承担的 \(F>0\)，
每家完整经济利润为 \(-F<0\)。
退出取得零，故至少无亏损进入这一必要条件失败。
\item 固定费或其他真实收入会改变利润公式；
成本或真实差异的改变会改变成本与需求。
这些安排不受保持原公式的否定结论排除。
\end{semanticsteps}
\end{proof}

\begin{remark}[价格下降、正常报酬与无效投入]
单位价格下降可以来自成本降低、匹配耗费减少或竞争加强。
这些行为是否浪费资源，不能由价格方向单独判断。
完整经济利润为零，意味着包括正常资本机会成本在内的成本得到覆盖；
它不同于营业现金为零，也不同于固定成本无法回收。
当 \(c=0\) 且固定成本已由其他明确来源覆盖时，
零单位使用价格也可以与有价值的正使用量相容。
本章只把满足定义~\ref{comp:effort-data}
和定义~\ref{comp:deletion} 的投入纳入可删除结论。
\end{remark}

\section{匹配、证明适配与资金机会成本的福利比较}

\begin{proposition}[固定两端模型的完整剩余]
\label{comp:hotelling-welfare}
在双活跃对称均衡中，
若 \(F\) 与 \(c\) 已涵盖模型内全部其他真实成本，
则总剩余为
\[
 W=M(v-c)-2F-\frac{M\theta}{4}.
\]
其中客户总匹配耗费为 \(M\theta/4\)。
价格支付作为客户与服务商之间的转移抵销。
\end{proposition}
\begin{certificate}
完全覆盖的客户量、真实价值、
匹配耗费与其他成本的互不重叠分解、
共同福利计量单位及完整固定成本。
分配权重若不同，须另列相应分配项。
\end{certificate}
\begin{proof}
\begin{semanticsteps}
\item 对称均衡中，
左半区由左端服务商服务，
右半区由右端服务商服务。
总匹配耗费为
\[
 M\theta\left(
 \int_0^{1/2}x\,dx+
 \int_{1/2}^1(1-x)\,dx
 \right)=\frac{M\theta}{4}.
\]
\item 所有 \(M\) 单位服务的总价值为 \(Mv\)，
增量真实成本为 \(Mc\)，
两家固定完整成本为 \(2F\)。
\item 客户支付总额为 \(M(c+\theta)\)。
它从客户剩余中扣除，
又等额进入两家收入，合并时相互抵销。
\item 将真实价值减去三个互不重复的真实成本项目，
得到所述公式。
也可直接相加客户剩余
\(M(v-c-5\theta/4)\)
与生产者利润 \(M\theta-2F\)，得到同一结果。
\end{semanticsteps}
\end{proof}

\begin{corollary}[固定成本复用的价格与剩余效应]
\label{comp:fixed-saving}
若经过证明的复用使每家完整固定成本从 \(F\)
降为 \(F-\eta\)，其中 \(0\le\eta\le F\)，
并保持 \(M,c,\theta,v\)、任务质量、
竞争结构及其他真实成本不变，
则对称均衡价格保持 \(c+\theta\)，
每家经济利润增加 \(\eta\)，总剩余增加 \(2\eta\)。
\end{corollary}
\begin{certificate}
命题~\ref{market:reuse-saving} 的同等任务成本节约，
固定成本的实际分摊、
所有保持不变的参数与覆盖条件，
以及两家活跃的退出条件。
\end{certificate}
\begin{proof}
\begin{semanticsteps}
\item 全域最佳反应的平方配方中，
固定成本只作为与价格无关的减项。
其下降不改变最佳价格 \(c+\theta\)。
\item 每家利润从 \(M\theta/2-F\)
变为 \(M\theta/2-(F-\eta)\)，
差额为 \(\eta\)。
\item 总剩余公式中的固定成本项
从 \(-2F\) 变为 \(-2(F-\eta)\)，
其他项不变，故总剩余增加 \(2\eta\)。
\end{semanticsteps}
因此，固定成本节约可以产生真实福利收益，
却不必在这个固定竞争结构中直接表现为单位价格下降。
\end{proof}

\begin{corollary}[真实匹配耗费下降的竞争效应]
\label{comp:matching-saving}
设 \(0<\theta'<\theta\)，
其他参数与完整成本均保持不变，
且 \(F\le M\theta'/2\)。
若这一变化只减少真实匹配耗费，
则双活跃均衡价格下降 \(\theta-\theta'\)，
每家经济利润减少 \(M(\theta-\theta')/2\)，
总剩余增加 \(M(\theta-\theta')/4\)。
\end{corollary}
\begin{certificate}
真实价值与质量保持证书、
匹配耗费的下降、完整成本不变、
新旧容量、覆盖及退出条件。
\end{certificate}
\begin{proof}
\begin{semanticsteps}
\item 原覆盖条件 \(v\ge c+2\theta\)
蕴涵 \(v\ge c+2\theta'\)；
给定的固定成本条件保证新模型可保持两家活跃。
\item 两个均衡价格分别为 \(c+\theta\)
和 \(c+\theta'\)，相减得到所述降幅。
\item 两个利润分别为 \(M\theta/2-F\)
和 \(M\theta'/2-F\)，相减得到利润减少额。
\item 总剩余中只改变匹配耗费项，
故增加额为 \(M(\theta-\theta')/4>0\)。
\end{semanticsteps}
如果新匹配差异不足以覆盖固定成本，
本推论的双活跃前提失败，
须重新研究进入、退出或收费安排。
\end{proof}

\begin{definition}[共同边界内的福利与成本分解]
\label{comp:welfare-data}
对可比较的方案 \(s\)，登记
\[
 \begin{aligned}
 W(s)={}&V(s)-C_{\rm prod}(s)-C_{\rm proof}(s)
          -C_{\rm match}(s)\\
        &-C_{\rm fund}(s)-C_{\rm gov}(s)-K(s).
 \end{aligned}
\]
各项分别为真实服务价值、生产成本、
证明与适配成本、客户真实匹配耗费、
资金的真实机会成本、治理成本和已定义的无效投入成本。
这些项目互不重复，并折算至同一时点。
若先前的 \(F\) 或 \(c\) 已经包含某项，
本分解从该总额中分出对应部分，不在原总额之外再次加计。
借款本金、采购付款、服务费和补助本身不是新增真实价值；
其资源机会成本和实施耗费按上述类别计入。
\end{definition}

\begin{theorem}[标准化与投入治理的精确福利条件]
\label{comp:welfare-condition}
对两个满足共同目标与权利约束的方案
\(s^0,s^1\)，令
\(\Delta X=X(s^1)-X(s^0)\)，并令
\(\mathcal C=\{{\rm prod,proof,match,fund,gov}\}\)。
则
\[
 W(s^1)-W(s^0)
 =\Delta V-\sum_{h\in\mathcal C}\Delta C_h
   +K(s^0)-K(s^1).
\]
因此严格福利改进当且仅当
\[
 \Delta V+K(s^0)-K(s^1)
 >\sum_{h\in\mathcal C}\Delta C_h.
\]
特别地，若真实服务和除治理外的真实成本均保持不变，
无效投入被全部删除，而新增治理成本为 \(g\ge0\)，
则严格改进当且仅当 \(K(s^0)>g\)。
\end{theorem}
\begin{certificate}
同一主体与期间的实际方案、
共同福利目标、全部直接和间接效应、
无重复成本分组、真实匹配耗费、
证明适配与资金机会成本，
以及被删除投入的非干扰证书。
\end{certificate}
\begin{proof}
\begin{semanticsteps}
\item 对两个方案分别写出
定义~\ref{comp:welfare-data} 的恒等式。
逐项相减并按同一类别合并，
得到第一式。
\item 严格福利改进的定义是左边大于零。
把真实成本差额移至右边，
得到所列充要条件；移项可逆。
\item 在特别情形中，
\(\Delta V=0\)，
除治理外的 \(\Delta C_h=0\)，
\(\Delta C_{\rm gov}=g\)，且 \(K(s^1)=0\)。
代入恒等式得到
\(W(s^1)-W(s^0)=K(s^0)-g\)。
\item 该差额严格为正等价于 \(K(s^0)>g\)。
因此投入被排除本身不替代治理成本比较。
\end{semanticsteps}
若价格和进入变化间接改变服务价值或资源耗用，
这些变化必须重新计入对应差额。
前述固定参数比较不能替代新均衡的全部福利核算。
\end{proof}

\section{合规市场与有效竞争的共同边界}

\begin{corollary}[存在正服务利润并排除无效投入的共同模型]
\label{comp:joint-market}
在固定两端模型中，
若合法用途与实际服务符合性成立，
治理满足定义~\ref{comp:deletion}，
零投入下全部非负价格仍为可行，
真实需求和价格收入保持本节给定的模型，
其完整成本已经计入 \(F,c\)，且
\[
 0\le F<\frac{M\theta}{2},
\]
则该模型存在
\[
 p_1=p_2=c+\theta>0,\qquad
 D_1=D_2=M/2>0,\qquad e_1=e_2=0
\]
的均衡，并且两家经济利润严格为正。
若另有逐前缀资金、实际交付和足额结算证书，
这些均衡销售可以形成实际正收入。
\end{corollary}
\begin{certificate}
共同任务和用途准入、
真实服务差异、需求及完整成本参数、
删除操作与收入非干扰、
全部价格与退出偏离的核查，
以及定理~\ref{market:prefix-cash}
要求的资金和实际支付证据。
\end{certificate}
\begin{proof}
\begin{semanticsteps}
\item 对任意价格和真实服务属性，
定理~\ref{comp:zero-effort}
使非零无效投入存在严格改进的删除偏离。
因此选择 \(e_i=0\) 不损失任何可得收入或准入资格。
\item 在零无效投入下，
需求、真实服务和其他成本正是固定两端模型。
价格全域最优反应给出
\(p_i=c+\theta\)，
每家数量为 \(M/2\)。
\item 任意同时改变价格和无效投入的偏离，
都可先删除其无效投入而提高或保持利润。
删除后的价格偏离又不能优于均衡价格，
故这种联合偏离也不能获利。
\item 严格固定成本条件给出
\(M\theta/2-F>0\)，
所以退出收益零也不能超过均衡利润。
均衡由此覆盖价格、投入和退出三类选择。
\item 逐前缀现金归纳使合同支付获得资金支持，
实际交付、验收和结算回执使合同收入成为实际收入。
这些支付前件不由静态均衡计算代替。
\end{semanticsteps}
\end{proof}

上一章的锁定合同组合与本章的竞争需求是不同市场构造。
本章先通过消费者选择和价格偏离确定成交量；
成交以后可以锁定合同并检查资金与结算。
不能在价格偏离时继续把 \(M/2\) 当作不变需求，
也不能把固定合同量定价的可行区间当作竞争均衡证明。

标准复用、真实服务差异和可删除无效投入的治理，
因而可以在共同证书下同时存在。
可证明的结论是有条件的非空合规市场、
选定服务商的完整成本回收、
已定义无效投入的消除及经比较确认的福利变化。
它们不保证任意服务商永久盈利，
也不把压低价格、获得零经济利润或转向更高效率的供给
归入应当排除的行为。
与劳动保护、任务商和权利修复的联合使用，
仍须通过第~\ref{joint:chapter}章的共同资源与语义核验。

\chapter{制度成本、最小兼容扩展与条件最优性}
\label{opt:chapter}

\section{目标约束与经济评价}

用途治理的经济比较必须预先规定保护何种权利、维持何种生产能力、允许何种风险，以及如何计入代价。
当人类参与被列为硬约束时，违反该约束的候选不属于同一可行比较集合；当其被列为加权福利项时，权重必须公开。
两种评价方式均可定义，但不能在比较中途更换。

\begin{definition}[共同可行候选族]
固定制度目标集合
\[
\mathcal G=\{\text{指定人类任务与权利，必要高保障能力，风险容限，资金与资源约束}\}.
\]
令 \(\mathfrak P_{\mathcal G}\) 为满足这些目标、使用同一主体和期间、并通过共同资源核验的制度候选族。
候选可以组合用途许可、事前核验、参与合同、社会保障、责任规则和不同融资安排。
名称不构成可行性或排除性的判据。
\end{definition}

\begin{definition}[完整成本与转移分账]
对候选 \(\pi\)，\(C_{\rm real}(\pi)\) 包括实际消耗的执行、检测、申诉、转型、维护、机会及环境资源。
转移支付 \(T_{ij}\) 单独记账。若评价边界同时包含付款人与收款人，并且没有交易成本，则净货币转移合计为零；
分配权重可以使同一转移产生福利效应，但不得再次作为相同口径的真实资源损失重复计入。
\end{definition}

\begin{lemma}[共同核算边界内的转移抵销]
\label{opt:transfer}
对有限主体集 \(\mathcal I\)，设 \(T_{ij}\) 为 \(i\) 向 \(j\) 的付款，则
\[
\sum_{i\in\mathcal I}\left(\sum_jT_{ji}-\sum_jT_{ij}\right)=0.
\]
\end{lemma}
\begin{certificate}
逐笔付款的唯一标识、双方归属、金额、实际到账与合并抵销表。
\end{certificate}
\begin{proof}
\begin{semanticsteps}
\item 有限双重求和交换指标后，\(\sum_{i,j}T_{ji}=\sum_{i,j}T_{ij}\)，相减为零。
\item 边界外转移、支付手续费和实施资源另列，不能由此抵销。
\end{semanticsteps}
\end{proof}

\section{相对于对照方案的净优势}

\begin{definition}[可比的净效应分解]
在预先固定的福利或净收益评价函数 \(V\) 下，比较用途治理方案 \(\pi_L\) 与对照 \(\pi\)。
将全部差异分为互不重复的三部分：避免的损害 \(H_{\rm avoid}\)、放弃的实际利益 \(B_{\rm forgone}\)，
以及新增真实成本 \(C_{\rm extra}\)，并规定
\[
V(\pi_L)-V(\pi)
=H_{\rm avoid}-B_{\rm forgone}-C_{\rm extra}.
\]
损害与利益包括被评价函数实际计入的劳动、收入、产出、风险及分配后果。跨期量先折算至同一时点。
\end{definition}

\begin{theorem}[条件净优势]
\label{opt:net}
若 \(\pi_L,\pi\in\mathfrak P_{\mathcal G}\)，并有可核验界
\[
H_{\rm avoid}\ge\underline h,
\quad B_{\rm forgone}\le\overline b,
\quad C_{\rm extra}\le\overline c,
\quad\underline h>\overline b+\overline c,
\]
则 \(V(\pi_L)>V(\pi)\)。
\end{theorem}
\begin{certificate}
共同评价函数、权重、反事实口径、效应分组、资金机会成本、误判损失、完整成本界和无重复计量表。
\end{certificate}
\begin{proof}
\begin{semanticsteps}
\item 由分解及界得
\[
V(\pi_L)-V(\pi)
\ge\underline h-\overline b-\overline c>0.
\]
该推导将经验估计与代数结论分离；不使用技术名称或历史融资方式替代任何一项界。
\end{semanticsteps}
\end{proof}

\begin{corollary}[可行候选族内的净收益最优性]
\label{opt:welfare}
若 \(\pi_L\in\mathfrak P_{\mathcal G}\)，且对每个 \(\pi\in\mathfrak P_{\mathcal G}\setminus\{\pi_L\}\) 均有
\(V(\pi_L)-V(\pi)\ge0\) 的比较证书，则 \(\pi_L\) 是该候选族中的最优解。
若所有不等式严格，则在该候选族内唯一最优。
\end{corollary}
\begin{certificate}
完整候选清单、逐候选的可行性和比较证书，以及所有并列值。
\end{certificate}
\begin{proof}
\begin{semanticsteps}
\item 最优性的定义为 \(V(\pi_L)\ge V(\pi)\) 对全部可行 \(\pi\) 成立，证书逐项给出该条件。
\item 严格性排除任何其他候选同值，故唯一。
\end{semanticsteps}
\end{proof}

\begin{remark}[成熟条件在比较中的作用]
运行独立和财务独立可以降低用途收缩导致的必要高保障能力损失，从而降低 \(B_{\rm forgone}\) 或部分转型成本。
它们并不自动给出 \(H_{\rm avoid}>B_{\rm forgone}+C_{\rm extra}\)。
高保障需求、普通用户利益、执法成本及其他可行制度仍须进入比较。
\end{remark}

\section{最小兼容制度变更}

\begin{definition}[制度基线与变更支集]
固定有限制度坐标集合 \(\mathcal U_0\)、基线 \(\Sigma_0=(\sigma_u^0)_{u\in\mathcal U_0}\)，
以及已确认失灵集合 \(\mathcal F\)。令
\[
\Delta(\Sigma;\Sigma_0)=\{u:\sigma_u\ne\sigma_u^0\},
\qquad
d_w(\Sigma,\Sigma_0)=\sum_{u\in\Delta(\Sigma;\Sigma_0)}w_u,
\quad w_u>0.
\]
对直接修复坐标 \(\mathcal R(\mathcal F)\) 及经证据确认的必要依赖 \(\operatorname{Dep}(\mathcal R(\mathcal F))\)，
兼容候选满足
\[
\Delta(\Sigma;\Sigma_0)
\subseteq\mathcal R(\mathcal F)\cup\operatorname{Dep}(\mathcal R(\mathcal F)).
\]
\end{definition}

市场价格、工资、产权、合同、结算、税收与社会保障中有效的坐标可以被列为保留基线。
新增用途规则、真实参与资格及履约修复只在已确认的失灵与必要依赖处改变制度。
这一接口采用参与经济形式系统中的最小兼容扩展结构\cite{InternalEconomics}。

\begin{theorem}[有限候选族的最小兼容解]
\label{opt:minimal}
设满足共同目标、修复要求和兼容条件的候选族 \(\mathfrak P_{\mathcal G}\) 有限非空。
按字典序最小化
\[
\left(\lvert\Delta(\Sigma;\Sigma_0)\rvert,
d_w(\Sigma,\Sigma_0)\right),
\]
则最小解集合非空。任一最小解保持所有与修复及必要依赖无关的基线坐标。
如果在全部并列最小解上再最小化有限的 \(C_{\rm real}\)，第三层最小解也存在。
\end{theorem}
\begin{certificate}
声明基线、失灵证据、依赖闭包、有限候选清单、逐候选共同可行性、预先固定权重、三层比较表和全部并列解。
\end{certificate}
\begin{proof}
\begin{semanticsteps}
\item 有限非空集合中的非负整数 \(\lvert\Delta\rvert\) 存在最小值。
\item 限制到取该值的非空子集，有限实数 \(d_w\) 也有最小值。
\item 再限制到前两层并列解，有限真实成本仍有最小值。
\item 兼容条件对全部候选成立，因此对每个最小解成立。
\end{semanticsteps}
\end{proof}

\begin{corollary}[支集不可严格缩小]
若 \(\Sigma_\star\) 为上述最小解，则不存在可行候选 \(\widehat\Sigma\) 满足
\(\Delta(\widehat\Sigma;\Sigma_0)\subsetneq\Delta(\Sigma_\star;\Sigma_0)\)。
\end{corollary}
\begin{certificate}
候选支集表和第一层最小值。
\end{certificate}
\begin{proof}
\begin{semanticsteps}
\item 反设存在可行 \(\widehat\Sigma\) 使
\(\Delta(\widehat\Sigma;\Sigma_0)\subsetneq\Delta(\Sigma_\star;\Sigma_0)\)。
\item 两支集有限，严格包含推出
\(\lvert\Delta(\widehat\Sigma;\Sigma_0)\rvert<\lvert\Delta(\Sigma_\star;\Sigma_0)\rvert\)。
\item 可行性使 \(\widehat\Sigma\) 属于同一比较族，第二步与第一层最小性矛盾。故不存在该候选。
\end{semanticsteps}
\end{proof}

\begin{proposition}[指定方案达到最小值的下界证书]
\label{opt:lowerbound}
设有两两不交的制度坐标集合 \(U_1,\ldots,U_k\)，且每个可行候选都必须改变每个 \(U_j\) 中至少一个坐标。
则所有可行候选满足 \(\lvert\Delta\rvert\ge k\)。
若指定候选 \(\Sigma_L\) 恰改变 \(k\) 个坐标，并在所有数量同为 \(k\) 的候选中达到最小 \(d_w\)，
则它是前两层的最小兼容解；若其真实成本在并列解中也最小，则达到第三层最小性。
\end{proposition}
\begin{certificate}
每一组坐标的必要性证明、组间不交检查、指定方案的可行实现，以及同层剩余候选的比较表。
\end{certificate}
\begin{proof}
\begin{semanticsteps}
\item 每个候选从 \(k\) 个不交集合中各改变至少一个坐标，故改变数量至少为 \(k\)。
\item 指定方案达到该下界，继而按后两项比较排除更优同层候选，得到字典序最小性。
\end{semanticsteps}
\end{proof}

\begin{remark}[评价口径的独立性]
制度变更最小性并不等于福利最大化。若要同时声称两者成立，需要同一方案分别通过定理~\ref{opt:net}
或推论~\ref{opt:welfare} 的净效应比较，以及定理~\ref{opt:minimal} 的变更比较。
最小兼容候选可以含有社会保障、公共供给或其他收入制度；只有明确的预算、资源或权利冲突限制这些组合。
\end{remark}

\section{可核验的制度比较边界}

\begin{proposition}[候选扩大时的最优值单调性]
若 \(\mathfrak P\subseteq\mathfrak P'\)，则对最大化净收益有
\[
\sup_{\pi\in\mathfrak P}V(\pi)
\le\sup_{\pi\in\mathfrak P'}V(\pi).
\]
原最优方案在扩大的候选族中仍最优，当且仅当新增候选均不取得更高值。
\end{proposition}
\begin{certificate}
原、新候选范围及新增候选的共同口径比较。
\end{certificate}
\begin{proof}
\begin{semanticsteps}
\item 对任意 \(\pi\in\mathfrak P\)，有 \(\pi\in\mathfrak P'\)，所以
\(V(\pi)\le\sup_{\rho\in\mathfrak P'}V(\rho)\)。对左侧取上确界得第一式。
\item 设原最优方案为 \(\pi_\star\)。它已满足
\(\forall\pi\in\mathfrak P,\ V(\pi)\le V(\pi_\star)\)。
\item 若新域仍最优，则对其子集 \(\mathfrak P'\setminus\mathfrak P\) 同样成立该不等式。
\item 反向，若新增子集均满足该界，将它与原域的界合取，并用
\(\mathfrak P'=\mathfrak P\cup(\mathfrak P'\setminus\mathfrak P)\)，得到新域全称最优条件。
\end{semanticsteps}
\end{proof}

\begin{remark}[制度结论的可修订性]
新的融资方式、较低成本的核验、不同的收入制度或新的需求证据可以改变候选边界。
制度比较证书随这些实质条件更新；有限候选内的最优性不扩张为未经比较的全球、永久或跨目标最优性。
\end{remark}

\chapter{联合可行性、有限构造与总定理}
\label{joint:chapter}

\section{共同证书的资源与语义接口}

前述各章分别处理运行、许可、劳动、收入、表示、修复与制度比较。
它们共同成立的关键，在于同一实施安排能否同时满足所有证书，而非各自找到一个彼此不相容的安排。
本章先声明联合验收条件，再给出有限模型中的非空见证。

\begin{definition}[联合证书]
\label{joint:certificate}
在同一声明域 \(\mathfrak D\) 下，联合证书记为
\[
\mathfrak C_\star=(\mathfrak C_M,\mathfrak C_U,\mathfrak C_Q,
\mathfrak C_B,\mathfrak C_R,\mathfrak C_F,\mathfrak C_J,\mathfrak C_{\rm joint}).
\]
各分量承担以下义务。
\begin{enumerate}[label=\textnormal{\arabic*.}]
  \item \(\mathfrak C_M\)：构造、运行、更新与财务四项成熟条件成立；存在不依赖受限常规替代扩张收入的期限内维持路径。
  \item \(\mathfrak C_U\)：用途准入覆盖所声明的实现载体与实际替代路径；实际覆盖及规避参数具有逐任务或劳动量加权界。
  \item \(\mathfrak C_Q\)：定义~\ref{q:three-institutions} 中三类制度的共同任务观察、动作、全部响应与目标精确对应；动态商有有限饱和或稳定性证书；共同任务像及共同链观察像一致；状态与链表示接口交换，并存在保持全部声明资格及逐人权利观察的合法读回。
  \item \(\mathfrak C_B\)：有效参与身份和过程、人数上限、正基本份额、到账资金、逐项义务、支付期限和再接入条件通过核验。
  \item \(\mathfrak C_R\)：指定缺陷满足 \(B\omega=0\)，并具有旧操作排除证书、新增修复解及实际非负执行；扩展复形采用逐人账本构造并有收缩 \(d_Ct+td_C=\id_{C^+}\)；\(E=P\oplus_\kappa C^+\) 的耦合满足 \(d_P\kappa+\kappa d_C=0\)。
  \item \(\mathfrak C_F\)：高保障任务存在旧不可达而新可达的真实目标；该任务域内每条旧允许路径在新制度中有相同初态、相同终态且保持共同约束的实现。
  \item \(\mathfrak C_J\)：净劳动效应中避免替代的下界超过其他反馈的净损失上界；劳动小时与就业人数采用各自证书。
  \item \(\mathfrak C_{\rm joint}\)：全部模块使用同一主体和规则索引；资金、设备、人员、时间、风险与既有权利不存在重复承诺或冲突。
\end{enumerate}
\end{definition}

\begin{lemma}[有限联合执行的合取保持]
\label{joint:preservation}
设 \(I_1,\ldots,I_k\) 为联合证书要求的不变量。同一执行安排满足各初态条件，且每个实际动作对全部登记响应同时保持各 \(I_j\)。
则任意允许有限执行保持 \(\bigwedge_{j=1}^kI_j\)。
\end{lemma}
\begin{certificate}
共同初态、全局调度、逐前缀未重复资源、单步合取证书及实现回执。
\end{certificate}
\begin{proof}
\begin{semanticsteps}
\item 令 \(I_\star(x)=\bigwedge_{j=1}^kI_j(x)\)。共同初态证书逐项给出 \(I_j(x_0)\)，合取引入得到 \(I_\star(x_0)\)。
\item 假定 \(I_\star(x_n)\)。合取消去得到每个 \(I_j(x_n)\)。
\item 同一实际动作 \(a_n\) 和响应 \(\xi_n\) 代入各模块单步证书，得到
\(I_j(\delta(x_n,a_n,\xi_n))\) 对全部 \(j\) 成立。
\item 以 \(x_{n+1}=\delta(x_n,a_n,\xi_n)\) 替换并合取，得到 \(I_\star(x_{n+1})\)。
\item 数学归纳完成；共同资源与调度证书确保所用动作属于各模块允许集合的交集。
\end{semanticsteps}
\end{proof}

\section{用途治理与劳动收入保障的总定理}

\begin{theorem}[阶段转换、任务保护与可执行修复的共同实现]
\label{joint:main}
若定义~\ref{joint:certificate} 的证书在共同见证上全部成立，令
\[
\epsilon_t=1-\underline\theta_t(1-\overline\varepsilon_t)<1,
\]
则在各自声明任务范围内有
\[
\boxed{\begin{aligned}
&D_t^G\le\epsilon_tD_t^0,\\
&H_t^G-H_t^0=A_t^{\rm dir}+G_t-C_t>0,\\
&y_{it}\ge\beta_t>0\qquad(i\in\mathcal E_t^\ast),\\
&\mathcal R_-(x_0)\subsetneq\mathcal R_+(x_0),\\
&Q_{\mathsf F}\cong Q_{\mathsf A}\cong Q_{\mathsf H}\quad\text{于已覆盖的共同任务},\\
&\overline C^{\mathsf F}_\bullet\cong\overline C^{\mathsf A}_\bullet\cong\overline C^{\mathsf H}_\bullet,\\
&[\omega]_{C^-}\ne0,\qquad[\omega]_{C^+}=0,\\
&H_\ast(E)\cong H_\ast(P).
\end{aligned}}
\]
其中 \(\mathsf F,\mathsf A,\mathsf H\) 均限于已覆盖的共同子任务；
\(y_{it}\) 为在约定期限内实际到账的净支付；\(C^-,C^+\) 为保留逐人权利的修复复形，\(P\) 为合法参与的阶段循环，
\(E\) 为相容耦合后的复形。已经证明的再接入条件在所覆盖的有限期间内逐期接续。
\end{theorem}
\begin{certificate}
联合证书全体、每项结论的任务索引、对偶排除与原始执行见证、逐人支付回执、商转移和链映射矩阵，
以及实际资源和期限的共同验收记录。
\end{certificate}
\begin{proof}
\begin{semanticsteps}
\item 定理~\ref{lab:thm-direct-bound} 给出直接替代界，定理~\ref{lab:thm-net-effects} 给出无重复的净劳动恒等式；
\(\mathfrak C_J\) 使其严格为正。
\item 定理~\ref{part:thm-allocation} 与命题~\ref{part:prop-timely-payment}
的资金、真实性和实际支付证书给出逐人到账下界，
偿付与接续归纳保证已覆盖期间的履约。
\item 定理~\ref{front:strict} 给出高保障任务的严格可达扩张。
\item 定理~\ref{q:dynamic} 与定理~\ref{q:common-image} 给出三侧状态商同构；
定理~\ref{ch:maximal-kernel} 与定理~\ref{ch:common-image} 给出链商同构，命题~\ref{ch:readback} 保证任务读回。
\item 旧缺陷的对偶证书排除 \(\omega\in\im A_-\)，新增合法修复给出 \(\omega\in\im A_+\) 及实际执行。
\item 引理~\ref{cyc:contraction} 的显式收缩和定理~\ref{cyc:retraction} 的相容耦合给出 \(E\) 到 \(P\) 的链同伦等价，从而同调同构。
\item 所有前件都由同一共同见证满足；引理~\ref{joint:preservation} 保持其执行条件。
\item 因此可对各结论作合取，而不发生主体、时点或资源的偷换。
\item 有限再接入由各期真实后继与下一期准入状态相同而逐期归纳。
\end{semanticsteps}
\end{proof}

\begin{corollary}[阶段转换解除特定扩张的必要性]
在总定理的成熟及共同资源条件下，存在期限 \(T\) 内的可行安排，能够维持指定高保障任务及必要研究辅助，
同时不采用被限制的常规替代扩张。故该扩张不是该期限内维持这些任务的必要条件。
\end{corollary}
\begin{certificate}
独立执行程序、研究与维护计划、可用资金来源、全部前缀收支和任务验收记录。
\end{certificate}
\begin{proof}
\begin{semanticsteps}
\item 共同见证给出一个实际安排 \(\pi_\star\)，满足
\(\operatorname{Maintain}_T(\pi_\star)=1\) 且 \(\operatorname{Expand}_C(\pi_\star)=0\)。
\item 若扩张为必要条件，则
\(\forall\pi,\operatorname{Maintain}_T(\pi)\Rightarrow\operatorname{Expand}_C(\pi)\)。
\item 令 \(\pi=\pi_\star\) 并肯定前件，将推出 \(\operatorname{Expand}_C(\pi_\star)=1\)，与第一步矛盾。
\item 因此该全称必要性命题为假；结论范围是证书覆盖的任务和期限。
\end{semanticsteps}
\end{proof}

\begin{corollary}[附加比较证书下的制度最优性]
若同一可行治理方案进一步通过声明候选族内的净收益比较，并且提交下界与实现证书，
确认其实际达到已声明的字典序最小值，则它在对应目标与评价顺序下达到各自最优性。
\end{corollary}
\begin{certificate}
除总证书外，增加候选全集、完整成本与收益、必要变更坐标的下界，以及全部并列方案。
\end{certificate}
\begin{proof}
\begin{semanticsteps}
\item 总证书确认同一方案 \(\pi_\star\in\mathfrak P_{\mathcal G}\)。
\item 净收益比较证书给出
\(\forall\pi\in\mathfrak P_{\mathcal G},V(\pi_\star)\ge V(\pi)\)，由最优性定义得到对应目标下最优。
\item 下界证书给出 \(\lvert\Delta(\pi)\rvert\ge k\)，指定方案达到 \(k\)，并在并列候选上分别达到后续最小值。
\item 逐层应用字典序定义得到最小兼容性。两类结论采用各自评价顺序，并由同一个 \(\pi_\star\) 联合实现。
\end{semanticsteps}
\end{proof}

\section{劳动权利、公共探索与可验证服务的联合实现}
\label{joint:service-section}

\begin{definition}[公共服务扩展证书]
\label{joint:service-certificate}
在定义~\ref{joint:certificate} 的共同见证上，事前登记一个非空已接纳服务请求域 \(\mathcal F_\star\)、
非空活跃提供者集合 \(\mathcal J_\star\)，以及下列附加证书。
对请求 \(r\)，以 \(H_r^{\rm keep}\) 表示在声明期限内本人持续要求交付、
相应许可及必要任务前件持续成立。
该谓词不读取是否已经按期完成的结果，不能据此事后删除未完成请求。
实际见证中至少有一个请求满足它。
\begin{enumerate}
\item 每个实际服务作用通过当前有效授权，证书与实际主体、用途、实现、输入、资源和规则绑定；
拒绝、补证、退款和复核保持其各自既有权利。
\item 对 \(r\in\mathcal F_\star\)，在 \(H_r^{\rm keep}\) 成立的执行范围中，
实际策略使用已验证的进展动作，有限秩、有界调度和期限相容；
服务验收指向合同约定结果，而非仅指流程终止。
本人合法撤回或许可条件改变时，改按相应关闭、退款和复核目标提交处置证书，
不把该处置记为原业务结果已经交付。
\item 真实劳动合同具有自愿进入退出、公开规则、既得债权、可达复核与按期付款的实际实现；
探索资格不以该劳动合同的评分或持续参加为条件。
\item 对每类同质可分探索服务，采用定义~\ref{access:divisible} 的正确资格、真实需求和费用口径，
登记 \(n_i\ge0\)、实际容量 \(K\ge0\)、合同价格 \(p>0\) 和到账预算 \(B_E\ge0\)，
并使用其零需求分支。配额的实际取得另满足推论~\ref{access:actual-delivery} 的资金、容量、调度和交付条件。
财富或劳动评分的比较采用定义~\ref{access:paired-witness} 的两份合法可达执行见证，
保持其全部必要准入、证据、资金和服务条件；比较实际交付时，两份执行中的自愿接受与撤回选择相同。
\item 具体标准实现与可验证任务实现具有相同的共同任务像，动作和全部响应覆盖、后继交换、状态读回及观察链商均通过验收。
\item 每个 \(j\in\mathcal J_\star\) 具有正实际交付量 \(Q_j\)、真实收入 \(R_j\)、完整经济成本 \(C_j\)，
满足 \(R_j\ge C_j\)；同时具有全部到期支出的逐前缀现金、设施和实际支付证书。
\item 已定义的无效投入可以被删除而不改变生产、准入和真实服务属性，且不直接或间接改变收入；
其正水平消耗严格正资源。所采用的定价或维护合同另有参与与成本回收证书。
\item 所有采购、支付、补助和清偿使用唯一交易标识；各模块共同占用的资金、设备、人员和时段没有重复承诺。
\end{enumerate}
\end{definition}

\begin{theorem}[公共服务、劳动保障和持续供给的共同实现]
\label{joint:service-main}
若定义~\ref{joint:certificate} 与定义~\ref{joint:service-certificate} 在同一个 \(z_\star\) 上成立，
且新增高保障目标另有定理~\ref{front:taskimage} 的结果像排除证书，
则原联合结论保持，并且在各自已覆盖任务上同时成立：
\[
\begin{aligned}
&\operatorname{ActualAct}(r,a)
 \Longrightarrow\operatorname{CurrentPermit}(r,a),\\
&r\in\mathcal F_\star\land H_r^{\rm keep}
 \Longrightarrow\operatorname{DeliveredByDeadline}(r),\\
&g(\mathcal R_-(x_0))\subsetneq g(\mathcal R_+(x_0)),\\
&\operatorname{ValidLabour}(i,t)
 \Longrightarrow y_{it}=b_{it}\ge\beta_t>0,\\
&N=\sum_i n_i>0
 \Longrightarrow a_i=\alpha n_i,\quad
 \alpha=\min\{1,K/N,B_E/(pN)\},\\
&n_i=n_j\Longrightarrow a_i=a_j,\\
&Q_W\cong Q_J,\qquad
 \overline C^W_\bullet\cong\overline C^J_\bullet,\\
&H_\ast(\overline C^W)\cong H_\ast(\overline C^J),\\
&\forall j\in\mathcal J_\star,\quad Q_j>0,\quad\Pi_j=R_j-C_j\ge0,\\
&e_j^\star=0\quad\text{于已定义的无效投入变量}.
\end{aligned}
\]
零总需求时，探索分配为零；零可用预算或容量时，不签发正交付承诺。
制度劳动权利、逐人余额及相应复核路径在声明有限执行内保持。
这里 \(e_j^\star\) 表示固定真实服务及其余行为后的利润最大化投入，
不单独声明全部生产与价格策略存在全局均衡。
\end{theorem}
\begin{certificate}
保存同一 \(z_\star\) 上的基础与附加证书，逐原子的前置与后置谓词，
实际选择的执行动作，全部响应及资源消耗，逐人义务和最终到账，需求分配与实际交付，
商映射及链矩阵，以及逐提供者成本和现金表。
理论上的最粗表示另须具备证据恢复和有限判定或可执行算法证书，才能进入实际签发流程。
\end{certificate}
\begin{proof}
\begin{semanticsteps}
\item 共同见证对基础证书的投影满足定义~\ref{joint:certificate}。
以肯定前件规则应用定理~\ref{joint:main}，得到直接替代、净劳动、原参与收入、分配修复和阶段循环结论。
\item 对每个服务作用，由入口覆盖取得当前检查记录；
由检查绑定和有效证据取得准入谓词；由实现交换取得实际作用符合性。
将三项合取，得到 \(\operatorname{ActualAct}\Rightarrow\operatorname{CurrentPermit}\)。
\item 固定 \(r\in\mathcal F_\star\) 且 \(H_r^{\rm keep}\) 成立的执行，
实际策略使正秩下降，其他允许步骤不增加秩。
有界调度与非停滞证书将定理~\ref{tra:liveness} 的步数、时间和资源界应用到该请求，
目标原像确认按期完成的是合同交付。
\item 新目标结果不属于旧结果像，实际执行使其属于新结果像，旧路径保持使旧像包含于新像。
三项谓词合取得到严格结果扩张。
\item 劳动真实性和真实贡献恢复应付合同，资金与结算记录给出
\(y_{it}=b_{it}\ge\beta_t\)。自愿退出、既得债权和复核的单步规则，
与金额规则共同满足定理~\ref{base:induction}，故所列制度权利逐前缀保持。
\item 对正总需求的探索类别，分别使用
\(\alpha\le1\)、\(\alpha N\le K\)、\(p\alpha N\le B_E\)，
得到需求上界、资源上界和支出上界。同一 \(n_i\) 取得同一配额。
两份合法执行见证保持所有有效输入，逐步交换使财富和劳动评分变化不改变该任务的配额及其交付。
\item 两种标准服务实现使用相同观察和动作响应语义。
共同像证书与定理~\ref{q:common-image} 给出状态商同构；
图链满射及边界交换给出观察链商同构，进而诱导同调同构。
\item 逐提供者的实际交易给出 \(Q_j>0\)。将同一交易的收入与完整成本相减，
以 \(R_j\ge C_j\) 得到 \(\Pi_j\ge0\)。逐前缀现金和实际支付证书另行确认期限内履约。
\item 对任意可行的正无效投入，删除映射保持收入和生产可行性，严格减少资源成本。
因此零投入严格优于任意正投入。
\item 上述每一步使用同一主体、请求、时点与资源见证。
引理~\ref{base:joint} 允许对这些已在同一 \(z_\star\) 上成立的谓词作合取。
步骤链的逻辑健全性由定理~\ref{atom:proof-induction} 保证。
\end{semanticsteps}
\end{proof}

\begin{corollary}[表示改善与经济净优势的分别验收]
\label{joint:service-comparison}
在定理~\ref{joint:service-main} 的条件上，若结果表示遗漏某项必须实施的治理决定，
且存在同结果而决定不同的有效证据对，则任务充分表示严格扩大可表示的该类决定集合。
若同一制度另有完整净收益比较证书，则进一步成立其声明比较中的经济净优势。
\end{corollary}
\begin{certificate}
结果纤维、真实治理类型、证据恢复、决定因子化、严格区分输入对，
以及对照制度、全部真实资源成本、分配权重和效应无重复表。
\end{certificate}
\begin{proof}
\begin{semanticsteps}
\item 设原结果表示为 \(q_0\)，真实类型为 \(\tau\)。投影给出
\(q_0=\operatorname{pr}_1(q_0,\tau)\)，故旧表示上的决定均可在联合表示上表达。
\item 严格见证满足 \(q_0(x)=q_0(y)\) 而 \(E(x)\ne E(y)\)。
引理~\ref{q:fiber} 排除 \(E\) 通过 \(q_0\) 因子化；联合类型证书使其通过新表示因子化。
\item 旧决定集合的包含与新增决定见证共同给出严格扩张。
可执行算法和证据条件把该表示结论落实到声明任务的实际判定。
\item 经济净优势另由定理~\ref{opt:net} 的同口径下界推出；其成本边界计入表示恢复、验证、维护和实施费用。
\end{semanticsteps}
\end{proof}


\section{一个联合条件非空的有限经济模型}

以下构造中的数值均为模型赋值，货币单位和劳动小时单位经过归一。
它展示条件的相容性，不作为任何地区、产业、政策效果或技术成熟状态的经验估计。
模型固定两期，所有状态坐标取有限整数值；所有动作均有明确前置条件和确定后继。

\subsection{资金、时间与技术承接}

\begin{definition}[两期共同配置]
各资金账户事先分别持有下列余额，账户用途互不重复。
\begin{center}
\begin{tabular}{p{0.58\textwidth}r}
\toprule
账户与用途 & 初始余额\\
\midrule
常规任务购买者，覆盖实际有酬劳动 & 250\\
高保障任务运行、更新、验证、治理及研究辅助 & 100\\
两期真实参与合同的支付资金 & 40\\
两期参与核验及结算成本 & 10\\
历史缺陷修复的担保余额 & 5\\
两名参与者原有现金，各为 5 & 10\\
无资格账户中的可依法追回金额 & 5\\
\midrule
全部账户合计 & 420\\
\bottomrule
\end{tabular}
\end{center}
高保障任务每期实际支出 40，其中运行 20，更新、验证、治理和研究辅助各 5。
两期后保留储备 20。所需实物投入与人员时段已分别预留。
程序在期初之前已经完成构造与验收，历史构造费用已经清付，不存在带入这两期的未付构造义务。
\end{definition}

\begin{proposition}[两期财务独立与运行成熟]
高保障任务具有独立存储并执行的确定性程序，更新版本保持原受保护观察，所有任务在一轮内验收。
其两期现金前缀均可履约，且不需要常规替代扩张收入。
\end{proposition}
\begin{certificate}
令能力状态 \(s=(s_1,s_2)\in\{0,1\}^2\)，初态 \((0,0)\)。旧动作仅能将 \(s_1\) 从 0 设为 1；
已验证新增程序可以在许可成立时将 \(s_2\) 从 0 设为 1。状态约束为 \(0\le s_j\le1\)。
不合法动作采用恒等回退，完成目标 \(s_2=1\) 的执行使用确定的合法新增动作。
新旧版本对这些状态及许可谓词采用相同转移。现金前缀为
\[
100\ge40+20,\qquad100=40+40+20.
\]
\end{certificate}
\begin{proof}
\begin{semanticsteps}
\item 有限状态的每项允许转移均保持 \(\{0,1\}^2\)，目标程序一次执行即将 \(s_2\) 设为 1，
既有 \(s_1\) 动作继续可用。
\item 执行只读取固定程序、状态和许可，不调用候选生成器。
\item 两个前缀收支不等式覆盖全部声明费用与储备，因此资金独立。
\item 旧路径始终满足 \(s_2=0\)，新路径达到 \(s_2=1\)，从而在这个抽象高保障任务中严格扩张可达能力。
\end{semanticsteps}
\end{proof}

\subsection{直接替代、需求与净劳动量}

\begin{proposition}[正的净劳动效应见证]
设共同期限内无替代时基础劳动量为 200，基线直接替代为 50，故实际基线劳动为 150。
制度覆盖全部声明替代任务，实际剩余替代为 5，新增需求及任务产生 6 个劳动小时，其他反馈减少 8 个劳动小时。
则新制度实际劳动量为 193，比基线多 43。
\end{proposition}
\begin{certificate}
\[
\underline\theta=1,\quad\overline\varepsilon=\tfrac1{10},
\quad D^0=50,\quad D^G=5,
\quad A^{\rm dir}=45,\quad G=6,\quad C=8.
\]
固定次序的中间劳动量为 \(150\to195\to201\to193\)。
假定单位工资为 1，购买者资金 250 覆盖新工资 193，实际可用劳动时间不少于 201，且各订单在共同调度中不重叠。
这里的 150、193 和基础量 200 均属于常规任务劳动域。
高保障运行、研究辅助、参与及治理模块的人类工时在两个比较情景中相同，并使用另行预留的时段，
不挤占上述常规任务容量；将这些不变工时合并以后，总劳动域的差额仍为 43。
\end{certificate}
\begin{proof}
\begin{semanticsteps}
\item 直接替代界为 \(D^G\le0.1D^0=5\)。
\item 中间差分分别为 \(45,6,-8\)，因此
\[
H^G-H^0=45+6-8=43>0,
\qquad H^G=200-5+6-8=193.
\]
规避引起的 5 已包含于 \(D^G\)，不重复扣除。
\item 工资需求小于实际购买资金；容量与调度证书将需求转为真实完成工时。
\end{semanticsteps}
\end{proof}

该构造是有资金支持的有限任务模型，未求解完整一般均衡。
它证明所列条件可以相容；实际政策效应仍需对需求、价格、劳动投入及反事实次序分别识别。

\subsection{正基本份额、履约与再接入}

\begin{proposition}[两期真实参与与实际收入]
每期接纳同两名真实参与者，\(N^{\max}=2\)、\(\beta=4\)、差异池 \(D=12\)，贡献权重为 \((1/3,2/3)\)。
两人每期按时实际得到 \((8,12)\)，两期总支付为 40。第一期支付足以保持第二期最低接入条件。
\end{proposition}
\begin{certificate}
每期已保留资金满足 \(2\cdot4+12=20\)。两期身份、真实履约、唯一申领及结算均通过。
各人每期必要现金成本为 3，下一期最低接入金额为 5，第二期名额、时间与执行路径均已保留。
参与治理支出每期 5，由独立的 10 覆盖。
\end{certificate}
\begin{proof}
\begin{semanticsteps}
\item \[
b_1=4+12/3=8,\qquad b_2=4+2\cdot12/3=12.
\]
逐人金额不小于 4，每期总额 20，两期总额 40，均在专用资金之内。
\item 第一期有 \(8-3=5\)、\(12-3=9\)，满足第二期最低接入要求，且第二期路径从第一期真实后继继续。
\item 两期末不据此承诺第三期；下一核算期限需要新的资金和容量证书。
\end{semanticsteps}
\end{proof}

\subsection{历史缺陷与新增修复方向}

\begin{proposition}[逐人权利、担保保持与完整追偿]
在账户顺序 \((H_1,H_2,O,X)\) 下，历史侵占额为 \((2,3)\)，要求修复后保持担保账户 \(O\) 原余额。
旧支付方向不能实现全部目标；增加追回方向后可以完成，且全部余额非负。
\end{proposition}
\begin{certificate}
\[
B=(1\ 1\ 1\ 1),\quad
A_-=\begin{pmatrix}1&0\\0&1\\-1&-1\\0&0\end{pmatrix},\quad
A_+=\begin{pmatrix}1&0&0\\0&1&0\\-1&-1&1\\0&0&-1\end{pmatrix},
\]
\[
\omega=(2,3,0,-5)^\top,
\quad\lambda_X=(0,0,0,1),
\quad A_+(2,3,5)^\top=\omega.
\]
修复子账户的状态顺序为
\[
(5,5,5,5)\longrightarrow(5,5,10,0)
\longrightarrow(7,5,8,0)\longrightarrow(7,8,5,0).
\]
\end{certificate}
\begin{proof}
\begin{semanticsteps}
\item \(BA_-=BA_+=0\)，\(\lambda_XA_-=0\)，而 \(\lambda_X\omega=-5\ne0\)，故旧操作无法实现该完整目标。
\item 新增解和给定顺序完成全部变化，所有中间余额非负，\(O\) 终值仍为 5。
\item \(\dim\ker B=3\)、\(\rank A_-=2\)、\(\rank A_+=3\)，所以指定缺陷由非零同调类变成边界。
\item 历史追偿的 5 单独核算，不再作为本期新劳动工资或参与合同收入计入。
\end{semanticsteps}
\end{proof}

\subsection{共同任务商与阶段同调}

\begin{proposition}[共同表示的非空构造]
将上述有限模型的全部任务相关状态组成 \(Q\)，保留主体、期次、资格、资源、应付、到账、时限及合法动作。
三种抽象实现取 \(X_j=Q\times\{0,1\}\)，每个实现只在无关记录位上不同；它们具有相同的任务商及兼容的链表示。
\end{proposition}
\begin{certificate}
取 \(q_j(x,z)=x\)，以及对任意合法登记标签 \(\ell\) 的转移
\[
\Phi_{j,\ell}(x,z)=\bigl(\overline\Phi_\ell(x),\,z\mathbin{\oplus}b_{j,\ell}\bigr),
\qquad b_{j,\ell}\in\{0,1\}.
\]
共同初态与目标分别登记为 \(Q_0,G_Q\subseteq Q\)，并取
\[
X_{j0}=Q_0\times\{0,1\},\qquad
G_j=G_Q\times\{0,1\}.
\]
实际动作与响应域采用
\(A_j(x,z)=A_Q(x)\)、\(\Xi_j((x,z),a)=\Xi_Q(x,a)\)。
观察族含所有 \(Q\) 类的指标函数；各实现另加入共同的未定义状态，非法后继送入该状态，并使该状态固定。
将共同状态及合法转移图的自由有理链表示作为三侧模板，采用共同顶点与共同标签边的逐生成元坐标指示，
保留源、标签、响应与目标观察。
\end{certificate}
\begin{proof}
\begin{semanticsteps}
\item 每个 \(q_j\) 满射且 \(q_j\Phi_{j,\ell}=\overline\Phi_\ell q_j\)。
\item 全部指标函数区分不同 \(Q\) 状态，故任何充分商不能进一步合并它们；同一 \(Q\) 状态的两个记录位则有相同观察及后继，
可以合并。
\item 因此共同最粗任务商恰为 \(Q\)。
\item 图链映射保持源与目标差，因而与边界交换；三侧共同链观察像相同，给出链商同构。
\item 这一构造只证明抽象共同任务表示的非空性；具体行业须另交全部语义对应证书。
\end{semanticsteps}
\end{proof}

\begin{proposition}[非零耦合下的阶段循环保持]
对两名参与者的扩展修复复形，设 \(z\) 为四阶段循环生成元。
取 \(\kappa_2(u_1,u_2,v)=u_1z\)，其余 \(\kappa_n=0\)。则 \(\kappa\) 相容且
\(H_\ast(P\oplus_\kappa C^+)\cong H_\ast(P)\)。
\end{proposition}
\begin{certificate}
\(d_Pz=0\)，修复复形的显式收缩，以及度数为 \(-1\) 的上述 \(\kappa\)。
这些系数表示规范化记录的耦合，不将货币与物理量直接相加。
\end{certificate}
\begin{proof}
\begin{semanticsteps}
\item 唯一可能非零的相容式为 \(d_P\kappa_2+\kappa_1d_C=d_P(u_1z)=0\)。
\item 故相容耦合定理的前件成立，显式回缩与链同伦给出同调同构。
\item 两期真实经济状态中的时间和余额发生变化；保留的是阶段循环类，并非全部状态回到原点。
\item 实际路径核验第一期至第二期的接续；第二期末为本期限终点，不由阶段环推导第三期资金或接纳义务。
\end{semanticsteps}
\end{proof}

\section{公共探索与成本回收的有限共同实例}
\label{joint:service-example}

以下为有限经济模型的参数证书。它与前述数值构造采用互不重复的资源子账户和执行时段，
用于证明公共申请、用途限制、按需探索与服务提供者持续供给的条件可以共同满足。
其中的有限任务表是数学实例，不作为现实企业已经完成标准认证或部署的经验证据。

\begin{definition}[实例的需求、预算与提供者]
设一个服务期间内有两名合格探索申请者，核验需求为 \((3,5)\) 个同类可分单位。
实际可交付容量为 \(K=8\)，探索专款为 \(B_E=32\)。两家服务商各锁定四个单位的合法订单，
所选单价 \(p=3\)，在整个 \([0,7/2]\) 报价区间内订单均有真实需求与付款覆盖。
每家单位交付成本为 \(c=1\)，固定经济成本为 \(F=6\)，期初工作资金为 10。
固定成本已包括该期间必要验证、维护、异常请求处理和正常资本成本。
独立预算不再将这些项目重复计入其他成本。
\end{definition}

\begin{proposition}[按需配额与正成本回收的同一见证]
\label{joint:service-numbers}
上述实例可同时实现全部探索需求、个人零账单支付、两家服务商正经济利润及全部交付成本的前缀覆盖。
\end{proposition}
\begin{certificate}
将八个实际请求分别登记，第一申请者的三个请求分配给第一家服务商，
第二申请者的一个请求分配给第一家、四个请求分配给第二家。
探索基金按每单位 3 支付。每家先用期初工作资金支付固定费用 6，
再按四个单位各支付交付费用 1；八份任务验收后分别取得 12 的最终服务收入。
价格、合同量、实际输入、资源时段和接收回执均使用相同请求索引。
\end{certificate}
\begin{proof}
\begin{semanticsteps}
\item 总需求为 \(N=3+5=8\)。代入分配规则得
\[
\alpha=\min\{1,8/8,32/(3\cdot8)\}=1,
\qquad(a_1,a_2)=(3,5).
\]
\item 实际交付合计为 8，不超过容量；总账单为 \(3\cdot8=24\le32\)，保留探索余额 8。
个人支付取零，基金支付分别为 9 与 15，逐笔满足 \(h_i+s_i=pa_i\)。
\item 两家各取得四个真实服务单位。
\(Q(\bar p-c)=4(7/2-1)=10>F=6\)，故各自的正利润价格区间非空。
实际单价 3 给出
\[
R_j=3\cdot4=12,\qquad
C_j=1\cdot4+6=10,\qquad\Pi_j=2>0.
\]
\item 支出前缀的最大累计额为 10；期初工作资金也为 10。
因此每个支出前缀非负。交付验收后到账 12，期末工作现金为 12。
初始资本回收与经济利润分别为 10 和 2。
\item 探索基金支付的 24 与两家收入总和 24 是同一组交易，不形成额外独立收益。
两家的实际生产成本合计 20 单独记录，资金转移与资源耗费保持不同核算性质。
\end{semanticsteps}
\end{proof}

\begin{proposition}[拒绝、退款与有效服务共存的有限状态实例]
\label{joint:service-state-model}
可在上述合同组合中加入一个不满足用途条件的请求和一个允许取消的预存请求，
同时保持八个有效请求的交付和上述成本回收结果。
\end{proposition}
\begin{certificate}
有效请求逐一携带身份、许可用途、版本、输入和配额绑定；实际受控执行对应完整有限任务表。
不合规请求只允许登记拒绝理由与复核，不发生受限业务作用。
取消请求预存 3 后按原标识退款 3，使用探索剩余专款中的暂时预留，不产生交付或额外服务收入。
相关检查与退款操作的资源费用已列于固定成本。
每个实际外部结算都有最终性记录；未完成状态仍保留对应债权。
\end{certificate}
\begin{proof}
\begin{semanticsteps}
\item 逐请求准入表对八个有效请求给出真，对指定越权用途给出假。
实际作用入口只接受当前准入为真的请求，因此拒绝请求不进入受限执行。
\item 取消请求的资金变化为 \(3a+3c=0\)，其中 \(a=(-1,1,0)^\top\)、\(c=-a\)。
预存时可用余额 8 足以支出 3，退款后恢复至 8；所有前缀余额非负。
\item 退款标识与结算标识互斥，该取消请求不进入八个有效请求的销量及收入分母。
所以 \(Q_j=4\)、\(R_j=12\) 和 \(\Pi_j=2\) 均保持。
\item 八个有效请求各采用有限下降策略；取消、拒绝和复核使用独立已预留的处理时段。
调度覆盖证明这些处置不会挤占有效请求的承诺期限。
\end{semanticsteps}
\end{proof}

\begin{proposition}[有限共同任务像与既有经济模块的积构造]
\label{joint:service-product}
将本节完整任务状态记为有限集 \(Q\)，采用有界余额精度、有限请求标识、有限规则版本和离散期限。
两种实现取
\[
X_W=Q\times\{0,1\},\qquad
X_J=Q\times\{0,1,2\},
\]
附加记录位不影响合同观察及后继。
两侧具有同一最粗任务商和观察图链商。
若本节模块与前述劳动、修复和能力模块使用互不重复的资源并具有保持各自期限的共同排程，
则二者的经济结论可以在积状态中共同实现。
\end{proposition}
\begin{certificate}
设 \(K_W=\{0,1\}\)、\(K_J=\{0,1,2\}\)，
共同初态和目标分别为 \(Q_0,G_Q\subseteq Q\)。登记
\[
X_{i0}=Q_0\times K_i,\qquad G_i=G_Q\times K_i,
\qquad A_i(x,z)=A_Q(x),
\]
以及 \(\Xi_i((x,z),a)=\Xi_Q(x,a)\)。
投影 \(q_i(x,z)=x\)，相同动作响应表，以及
\[
\Phi_{i,\ell}(x,z)=(\bar\Phi_\ell(x),z).
\]
共同观察含 \(Q\) 中状态及其带标签边的指标函数。
资金、设备、劳动者时段和预算分别登记在不交资源坐标；
规则读取使用相同版本。积状态另保存同一个真实时间 \(t\)。
某模块执行业务时，另一模块采用该时段内已获验证的等待转移；
等待只保留不应改变的非时间字段，时间及相应派生观察按共同 \(t\) 更新。
共同排程逐项保存两侧等待的允许性、费用、未完任务可实现性与到期目标；
各时段中实际发生的义务和资源耗费均进入相应账本。
无效投入的删除映射固定全部任务和合同收入，且其正值另耗严格正资源。
\end{certificate}
\begin{proof}
\begin{semanticsteps}
\item 对每个 \(i\)，投影满射且满足
\(q_i\Phi_{i,\ell}=\bar\Phi_\ell q_i\)。所有合同观察通过该投影读回。
\item 指标函数区分不同 \(Q\) 状态；同一 \(Q\) 状态的附加记录位不影响任何有限后继观察。
因此动态等价类恰为投影纤维，最粗商为 \(Q\)。
\item 顶点和标签边的线性延拓与边界交换，两侧图链满射至同一模板。
以这些映射为观察，最大不可观察核即其链核，第一同构定理给出共同观察链商。
\item 在积状态 \((x_{\rm old},x_{\rm srv},t)\) 上，原模块作用采用
\[
(x_{\rm old},x_{\rm srv},t)
\longmapsto
(\delta_{\rm old}x_{\rm old},W_{\rm srv}^{t,t'}x_{\rm srv},t'),
\]
其中 \(W_{\rm srv}^{t,t'}\) 是同一时段上的已核验等待转移；服务模块作用对称定义。
主动作用和等待转移分别保持各自 \(I_{\rm old}(x_{\rm old},t)\) 与
\(I_{\rm srv}(x_{\rm srv},t)\)，因此它们的合取在同一 \(t'\) 成立。
\item 对共同排程中的原子次数归纳，双方资源与权利保持。
各自的进展及等待上界再由共同排程保证，使期限与未完任务可实现性同时成立。
资金来源相加只使用不交子账户，因此不会把同一单位资金承诺两次。
\item 删除无效投入保持生产和收入，严格降低成本，因此其最优值为零。
该实例形成一个共同见证，有限经济结论可以合取。
\end{semanticsteps}
\end{proof}


\section{经验验收与结论的适用域}

\begin{definition}[分项经验验收]
每项现实应用分别登记下列可观测量及允许误差：任务覆盖与实际替代；独立运行、期限和成本；
真实需求与人类份额；资格错判、违规漏检及复核结果；逐人应付和到账；资金前缀与剩余承诺；
新增目标的旧排除与新执行；以及共同候选上的成本与收益。
\end{definition}

\begin{proposition}[有误差界时的保守验收]
对欲确认的净量 \(Z=z_1+z_2-z_3\)，若各分量有区间
\(z_1\in[\underline z_1,\overline z_1]\)、\(z_2\in[\underline z_2,\overline z_2]\)、\(z_3\in[\underline z_3,\overline z_3]\)，
则
\[
Z\ge\underline z_1+\underline z_2-\overline z_3.
\]
只有该下界严格为正，才能由这些区间确认 \(Z>0\)。
\end{proposition}
\begin{certificate}
区间来源、共同总体、估计期限、误差口径及联合覆盖条件。若是统计区间，其置信陈述须保持原有联合概率含义。
\end{certificate}
\begin{proof}
\begin{semanticsteps}
\item 以有利项下界和不利项上界逐项替换，得最坏情形下界。
\item 下界不大于零时，区间信息允许 \(Z\le0\)，因而不能由这些区间单独确认严格正号。
\end{semanticsteps}
\end{proof}

\begin{remark}[最终经济含义]
该制度构造把技术承接、用途准入、劳动需求、真实参与和收益修复组织为相容的经济任务。
技术成熟提供阶段转换的可行性；用途覆盖限制指定替代；需求与净反馈决定实际劳动效果；资金和履约把资格转为收入；
任务商保留必要区分；修复操作纠正指定的分配障碍；制度比较确定在声明目标下是否值得采用。
每项结论由自己的证书承担。联合实施要求这些证书在共同资源和真实后继状态上同时成立。
\end{remark}

\appendix
\chapter{证书接口与符号约定}
\label{app:certificates}

\section{联合验收的最小信息字段}

下表列出每类结论必须保留的证书接口。
这些字段不规定唯一的软件实现；相同任务可以采用不同登记载体，只要能够独立恢复必要事实并核验后继保持。
\begingroup
\small
\renewcommand{\arraystretch}{1.18}
\begin{longtable}{p{0.17\textwidth}p{0.41\textwidth}p{0.33\textwidth}}
\toprule
结论对象 & 必须登记的见证 & 可确认的结果\\
\midrule
\endfirsthead
\toprule
结论对象 & 必须登记的见证 & 可确认的结果\\
\midrule
\endhead
阶段成熟 & 构造成果；运行性能；版本迁移；全部维护与研究辅助费用 & 期限内独立承接及特定扩张非必要性\\
用途准入 & 实际任务；控制权限；责任主体；许可；全部执行入口 & 所声明部署用途边界保持\\
劳动保护 & 同口径基线；逐任务劳动量；覆盖；实际规避 & 直接替代上界\\
净劳动变化 & 完整反事实路径；需求、新任务、系数及资源变化；去重归属 & 指定核算域净劳动效应\\
真实参与 & 接纳；实际人类过程；完成；唯一请求权；错误复核 & 真实资格与最终确认的对应\\
参与份额 & 人数上限；基本份额；贡献；差异池；专用资金 & 逐人正应付与总额覆盖\\
实际履约 & 真实义务；到账时点；支付设施；队列；最终性回执 & 按期实际支付与权利余额保持\\
再次接入 & 同一真实末态；逐前缀现金；席位；时间；联合容量 & 有限期自愿参与的可行接续\\
法律激励 & 登记偏离；合法基准；发现与执行概率；可追偿金额 & 指定行为模型中的严格激励约束\\
动态任务商 & 初始观察；全部动作响应；分块见证；稳定转移 & 相对于该任务的最粗充分表示\\
观察链商 & 边界矩阵；观察矩阵；核与像；共同计量；状态接口 & 最粗线性链商与任务读回\\
分配修复 & 逐人缺陷；旧操作分离函数；新增解；合法余额路径 & 指定障碍成为实际可执行边界\\
循环保持 & 阶段图；链收缩；耦合相容；真实路径提升 & 阶段同调保持及有限执行接续\\
能力扩张 & 旧效果全覆盖；新目标排除；新执行；旧路径端点保持 & 指定任务可达集合严格扩张\\
制度最优 & 共同目标；有限候选；完整成本；下界；并列解 & 声明候选与评价口径下的最优性\\
公共签发 & 主体委托；实际用途；版本输入绑定；完整作用入口；拒绝与救济 & 精确准入及有效请求的有限履约\\
劳动权利 & 自愿进退；既得债权；可达复核；可比较贡献；最终支付 & 可兑现报酬与制度权利保持\\
探索机会 & 完整资格路径；核验需求；预算容量；真实交付；对照见证 & 同需同待遇及声明变量的非干扰\\
标准实现 & 合同语义；共同像；动作响应覆盖；链交换与状态读回 & 共同任务商及观察链商\\
服务市场 & 真实订单；支付能力；全部固定与可变成本；逐前缀现金 & 正交付及指定提供者成本回收\\
有效竞争 & 收入对无效投入的完整非干扰；删除可行性；真实服务差异 & 无效投入为零及给定模型的均衡\\
联合实现 & 统一索引；资源总账；排程；权利冲突检查 & 局部证书的共同可行合取\\
\bottomrule
\end{longtable}
\endgroup

\section{主要符号及其局部语义}

同一字母在不同模型中可以承担不同的标准数学角色；每个章节先声明其类型。
跨章联用时，以带类型的对象及证书索引对应，不把字形相同当作对象相同。
\begingroup
\small
\renewcommand{\arraystretch}{1.15}
\begin{longtable}{p{0.28\textwidth}p{0.63\textwidth}}
\toprule
符号 & 含义\\
\midrule
\endfirsthead
\toprule
符号 & 含义\\
\midrule
\endhead
\(\mathfrak D,\nu,T\) & 共同声明域、规则版本、有限期限。\\
\(\mathcal M,X,A,\Xi,\delta\) & 任务系统、状态、允许动作、全部登记响应和后继。\\
\(P,W,V\)（阶段转换章） & 候选生成器、可验证执行器、独立验证器。\\
\(I,G,\rho\) & 不变量或安全域、目标集合、有限完成秩。\\
\(\mathcal U_{\rm s},\mathcal U_{\rm o}\) & 高保障用途集合与受保护常规用途集合。\\
\(D_t^0,D_t^G,A_t^{\rm dir}\) & 基线直接替代、治理后直接替代、已避免的直接替代劳动量。\\
\(\underline\theta_t,\overline\varepsilon_t\) & 有效覆盖比例下界、覆盖部分的实际规避比例上界。\\
\(H_t,J_t^F,N_t^{\rm emp}\) & 实际有酬劳动小时、全时当量、实际就业人数。\\
\(G_t,C_t\)（净劳动分解） & 固定比较路径中的其他正差分与负差分合计。\\
\(\mathcal E_t^\ast,\widehat{\mathcal E}_t\) & 真实有效参与集合与核验程序最终确认集合。\\
\(\beta_t,D_t,\lambda_{it}\)（分配合同） & 基本份额、差异池、归一贡献权重。\\
\(b_{it},y_{it}\) & 逐人应付金额与最终到账金额。\\
\(B_n,L_{ik,n},Y_{ik,n},\Gamma_{ik,n}\) & 可履约资产、未付义务、累计到账、累计成立权利。\\
\(q_0,\tau,q^+,q_\ast\) & 旧结果表示、真实类型、最粗静态细化、最粗动态充分商。\\
\(C_\bullet,d,O_n,N_n\) & 链复形、边界、线性观察、最大不可观察子复形。\\
\(A_-,A_+,B\)（修复复形） & 旧、新修复操作矩阵及资源相容矩阵。\\
\(\omega,\lambda_X,r\) & 指定配置缺陷、读取无资格账户的分离函数、追回方向。\\
\(P,z\)（循环章节） & 四阶段图链复形及其一维循环生成元。\\
\(t,\kappa,E,j,r_E,T_E\) & 修复收缩、耦合、扩展复形、嵌入、回缩、链同伦。\\
\(\partial_a O_k,\mathfrak a,\mathfrak p_j\) & 类型化语义微分、经济作用原子、单步逻辑证书。\\
\(n_i,N,K,B_E,\alpha,a_i\)（探索服务） & 核验需求、需求总量、容量、预算、配给系数、实际配额。\\
\(Q_j,R_j,C_j,\Pi_j\)（服务市场） & 真实交付量、实际收入、完整经济成本、经济利润。\\
\(M,c,\theta,F,v\)（差异化定价） & 连续需求规模、可变成本、真实匹配系数、固定成本、共同服务价值。\\
\(\Sigma_0,\Delta,d_w\) & 制度基线、变更支集、加权变更距离。\\
\(\mathfrak P_{\mathcal G},V,C_{\rm real}\) & 共同目标下候选族、评价函数、真实资源成本。\\
\bottomrule
\end{longtable}
\endgroup

\section{证明依赖的读法}

有限运行安全、版本切换、跨境运输、支付余额和参与接续均由初态及单步条件按执行次数归纳。
任务商的最粗性由因子化的普遍性质证明；有限饱和由分区块数的上界证明。
净劳动与实际收入比较使用同口径恒等式和有证据的参数界。
分配修复使用核、像、秩与对偶见证；合法循环保持使用显式链同伦。
制度最优性使用共同目标下的比较和有限候选上的下界。
公共签发与劳动权利采用逐原子保持；探索机会使用完整路径非干扰及资源约束；
服务市场使用需求、付款、成本和现金前缀；竞争结论使用明确策略空间中的全局最佳反应。

因此，论文中的推理可以按下列类型化顺序核验：
\[
\begin{aligned}
&\text{对象、规则、单位与时间登记}\\
&\quad\Longrightarrow\text{经验前件与规范目标明确}\\
&\quad\Longrightarrow\text{局部形式证书与实际实现对应}\\
&\quad\Longrightarrow\text{共同资源和路径接续检查}\\
&\quad\Longrightarrow\text{声明范围内的联合经济结论}.
\end{aligned}
\]
这里每个箭头都指向文中明确列出的证明义务。对未覆盖任务、未来期限或新增候选，重新登记相应证书后再作判断。

\cleardoublepage
\addcontentsline{toc}{chapter}{参考文献}
\begin{thebibliography}{99}

\bibitem{InternalEconomics}
经济学形式系统汇编（限定内部资料）.
本文采用其中市场参与、参与治理、偏好与权利、规范竞争、生产条件修复及真人参与防作弊的数学接口。
相关定位为基本份额与差异份额覆盖、逐人义务账本、真实退出、最小兼容扩展及有限线性修复的解与分离证书。
各项实际使用的定义和证明在正文中完整重述。

\bibitem{InternalTheory}
任务保持与演化治理理论汇编（限定内部资料）.
本文采用其中《率与商》的精确任务商、有限可达性、最粗动态充分商、联合分型细化及任务保持复用接口。
相关定位包括 \texttt{eq:exact}、\texttt{thm:reach}、
\texttt{thm:saturated-quotient} 和 \texttt{thm:joint-typing-refinement}。
正文按经济任务语义给出相应实例及独立证明。

\bibitem{AcemogluRestrepo2019}
Acemoglu, D., and P. Restrepo.
Automation and New Tasks: How Technology Displaces and Reinstates Labor.
\emph{Journal of Economic Perspectives}, 2019, 33(2): 3--30.
\url{https://doi.org/10.1257/jep.33.2.3}.

\bibitem{Autor2015}
Autor, D. H.
Why Are There Still So Many Jobs? The History and Future of Workplace Automation.
\emph{Journal of Economic Perspectives}, 2015, 29(3): 3--30.
\url{https://doi.org/10.1257/jep.29.3.3}.

\bibitem{Becker1968}
Becker, G. S.
Crime and Punishment: An Economic Approach.
\emph{Journal of Political Economy}, 1968, 76(2): 169--217.
\url{https://doi.org/10.1086/259394}.

\bibitem{Necula1997}
Necula, G. C.
Proof-Carrying Code.
In \emph{Proceedings of the 24th ACM SIGPLAN-SIGACT Symposium on Principles of Programming Languages},
1997: 106--119. ACM.
\url{https://doi.org/10.1145/263699.263712}.
作者公开论文目录：\url{https://people.eecs.berkeley.edu/~necula/papers.html}.

\bibitem{IAEACNS}
International Atomic Energy Agency.
\emph{Convention on Nuclear Safety}.
INFCIRC/449. Vienna: IAEA, 1994.
本文采用第 2--3、5、7--10、14、20 条的对象范围与责任结构。
\url{https://www.iaea.org/sites/default/files/infcirc449.pdf}.

\bibitem{Hatcher2002}
Hatcher, A.
\emph{Algebraic Topology}.
Cambridge University Press, 2002.
第 2 章：链复形、同调、链同伦及正合序列。
\url{https://pi.math.cornell.edu/~hatcher/AT/ATpage.html}.

\bibitem{JakartaPlatform11}
Jakarta EE Platform Team.
\emph{Jakarta EE Platform, Version 11.0}.
规范正文署期：2025-03-10。
采用平台合同、部署授权、事务边界与服务质量规定。
\url{https://jakarta.ee/specifications/platform/11/jakarta-platform-spec-11.0.html}.

\bibitem{JakartaSecurity4}
Jakarta Security Team.
\emph{Jakarta Security, Version 4.0}.
规范正文署期：2024-05-09。
\url{https://jakarta.ee/specifications/security/4.0/jakarta-security-spec-4.0.html}.

\bibitem{JakartaAuthorization3}
Jakarta Authorization Team.
\emph{Jakarta Authorization, Version 3.0}.
规范正文署期：2024-04-07。
采用政策配置、权限判定与执行合同。
\url{https://jakarta.ee/specifications/authorization/3.0/jakarta-authorization-spec-3.0.html}.

\bibitem{JakartaREST4}
Jakarta RESTful Web Services Team.
\emph{Jakarta RESTful Web Services, Version 4.0.0}.
规范正文署期：2024-04-03。
\url{https://jakarta.ee/specifications/restful-ws/4.0/jakarta-restful-ws-spec-4.0.html}.

\bibitem{JakartaTransactions2}
Jakarta Transactions Team.
\emph{Jakarta Transactions, Version 2.0}.
规范正文署期：2020-10-08。
采用资源登记、事务完成及最终状态通知。
\url{https://jakarta.ee/specifications/transactions/2.0/jakarta-transactions-spec-2.0.html}.

\bibitem{JakartaConcurrency31}
Jakarta Concurrency Team.
\emph{Jakarta Concurrency, Version 3.1}.
规范正文署期：2024-06-21。
采用上下文、至多一次任务及事务范围规定。
\url{https://jakarta.ee/specifications/concurrency/3.1/jakarta-concurrency-spec-3.1.html}.

\bibitem{JakartaConnectors21}
Jakarta Connectors Team.
\emph{Jakarta Connectors, Version 2.1}.
规范正文署期：2022-04-17。
采用外部连接、安全与声明事务支持级别。
\url{https://jakarta.ee/specifications/connectors/2.1/jakarta-connectors-spec-2.1.html}.

\bibitem{KatzShapiro1994}
Katz, M. L., and C. Shapiro.
Systems Competition and Network Effects.
\emph{Journal of Economic Perspectives}, 1994, 8(2): 93--115.
\url{https://doi.org/10.1257/jep.8.2.93}.

\bibitem{Hotelling1929}
Hotelling, H.
Stability in Competition.
\emph{The Economic Journal}, 1929, 39(153): 41--57.
\url{https://doi.org/10.2307/2224214}.
本文固定服务商位置，独立证明所列价格子博弈的最佳反应。

\end{thebibliography}
\end{document}
